\section{Foundational literature spine}
\label{sec:literature-spine}

The same phrase ``the Collatz map'' is used for maps with different state
spaces and different clocks.  This section does not normalize those
presentations.  It first records the source symbols and then introduces
project-local subscripts only where two source symbols would otherwise
collide.

\subsection{Three clocks and their exact first-return relation}

Let $\Npos=\{1,2,\ldots\}$ and let $\Opos$ be its odd elements.  The current
Tao source writes $\operatorname{Col}:\Npos\to\Npos$ for the unshortened map
\cite[Sec.~1.1]{Tao2026v7}; here it is displayed as
\begin{equation}
\label{eq:unshortened-U}
 U(n)=
 \begin{cases}
  3n+1,&n\text{ odd},\\
  n/2,&n\text{ even}.
 \end{cases}
\end{equation}
B\"ohm--Sontacchi write $u:\Z\to\Z$, and Bernstein--Lagarias write $T$, for
the shortened map
\begin{equation}
\label{eq:shortened-T}
 T(n)=
 \begin{cases}
  (3n+1)/2,&n\text{ odd},\\
  n/2,&n\text{ even}.
 \end{cases}
\end{equation}
The same branch formula extends continuously to $T:\Ztwo\to\Ztwo$
\cite[Eqs.~(1.1)--(1.3)]{BernsteinLagarias1996}.  Crandall instead writes
$C:D^+\to D^+$, where $D^+=\Opos$, for the odd-to-odd map
\begin{equation}
\label{eq:crandall-C}
 e(x)=\vtwo(3x+1),
 \qquad
 C_{\mathrm{Cr}}(x)=\frac{3x+1}{2^{e(x)}}.
\end{equation}
The subscript in $C_{\mathrm{Cr}}$ is editorial disambiguation; the source
symbol is $C$ \cite[p.~1281, Eq.~(1.1)]{Crandall1978}.  Tao calls this same
odd-to-odd map $\operatorname{Syr}$ \cite[Sec.~1.2,
Eq.~(S-def)]{Tao2026v7}.

\begin{proposition}[the two exact first-return clocks]
\label{prop:first-return-clocks}
For $x\in\Opos$, put $a=e(x)$.  Then $a\geq1$ and
\begin{align}
 T^j(x)&=\frac{3x+1}{2^j} &&(1\leq j\leq a),
 \label{eq:T-first-return-intermediates}\\
 U^{j+1}(x)&=\frac{3x+1}{2^j} &&(0\leq j\leq a).
 \label{eq:U-first-return-intermediates}
\end{align}
The values in \eqref{eq:T-first-return-intermediates} are even for
$1\leq j<a$ and odd for $j=a$.  Consequently
\begin{equation}
\label{eq:exact-first-return}
 C_{\mathrm{Cr}}(x)=T^{e(x)}(x)=U^{e(x)+1}(x),
\end{equation}
and these are respectively the first returns of the shortened and
unshortened orbits to $\Opos$.
\end{proposition}

\begin{proof}
Because $x$ is odd, $3x+1$ is a positive even integer, so
$a=\vtwo(3x+1)\geq1$.  The odd branch of $T$ gives the case $j=1$ of
\eqref{eq:T-first-return-intermediates}.  As long as $j<a$, the quotient has
$2$-adic valuation $a-j>0$, so the next branch of $T$ divides it by $2$.
Induction proves the first identity and shows that no earlier value is odd.
For $U$, its first branch sends $x$ to $3x+1$ and each of the next $a$
branches divides by $2$, proving the second identity and its first-return
claim.  The terminal quotient is exactly \eqref{eq:crandall-C}.
\end{proof}

\begin{nonclaim}
Equation \eqref{eq:exact-first-return} is a variable-time first-return
identity.  It does not identify one iterate of $C_{\mathrm{Cr}}$ with one
iterate of $T$ or $U$, and it does not permit stopping times, periods, or
density statements to be transferred without translating the clock.
\end{nonclaim}

\subsection{Krasikov--Lagarias predecessor counts and the whole-orbit clock}
\label{subsec:KL-predecessor-clock}

Krasikov--Lagarias use exactly the shortened map \(T\) in
\eqref{eq:shortened-T}.  For a fixed positive target
\(a\not\equiv0\pmod3\), they define
\begin{align}
 \pi_a(x)
  &=
  \#\{n\in\Npos: n\le x,\ \exists j\in\Nzero,\ T^j(n)=a\},
  \label{eq:KL-pi}\\
 \pi_a^*(x)
  &=
  \#\{n\in\Npos:n\le x,\ \exists j\in\Nzero,\ T^j(n)=a,\
       T^i(n)\le x\ (0\le i\le j)\}.
  \label{eq:KL-pi-star}
\end{align}
Thus \(\pi_a^*(x)\le\pi_a(x)\).  For \(k\ge2\), a residue
\(m\bmod3^k\) not divisible by \(3\), and \(y\ge0\), the typed auxiliary
function is
\begin{equation}
\label{eq:KL-phi}
 \phi_k^m(y)=
 \inf_{\substack{a\equiv m\pmod{3^k}\\
                  a\text{ not contained in a finite }T\text{-cycle}}}
 \pi_a^*(2^ya).
\end{equation}
The printed exponent \(3^j\) in the congruence is an unbound-index source
defect; Source Issue~\ref{sourceissue:KL-2003-defects} records the
manifestation and the forced repair.

\begin{proposition}[shortened-to-unshortened clock embedding]
\label{prop:KL-Tao-clock-embedding}
Fix \(N\in\Npos\), put \(n_j=T^j(N)\), and define
\(\iota_N:\Nzero\to\Nzero\) by
\begin{equation}
\label{eq:KL-Tao-clock}
 s_0=0,\qquad
 s_{j+1}=s_j+
 \begin{cases}
  1,&n_j\text{ even},\\
  2,&n_j\text{ odd},
 \end{cases}
 \qquad \iota_N(j)=s_j.
\end{equation}
Then
\begin{equation}
\label{eq:KL-Tao-clock-intertwining}
 U^{\iota_N(j)}(N)=T^j(N)\qquad(j\in\Nzero).
\end{equation}
The map \(\iota_N\) is strictly increasing.  Its fibres are singletons on
its image and its equality kernel is the diagonal.  The times outside its
image are exactly \(s_j+1\) for those \(j\) with \(n_j\) odd, and the omitted
state is
\begin{equation}
\label{eq:KL-Tao-omitted-state}
 U^{s_j+1}(N)=3n_j+1\ge4.
\end{equation}
Consequently
\begin{equation}
\label{eq:KL-Tao-reach-one}
 \exists j\in\Nzero,\ T^j(N)=1
 \quad\Longleftrightarrow\quad
 \exists r\in\Nzero,\ U^r(N)=1.
\end{equation}
\end{proposition}

\begin{proof}
For even \(n\), \(T(n)=U(n)\).  For odd \(n\),
\[
 U(n)=3n+1,\qquad U^2(n)=\frac{3n+1}{2}=T(n).
\]
Induction using the two increments in \eqref{eq:KL-Tao-clock} proves
\eqref{eq:KL-Tao-clock-intertwining}.  Every increment is positive, which
proves strict increase, singleton fibres, and the diagonal equality kernel.
An increment of two contributes precisely the intermediate time \(s_j+1\);
an increment of one contributes none.  This proves the asserted complement
of the image and \eqref{eq:KL-Tao-omitted-state}.  Every \(U\)-time is
therefore either an image time or one of those intermediate times.  The
latter cannot carry the value \(1\), while image times carry exactly the
\(T\)-orbit values, proving \eqref{eq:KL-Tao-reach-one}.
\end{proof}

Since every \(U\)-orbit value is a positive integer,
\(\operatorname{Col}_{\min}(N)=1\) if and only if the \(U\)-orbit reaches
\(1\).  Proposition~\ref{prop:KL-Tao-clock-embedding} therefore gives the
exact set identity
\begin{equation}
\label{eq:KL-Tao-count-identity}
 \{N\le x:\operatorname{Col}_{\min}(N)=1\}
 =
 \{N\le x:\exists j\in\Nzero,\ T^j(N)=1\},
\end{equation}
where both sets are subsets of \(\Npos\).  The cardinality on the right is
\(\pi_1(x)\).  The clock map loses the intermediate state after an odd
input, but \eqref{eq:KL-Tao-omitted-state} proves that this information loss
does not affect target-\(1\) reachability.

\begin{theorem}[Krasikov--Lagarias, source theorem]
\label{thm:KL-2003-predecessor-bound}
For every fixed \(a\in\Npos\) with \(a\not\equiv0\pmod3\), there exists
\(x_0(a)\) such that
\begin{equation}
\label{eq:KL-084}
 \pi_a(x)\ge x^{0.84}\qquad(x\ge x_0(a)).
\end{equation}
\end{theorem}

This is Theorem~6.1 of
\cite[p.~254]{KrasikovLagarias2003}.  Its exponent is exactly
\(0.84=21/25\).  The paper reports a computer-found positive feasible
solution of \(L_{11}^{NT}(\lambda)\) at
\(\lambda=1.7922310\).  The exact comparison
\begin{equation}
\label{eq:KL-exponent-margin}
 1792231^{25}-2^{21}10^{150}>0
\end{equation}
proves \(\log_2(1.7922310)>21/25\) without floating-point arithmetic.
The paper does not print the \(k=11\) coordinate vector or an independently
checkable feasibility certificate.  Thus
\eqref{eq:KL-exponent-margin} is independently certified, while the
feasibility premise and Theorem~\ref{thm:KL-2003-predecessor-bound} retain
their status as a published computer-assisted result.

\begin{proposition}[typed final-target handoff]
\label{prop:KL-target-handoff}
Let \(k\ge2\), let \(1<\lambda\le2\), and suppose the explicitly defined
linear program \(L_k^{NT}(\lambda)\) has a feasible solution with principal
coordinates \(c_k^m\).  Put \(\gamma=\log_2\lambda\).  For every
\(0<\beta<\gamma\) and every fixed \(a\in\Npos\) with
\(a\not\equiv0\pmod3\), there exists \(x_0(a)\) such that
\[
 \pi_a(x)\ge x^\beta\qquad(x\ge x_0(a)).
\]
\end{proposition}

\begin{proof}
Theorem~2.2 of the source gives, for \(m\in[3^k]\),
\[
 \phi_k^m(y)\ge
 \Delta_1c_k^m\lambda^y,\qquad
 \Delta_1=\frac{1}{4\max_r c_k^r},\qquad y\ge0.
\]
First suppose that \(a\) is not in a finite \(T\)-cycle and
\(a\equiv2\pmod3\).  With \(m=a\bmod3^k\) and
\(y=\log_2(x/a)\), equations \eqref{eq:KL-pi-star} and
\eqref{eq:KL-phi} give, for \(x\ge a\),
\[
 \pi_a(x)\ge\Delta_1c_k^m a^{-\gamma}x^\gamma.
\]
If instead \(a\equiv1\pmod3\), source equation~(2.1),
\(\phi_k^m(y)=\phi_k^{2m}(y-1)\) for \(y\ge1\), moves to the residue
\(2m\equiv2\pmod3\) and contributes the further fixed factor
\(\lambda^{-1}\).  In either case the positive fixed prefactor dominates
\(x^{\beta-\gamma}\) once \(x\) exceeds an \(a\)-dependent threshold.

Now suppose \(a\) lies in a finite \(T\)-cycle, and let \(M\) be the maximum
on that cycle.  Choose \(r\ge1\) with \(b=2^ra>M\).  Then \(T^r(b)=a\).
If \(b\) lay in a finite cycle, its orbit would reach \(a\), so it would lie
in the same cycle and satisfy \(b\le M\), a contradiction.  The noncyclic
case applies to \(b\), and every orbit reaching \(b\) reaches \(a\).
Therefore \(\pi_a(x)\ge\pi_b(x)\ge x^\beta\) beyond \(x_0(b)\).
\end{proof}

\begin{sourceissue}[Krasikov--Lagarias manifestation and proof-text defects]
\label{sourceissue:KL-2003-defects}
The published paper prints \(3^j\), with \(j\) unbound, in
\eqref{eq:KL-phi}; prints the real variable \(y\), rather than the residue
\(m\), in a congruence modulo \(3\); prints
\(\delta=\beta_2-\beta_1>0\) immediately after
\(\beta_1>\beta_2>\cdots\), although the halting contradiction requires
\(\delta<0\); closes the Theorem~5.1 induction with
\(k\in[3^m]\) rather than \(m\in[3^k]\); reverses the indices in
\(\Delta c_m^k\lambda^y\); and omits a parenthesis from \(T^{(j)}\) in
Theorem~6.1.  The arXiv source additionally prints the two three-lift
exponents in (L4) as \(3^k\), while the journal and the source's own
subsequent minimum use \(3^{k-1}\), and substitutes a tree-vertex macro for
\(\nu\) in the initial interval of Theorem~5.1.  The exact manifestations,
typed repairs, and dependency reasons are recorded in
\texttt{qa/KRASIKOV-LAGARIAS-2003-F029-difference-inequality-audit.md}.
\end{sourceissue}

\begin{nonclaim}
Theorem~\ref{thm:KL-2003-predecessor-bound} does not prove the Collatz
conjecture, positive density, density one, a target-uniform threshold, or the
claim for targets divisible by \(3\).  Since \(x^{0.84}/x\to0\), it is not an
\emph{almost all} theorem.  The paper does not prove the stronger rounded
exponent \(0.8417560\), strict increase or convergence to \(2\) of its
optimal \(\lambda_k\), or optimality of the difference-inequality method.
\end{nonclaim}

\subsection{Terras--Allouche--Korec: exact finite coordinates and
density-one power descent}
\label{subsec:Terras-Allouche-Korec-density}

The four sources in this subsection use the shortened map
\[
 T(y)=
 \begin{cases}
  (3y+1)/2,&y\equiv1\pmod2,\\
  y/2,&y\equiv0\pmod2,
 \end{cases}
\]
which is exactly \eqref{eq:shortened-T}.  It is not one iterate of the
unshortened map \(U\) on an odd input.  The transfer to the \(U\)-orbit is
made only after the source results have been stated and audited in their own
clock.

\subsubsection{Terras's parity coordinate}

Terras works on \(\Nzero\).  Put
\begin{align}
 X(y)&=
 \begin{cases}1,&y\text{ odd},\\0,&y\text{ even},\end{cases}
 &
 X_i(y)&=X(T^i(y)),                                      \label{eq:Terras-X}\\
 S_k(y)&=\sum_{i=0}^{k-1}X_i(y),
 &
 \lambda_k(y)&=2^{-k}3^{S_k(y)},                         \label{eq:Terras-lambda}\\
 E_k(y)&=(X_0(y),\ldots,X_{k-1}(y))\in\{0,1\}^k.         \label{eq:Terras-E}
\end{align}
Terras's Theorem~1.1 gives the exact affine expression
\begin{equation}
\label{eq:Terras-affine}
 T^k(y)=\lambda_k(y)y+\rho_k(y),
\qquad
 \rho_k(y)=\frac{\lambda_k(y)}2
 \sum_{i=0}^{k-1}X_i(y)\lambda_{i+1}(y)^{-1}.
\end{equation}
In particular \(\rho_k(y)\ge0\); it is not an unspecified error term
\cite[pp.~241--242, Thm.~1.1]{Terras1976}.

\begin{proposition}[Terras's exact parity coordinate]
\label{prop:Terras-parity-coordinate}
For every \(k\in\Nzero\) and \(x,y\in\Nzero\),
\begin{equation}
\label{eq:Terras-fibre}
 E_k(x)=E_k(y)
 \quad\Longleftrightarrow\quad
 x\equiv y\pmod{2^k}.
\end{equation}
Consequently \(E_k\) induces a bijection
\[
 \Z/2^k\Z\xrightarrow{\;\sim\;}\{0,1\}^k.
\]
Thus, on every complete interval of \(2^k\) consecutive integers, each
chronological parity word occurs exactly once.
\end{proposition}

\begin{proof}
The assertion is trivial for \(k=0\).  Suppose it is known for \(k\), and
write \(\epsilon=X(x)=X(y)\).  On the common parity branch,
\begin{equation}
\label{eq:Terras-one-step-difference}
 T(x)-T(y)=\frac{3^\epsilon(x-y)}2.
\end{equation}
If \(x\equiv y\pmod{2^{k+1}}\), then \(x,y\) have the same parity and
\eqref{eq:Terras-one-step-difference} gives
\(T(x)\equiv T(y)\pmod{2^k}\).  The inductive hypothesis applied to their
next \(k\) iterates proves \(E_{k+1}(x)=E_{k+1}(y)\).

Conversely, if \(E_{k+1}(x)=E_{k+1}(y)\), then \(x,y\) have the same
\(\epsilon\), and the final \(k\) coordinates give
\(E_k(T(x))=E_k(T(y))\).  Hence \(T(x)\equiv T(y)\pmod{2^k}\).
Equation~\eqref{eq:Terras-one-step-difference} and the oddness of
\(3^\epsilon\) imply \(2^{k+1}\mid x-y\).  This proves
\eqref{eq:Terras-fibre} by induction.  Both finite sets in the induced map
have \(2^k\) elements, so the injection is a bijection.  Translating any
complete interval to a complete residue system proves the last assertion.
\end{proof}

This proposition is the arithmetic content behind the independent
Bernoulli coordinates in Terras's Corollaries~1.3--1.4.  No random surrogate
has replaced the orbit: probabilities below are normalized, exact counts on
\(\Z/2^k\Z\).

\subsubsection{Coefficient stopping and actual stopping}

Terras defines the two positive-time stopping functions
\begin{align}
 \chi(y)&=\min\{j\in\Npos:T^j(y)<y\},                    \label{eq:Terras-chi}\\
 \tau(y)&=\min\{j\in\Npos:\lambda_j(y)<1\},              \label{eq:Terras-tau}
\end{align}
with value \(\infty\) if the relevant set is empty.  Time zero is excluded.
Since \(\rho_j(y)\ge0\), \eqref{eq:Terras-affine} gives
\(\tau(y)\le\chi(y)\).  If \(\tau(y)=j\), Terras's Theorem~1.9 proves that
all sufficiently large \(z\equiv y\pmod{2^j}\) satisfy \(\chi(z)=j\).
It follows that, for each fixed \(j\), the two sets
\(\{\tau=j\}\) and \(\{\chi=j\}\) differ by only finitely many integers
\cite[pp.~243--244]{Terras1976}.

For \(k\ge1\), define the finite-coordinate tail
\begin{equation}
\label{eq:Terras-Fk}
 F(k)=\mathbb P[\tau\ge k]
 =2^{-k}\#\{1\le y\le2^k:\tau(y)\ge k\}.
\end{equation}
The final equality is independent of the chosen complete residue block by
Proposition~\ref{prop:Terras-parity-coordinate}.

\begin{theorem}[Terras 1976 coefficient-tail theorem, with the printed
support bound repaired]
\label{thm:Terras-1976-coefficient-tail}
The sequence \(F(k)\) is nonincreasing and
\[
 \lim_{k\to\infty}F(k)=0.
\]
More exactly, if
\(\gamma=\log 2/\log 3\), then every length-\(k\) parity word counted by
\(F(k)\) has at most
\begin{equation}
\label{eq:Terras-repaired-zero-cutoff}
 q_k=\left\lfloor (k-1)(1-\gamma)\right\rfloor+1
\end{equation}
zero coordinates.  Hence
\begin{equation}
\label{eq:Terras-binomial-tail}
 F(k)\le 2^{-k}\sum_{a=0}^{q_k}\binom{k}{a}
 \xrightarrow[k\to\infty]{}0.
\end{equation}
\end{theorem}

\begin{proof}
The tail events \(\{\tau\ge k+1\}\subseteq\{\tau\ge k\}\) prove
monotonicity.  Consider a word counted by \(F(k)\), and let \(b'\) be the
number of ones among its first \(k-1\) coordinates.  Since no coefficient
drop has occurred before time \(k\),
\[
 \lambda_{k-1}=\frac{3^{b'}}{2^{k-1}}\ge1.
\]
For \(k>1\) equality is impossible because a positive power of \(2\) is not
a power of \(3\); thus \(b'>(k-1)\gamma\).  The first \(k-1\) positions
therefore contain fewer than \((k-1)(1-\gamma)\) zeros, and the last
coordinate contributes at most one more.  Integrality gives
\eqref{eq:Terras-repaired-zero-cutoff}.

Proposition~\ref{prop:Terras-parity-coordinate} turns the word count into
\eqref{eq:Terras-binomial-tail}.  Finally
\[
 3^2>2^3\quad\Longrightarrow\quad
 \gamma>\frac12
 \quad\Longrightarrow\quad
 \frac{q_k}{k}\longrightarrow1-\gamma<\frac12.
\]
The lower binomial tail in \eqref{eq:Terras-binomial-tail} therefore tends
to zero, by the standard de Moivre--Laplace estimate used in the source (or
by any equivalent fixed-distance binomial tail bound).
\end{proof}

\begin{sourceissue}[Terras 1976 proof-text defects and forced local repairs]
\label{sourceissue:Terras-1976-defects}
Five source distinctions are required.

\begin{enumerate}
\item Proposition~1.7 prints, without a sign hypothesis,
\[
 T^k(y)<y
 \quad\Longleftrightarrow\quad
 \frac{\rho_k(y)}{1-\lambda_k(y)}<y.
\]
It is false as printed: for \(y=1,k=1\),
\[
 T(1)=2\not<1,\qquad
 \lambda_1(1)=\frac32,\qquad
 \rho_1(1)=\frac12,
\]
but \(\rho_1(1)/(1-\lambda_1(1))=-1<1\).  Division is valid only under
\(\lambda_k(y)<1\).  Theorem~1.9 invokes it under \(\tau(y)=k\), which
supplies precisely that missing hypothesis.

\item The proof of Theorem~1.9 cites Proposition~1.9 for
\(\tau\le\chi\).  The relevant result is Proposition~1.8.

\item Definition~1.12 prints the proper-prefix range \(1\le i\le k-2\).
Formula~(9), the stopping event, and the subsequent recurrence require every
proper nonempty prefix, \(1\le i\le k-1\).  The literal range already loses
the decisive prefix when \(k=2\).

\item In Theorem~1.17, a terminal word of length \(k\), with \(a\) zeros and
\(b=k-a\) ones, has last digit zero and an active length-\(k-1\) prefix.
Therefore
\[
 \frac{3^b}{2^{k-1}}>1,\qquad
 a<1+(k-1)(1-\gamma).
\]
The source drops \(+1\), and then uses
\(a\le\lfloor k(1-\gamma)\rfloor\).  The terminal word \((1,0)\) at
\(k=2\) has \(a=1\), while the printed cutoff is \(0\).  The repaired
cutoff is exactly \eqref{eq:Terras-repaired-zero-cutoff}; its ratio to \(k\)
has the same strict limit below \(1/2\), so the theorem's limiting conclusion
survives.

\item Definition~1.10 introduces point masses and lower-tail sums before
Theorem~1.11 uses \(\mathbb P[\tau\ge k]\); formula~(6) supplies the direct
complete-block definition only afterwards.  Equation~\eqref{eq:Terras-Fk}
states the tail on its finite block before use.
\end{enumerate}
These are local repairs.  None is treated as text literally printed by
Terras \cite[pp.~243--247]{Terras1976}.
\end{sourceissue}

\begin{theorem}[Terras 1979 finite-complement density identity]
\label{thm:Terras-1979-density-identity}
For every \(k\in\Npos\), the natural density exists and satisfies
\begin{equation}
\label{eq:Terras-1979-density}
 \delta\{y\in\Npos:\chi(y)\ge k\}
 =\mathbb P[\tau\ge k]=F(k).
\end{equation}
Both source inequalities are non-strict \(\ge\).
\end{theorem}

\begin{proof}
For each fixed \(j\), the 1976 coset theorem recalled above says that
\(\{\chi=j\}\) and \(\{\tau=j\}\) differ by a finite set.  The second set is
a union of residue classes modulo \(2^j\); hence
\[
 \delta\{\chi=j\}=\mathbb P[\tau=j].
\]
Only the finite set of indices \(1\le j<k\) is now summed:
\[
 \delta\{\chi<k\}
 =\sum_{j=1}^{k-1}\mathbb P[\tau=j]
 =\mathbb P[\tau<k].
\]
Taking the complements in \(\Npos\) and in the finite probability space
\(\Z/2^k\Z\) gives \eqref{eq:Terras-1979-density}.  This is the explicit
finite-complement closure supplied by the 1979 note; it does not interchange
natural density with an infinite sum
\cite[pp.~101--102]{Terras1979}.
\end{proof}

\begin{corollary}[Terras density-one finite stopping]
\label{cor:Terras-density-one-stopping}
The set \(\{y\in\Npos:\chi(y)<\infty\}\) has natural density one.
\end{corollary}

\begin{proof}
For every \(k\),
\[
 \{\chi=\infty\}\subseteq\{\chi\ge k\}.
\]
Theorem~\ref{thm:Terras-1979-density-identity} and
Theorem~\ref{thm:Terras-1976-coefficient-tail} therefore give
\[
 \overline\delta\{\chi=\infty\}
 \le F(k)\xrightarrow[k\to\infty]{}0.
\]
The exceptional set has upper density zero, so its complement has natural
density one.
\end{proof}

\begin{nonclaim}
Corollary~\ref{cor:Terras-density-one-stopping} leaves open a zero-density
exceptional set and does not prove the Collatz conjecture.  It is a theorem
about first descent below the starting value, not eventual arrival at \(1\).
The 1979 argument proves existence of the displayed natural densities by a
finite complement; it supplies no license to sum arbitrary density limits.
\end{nonclaim}

\subsubsection{Allouche's unnormalized residue system}

Allouche begins with positive integers \(n,d\), with \(d\ge2\) and
\(\gcd(n,d)=1\), and a chosen complete representative system
\begin{equation}
\label{eq:Allouche-Ad}
 A_d=\{0,r_1,\ldots,r_{d-1}\}\subset\Z
\end{equation}
modulo \(d\).  Let \(\varphi:\Z\to A_d\) send an integer to its chosen
representative.  The map is
\begin{equation}
\label{eq:Allouche-g}
 g(\ell)=
 \begin{cases}
  \ell/d,&\ell\equiv0\pmod d,\\[1mm]
  (n\ell-\varphi(n\ell))/d,&\ell\not\equiv0\pmod d.
 \end{cases}
\end{equation}
After the introduction, the source assumes \(n>d\ge2\).  The data
\((n,d,A_d,\varphi)\) are retained: in the classical specialization
\[
 n=3,\qquad d=2,\qquad A_2=\{0,-1\},
\]
\eqref{eq:Allouche-g} is the shortened Collatz map on \(\Z\).  Replacing the
source representative \(-1\) by \(1\) would alter the affine remainder and
is not done.

For \(k\ge1\), set
\begin{align}
 \alpha(k,j)
  &=\#\{0\le v<k:g^v(j)\not\equiv0\pmod d\},              \label{eq:Allouche-alpha}\\
 P_k(j)
  &=g^k(j)-\frac{n^{\alpha(k,j)}}{d^k}j.                 \label{eq:Allouche-P}
\end{align}
Also put
\begin{equation}
\label{eq:Allouche-M}
 R=\max_{a\in A_d}|a|,
 \qquad
 M=\frac{R}{n-d}.
\end{equation}

\begin{theorem}[Allouche's affine coordinate and exact branch counts]
\label{thm:Allouche-affine-counting}
For \(k\ge1\), \(j,\lambda\in\Z\), and \(0\le i\le k\),
\begin{align}
 g^k(d^k\lambda+j)
  &=n^{\alpha(k,j)}\lambda+g^k(j),                        \label{eq:Allouche-affine}\\
 \alpha(k,d^k\lambda+j)&=\alpha(k,j),                     \label{eq:Allouche-alpha-fibre}\\
 |P_k(j)|&\le M\frac{n^k}{d^k},                           \label{eq:Allouche-remainder}\\
 \#\{0\le j<d^k:\alpha(k,j)=i\}
  &=\binom{k}{i}(d-1)^i.                                 \label{eq:Allouche-count}
\end{align}
Thus \(j\bmod d^k\) is the residue-fibre coordinate, \(\lambda\) is its
integral quotient coordinate, and the affine multiplier is constant on the
fibre.

More quantitatively, let \(A<B\) be rational numbers, neither equal to any
\(n^\mu/d^\nu\) with \(\mu\in\Nzero\), \(\nu\in\Npos\).  Let \(C\) be the
maximum of the absolute denominators of \(A,B\), and
\(\vartheta=\mathbf 1_{(A,B)}\).  Then the finite-coordinate comparison used
in Allouche's Theorem~1 is
\begin{multline}
\label{eq:Allouche-quantitative-handoff}
 \left|
 \sum_{1\le\ell\le x}
  \vartheta\left(\frac{g^k(\ell)}{\ell}\right)
 -\frac{x}{d^k}\sum_{i=0}^k
  \binom{k}{i}(d-1)^i
  \vartheta\left(\frac{n^i}{d^k}\right)
 \right|\\
 \le \frac52d^k+MCn^k.
\end{multline}
For every \(\varepsilon>0\), there are \(k_0\) and
\(\eta\in(0,1)\), depending only on the fixed source data and
\(\varepsilon\), such that
\begin{equation}
\label{eq:Allouche-binomial-concentration}
 \frac1{d^k}
 \sum_{\left|i-\frac{d-1}{d}k\right|>\varepsilon k}
 \binom{k}{i}(d-1)^i
 \le\eta^k
 \qquad(k\ge k_0).
\end{equation}
\end{theorem}

\begin{proof}
For \(a\in\Z\), \(t\in\Z\), and
\(\epsilon=1_{\{a\not\equiv0\pmod d\}}\), the representative is unchanged
after adding \(dt\), and direct substitution in \eqref{eq:Allouche-g}
gives
\begin{equation}
\label{eq:Allouche-triangular-step}
 g(a+dt)=g(a)+n^\epsilon t.
\end{equation}
Apply this first with \(t=d^{k-1}\lambda\), and then induct on \(k\).
The first branch symbol is unchanged, the next \(k-1\) symbols are those of
\(g(j)\), and their multipliers multiply.  This proves
\eqref{eq:Allouche-affine} and \eqref{eq:Allouche-alpha-fibre}.

At a nondivisible step the additive correction in
\eqref{eq:Allouche-g} has absolute value at most \(R/d\); at a divisible
step it is zero.  Expanding \(k\) iterates and bounding every later branch
multiplier by \(n/d\) gives
\[
 |P_k(j)|
 \le\frac Rd\sum_{r=0}^{k-1}\left(\frac nd\right)^r
 =\frac{R}{n-d}\left(\frac{n^k}{d^k}-1\right)
 \le M\frac{n^k}{d^k},
\]
which proves \eqref{eq:Allouche-remainder}.

It remains to count branches without replacing them by a stochastic model.
The map
\[
 j\bmod d^k
 \longmapsto
 \bigl(j\bmod d,\ g(j)\bmod d^{k-1}\bigr)
\]
is a bijection.  Indeed, for first residue \(a=0\), write \(j=dq\), so
\(g(j)=q\), which determines \(q\bmod d^{k-1}\).  For \(a\ne0\), write
\(j=a+dq\); then
\[
 g(j)=\frac{na-\varphi(na)}d+nq,
\]
and \(\gcd(n,d)=1\) makes \(q\bmod d^{k-1}\) unique.  There is one divisible
first residue and \(d-1\) nondivisible ones.  Induction therefore gives
\[
 \binom{k-1}{i}(d-1)^i+
 (d-1)\binom{k-1}{i-1}(d-1)^{i-1}
 =\binom{k}{i}(d-1)^i,
\]
proving \eqref{eq:Allouche-count}.  Finally,
\eqref{eq:Allouche-quantitative-handoff} is the exact proposition on
pp.~9-08--9-10 of the source, obtained by partitioning the summation into
residue fibres, using \eqref{eq:Allouche-affine}--\eqref{eq:Allouche-remainder},
and bounding the two rational interval endpoints.  Equation
\eqref{eq:Allouche-binomial-concentration} is source Lemma~4
\cite[pp.~9-04--9-11]{Allouche1979}.
\end{proof}

For completeness, the hypotheses actually consumed by Allouche's
Theorem~1 are recorded rather than compressed into the phrase
\emph{a large-deviation argument}.  Given \(a>1\), put
\[
 k=\left\lfloor\frac{\log x}{a\log n}\right\rfloor,
 \qquad
 \mathcal F(x)=
 \sum_{1\le\ell\le x}
 \vartheta\left(\frac{g^k(\ell)}{\ell}\right).
\]
The theorem has three separate cases:
\begin{enumerate}
\item if \(A<B<0\), then \(\mathcal F(x)=O(Cx^{1/a})\);
\item if \(A<B\), \(B>0\), and
\[
 \frac{d-1}{d}-\frac{\log d}{\log n}
 -\frac{\log B}{k\log n}\ge\varepsilon,
\]
then, for \(k\ge k_0(\varepsilon)\),
\[
 \mathcal F(x)=
 O\left(Cx^{1/a}+x^{1-\frac{|\log\eta|}{a\log n}}\right);
\]
\item if \(0<A<B\) and
\[
 \frac{\log d}{\log n}+\frac{\log A}{k\log n}
 -\frac{d-1}{d}\ge\varepsilon,
\]
the same estimate holds.
\end{enumerate}
The implied constants depend on \(n,d,A_d\).  These are distinct
hypothesis packages; no one of them is silently substituted for another.

\begin{theorem}[Allouche's branch-slope obstruction, with the displayed
domain completed]
\label{thm:Allouche-slope-obstruction}
Let \(d<n\) be coprime positive integers.  After the local completion below,
there are functions \(F,G:\Z\to\Z\) satisfying
\begin{align*}
 F(\ell)&=\ell/d+O(1),&G(\ell)&=\ell/d+O(1)
 &&(\ell\equiv0\pmod d),\\
 F(\ell)&=n\ell/d+O(1),&G(\ell)&=n\ell/d+O(1)
 &&(\ell\not\equiv0\pmod d),
\end{align*}
such that every \(F\)-orbit is eventually periodic and only finitely many
periods occur, whereas
\[
 |G^k(\ell)|\longrightarrow\infty
\]
for every \(\ell\) outside a finite set.
\end{theorem}

\begin{proof}[Exact source audit and independent domain repair]
Allouche chooses \(u\) with \(d<n<d^u\).  The displayed divisible branch
of \(F\) is specified only for
\[
 \ell\equiv d^i\pmod{d^u},\qquad1\le i\le u.
\]
These classes do not exhaust the multiples of \(d\): for \(d=2,u=3\), the
class \(\ell\equiv6\pmod8\) is missing.  Thus the printed cases do not
literally define the announced function \(F:\Z\to\Z\).

For every \(\ell\equiv0\pmod d\), let \(\rho_\ell\) be the unique integer
with
\begin{equation}
\label{eq:Allouche-rho-repair}
 0\le\rho_\ell<d^{u-1},
 \qquad
 \frac{\ell}{d}+\rho_\ell\equiv0\pmod{d^{u-1}},
\end{equation}
and define the completed divisible branch by
\begin{equation}
\label{eq:Allouche-F-repair}
 F(\ell)=\frac{\ell}{d}+\rho_\ell.
\end{equation}
Existence and uniqueness are the unique representative of
\(-\ell/d\bmod d^{u-1}\) in \([0,d^{u-1})\).  On a source class
\(\ell\equiv d^i\pmod{d^u}\), it gives
\[
 \rho_\ell=
 \begin{cases}
  d^{u-1}-d^{i-1},&i<u,\\
  0,&i=u,
 \end{cases}
\]
so \eqref{eq:Allouche-F-repair} agrees exactly with every printed divisible
case.  It also proves, on every formerly missing divisible residue,
\begin{equation}
\label{eq:Allouche-repair-properties}
 F(\ell)\equiv0\pmod{d^{u-1}},
 \qquad
 \left|F(\ell)-\frac{\ell}{d}\right|<d^{u-1}.
\end{equation}
Retain the printed nondivisible branches unchanged.  The remainder of the
source proof uses the divisible branch only through
\eqref{eq:Allouche-repair-properties}: pp.~9-12--9-14 obtain a uniform bound
for the \(F\)-iterates, set \(G=F+1\), obtain an expanding lower bound for
the \(G\)-iterates, and use finite predecessor fibres to isolate the finite
exceptional set.  Therefore that proof applies to the completed map
\cite[pp.~9-04, 9-12--9-14]{Allouche1979}.  Formula
\eqref{eq:Allouche-F-repair} is an edition reconstruction, not a formula
silently attributed to the printed source.
\end{proof}

\begin{proposition}[Allouche's exact classical threshold extracted from
the affine and counting lemmas]
\label{prop:Allouche-threshold-extraction}
Put
\begin{equation}
\label{eq:Allouche-cA}
 c_A=\frac32-\log_3 2.
\end{equation}
For every real \(c>c_A\), the set
\begin{equation}
\label{eq:Allouche-Mc}
 \mathcal M_c^{T}
 =\{y\in\Npos:\exists r\in\Nzero,\ T^r(y)<y^c\}
\end{equation}
has natural density one.
\end{proposition}

\begin{proof}
In the classical specialization of
Theorem~\ref{thm:Allouche-affine-counting}, write
\[
 A_m(y)=\alpha(m,y)=\sum_{i=0}^{m-1}X_i(y).
\]
Equations \eqref{eq:Allouche-P} and \eqref{eq:Allouche-remainder}, with
\(M=1\), give
\begin{equation}
\label{eq:Allouche-classical-affine}
 T^m(y)=\frac{3^{A_m(y)}}{2^m}y+P_m(y),
 \qquad
 |P_m(y)|\le\left(\frac32\right)^m.
\end{equation}

Choose \(\varepsilon>0\) with \(c_A+\varepsilon<c\).  On the
ternary-height block \(3^m\le y<3^{m+1}\), call \(y\) good when
\[
 A_m(y)\le\left(\frac12+\varepsilon\right)m.
\]
For a good \(y\),
\[
 \frac{3^{A_m(y)}}{2^m}y
 \le y\left(\frac{3^{1/2+\varepsilon}}2\right)^m.
\]
If \(q=3^{1/2+\varepsilon}/2\), the ratio
\(y/3^m\in[1,3)\) gives the uniform comparison
\[
 q^m\le 3^{|\log_3q|}y^{\log_3q}.
\]
Since \(1+\log_3q=c_A+\varepsilon\), the main term is at most
\[
 C_\varepsilon y^{c_A+\varepsilon},
 \qquad C_\varepsilon=3^{|\log_3q|}.
\]
The remainder has strictly smaller exponent:
\[
 \left(\frac32\right)^m
 \le y^{\log_3(3/2)},
 \qquad
 \log_3(3/2)=c_A-\frac12<c_A.
\]
Hence \(T^m(y)<y^c\) for all sufficiently large good \(y\).

By \eqref{eq:Allouche-count}, every complete residue block of length \(2^m\)
contains exactly
\[
 \sum_{i>(1/2+\varepsilon)m}\binom mi
\]
bad residues.  Its normalized value tends to zero by
\eqref{eq:Allouche-binomial-concentration}.  The interval
\([3^m,3^{m+1})\) is a union of complete \(2^m\)-blocks plus at most two end
fragments, of total length below \(2^{m+1}\).  Their relative contribution is
\[
 O\left(\frac{2^m}{3^m}\right)=O((2/3)^m).
\]
Thus the proportion outside \(\mathcal M_c^T\) in each ternary-height block
tends to zero.  Summing these geometrically increasing blocks proves that
the exceptional proportion up to an arbitrary \(x\) tends to zero: after a
fixed initial segment, each block error is at most
\(\epsilon_0 3^m\), and \(\sum_{m\le M}3^m<\tfrac32\,3^M\), while
\(x\ge3^M\) on the last block.  Since \(\epsilon_0\) is arbitrary,
\(\delta(\mathcal M_c^T)=1\).
\end{proof}

\begin{nonclaim}
Proposition~\ref{prop:Allouche-threshold-extraction} is an edition
reconstruction from Allouche's printed Lemmas~1--4 and quantitative
proposition; the decimal exponent is not printed there as a standalone
theorem.  Korec explicitly records
\(3/2-\log_3 2=0.86907\ldots\) as the Allouche threshold
\cite[p.~86]{Korec1994}.  The endpoint \(c=c_A\) is not included.
Theorem~\ref{thm:Allouche-slope-obstruction} shows, after its explicit domain
repair, that the two leading branch slopes and bounded additive errors do
not determine global orbit behaviour; it is not a result about the endpoint
of the classical Collatz map.
\end{nonclaim}

\subsubsection{Korec's improved power threshold}

For real \(c\), Korec defines
\begin{equation}
\label{eq:Korec-Mc}
 M_c=\{y\in\Nzero:\exists r\in\Nzero,\ T^r(y)<y^c\}.
\end{equation}
With \(X_i,E_m,S_m\) as in
\eqref{eq:Terras-X}--\eqref{eq:Terras-E}, he also defines, with a
non-strict boundary,
\begin{equation}
\label{eq:Korec-Umd}
 U(m,d)=\#\{0\le y<2^m:S_m(y)\le md\}.
\end{equation}
Proposition~\ref{prop:Terras-parity-coordinate} gives the exact complete-block
identity
\[
 U(m,d)=\sum_{0\le i\le md}\binom mi,
\]
where an empty or nonintegral endpoint is interpreted by the displayed
integer condition.  Korec's Lemma~2 consequently gives
\begin{equation}
\label{eq:Korec-CLT}
 d>\frac12
 \quad\Longrightarrow\quad
 \frac{U(m,d)}{2^m}\longrightarrow1.
\end{equation}

\begin{theorem}[Korec 1994 density estimate, with three local proof-text
repairs]
\label{thm:Korec-density-estimate}
If
\[
 c>\log_4 3,
\]
then \(M_c\) has natural density one.  The inequality is strict.
\end{theorem}

\begin{proof}
It suffices to prove the assertion for
\(\log_4 3<c<1\): if \(c\ge1\), choose
\(\log_4 3<c'<1\), and \(M_{c'}\subseteq M_c\) away from the immaterial
inputs \(0,1\).

For a large cutoff \(a\), choose the least positive integer \(m\) such that
\begin{equation}
\label{eq:Korec-m-choice}
 a\le m^2 2^m.
\end{equation}
Define the repaired parameter
\begin{equation}
\label{eq:Korec-dstar}
 d_*=\frac12\left(\frac{c}{\log_2 3}+\frac12\right).
\end{equation}
Since
\[
 c>\log_4 3=\frac12\log_2 3,
\]
the two endpoints whose arithmetic mean occurs in
\eqref{eq:Korec-dstar} satisfy
\[
 \frac12<\frac{c}{\log_2 3};
\]
hence
\begin{equation}
\label{eq:Korec-dstar-interval}
 \frac12<d_*<\frac{c}{\log_2 3}.
\end{equation}

Consider \(m2^m\le y<a\), and put \(k=S_m(y)\).  For \(0\le p<m\),
every branch of \(T\) divides its input by at most \(2\), so
\[
 T^p(y)\ge\frac{y}{2^p}>m.
\]
At each of the \(k\) odd steps the exact ratio is
\((3T^p(y)+1)/(2T^p(y))\), while at each even step it is \(1/2\).
Therefore
\begin{align}
 T^m(y)
 &<y\left(\frac{3m+1}{2m}\right)^k
       \left(\frac12\right)^{m-k} \notag\\
 &=y\frac{3^k}{2^m}
       \left(1+\frac1{3m}\right)^k \notag\\
 &\le y\frac{3^k}{2^m}
       \left(1+\frac1{3m}\right)^m
 <y\frac{3^k}{2^{m-1}}.                                \label{eq:Korec-product}
\end{align}
The final inequality uses
\((1+1/(3m))^m<e^{1/3}<2\).

Because \(c<1\), equations
\eqref{eq:Korec-m-choice} and \eqref{eq:Korec-product} show that
\(T^m(y)<y^c\) whenever
\begin{equation}
\label{eq:Korec-sufficient-product}
 (m^2 2^m)^{1-c}3^k\le2^{m-1}.
\end{equation}
Taking base-two logarithms and preserving every denominator proves the exact
equivalence
\begin{equation}
\label{eq:Korec-exact-k-condition}
 \frac{k}{m}
 \le
 \frac{c}{\log_2 3}
 -\frac{1+2(1-c)\log_2m}{m\log_2 3}.
\end{equation}
The right side tends to \(c/\log_2 3>d_*\).  Hence, for all sufficiently
large \(m\),
\[
 S_m(y)\le md_*
 \quad\Longrightarrow\quad
 T^m(y)<y^c.
\]
The non-strict boundary is the one counted in \eqref{eq:Korec-Umd}.

By \eqref{eq:Korec-CLT} and \eqref{eq:Korec-dstar-interval}, the proportion
of such \(y\) in every complete interval of \(2^m\) consecutive integers
tends to one.  The minimality in \eqref{eq:Korec-m-choice} gives
\[
 a>(m-1)^2 2^{m-1}.
\]
Consequently the discarded initial interval below \(m2^m\) and the single
terminal fragment shorter than \(2^m\) have relative sizes bounded by
\[
 \frac{m2^m}{a}<\frac{2m}{(m-1)^2}\longrightarrow0,
 \qquad
 \frac{2^m}{a}<\frac{2}{(m-1)^2}\longrightarrow0.
\]
All remaining points lie in complete \(2^m\)-blocks, where the exceptional
proportion is
\[
 1-\frac{U(m,d_*)}{2^m}\longrightarrow0.
\]
Thus \(\#(M_c\cap[0,a))/a\to1\), proving the theorem.
\end{proof}

\begin{sourceissue}[Korec 1994: the three printed proof defects]
\label{sourceissue:Korec-1994-defects}
The theorem above is the published conclusion, but its printed proof
requires three visible repairs \cite[pp.~86--88]{Korec1994}.
\begin{enumerate}
\item The source sets
\[
 d=\frac12\left(\frac{c}{\log_4 3}+\frac12\right)
\]
and immediately asserts \(d<c/\log_2 3\).  Since
\[
 \frac12\left(\frac{c}{\log_4 3}+\frac12\right)
 =\frac{c}{\log_2 3}+\frac14,
\]
the assertion is false.  Equation~\eqref{eq:Korec-dstar} replaces the
denominator \(\log_4 3\) by the \(\log_2 3\) required by the ensuing
comparison and proves the claimed open interval.

\item The source calls
\[
 \frac{k}{m}\le
 \frac{c}{\log_2 3}
 -\frac{1+2(1-c)\log_2m}{m}
\]
equivalent to \eqref{eq:Korec-sufficient-product}.  Exact rearrangement gives
\eqref{eq:Korec-exact-k-condition}, with an additional
\(\log_2 3\) in the final denominator.  Because the numerator is positive
and \(\log_2 3>1\), the printed condition subtracts more; it remains a
stronger sufficient condition, but it is not equivalent.

\item The proof text uses \(S_m(y)<md\), whereas the definition
\eqref{eq:Korec-Umd} and the final block count use \(S_m(y)\le md\).
The strict margin in \eqref{eq:Korec-dstar-interval} proves the implication
with \(S_m(y)\le md_*\) for all sufficiently large \(m\), so the defined
block count applies including its boundary.
\end{enumerate}
The witness time is \(r=m\), and \(m\le\log_2y\) on the retained interval.
It is controlled by the starting value, not by a constant independent of
that value.
\end{sourceissue}

\begin{sourceissue}[Tao line 165: sign-inconsistent Allouche threshold]
\label{sourceissue:Tao-line165-threshold}
Tao's v5 and v7 source line~165, and the 2022 journal p.~2, print
\begin{equation}
\label{eq:Tao-line165-printed}
 \theta>\frac32-\frac{\log3}{\log2}\approx0.869.
\end{equation}
The symbolic expression and decimal cannot both be correct.  Indeed,
\[
 3^2>2^3
 \quad\Longrightarrow\quad
 \frac{\log3}{\log2}>\frac32,
\]
so the printed symbolic right side is negative.  The exact correction,
proved independently in Proposition~\ref{prop:Allouche-threshold-extraction}
and explicitly reported by Korec, is
\begin{equation}
\label{eq:Tao-line165-repaired}
 \boxed{\theta>
 \frac32-\log_3 2
 =\frac32-\frac{\log2}{\log3}}
 \approx0.8690702464.
\end{equation}
Tao's subsequent Korec threshold
\[
 \theta>\log_4 3\approx0.7924812504
\]
is printed correctly \cite[Sec.~1.1, source line~165]{Tao2026v7}.  This
repair is not inferred from the decimal alone: the sign obstruction,
Allouche-coordinate derivation, and Korec's exact retrospective formula
agree.
\end{sourceissue}

\begin{proposition}[exact Tao--Korec threshold and clock crosswalk]
\label{prop:Tao-Korec-threshold-crosswalk}
Let
\[
 c_K=\log_4 3,
 \qquad
 c_A=\frac32-\log_3 2.
\]
Then:
\begin{enumerate}
\item \(c_K<c_A\), with exact difference
\begin{equation}
\label{eq:Tao-Korec-threshold-gap}
 c_A-c_K=
 \frac{(\log_2 3-1)(2-\log_2 3)}{2\log_2 3}>0.
\end{equation}
\item For every \(y\in\Npos\) and real \(c\),
\begin{equation}
\label{eq:Tao-Korec-event-equality}
 \exists j\in\Nzero,\ T^j(y)<y^c
 \quad\Longleftrightarrow\quad
 \exists r\in\Nzero,\ U^r(y)<y^c.
\end{equation}
Consequently the density statements of
Proposition~\ref{prop:Allouche-threshold-extraction} and
Theorem~\ref{thm:Korec-density-estimate} transfer to Tao's whole-orbit
minimum without any change in their sets of starting values.
\item For \(N\in\Opos\), define
\[
 a_j=\vtwo(3\operatorname{Syr}^j(N)+1),
 \qquad
 A_r=\sum_{j=0}^{r-1}a_j.
\]
Then the exact accelerated affine coordinate is
\begin{equation}
\label{eq:Syr-exact-affine-r}
 \operatorname{Syr}^r(N)
 =\frac{3^r}{2^{A_r}}N+
 \sum_{j=0}^{r-1}\frac{3^{r-1-j}}{2^{A_r-A_j}},
 \qquad A_0=0.
\end{equation}
For \(N>1,r\ge1\), the exponent of its multiplicative term is exactly
\begin{equation}
\label{eq:Syr-exponent-coordinate}
 \log_N\left(\frac{3^r}{2^{A_r}}N\right)
 =1+\frac{r}{\log N}
 \left(\log3-\frac{A_r}{r}\log2\right).
\end{equation}
At the parameter coordinate
\[
 \frac{A_r}{r}=2,\qquad
 \frac{r}{\log N}=\frac1{\log4},
\]
the right side of \eqref{eq:Syr-exponent-coordinate} is exactly
\(\log_4 3\).
\end{enumerate}
\end{proposition}

\begin{proof}
Put \(L=\log_2 3\).  Since \(2<3<4\), one has \(1<L<2\), and
\[
 c_A-c_K
 =\frac32-\frac1L-\frac L2
 =\frac{(L-1)(2-L)}{2L}>0,
\]
proving \eqref{eq:Tao-Korec-threshold-gap}.

Proposition~\ref{prop:KL-Tao-clock-embedding} proves that every shortened
state \(T^j(y)\) is the unshortened state \(U^{\iota_y(j)}(y)\), giving the
forward implication in \eqref{eq:Tao-Korec-event-equality}.  Conversely,
every unshortened time outside the image of \(\iota_y\) is the intermediate
state
\[
 U^{s_j+1}(y)=3T^j(y)+1
\]
following an odd shortened state.  Its next unshortened state is
\[
 U^{s_j+2}(y)=\frac{3T^j(y)+1}{2}=T^{j+1}(y).
\]
If the intermediate state is below \(y^c\), then its half is also below
\(y^c\).  Thus every unshortened witness produces a shortened witness,
which proves the reverse implication and the exact equality of starting
sets.

Finally, the first-return identity gives
\[
 \operatorname{Syr}^{j+1}(N)
 =\frac{3\operatorname{Syr}^j(N)+1}{2^{a_j}}.
\]
Induction expands this recurrence to \eqref{eq:Syr-exact-affine-r}; every
summand in its remainder is positive and its denominator records the exact
remaining valuation sum.  Taking logarithms of the multiplicative term
gives \eqref{eq:Syr-exponent-coordinate}.  Substitution of the two displayed
parameter coordinates gives
\[
 1+\frac{\log3-2\log2}{\log4}
 =\frac{\log3}{\log4}
 =\log_4 3.
\]
\end{proof}

\begin{nonclaim}
The exact coordinate calculation in
\eqref{eq:Syr-exponent-coordinate} explains Tao's statement that optimizing
the descent parameters recovers Korec
\cite[source line~695]{Tao2026v7}; it is not a substitute for Korec's
natural-density proof.  Tao's locally displayed loose choices
\(n_0=0.1\log_2x\) and \(A_{n_0}>1.9n_0\) do not themselves yield the
optimal exponent.  The separate \(1.9n\) versus \(1.9n_0\) source defect is
audited elsewhere and is not conflated with
Source Issue~\ref{sourceissue:Tao-line165-threshold}.

None of the Terras, Allouche, or Korec density-one theorems includes its
strict endpoint or removes a possible zero-density exceptional set.  None
proves the Collatz conjecture.  Equation~\eqref{eq:Tao-Korec-event-equality}
identifies the existence of a descent witness, not one iterate of \(T\) with
one iterate of \(U\); their clocks remain related by the variable-time map
\(\iota_y\).
\end{nonclaim}

\subsection{Run lengths and chronological parity words}

Put
\[
 \mathcal A=(\Z_{\geq1})^{\Nzero},
 \qquad
 \mathcal B=\{(b_k)_{k\geq0}\in\{0,1\}^{\Nzero}:
 b_0=1\text{ and infinitely many }b_k\text{ equal }1\}.
\]
Let $\sigma_{\mathcal A}$ and $\sigma_{\mathcal B}$ denote the left shifts.
For $\mathbf a=(a_0,a_1,\ldots)\in\mathcal A$, set
\[
 s_0=0,
 \qquad
 s_j=\sum_{i=0}^{j-1}a_i\quad(j\geq1),
\]
and define $B:\mathcal A\to\mathcal B$ by
\begin{equation}
\label{eq:run-to-parity-B}
 B(\mathbf a)_k=1
 \quad\Longleftrightarrow\quad
 k=s_j\text{ for some }j\geq0.
\end{equation}

\begin{proposition}[run-length/parity-word isomorphism]
\label{prop:run-parity-isomorphism}
The map $B:\mathcal A\to\mathcal B$ is a bijection.  If the one positions of
$\mathbf b\in\mathcal B$ are
$0=s_0<s_1<s_2<\cdots$, then
\begin{equation}
\label{eq:parity-to-run-B-inverse}
 B^{-1}(\mathbf b)_j=s_{j+1}-s_j.
\end{equation}
Moreover, for every $\mathbf a\in\mathcal A$,
\begin{equation}
\label{eq:variable-shift-B}
 B(\sigma_{\mathcal A}\mathbf a)
 =\sigma_{\mathcal B}^{a_0}B(\mathbf a).
\end{equation}
\end{proposition}

\begin{proof}
Every $a_i$ is positive, so the partial sums $s_j$ are strictly increasing,
start at zero, and are unbounded.  Thus \eqref{eq:run-to-parity-B} defines an
element of $\mathcal B$.  Conversely, the one positions of an element of
$\mathcal B$ form a strictly increasing unbounded sequence beginning at zero;
their successive positive gaps define \eqref{eq:parity-to-run-B-inverse}.
Taking partial sums of those gaps recovers the original one positions, which
proves both inverse identities.  Removing the first gap translates every
later one position by $-a_0$, exactly as the right side of
\eqref{eq:variable-shift-B} does.
\end{proof}

For $x\in\Opos$, define the accelerated exponent code and shortened parity
code by
\begin{align}
 E(x)&=\bigl(e(C_{\mathrm{Cr}}^j(x))\bigr)_{j\geq0}
       \in\mathcal A,
 \label{eq:E-accelerated-code}\\
 P(x)&=\bigl(T^k(x)\bmod2\bigr)_{k\geq0}
       \in\mathcal B.
 \label{eq:P-parity-code}
\end{align}
Both maps are total on $\Opos$: every accelerated value remains a positive odd
integer and each return time is finite.

\begin{theorem}[typed first-return coding morphism]
\label{thm:first-return-coding}
For every $x\in\Opos$,
\begin{equation}
\label{eq:PBE-factorization}
 P(x)=B(E(x)).
\end{equation}
The dynamics satisfy
\begin{equation}
\label{eq:E-shift}
 E(C_{\mathrm{Cr}}x)=\sigma_{\mathcal A}E(x),
 \qquad
 P(C_{\mathrm{Cr}}x)=\sigma_{\mathcal B}^{e(x)}P(x).
\end{equation}
The fibres of $E$ have cardinality at most one.  No surjectivity of
$E:\Opos\to\mathcal A$ is asserted.
\end{theorem}

\begin{proof}
Proposition~\ref{prop:first-return-clocks} shows that the shortened parity
block beginning at $C_{\mathrm{Cr}}^j(x)$ is
\[
 1\,0^{\,e(C_{\mathrm{Cr}}^j(x))-1}.
\]
Concatenating these blocks gives exactly \eqref{eq:PBE-factorization}.  The
first equality in \eqref{eq:E-shift} follows directly from the definition of
$E$; the second follows either from the first equality,
Proposition~\ref{prop:run-parity-isomorphism}, and
\eqref{eq:PBE-factorization}, or directly from
$C_{\mathrm{Cr}}(x)=T^{e(x)}(x)$.

For injectivity, let
\[
 Q_3(z)=\sum_{k\geq0}(T^k(z)\bmod2)2^k
 \qquad(z\in\Ztwo)
\]
be the Bernstein--Lagarias chronological parity encoder.  Their theorem,
reproved below as Proposition~\ref{prop:BL-Q}, makes $Q_3$ injective.  If
$E(x)=E(y)$, then \eqref{eq:PBE-factorization} gives $P(x)=P(y)$, hence
$Q_3(x)=Q_3(y)$ and $x=y$.  The proof identifies the image only as
$E(\Opos)$; it supplies no realization of an arbitrary positive-integer run
sequence by a positive orbit.
\end{proof}

\subsection{The accelerated odd \texorpdfstring{$2$}{2}-adic survivor and its
countable shift}

Urata writes
\[
 \mathbb U_2=1+2\Ztwo=\Ztwo^{\times}
\]
for the closure of the odd-type rationals.  His continuous extension is the
map
\[
 \varphi:\Ztwo\setminus\{0,-1/3\}\longrightarrow\mathbb U_2
\]
given, without merging its two branches, by
\begin{equation}
\label{eq:Urata-full-map}
 \varphi(x)=
 \begin{cases}
  (3x+1)/2^{\vtwo(3x+1)},
    &x\in\mathbb U_2\setminus\{-1/3\},\\
  x/2^{\vtwo(x)},
    &x\in\Ztwo\setminus(\mathbb U_2\cup\{0\}).
 \end{cases}
\end{equation}
The restriction to odd units is the accelerated map used below; on positive
odd integers it is exactly $C_{\mathrm{Cr}}$.  The full even-input branch in
\eqref{eq:Urata-full-map} is retained because it belongs to the source
definition, but it is not smuggled into the odd first-return system
\cite[printed pp.~5--8]{Urata2003}.

\begin{sourceissue}[typing Urata's valuation-continuity sentence]
After defining the exponent of zero to be infinity, the text calls the
exponent continuous on $\mathbb Q_2$ without specifying a topology on
$\Z\cup\{\infty\}$.  Continuity at zero is therefore not yet a typed
statement.  What is used here is the exact local constancy of $\vtwo$ on
$\mathbb Q_2^{\times}$; compactness of $\Ztwo$ is invoked independently.
\end{sourceissue}

For $p\geq1$, define
\begin{equation}
\label{eq:Urata-inverse-branch}
 F_p:\mathbb U_2\longrightarrow\mathbb U_2,
 \qquad F_p(y)=\frac{2^p y-1}{3}.
\end{equation}
The literal identity
\[
 3F_p(y)+1=2^p y
\]
shows that $\vtwo(3F_p(y)+1)=p$ and $\varphi(F_p(y))=y$.  Conversely every
odd-unit preimage has this form with its unique valuation.  Thus
\begin{equation}
\label{eq:Urata-fibre}
 \varphi^{-1}(y)\cap\mathbb U_2
 =\{F_p(y):p\geq1\},
\end{equation}
an explicitly labelled countably infinite fibre
\cite[printed p.~8, Example~2]{Urata2003}.

For an actual odd-unit orbit $x_k=\varphi^k(x_0)$, put
\[
 p_k=\vtwo(3x_{k-1}+1),\qquad
 S_0=0,\qquad S_n=\sum_{k=1}^{n}p_k,
\]
and retain the coordinate polynomial
\begin{equation}
\label{eq:Urata-Cn}
 C_n(p_1,\ldots,p_n)
 =\sum_{j=1}^{n}3^{n-j}2^{S_{j-1}}.
\end{equation}
This is a reindexing of Urata's displayed $C(p_1,\ldots,p_{n-1})$, not a
replacement source formula.  Successive substitution gives
\begin{equation}
\label{eq:Urata-finite-recurrence}
 2^{S_n}x_n=C_n+3^n x_0.
\end{equation}
Urata Proposition~1 records this relation and the resulting complete
$n$-step inverse set.  Propositions~2--3 show that common first $n$ exponents
give $|x-y|_2\leq2^{-S_n}$ and that a complete exponent sequence determines at
most one survivor.  In fact subtraction of
\eqref{eq:Urata-finite-recurrence}, together with the even difference of two
odd $n$th iterates, gives the sharper bound
\begin{equation}
\label{eq:Urata-sharp-prefix-bound}
 |x-y|_2\leq2^{-(S_n+1)}.
\end{equation}

\begin{sourceissue}[the Cauchy bound in Urata Proposition 4]
\label{issue:Urata-prop4-bound}
Proposition~4 constructs approximants $x_n$ whose first $n$ exponents are
prescribed.  For $m\leq n$, the pair $x_m,x_n$ shares only the first $m$
prescribed exponents, but the printed comparison uses $S_n$ rather than
$S_m$.  The stronger printed bound is false.  For $p_1=1,p_2=3$, the first
two approximants are
\[
 x_1=\frac13,\qquad x_2=\frac{11}{9},\qquad
 |x_2-x_1|_2=\left|\frac89\right|_2=2^{-3}>2^{-4}=2^{-S_2}.
\]
This is a proof defect in the displayed Cauchy estimate, not a counterexample
to the proposition.  The official formula \eqref{eq:Urata-Cn} itself is
consistent; no correction from the later unauthenticated re-typesetting is
attributed to Urata.
\end{sourceissue}

Define the survivor domain
\begin{equation}
\label{eq:Urata-survivor}
 X=\mathbb U_2\setminus
 \bigcup_{j\geq0}\varphi^{-j}\{-1/3\}.
\end{equation}

\begin{theorem}[repaired countable-alphabet conjugacy]
\label{thm:Urata-countable-conjugacy}
The exponent encoder
\begin{equation}
\label{eq:Urata-exponent-encoder}
 E_U:X\longrightarrow(\Z_{\geq1})^{\Nzero},
 \qquad
 E_U(x)_j=\vtwo(3\varphi^j(x)+1),
\end{equation}
is a homeomorphism.  If
$\mathbf p=(p_1,p_2,\ldots)$ and
$S_j=p_1+\cdots+p_j$, its inverse is the unnormalized coordinate series
\begin{equation}
\label{eq:Urata-explicit-inverse}
 E_U^{-1}(\mathbf p)
 =-\sum_{j=1}^{\infty}\frac{2^{S_{j-1}}}{3^j}.
\end{equation}
Moreover
\begin{equation}
\label{eq:Urata-shift-conjugacy}
 E_U\circ\varphi
 =\sigma_{\mathcal A}\circ E_U.
\end{equation}
Urata Propositions~3--4 supply the uniqueness and existence statements; the
explicit inverse, repaired proof, and conjugacy formulation here are separate
constructive reconstructions with no novelty claim.
\end{theorem}

\begin{proof}
The series in \eqref{eq:Urata-explicit-inverse} converges because
$S_{j-1}\geq j-1$.  More generally define, for $k\geq0$,
\[
 x_k=-\sum_{r=1}^{\infty}
 \frac{2^{p_{k+1}+\cdots+p_{k+r-1}}}{3^r},
\]
where the empty exponent sum is zero.  Its first term is $-1/3$ and every
later term is even, so $x_k$ is an odd unit.  Direct cancellation gives
\begin{equation}
\label{eq:Urata-suffix-identity}
 3x_k+1=2^{p_{k+1}}x_{k+1}.
\end{equation}
Therefore the valuation is exactly $p_{k+1}$,
$\varphi(x_k)=x_{k+1}$, and no $x_k$ is $-1/3$.  This proves existence in
$X$ without the defective Cauchy estimate.  Urata Proposition~3, or
\eqref{eq:Urata-sharp-prefix-bound} with $n\to\infty$, proves uniqueness.

Every finite exponent prefix is locally constant on $X$.  Conversely,
\eqref{eq:Urata-sharp-prefix-bound} shows that common prefixes tending to
infinite length force the inverse points together.  This proves continuity in
both directions.  Dropping the first valuation from the definition proves
\eqref{eq:Urata-shift-conjugacy}.
\end{proof}

\begin{corollary}[exact exponent cylinders and Haar law]
\label{cor:Urata-cylinders-Haar}
For $w=(p_1,\ldots,p_n)$, put
$F_w=F_{p_1}\circ\cdots\circ F_{p_n}$.  Then
\begin{align}
 F_w(y)&=\frac{2^{S_n}y-C_n}{3^n},
 \label{eq:Urata-Fw-coordinate}\\
 F_w(\mathbb U_2)&=F_w(1)+2^{S_n+1}\Ztwo.
 \label{eq:Urata-cylinder-ball}
\end{align}
The exponent cylinder in $X$ is $F_w(X)$; its ambient closure is the ball
\eqref{eq:Urata-cylinder-ball}.  For normalized Haar measure $\mu_U$
on $\mathbb U_2$,
\begin{equation}
\label{eq:Urata-cylinder-measure}
 \mu_U(F_w(X))
 =2^{-S_n}=\prod_{j=1}^{n}2^{-p_j}.
\end{equation}
Consequently $E_U$ carries Haar measure to the Bernoulli product law
$\Pr(p_j=k)=2^{-k}$ for $k\geq1$, and the accelerated survivor system is
Bernoulli and strongly mixing modulo its countable null exceptional set.
\end{corollary}

\begin{proof}
Composition of \eqref{eq:Urata-inverse-branch} gives
\eqref{eq:Urata-Fw-coordinate}.  Since
$\mathbb U_2=1+2\Ztwo$ and $3^n$ is a unit, its difference set is exactly
$2^{S_n+1}\Ztwo$, proving \eqref{eq:Urata-cylinder-ball}.  This ball has
ambient $\Ztwo$-Haar measure $2^{-(S_n+1)}$ and hence normalized
$\mathbb U_2$-measure $2^{-S_n}$.  The removed preimage tree of $-1/3$ is
countable, so it has measure zero.  Cylinder products prove the product law;
the Bernoulli shift consequences follow through
\eqref{eq:Urata-shift-conjugacy}.
\end{proof}

\begin{proposition}[exact induced-system bridge]
\label{prop:Urata-binary-bridge}
For $x\in X$ and $S_j=p_1+\cdots+p_j$ from
$E_U(x)$,
\begin{equation}
\label{eq:Urata-induced-times}
 \varphi^j(x)=T^{S_j}(x),
 \qquad
 (T^k(x)\bmod2)_{k\geq0}=B(E_U(x)).
\end{equation}
Thus $X$ is the first-return system of the shortened binary shift on odd
inputs whose parity sequence contains infinitely many ones.  The excluded
word $1000\ldots$ is the parity word of $-1/3$, and words with only finitely
many ones are its accelerated preimages.
\end{proposition}

\begin{proof}
Equation \eqref{eq:Urata-suffix-identity} says that one odd branch followed by
$p_j-1$ even divisions reaches the next odd unit, so the elapsed shortened
time is exactly $p_j$.  Summing the retained times proves the first identity;
the intervening parity block is $1\,0^{p_j-1}$, which proves the second through
Proposition~\ref{prop:run-parity-isomorphism}.  The final exceptional-set
description follows from the total Bernstein--Lagarias parity conjugacy.
\end{proof}

\begin{proposition}[periodic exponent words]
\label{prop:Urata-periodic-exponents}
For a genuine survivor orbit with first $n$ exponents,
\begin{equation}
\label{eq:Urata-periodic-fraction}
 \varphi^n(x)=x
 \quad\Longleftrightarrow\quad
 (2^{S_n}-3^n)x=C_n
 \quad\Longleftrightarrow\quad
 x=\frac{C_n}{2^{S_n}-3^n}.
\end{equation}
The orbit returns after $n$ accelerated steps exactly when its full exponent
sequence is $n$-periodic.  Its least accelerated period equals the least shift
period; neither need equal a displayed nonleast $n$.
\end{proposition}

\begin{proof}
The first equivalence is \eqref{eq:Urata-finite-recurrence} with $x_n=x_0$;
$2^{S_n}\neq3^n$ gives the second.  Repetition of a returning orbit repeats
the valuations.  Conversely a periodic exponent sequence determines by
Theorem~\ref{thm:Urata-countable-conjugacy} the same point after one word.
These are Urata Propositions~5--6 with the least-period qualification made
explicit \cite[printed pp.~10--11]{Urata2003}.
\end{proof}

\begin{nonclaim}
Urata's arbitrary exponent realization holds on the enlarged domain
$X\subset\Ztwo^{\times}$, not on $\Opos$.  Thus the earlier
$E:\Opos\to\mathcal A$ remains only injective.  Rationality of
\eqref{eq:Urata-periodic-fraction} does not imply integrality or a new positive
cycle; printed p.~11 asks when that fraction is integral and leaves the
question open.  The rational eventual-periodicity paragraph on printed p.~6
is conjectural.  The variable cylinder modulus is $2^{S_n+1}$, not $2^n$, and
the Haar statements do not transfer to the countable Haar-null ordinary
integers.
\end{nonclaim}

\subsection{Chronological affine words and fixed returns}

On $\Q$ put
\[
 h_0(z)=\frac z2,
 \qquad
 h_1(z)=\frac{3z+1}{2}.
\]
Let $\boldsymbol\varepsilon=(\varepsilon_0,\ldots,\varepsilon_{L-1})$ be a
chronological binary word of length $L\geq1$.  Suppose it has $m$ ones, at
positions
\[
 0\leq v_0<\cdots<v_{m-1}<L.
\]

\begin{proposition}[formal chronological composition]
\label{prop:formal-chronological-composition}
For every $z\in\Q$,
\begin{equation}
\label{eq:formal-affine-word}
 h_{\varepsilon_{L-1}}\circ\cdots\circ h_{\varepsilon_0}(z)
 =\frac{3^m z+
 \displaystyle\sum_{r=0}^{m-1}3^{m-r-1}2^{v_r}}{2^L}.
\end{equation}
The sum is empty when $m=0$.
\end{proposition}

\begin{proof}
After the first $j$ chronological branches, let
\[
 m_j=\sum_{i=0}^{j-1}\varepsilon_i,
 \qquad
 A_j=\sum_{\substack{0\leq r<m\\v_r<j}}
       3^{m_j-r-1}2^{v_r}.
\]
The induction invariant is
\[
 h_{\varepsilon_{j-1}}\circ\cdots\circ h_{\varepsilon_0}(z)
 =\frac{3^{m_j}z+A_j}{2^j}.
\]
It is the identity at $j=0$.  A zero branch only adds one power of $2$ to
the denominator.  A one branch multiplies the old numerator by $3$ and adds
$2^j$; this raises every old exponent of $3$ by one and inserts the new term
$3^0 2^j$.  The invariant at $j=L$ is \eqref{eq:formal-affine-word}.
\end{proof}

The formal word equals $T^L(z)$ exactly when every prefix value has the parity
prescribed by the next $\varepsilon_j$.  Algebra alone does not supply that
itinerary condition.

\begin{sourceissue}[B\"ohm--Sontacchi Proposition 4]
\label{sourceissue:BS-prop4}
B\"ohm--Sontacchi Proposition~1 supplies the formal affine-composition
algebra.  Their printed Proposition~4 displays its specialization as $u^n(x)$
while quantifying over every $x\in\Z$ and every run tuple
\cite[printed p.~263, Props.~1 and~4]{BohmSontacchi1978}.  That unrestricted
iterate statement is false.  With $x=0$, $m=1$, and $n_0=n_1=1$, the formal
chronological word is $(0,1,0)$ and
\[
 h_0\circ h_1\circ h_0(0)=\frac14,
 \qquad
 T^3(0)=0.
\]
Thus the right side survives as a formal composition; equality with an
iterate requires a separate branch-correctness proof.
\end{sourceissue}

\begin{sourceissue}[B\"ohm--Sontacchi Proposition 5 position range]
\label{sourceissue:BS-prop5}
The printed statement gives $v_{i-1}<v_i$ only for $1\leq i<m$ and then
$v_m=n$; it does not visibly print the necessary final inequality
$v_{m-1}<v_m$ \cite[printed p.~263, Prop.~5]{BohmSontacchi1978}.  If the
display is read literally, $m=n=1$ and $(v_0,v_1)=(1,1)$ give the candidate
$x=-2$, while $T(-2)=-1$.  The formal run construction does impose
$v_m-v_{m-1}=1+n_m>0$.  This edition therefore states positions directly as
$0\leq v_0<\cdots<v_{m-1}<L$; that final inequality is an explicit
reconstruction, not silently attributed printed text.
\end{sourceissue}

\begin{lemma}[wrong branches cannot close in $\Ztwo$]
\label{lem:wrong-branch}
For $z\in\Ztwo$, exactly one of $h_0(z),h_1(z)$ lies in $\Ztwo$, namely
$h_{z\bmod2}(z)$.  If $v_2(y)<0$, then
\[
 v_2(h_0(y))=v_2(h_1(y))=v_2(y)-1.
\]
\end{lemma}

\begin{proof}
If $z$ is even, $z/2$ is integral and $(3z+1)/2$ is not; if $z$ is odd, the
roles reverse.  When $v_2(y)<0$, the valuations of $3y$ and $1$ are unequal,
so the ultrametric equality gives $v_2(3y+1)=v_2(y)$; division by $2$ proves
the second assertion for both branches.
\end{proof}

\begin{theorem}[recovered fixed-return theorem]
\label{thm:recovered-fixed-return}
Define
\begin{equation}
\label{eq:fixed-return-candidate}
 R(\boldsymbol\varepsilon)=
 \frac{\displaystyle\sum_{r=0}^{m-1}3^{m-r-1}2^{v_r}}
      {2^L-3^m}\in\Q.
\end{equation}
If $R(\boldsymbol\varepsilon)\in\Ztwo$, every branch of the formal word is
correct and
\[
 T^L(R(\boldsymbol\varepsilon))=R(\boldsymbol\varepsilon).
\]
Conversely, if $z\in\Ztwo$, $T^L(z)=z$, and
$\boldsymbol\varepsilon$ is its chronological parity word, then
$z=R(\boldsymbol\varepsilon)$.  If $z\in\Z$, this is an integer return.  Its
exact period divides $L$ and need not equal $L$.
\end{theorem}

\begin{proof}
The denominator in \eqref{eq:fixed-return-candidate} is nonzero: a positive
power of $2$ cannot equal a nonnegative power of the odd integer $3$.
Solving the fixed-point equation for \eqref{eq:formal-affine-word} gives
exactly \eqref{eq:fixed-return-candidate}.  Start its right-to-left evaluation
at $R\in\Ztwo$.  If a wrong branch occurs, Lemma~\ref{lem:wrong-branch} sends
the current value outside $\Ztwo$ with negative valuation, and every later
formal branch decreases that valuation.  The final value cannot then equal
the initial integral value.  Hence all branches are correct and the formal
fixed point is a genuine $T$-return.  The converse substitutes the actual
itinerary into Proposition~\ref{prop:formal-chronological-composition} and
solves the same affine equation.  A return after $L$ iterates says only that
the least positive period divides $L$.  This last scope agrees with the
unnumbered paragraph following B\"ohm--Sontacchi Proposition~5
\cite[printed p.~264]{BohmSontacchi1978}.
\end{proof}

\begin{corollary}[Crandall's accelerated cycle identity]
\label{cor:crandall-cycle-identity}
Let $x\in\Opos$, let
$a_j=e(C_{\mathrm{Cr}}^j(x))$, and put
\[
 A_0=0,
 \qquad
 A_j=\sum_{r=0}^{j-1}a_r\quad(1\leq j\leq k).
\]
If $C_{\mathrm{Cr}}^k(x)=x$, then
\begin{equation}
\label{eq:crandall-cycle-identity}
 x(2^{A_k}-3^k)
 =\sum_{j=0}^{k-1}2^{A_j}3^{k-1-j}.
\end{equation}
Conversely, if $0=A_0<A_1<\cdots<A_k$ and an odd positive integer $x$
satisfies \eqref{eq:crandall-cycle-identity}, then the corresponding formal
shortened word fixes $x$; by Theorem~\ref{thm:recovered-fixed-return} it is a
genuine $T$-return, and its odd returns give
$C_{\mathrm{Cr}}^k(x)=x$.
\end{corollary}

\begin{proof}
Theorem~\ref{thm:first-return-coding} places the $k$ ones of the shortened
itinerary at $A_0,\ldots,A_{k-1}$ and gives total shortened length $A_k$.
Equation \eqref{eq:fixed-return-candidate}, with $L=A_k$ and $m=k$, is
equivalent to \eqref{eq:crandall-cycle-identity}.  Conversely, the same
identity makes $x$ the formal fixed point of that word.  Its membership in
$\Z\subset\Ztwo$ supplies branch correctness, and the one positions split the
genuine itinerary into the prescribed $k$ first-return blocks.  This is
Crandall's Eq.~(7.3) and Corollary~7.1, including its return-time rather than
exact-period direction \cite[pp.~1287--1288]{Crandall1978}.
\end{proof}

\subsection{Crandall's complete accelerated programme}
\label{subsec:Crandall-complete-audit}

This subsection reconstructs all twelve printed pages of Crandall's 1978
article from the published witness \cite[pp.~1281--1292]{Crandall1978}.  The
source uses only the accelerated odd map \eqref{eq:crandall-C}.  Its
trajectory is the indexed sequence
\[
 \mathcal T_m=\bigl(C_{\mathrm{Cr}}^j(m)\bigr)_{j\geq1},
\]
stopped after the least index $j$ for which $C_{\mathrm{Cr}}^j(m)=1$, if
such an index exists, and left infinite otherwise.  Its height is that least
index, or $\infty$ when there is none.  Thus the braces used in the source's
display denote an ordered sequence, not its underlying set; repeated values
in a nonterminating periodic trajectory do not make its height finite.  In
particular, the source table has $\mathcal T_1=(1)$ and $h(1)=1$; this is not
the zero-step hitting-time convention.

For a segment $x_i=C_{\mathrm{Cr}}^i(x_0)$, put
\begin{equation}
\label{eq:Crandall-cumulative-time}
 e_i=\vtwo(3x_i+1),
 \qquad
 A_0=0,
 \qquad
 A_j=\sum_{i=0}^{j-1}e_i.
\end{equation}
Repeated use of Proposition~\ref{prop:first-return-clocks} gives
\begin{equation}
\label{eq:Crandall-three-clocks}
 C_{\mathrm{Cr}}^j(x_0)=T^{A_j}(x_0)=U^{A_j+j}(x_0).
\end{equation}
Thus every height, trajectory position, and period in Crandall's paper is
accelerated time $j$.  Neither $A_j$ nor $A_j+j$ may be substituted for it.

\subsubsection{Large finite excursions and the random-walk nonclaim}

\begin{proposition}[Crandall, Theorem 2.2]
\label{prop:Crandall-excursions}
For $k\geq2$, let $m_k=2^k-1$.  Then
\begin{equation}
\label{eq:Crandall-excursion-orbit}
 C_{\mathrm{Cr}}^j(m_k)=3^j2^{k-j}-1
 \qquad(0\leq j<k),
\end{equation}
and hence
\begin{equation}
\label{eq:Crandall-excursion-ratio}
 \frac{\sup\mathcal T_{m_k}}{m_k}>
 \left(\frac32\right)^{k-1}.
\end{equation}
Consequently the family of ratios
\[
 \left\{\frac{\sup\mathcal T_m}{m}:m\in\Opos\right\}
\]
is unbounded.  For the ratios with denominator $m^\alpha$, this construction
proves unboundedness only for $\alpha<\log_2 3$.
\end{proposition}

\begin{proof}
Equation~\eqref{eq:Crandall-excursion-orbit} is immediate at $j=0$.  If it
holds with $j<k-1$, then
\[
 3(3^j2^{k-j}-1)+1
 =2(3^{j+1}2^{k-j-1}-1),
\]
whose second factor is odd; the next accelerated exponent is therefore one.
At $j=k-1$, the trajectory contains $2\cdot3^{k-1}-1$.  Direct
cross-multiplication gives
\[
 \frac{2\cdot3^{k-1}-1}{2^k-1}>
 \frac{3^{k-1}}{2^{k-1}}
\]
because $3^{k-1}>2^{k-1}$.  This proves the stated bound.  Dividing the same
member by $m_k^\alpha$ yields an exponential factor
$3^{k-1}/2^{\alpha k}$, which this construction forces upward exactly when
$\alpha<\log_2 3$ \cite[pp.~1282--1283, Thm.~2.2 and following paragraph]
{Crandall1978}.
\end{proof}

Section~3 assigns the heuristic probabilities
$\Pr(e(m)=a)=2^{-a}$ for $a\geq1$.  Its model logarithmic drift is
\begin{equation}
\label{eq:Crandall-heuristic-drift}
 \sum_{a\geq1}2^{-a}\log(3/2^a)=-\log(4/3).
\end{equation}
It then proposes the individual estimate
$h(m)\sim\log m/\log(4/3)$.  For the source average
\begin{equation}
\label{eq:Crandall-average-height}
 H(x)=\frac2x\sum_{\substack{m\leq x\\m\in\Opos}}h(m),
\end{equation}
Conjecture~3.1 is
\begin{equation}
\label{eq:Crandall-average-conjecture}
 H(x)\sim\frac{2\log x}{\log(16/9)}.
\end{equation}

\begin{nonclaim}
\label{nonclaim:Crandall-random-walk}
Equations \eqref{eq:Crandall-heuristic-drift}--
\eqref{eq:Crandall-average-conjecture} are explicitly heuristic or
conjectural.  They prove neither termination nor unboundedness.  The source
also reports Everett's almost-everywhere descent theorem, but does not prove
it.  Everett's primary proof and the exact variable-clock transfer are now
reconstructed in Theorem~\ref{thm:Everett-density-one-descent} and
Proposition~\ref{prop:Everett-Crandall-crosswalk}; neither follows from the
heuristic drift displayed here.  The sentence after Conjecture~3.1 calls it
``(3.1)'', duplicating the number of the individual-height heuristic
\cite[pp.~1282--1284]{Crandall1978}.
\end{nonclaim}

\subsubsection{Backward coordinates and the exact finite-height address}

Crandall defines, for $a\in\Z_{>0}$ and $n\in\Q$,
\begin{equation}
\label{eq:Crandall-backward-branch}
 B_a(n)=\frac{2^a n-1}{3},
 \qquad
 B_{a_j\cdots a_1}=B_{a_j}\circ\cdots\circ B_{a_1}.
\end{equation}
The subscript order is backward: when all displayed values are integral,
$C_{\mathrm{Cr}}$ removes the leftmost applied branch $B_{a_j}$.

For a natural-order address $\mathbf a=(a_1,\ldots,a_j)$, define
\begin{equation}
\label{eq:Crandall-address-recursion}
 n_0=1,
 \qquad
 n_i=B_{a_i}(n_{i-1})\quad(1\leq i\leq j).
\end{equation}
Let $G^\sharp$ be the set of such finite addresses satisfying
\begin{align}
 a_1&>2,
 \label{eq:Crandall-address-first-gate}\\
 2^{a_i}n_{i-1}&\equiv1\pmod3
 &&(1\leq i\leq j),
 \label{eq:Crandall-address-integrality-gate}\\
 2^{a_i}n_{i-1}&\not\equiv1\pmod9
 &&(1\leq i<j).
 \label{eq:Crandall-address-continuation-gate}
\end{align}

\begin{theorem}[Crandall's finite-height address theorem, repaired]
\label{thm:Crandall-address-bijection}
Let
\[
 \mathcal M=\{m\in\Opos:m>1,\ h(m)<\infty\}.
\]
The map
\begin{equation}
\label{eq:Crandall-address-evaluation}
 \operatorname{ev}:G^\sharp\longrightarrow\mathcal M,
 \qquad
 (a_1,\ldots,a_j)\longmapsto B_{a_j\cdots a_1}(1)
\end{equation}
is a bijection.  Its inverse sends $m\in\mathcal M$ to the length-$h(m)$
address
\begin{equation}
\label{eq:Crandall-address-inverse}
 a_i=e(C_{\mathrm{Cr}}^{h(m)-i}(m))
 \qquad(1\leq i\leq h(m)).
\end{equation}
The exact trajectory of $n_j=\operatorname{ev}(\mathbf a)$ is
\[
 n_j\longmapsto n_{j-1}\longmapsto\cdots\longmapsto n_1
 \longmapsto1,
\]
so its height is exactly $j$.
\end{theorem}

\begin{proof}
Suppose first that \eqref{eq:Crandall-address-integrality-gate} holds through
index $i$.  Then $n_i=(2^{a_i}n_{i-1}-1)/3$ is integral.  Since the numerator
is odd and the denominator is odd, $n_i$ is odd.  Moreover
\[
 n_i\equiv0\pmod3
 \quad\Longleftrightarrow\quad
 2^{a_i}n_{i-1}\equiv1\pmod9.
\]
Thus \eqref{eq:Crandall-address-continuation-gate} says exactly that every
nonterminal $n_i$ is prime to $3$, which is necessary for another inverse
branch to be integral.  From
$3n_i+1=2^{a_i}n_{i-1}$ and oddness of $n_{i-1}$, one has
$\vtwo(3n_i+1)=a_i$ and hence $C_{\mathrm{Cr}}(n_i)=n_{i-1}$.  Finally,
$a_1>2$ ensures $n_1\ne1$, so the first occurrence of $1$ is terminal and
the height is $j$.

Conversely, a finite trajectory from $m>1$ supplies the valuations in
\eqref{eq:Crandall-address-inverse}.  Reversing its branch equations gives
\eqref{eq:Crandall-address-recursion}.  Every nonterminal orbit value must be
prime to $3$: if $3\mid n_i$, then no equation
$3u+1=2^a n_i$ is possible modulo $3$.  Hence all three gates hold.  The
determinism of $C_{\mathrm{Cr}}$ recovers every $a_i$ uniquely, so evaluation
and \eqref{eq:Crandall-address-inverse} are literal two-sided inverses.  This
is the corrected content of Lemmas~4.1--4.3 and Theorem~4.1
\cite[pp.~1284--1285]{Crandall1978}.
\end{proof}

\begin{sourceissue}[two finite-address endpoint defects]
\label{sourceissue:Crandall-address-endpoints}
The printed definition of $G$ imposes a non-$1\pmod9$ condition on $a_1$
and a terminal mod-$3$ condition on $a_j$.  At length $j=1$, both apply to
the same letter.  It thereby excludes the valid height-one example
\[
 B_6(1)=21,
 \qquad C_{\mathrm{Cr}}(21)=1,
 \qquad 2^6\equiv1\pmod9.
\]
The mod-$9$ gate belongs only to nonterminal indices, as in
\eqref{eq:Crandall-address-continuation-gate}.  Independently, the printed
Lemma~4.3 includes ``an integer of height $j$'' while $a_1>2$ excludes
$B_2(1)=1$; the source table nevertheless has $h(1)=1$.  Restricting the
lemma to the $m>1$ set used by Theorem~4.1 repairs this endpoint.  Likewise,
inside $\mathcal M$ the height-one family $(4^s-1)/3$ has $s\geq2$, not
merely $s\geq1$ \cite[pp.~1282, 1284--1285]{Crandall1978}.
\end{sourceissue}

The address theorem is genuinely a bijection, not a merely suggestive
parameterization: its domain, codomain, evaluation, inverse, and exact clock
are all displayed.  It is also a finite-orbit theorem.  It provides no
infinite address space and no conclusion about nonterminating positive
orbits.

\subsubsection{Fixed-height and total termination counts}

Let
\[
 \pi_h(x)=\#\{m\in\mathcal M:h(m)=h,\ m\leq x\}.
\]
The elementary size estimate is
\begin{equation}
\label{eq:Crandall-backward-size}
 B_{a_j\cdots a_1}(1)
 <\frac{2^{a_1+\cdots+a_j}}{3^j}.
\end{equation}
It follows by induction from \eqref{eq:Crandall-backward-branch}, retaining
the discarded negative translation rather than replacing the branch by an
equality.

\begin{sourceissue}[domain of the address-counting lemma]
\label{sourceissue:Crandall-lemma52-domain}
Printed Lemma~5.2 says that for every real $z>0$, the number of length-$j$
addresses in $G$ with coordinate sum at most $z$ is at least
\begin{equation}
\label{eq:Crandall-address-count}
 \left(2\left\lfloor\frac{z-2}{6j}\right\rfloor\right)^j.
\end{equation}
At $z=1,j=2$, the right side is $4$, while no two positive coordinates have
sum at most $1$.  The proof's blocks of six require a nonnegative number of
blocks.  The statement is correct at the scope used later after requiring
$z\geq2$, or globally after replacing the floor by its nonnegative part.
\end{sourceissue}

\begin{theorem}[Crandall, fixed accelerated height]
\label{thm:Crandall-fixed-height-count}
There are real constants $r,x_0>0$, independent of $h$, such that for every
$h\in\Z_{>0}$,
\begin{equation}
\label{eq:Crandall-fixed-height-hypothesis}
 x>\max\{x_0,2^{h/r}\}
\end{equation}
implies
\begin{equation}
\label{eq:Crandall-fixed-height-bound}
 \pi_h(x)>\frac{(\log_2(x^r))^h}{h!}.
\end{equation}
In particular, infinitely many positive odd integers have each prescribed
positive accelerated height $h$.
\end{theorem}

\begin{proof}
Theorem~\ref{thm:Crandall-address-bijection},
\eqref{eq:Crandall-backward-size}, and the repaired form of
\eqref{eq:Crandall-address-count} give the source's explicit preliminary
bound
\[
 \pi_h(x)\geq
 \left(\frac{\log_2(3^h x/4)}{3h}-2\right)^h
\]
once the displayed base is positive.  Choose $x_0$ so that $x/4>x^{1/2}$
for $x>x_0$.  Choose $0<r<1$ with
\begin{equation}
\label{eq:Crandall-r-choice}
 r\log_2(3/64)+\frac12>3\mathrm e r>0.
\end{equation}
If also $x>2^{h/r}$, then $h<r\log_2x$.  Substituting this strict inequality
in the preliminary bound and using $h!\mathrm e^h>h^h$ yields
\eqref{eq:Crandall-fixed-height-bound}.  This follows the source proof while
making its strict threshold and the repaired Lemma~5.2 domain explicit
\cite[p.~1286, Thm.~5.1]{Crandall1978}.
\end{proof}

Let
\begin{equation}
\label{eq:Crandall-pi-total}
 \pi(x)=\#\{m\in\mathcal M:m\leq x\}
       =\sum_{h\geq1}\pi_h(x).
\end{equation}

\paragraph{The exact handbook route.}
Crandall's remark after Lemma~6.1 says only that ``various formulas''
pertaining to the lemma occur in reference~4; it gives no handbook section,
page, or formula number \cite[p.~1287]{Crandall1978}.  The relevant formulas
can nevertheless be identified exactly.  Abramowitz--Stegun use $P(z)$ for
the standard normal lower tail, but also use $P(\chi^2\mid\nu)$ and
$Q(\chi^2\mid\nu)=1-P(\chi^2\mid\nu)$ for the lower and upper chi-square
probabilities.  Their formulas 26.2.2--26.2.6 give $P(0)=Q(0)=1/2$; formula
26.2.29 relates this $P$ to $\operatorname{erf}$; formula 26.4.11 gives the
large-$\nu$ normal coordinate
\[
 \frac{\chi^2-\nu}{\sqrt{2\nu}};
\]
and formula 26.4.21 gives the exact finite Poisson sum below
\cite[printed pp.~931, 934, 940--941]{AbramowitzStegun1964}.  These are
identified relevant formulas, not locators supplied by Crandall.

\begin{proposition}[exact Poisson--gamma--chi-square coordinate]
\label{prop:AS-Poisson-gamma-morphism}
For $n\in\Nzero$ and $t\geq0$, set
\begin{equation}
\label{eq:AS-Poisson-gamma-morphism}
 F_n(t)=e^{-t}\sum_{j=0}^{n}\frac{t^j}{j!}.
\end{equation}
Then, with $\Gamma(a,t)$ denoting the upper incomplete gamma function,
\begin{equation}
\label{eq:AS-Poisson-gamma-identities}
 F_n(t)=\frac{\Gamma(n+1,t)}{\Gamma(n+1)}
       =Q(2t\mid2n+2).
\end{equation}
For fixed $n$,
\begin{equation}
\label{eq:AS-Poisson-gamma-derivative}
 F_n'(t)=-e^{-t}\frac{t^n}{n!}\qquad(t>0),
\end{equation}
so $F_n$ is a continuous strictly decreasing bijection
$[0,\infty)\to(0,1]$.  Its inverse is the unique upper-tail quantile on this
fixed-$n$ fibre; no elementary inverse formula is asserted.
\end{proposition}

\begin{proof}
Integration by parts gives
\[
 \Gamma(n+1,t)=n!e^{-t}\sum_{j=0}^{n}\frac{t^j}{j!}
\]
by induction from $\Gamma(1,t)=e^{-t}$, proving the first equality in
\eqref{eq:AS-Poisson-gamma-identities}.  In Abramowitz--Stegun 26.4.21 take
$c=n+1$, $\nu=2c$, $m=t$, and $\chi^2=2t$ to obtain the second equality
literally \cite[printed p.~941, Eq.~26.4.21]{AbramowitzStegun1964}.
Differentiating the finite sum makes every term cancel except the displayed
last term in \eqref{eq:AS-Poisson-gamma-derivative}.  Finally $F_n(0)=1$,
while $F_n(t)\to0$ because an exponential dominates the fixed polynomial;
strict monotonicity supplies the stated fibres and inverse.
\end{proof}

\begin{lemma}[completion of Crandall's Poisson lemma]
\label{lem:Crandall-Poisson-half-mass}
Let $k(t)$ be integer-valued and suppose
\begin{equation}
\label{eq:Crandall-Poisson-moving-cutoff}
 \frac{k(t)-t}{\sqrt t}\longrightarrow0.
\end{equation}
Then
\begin{equation}
\label{eq:Crandall-Poisson-half-mass}
 \lim_{t\to\infty}e^{-t}
 \sum_{u=0}^{k(t)}
 \frac{t^u}{u!}=\frac12.
\end{equation}
Consequently the same limit holds after deleting $u=0$ for each of the
following three exact upper cutoffs:
\begin{equation}
\label{eq:Crandall-Poisson-three-cutoffs}
 k(t)=\lfloor t\rfloor-1,
 \qquad
 k(t)=\lfloor t\rfloor,
 \qquad
 k(t)=\lceil t\rceil-1.
\end{equation}
The first is Crandall's printed strict cutoff
$u<\lfloor t\rfloor$; the third is the maximal integer cutoff uniformly
authorized by the strict height hypothesis in Theorem~5.1.
\end{lemma}

\begin{proof}
The handbook route to the source case is now literal.  Put
$n=\lfloor t\rfloor$, $\nu=2n$, and $\chi^2=2t$ in
Proposition~\ref{prop:AS-Poisson-gamma-morphism}.  Formula 26.4.21 identifies
the sum $u=0,\ldots,n-1$ with $Q(2t\mid2n)$, while the normal coordinate in
26.4.11 is
\[
 \frac{2t-2n}{\sqrt{4n}}=\frac{t-n}{\sqrt n}\longrightarrow0.
\]
Together with $P(0)=Q(0)=1/2$, this is the exact classical formula path to
Crandall's limit.  Formula 26.4.11 prints neither an error term nor a proof on
that page, however, so it is not used as an unspoken uniformity assertion.
The following argument controls all three real-$t$ endpoints directly.

Let $X_t$ have the Poisson distribution with mean $t$.  The characteristic
function of $Y_t=(X_t-t)/\sqrt t$ is
\[
 \mathbf E(e^{isY_t})
 =\exp\!\left(t\left(e^{is/\sqrt t}-1-is/\sqrt t\right)\right).
\]
The Taylor remainder is $O_s(t^{-1/2})$, so this converges to
$e^{-s^2/2}$.  The L\'evy continuity theorem gives
$Y_t\Rightarrow N(0,1)$.  Write $a_t=(k(t)-t)/\sqrt t$.  For every fixed
$\delta>0$, condition \eqref{eq:Crandall-Poisson-moving-cutoff} gives
\[
 \Pr(Y_t\leq-\delta)
 \leq \Pr(Y_t\leq a_t)
 \leq \Pr(Y_t\leq\delta)
\]
for all sufficiently large $t$.  Weak convergence at the two fixed continuity
points gives lower and upper limits between $\Phi(-\delta)$ and
$\Phi(\delta)$; letting $\delta\downarrow0$ proves
\eqref{eq:Crandall-Poisson-half-mass}.  Every cutoff in
\eqref{eq:Crandall-Poisson-three-cutoffs} differs from $t$ by at most a fixed
constant.  Removing the $u=0$ term changes the probability by exactly
$e^{-t}\to0$.

The high-resolution source page prints
$u<\lfloor t\rfloor$, not $u\leq\lfloor t\rfloor$.  Crandall points to
error-function asymptotics but prints no proof
\cite[p.~1287, Lem.~6.1]{Crandall1978}.  The argument above supplies the
all-real analytic completion while explicitly retaining its invocation of
L\'evy's continuity theorem.
\end{proof}

\begin{theorem}[Crandall's power-law theorem, with an explicit witness]
\label{thm:Crandall-power-law-count}
There are $c>0$ and $X_0>0$ such that
\begin{equation}
\label{eq:Crandall-power-law-count}
 \pi(x)>x^c\qquad(x>X_0).
\end{equation}
The source leaves $c$ unspecified.  Its proof, with the endpoint repaired,
permits the concrete edition choice $c=1/100$.
\end{theorem}

\begin{proof}
\phantomsection\label{sourceissue:Crandall-theorem61-endpoint}
For $x>x_0$, Theorem~\ref{thm:Crandall-fixed-height-count} may be summed only
over $h<r\log_2x$.  The printed proof sums through
$\lfloor r\log_2x\rfloor$, which includes one unauthorized equality term
when $r\log_2x$ is integral.  Summing instead through
$\lceil r\log_2x\rceil-1$ and applying
Lemma~\ref{lem:Crandall-Poisson-half-mass} gives, for every fixed
$d\in(0,1/2)$ and all sufficiently large $x$,
\begin{equation}
\label{eq:Crandall-power-intermediate}
 \pi(x)>(1/2-d)e^{r\log_2x}
       =(1/2-d)x^{r/\log2}.
\end{equation}
This proves the source theorem for every $c<r/\log2$ after increasing the
threshold.

For an exact explicit witness, take $r=1/100$.  Since
$3/64>2^{-5}$ and $\mathrm e<3$, the left side of
\eqref{eq:Crandall-r-choice} is greater than $45/100$, while its middle term
$3\mathrm e/100$ is less than $9/100$.  Thus this $r$ is admissible.  Since
$\log2<1$, one has $r/\log2>1/100$; the fixed positive factor in
\eqref{eq:Crandall-power-intermediate} is eventually absorbed by the excess
power.  Hence $c=1/100$ works.  This is an edition-supplied sharpening of the
source's existential conclusion, not a number printed by Crandall.
\end{proof}

\begin{nonclaim}
\label{nonclaim:Crandall-no-positive-density}
Theorem~\ref{thm:Crandall-power-law-count} is a lower bound of order $x^c$
for some $c<1$.  It is compatible with natural density zero and is not a
positive-density theorem.  It proves neither $\pi(x)\gg x$ nor convergence
for almost every odd starting value.  The frozen route text that calls this
positive density is routing metadata, not mathematical evidence.
\end{nonclaim}

\subsubsection{Membership coordinates and the return obstruction}

For a backward word $m=B_{b_k\cdots b_1}(n)$, Crandall uses the reverse
cumulative sums
\begin{equation}
\label{eq:Crandall-source-A-coordinates}
 A_0=0,
 \qquad
 A_i=\sum_{j=k-i+1}^{k}b_j\quad(1\leq i\leq k).
\end{equation}
Expanding the affine composition gives
\begin{equation}
\label{eq:Crandall-membership-identity}
 2^{A_k}n-3^km
 =\sum_{j=0}^{k-1}2^{A_j}3^{k-1-j}.
\end{equation}

\begin{theorem}[Crandall, exact trajectory-membership coordinates]
\label{thm:Crandall-membership}
For $m,n\in\Opos$, one has $n\in\mathcal T_m$ if and only if there are an
integer $k\geq1$ and integers
\[
 0=A_0<A_1<\cdots<A_k
\]
satisfying \eqref{eq:Crandall-membership-identity}.  If such data exist, then
$n=1$ or $n$ is the $k$-th displayed member of $\mathcal T_m$.
\end{theorem}

\begin{proof}
If $n=C_{\mathrm{Cr}}^k(m)$, reverse its $k$ branch equations and take
$b_i=e(C_{\mathrm{Cr}}^{k-i}(m))$; expansion gives
\eqref{eq:Crandall-membership-identity}.  Conversely, define
\begin{equation}
\label{eq:Crandall-recovered-differences}
 d_i=A_{k-i+1}-A_{k-i}\quad(1\leq i\leq k).
\end{equation}
The strict increase makes every $d_i$ positive, and solving
\eqref{eq:Crandall-membership-identity} gives
$m=B_{d_k\cdots d_1}(n)$.  Crandall's Lemma~4.1 then proves that every
formal inverse branch is integral, odd, and removed by the actual accelerated
map.  Hence $C_{\mathrm{Cr}}^k(m)=n$, with $n=1$ treated as the terminating
case \cite[pp.~1287--1288, Eq.~(7.2), Thm.~7.1]{Crandall1978}.
\end{proof}

\begin{sourceissue}[Theorem 7.1 difference index]
\label{sourceissue:Crandall-theorem71-index}
The printed converse assigns
$d_i=A_{k-i+1}-A_{k-1}$.  For $i=2$ this can give zero and it does not recover
successive coordinate gaps.  Equation~\eqref{eq:Crandall-recovered-differences}
is the forced correction.  It is exactly the difference needed to recover
$A_i$ from the $d$-word and makes the printed conclusion
$m=B_{d_k\cdots d_1}(n)$ valid.
\end{sourceissue}

Taking $n=m$ in \eqref{eq:Crandall-membership-identity} gives precisely
\eqref{eq:crandall-cycle-identity}.  The forward clause of Crandall's printed
Corollary~7.1 assumes that the trajectory has period $k$; its converse says
that the displayed equation produces a return at the $k$-th element.
Combining that converse with Theorem~\ref{thm:Crandall-membership} yields the
edition's slightly more general equivalence: the equation with strict
$A$-coordinates characterizes a return after $k$ accelerated steps.  It does
not assert that $k$ is the least period.  Positivity of the right side forces
\begin{equation}
\label{eq:Crandall-positive-cycle-denominator}
 2^{A_k}>3^k.
\end{equation}

\begin{proposition}[Crandall's explicit Diophantine obstruction]
\label{prop:Crandall-cycle-obstruction}
Suppose integers $a,b$ satisfy
\begin{equation}
\label{eq:Crandall-obstruction-hypothesis}
 b>1,
 \qquad
 0<2^a-3^b\mid 2^{a-b}-1.
\end{equation}
Put
\begin{equation}
\label{eq:Crandall-obstruction-construction}
 d=\frac{2^{a-b}-1}{2^a-3^b},
 \qquad
 m=2^bd-1.
\end{equation}
Then $m\in\Opos$, $m>1$, and $C_{\mathrm{Cr}}^b(m)=m$.  Therefore the
Collatz conjecture would be false.  The least period is allowed to divide
$b$.
\end{proposition}

\begin{proof}
The positive divisor relation forces $a>b$, so $d$ is a positive integer and
$m$ is positive odd.  If $m=1$, then $d=2^{1-b}$ is not an integer because
$b>1$; hence $m>1$.  Multiplying the definition of $d$ by $2^b$ gives
\[
 m(2^a-3^b)=3^b-2^b
 =\sum_{j=0}^{b-1}2^j3^{b-1-j}.
\]
This is the return equation with $k=b$, $A_j=j$ for $0\leq j<b$, and
$A_b=a$.  The strictness $A_{b-1}<A_b$ follows from $a>b$.  Corollary~
\ref{cor:crandall-cycle-identity} supplies the genuine return
\cite[p.~1288]{Crandall1978}.
\end{proof}

\subsubsection{Pillai's fixed-difference citation and the exact P\'olya route}

Crandall reports two external Diophantine results.  The first says that, for
each fixed integer $z$, the equation
\begin{equation}
\label{eq:Crandall-fixed-difference}
 2^x-3^y=z
\end{equation}
has only finitely many integer solutions $(x,y)$
\cite[ref.~9 and p.~1288]{Crandall1978}.  Pillai's cited primary paper has now
been read in full.  Its Corollary~1 states the corresponding positive-
exponent, positive-difference conclusion
\cite[p.~4]{Pillai1931}.  The paper also identifies the qualitative result as
a special case of P\'olya's earlier fixed-prime gap theorem
\cite[p.~2]{Pillai1931}.  Using P\'olya's sourced gap theorem yields exactly
the finiteness required by Crandall and bypasses Pillai's defective
quantitative derivation.  The proof-dependency audit of P\'olya's explicit
Thue citations is recorded separately below.

\begin{sourceissue}[Pillai's literal quantitative proof]
\label{sourceissue:Pillai-degree-gap}
Pillai quotes Siegel's approximation exponent correctly but calls its
constant a positive \emph{integer}; Siegel asserts only a positive real
constant \cite[p.~174]{Siegel1921}.  More seriously, Pillai puts
$\alpha=(A/B)^{1/r}$ and asserts that $\alpha$ has degree $r$ whenever $A$
and $B$ are not both perfect $r$th powers \cite[p.~3]{Pillai1931}.  This is
false: $(4/1)^{1/4}=\sqrt2$ has degree two, and
$(2/2)^{1/2}=1$ has degree one although neither displayed coefficient is a
square.  His next comparison bounds $Y/X$ only from above, although the step
which replaces $Y^{r-1}$ by a constant multiple of $X^{r-1}$ needs a lower
bound.  Finally, the proof of Theorem~I applies the lemma to every residue
pair without checking its ``not both powers'' hypothesis; the residue pair
$h=\ell=0$ with outer coefficients one violates it.  Thus the printed
derivation is not admitted literally.
\end{sourceissue}

The defect is repairable on the exact source theorem rather than by changing
Pillai's conclusion.  For $q\geq1$ write
\begin{equation}
\label{eq:Pillai-mq}
 \mu(q)=\min_{1\leq\lambda\leq q}
 \left(\frac{q}{\lambda+1}+\lambda\right).
\end{equation}

\begin{theorem}[Siegel's rational-approximation theorem]
\label{thm:Siegel-rational-approximation}
Let $\xi$ be algebraic of degree $q\geq2$ and let $\epsilon>0$.
There is a positive real constant $c=c(\xi,\epsilon)$ such that, for all
rational integers $u,v$ with $v>0$,
\[
 \left|\xi-\frac uv\right|
 >\frac{c}{v^{\mu(q)+\epsilon}}.
\]
The constant is not asserted to be an integer or numerically effective.
\end{theorem}

\begin{proof}
This is the denominator form printed on Siegel's p.~174
\cite[p.~174]{Siegel1921}.  His Hauptsatz is stated in height form on
printed p.~178 and proved through p.~191; Zus\"atze~1--3 give a uniform
positive constant and extend the conclusion to nonintegral algebraic
$\xi$ \cite[pp.~178, 191]{Siegel1921}.  Although the Hauptsatz minimizes over
$1\leq s<q$, this is the same minimum as in \eqref{eq:Pillai-mq}: the
excluded endpoint has value $q+q/(q+1)>q$, whereas the candidate $s=1$ has
value $q/2+1\leq q$.  If $u/v=p/s$ is reduced, then
$s\leq v$.  In the error-$<1$ case,
$H(p/s)=\max\{|p|,s\}\leq\max\{1,|\xi|+1\}s$; the complementary error
case is absorbed by decreasing the positive constant.  This is the exact
height-to-denominator map, with no assertion of an effective constant.
\end{proof}

\begin{lemma}[Monotonicity of Siegel's exponent]
\label{lem:Siegel-exponent-monotone}
The sequence $\mu(q)$ is nondecreasing on $\Npos$.
\end{lemma}

\begin{proof}
For an old candidate $1\leq\lambda\leq q$, replacing $q$ by $q+1$
increases $q/(\lambda+1)+\lambda$ by $1/(\lambda+1)$.  The only new
candidate is $\lambda=q+1$, whose value
$(q+1)/(q+2)+q+1$ is greater than $q+1$.  But
$\mu(q)\leq q/2+1\leq q+1$, using the old candidate $\lambda=1$.
Every candidate for $\mu(q+1)$ is therefore at least $\mu(q)$.
\end{proof}

\begin{proposition}[Actual-degree repair of Pillai's lemma]
\label{prop:Pillai-actual-degree-repair}
Fix $A,B,r\in\Npos$ with $r\geq2$ and fix $\eta>0$.  There is a constant
$C=C(A,B,r,\eta)>0$ such that, for all $X,Y\in\Npos$ with
$AX^r>BY^r$,
\begin{equation}
\label{eq:Pillai-repaired-power-separation}
 AX^r-BY^r>CX^{r-\mu(r)-\eta}.
\end{equation}
Consequently, for every $\delta>0$ the coefficient-free inequality
$AX^r-BY^r>X^{r-\mu(r)-\delta}$ holds for all sufficiently large $X$.
No hypothesis that $A$ and $B$ are not both $r$th powers is required.
\end{proposition}

\begin{proof}
Put $\alpha=(A/B)^{1/r}$ and
$d=[\Q(\alpha):\Q]\leq r$.  The exact factorization is
\begin{equation}
\label{eq:Pillai-exact-factorization}
 AX^r-BY^r
 =B(\alpha X-Y)
   \sum_{j=0}^{r-1}(\alpha X)^{r-1-j}Y^j.
\end{equation}
If $d\geq2$, Siegel's theorem at the \emph{actual} degree gives a constant
$c>0$ for which
\[
 \left|\alpha-\frac YX\right|>cX^{-\mu(d)-\eta}.
\]
The difference is positive, and Lemma~\ref{lem:Siegel-exponent-monotone}
gives $\mu(d)\leq\mu(r)$.  Hence
$\alpha X-Y>cX^{1-\mu(r)-\eta}$.  Retaining the positive $j=0$ term in
the sum in \eqref{eq:Pillai-exact-factorization} gives
\[
 AX^r-BY^r
 >Bc\alpha^{r-1}X^{r-\mu(r)-\eta}.
\]
If $d=1$, write $\alpha=P/Q$ in lowest positive terms.  Strict positivity
excludes $PX=QY$, so $PX-QY\geq1$ and $\alpha X-Y\geq1/Q$.  Since $r\geq2$
and $Y>0$, the factor sum in \eqref{eq:Pillai-exact-factorization} is strictly
larger than its $j=0$ term.  Hence
\[
 AX^r-BY^r>\frac{B\alpha^{r-1}}QX^{r-1}.
\]
Moreover $\mu(r)>1$, because every candidate in \eqref{eq:Pillai-mq} has
$\lambda\geq1$ and a strictly positive first summand.  The displayed bound
therefore implies \eqref{eq:Pillai-repaired-power-separation} for $X\geq1$.
Finally, with $\eta=\delta/2$, the ratio of the established lower bound to
$X^{r-\mu(r)-\delta}$ is $CX^{\delta/2}$, which exceeds one for
$X>C^{-2/\delta}$.
\end{proof}

\begin{theorem}[Repaired quantitative Pillai theorem]
\label{thm:Pillai-repaired-quantitative}
Fix $A,B,M,N\in\Npos$ with $M,N>1$ and $\delta>0$, and assume, as Pillai
does, that $\log M/\log N\notin\Q$.  Then there is an integer
$x_0=x_0(A,B,M,N,\delta)\geq0$ such that
\begin{equation}
\label{eq:Pillai-repaired-theorem-I}
 AM^x-BN^y>M^{x(1-\delta)}
\end{equation}
whenever $x,y\in\Npos$, $x>x_0$, and $AM^x>BN^y$.
\end{theorem}

\begin{proof}
Replace $\delta$ by $\delta_0=\min(\delta,1/2)$.  The required large $r$
exists constructively: for $k=\lceil\sqrt r\rceil$ and
$\lambda=k-1\geq1$,
\[
 \mu(r)\leq\frac r{k}+k-1<2\sqrt r,
\]
so $\mu(r)/r<2/\sqrt r$.  Choose, for example, an integer
$r>(4/\delta_0)^2$, increasing it to at least two if necessary.  Then
$\mu(r)/r<\delta_0/2$.  Choose $\eta>0$ with
$\eta/r<\delta_0/4$.  Write uniquely
\[
 x=rs+h,\qquad y=rt+\ell,
 \qquad 0\leq h,\ell<r.
\]
For the fixed residue pair $(h,\ell)$, put
$A_h=AM^h$, $B_\ell=BN^\ell$, $X=M^s$, and $Y=N^t$.
Proposition~\ref{prop:Pillai-actual-degree-repair} gives
\[
 AM^x-BN^y
 >C_{h,\ell}M^{(x-h)(1-(\mu(r)+\eta)/r)}.
\]
Put
\[
 C_*=\min_{0\leq h,\ell<r}C_{h,\ell}>0,
 \qquad
 \gamma=1-\frac{\mu(r)+\eta}{r},
 \qquad
 \varepsilon=\delta_0-\frac{\mu(r)+\eta}{r}>\frac{\delta_0}{4}.
\]
Here $\gamma>1-3\delta_0/4>0$.  For every residue pair,
\[
 \frac{C_{h,\ell}M^{(x-h)\gamma}}{M^{x(1-\delta_0)}}
 =C_{h,\ell}M^{x\varepsilon-h\gamma}
 \geq C_*M^{x\varepsilon-(r-1)\gamma}.
\]
The last quantity exceeds one whenever
\[
 x>\frac{(r-1)\gamma-\log_M C_*}{\varepsilon},
\]
a threshold independent of $y,h,\ell$.  This proves
$AM^x-BN^y>M^{x(1-\delta_0)}\geq M^{x(1-\delta)}$.
The preceding proof establishes the same quantified conclusion for all
$M,N>1$, without the irrational-log hypothesis.  This stronger statement is
an edition result; Pillai's sourced theorem remains recorded with his
hypothesis.
\end{proof}

We now type the independent qualitative route used for Crandall.  For a
finite ordered tuple $S=(p_1,\ldots,p_r)$ of pairwise distinct primes define
\[
 \mathcal A_S
 =\{p_1^{e_1}\cdots p_r^{e_r}:e_i\in\Nzero\}.
\]

\begin{proposition}[Prime-exponent coordinate morphism]
\label{prop:Polya-smooth-coordinate}
The map
\[
 E_S:(\Nzero^r,+)\longrightarrow(\mathcal A_S,\mathord\cdot),
 \qquad
 E_S(e_1,\ldots,e_r)=\prod_{i=1}^r p_i^{e_i},
\]
is a commutative-monoid isomorphism.  Its inverse is
$a\mapsto(\nu_{p_1}(a),\ldots,\nu_{p_r}(a))$; every fibre is a singleton.
Thus the coordinate loses no exponent information while the prime set $S$
is retained.
\end{proposition}

\begin{proof}
Multiplication in $\mathcal A_S$ adds every valuation, so $E_S$ preserves
the monoid law.  Existence and uniqueness of prime factorization prove both
surjectivity and the displayed two-sided inverse.
\end{proof}

\begin{theorem}[Thue's fixed-value theorem at P\'olya's exact scopes]
\label{thm:Thue-fixed-binary-form}
Let $c\in\Z\setminus\{0\}$.
\begin{enumerate}
\item If $U\in\Z[X,Y]$ is homogeneous and irreducible of degree greater
than two, then $U(x,y)=c$ has only finitely many solutions in $\Z^2$.
\item More generally, if a homogeneous integral binary form $F$ has
infinitely many integer solutions of $F(x,y)=c$, then, up to multiplication
by a nonzero scalar, $F$ is a power of a linear form or an indefinite
quadratic form.
\end{enumerate}
\end{theorem}

\begin{proof}
Thue's Theorem~IV is printed for positive $x,y$
\cite[pp.~303--304]{Thue1909}.  On the off-axis domain, the map
\[
 (x,y)\longmapsto
 \bigl(\operatorname{sgn}x,\operatorname{sgn}y,|x|,|y|\bigr)
\]
is a bijection
\[
 (\Z\setminus\{0\})^2
 \longrightarrow
 \{\pm1\}^2\times\Npos^2,
\]
with inverse $(\epsilon,\delta,X,Y)\mapsto(\epsilon X,\delta Y)$.
For each of the four sign pairs, the transformed form
$U_{\epsilon,\delta}(X,Y)=U(\epsilon X,\delta Y)$ remains integral,
homogeneous, irreducible, and of the same degree, so Thue's positive-quadrant
theorem gives a finite set.  If $y=0$, irreducibility in degree greater than
one forces the coefficient $u_d$ of $X^d$ to be nonzero, and the remaining
equation is $u_dx^d=c$; the other axis is identical.  This proves the exact
all-integer extension in part~1.

Part~2 is the exceptional-form formulation quoted by P\'olya
\cite[pp.~144--145]{Polya1918}.  His locator is Thue II, printed p.~3.  Both
available primary scans omit printed pp.~2--3, so that page is not represented
as visually verified.  Thue II printed p.~1 gives the equivalent
dehomogenized theorem and its degree-one/degree-two power exceptions
\cite[p.~1]{Thue1908II}.  The manifestation gap, P\'olya's exact quotation,
and the relation between the two formulations are recorded in
\path{qa/THUE-1908-1909-F017-binary-form-audit.md}.
\end{proof}

\begin{proposition}[P\'olya's residue-cell cubic morphism]
\label{prop:Polya-residue-cell-cubic}
Let
\[
 L_1(n)=an+b,\qquad L_2(n)=cn+d,
\]
where $a,b,c,d\in\Z$, $a,c\ne0$, and
$\Delta=ad-bc\ne0$.  Fix pairwise distinct primes
$p_1,\ldots,p_t$ and residue vectors
$\mathbf r,\mathbf s\in\{0,1,2\}^t$.  Let
$D_{\mathbf r,\mathbf s}$ be the set of integers $n$ for which
$L_1(n),L_2(n)>0$, both values have prime support in
$\{p_1,\ldots,p_t\}$, and
\[
 \nu_{p_i}(L_1(n))\equiv r_i\pmod3,\qquad
 \nu_{p_i}(L_2(n))\equiv s_i\pmod3
 \quad(1\leq i\leq t).
\]
Put
\[
 R=\prod_{i=1}^t p_i^{r_i},\qquad
 S_0=\prod_{i=1}^t p_i^{s_i},\qquad
 G_{\mathbf r,\mathbf s}(X,Y)=aS_0Y^3-cRX^3.
\]
Then
\[
 \Phi_{\mathbf r,\mathbf s}:D_{\mathbf r,\mathbf s}
 \longrightarrow
 \mathcal Z_{\mathbf r,\mathbf s}
 =\{(X,Y)\in\Npos^2:G_{\mathbf r,\mathbf s}(X,Y)=\Delta\},
\]
defined by
\[
 X=\prod_{i=1}^t
 p_i^{(\nu_{p_i}(L_1(n))-r_i)/3},
 \qquad
 Y=\prod_{i=1}^t
 p_i^{(\nu_{p_i}(L_2(n))-s_i)/3},
\]
is well defined and injective.  Its inverse on its image is
\[
 n=\frac{RX^3-b}{a}=\frac{S_0Y^3-d}{c}.
\]
The target $\mathcal Z_{\mathbf r,\mathbf s}$ is finite.
\end{proposition}

\begin{proof}
The residue hypotheses make the displayed exponents nonnegative integers and
give the exact identities
\[
 an+b=RX^3,\qquad cn+d=S_0Y^3.
\]
Multiplying the second identity by $a$, the first by $c$, and subtracting
gives
\[
 aS_0Y^3-cRX^3=ad-bc=\Delta.
\]
The two displayed recovery formulas agree by this equation and recover the
original $n$, proving injectivity and singleton fibres on the image.

The cubic has
\[
 \operatorname{Disc}(G_{\mathbf r,\mathbf s})
 =-27(aS_0)^2(cR)^2\ne0.
\]
It need not be irreducible over $\Q$, so part~1 of
Theorem~\ref{thm:Thue-fixed-binary-form} cannot be invoked alone.  Suppose
instead that, over a characteristic-zero field,
\[
 G_{\mathbf r,\mathbf s}(X,Y)=\lambda(uX+vY)^3.
\]
Its nonzero $X^3$ and $Y^3$ coefficients force
$\lambda,u,v\ne0$, whereas its zero $X^2Y$ coefficient would give
$3\lambda u^2v=0$, a contradiction.  It is therefore not a scalar linear
cube.  A quadratic power cannot have degree three.  Part~2 of
Theorem~\ref{thm:Thue-fixed-binary-form} now proves finiteness.  If one uses
P\'olya's positive-right-side phrasing and $\Delta<0$, the exact involution
\[
 (G_{\mathbf r,\mathbf s},\Delta,X,Y)
 \longmapsto
 (-G_{\mathbf r,\mathbf s},-\Delta,X,Y)
\]
preserves the solution set and exceptional-form status.
\end{proof}

\begin{theorem}[P\'olya's fixed-prime gap theorem]
\label{thm:Polya-smooth-gaps}
Let $r\geq2$ and enumerate $\mathcal A_S$ increasingly as
$a_0<a_1<a_2<\cdots$.  Then
\begin{equation}
\label{eq:Polya-smooth-gaps}
 \lim_{j\to\infty}(a_{j+1}-a_j)=\infty.
\end{equation}
\end{theorem}

\begin{proof}
This is the exact statement in P\'olya's Section~4, equations (12)--(14)
\cite[pp.~147--148]{Polya1918}.  P\'olya proves it equivalent to his Satz~I
for products of two essentially distinct rational linear factors.  His proof
of Satz~I reduces finitely supported prime-exponent vectors modulo three,
fixes a residue vector, and invokes Thue's finiteness theorem for the resulting
binary cubic \cite[pp.~144--145]{Polya1918}.  Proposition~
\ref{prop:Polya-residue-cell-cubic} verifies the exact dependency, including
the reducible-cubic cells not covered by Thue's irreducible-form theorem
alone.  Thue I, Satz~12, records the same modulo-three proof template
\cite[pp.~30--31]{Thue1908I}; P\'olya supplies the residue-cell reduction used
here.  No effective gap threshold is imported.
\end{proof}

\begin{corollary}[Finite fixed-difference fibres]
\label{cor:Polya-fixed-difference-fibres}
For every nonempty finite prime set $S$ and every fixed
$z\in\Z\setminus\{0\}$, the
fibre
\[
 \Delta_S^{-1}(z)
 =\{(u,v)\in\mathcal A_S^2:u-v=z\},
 \qquad \Delta_S(u,v)=u-v,
\]
is finite.
\end{corollary}

\begin{proof}
If $|S|\geq2$, choose an index after which every adjacent gap in
\eqref{eq:Polya-smooth-gaps} exceeds $|z|$.  The difference between any two
distinct later terms is a sum of one or more adjacent positive gaps, hence
also exceeds $|z|$.  Only pairs with at least one earlier coordinate remain.
If $a_0,\ldots,a_{J-1}$ are the earlier terms, then choosing one of these
values and which coordinate of $(u,v)$ it occupies leaves the other
coordinate uniquely determined by $u-v=z$; hence at most $2J$ pairs remain.
If
$S=\{p\}$, use directly $p^{j+1}-p^j=(p-1)p^j\to\infty$.
\end{proof}

\begin{proposition}[Exact fixed-$z$ consequence of P\'olya's sourced theorem]
\label{prop:Crandall-fixed-z-finiteness}
For every $z\in\Z$, equation \eqref{eq:Crandall-fixed-difference} has only
finitely many solutions $(x,y)\in\Z^2$.
\end{proposition}

\begin{proof}
First prove the exponent-domain reduction.  If $x=-a<0$ and $y\geq0$, then
\[
 2^x-3^y=\frac{1-2^a3^y}{2^a}
\]
has odd numerator and is not an integer.  If $x\geq0$ and $y=-b<0$, the
numerator of $(3^b2^x-1)/3^b$ is $-1$ modulo $3$.  If both exponents are
negative, then
\[
 2^{-a}-3^{-b}=\frac{3^b-2^a}{2^a3^b}
\]
has absolute value strictly between zero and one: its numerator is nonzero
by unique factorization and has absolute value smaller than the denominator.
Thus an integer-valued difference forces $x,y\in\Nzero$.

Define
\[
 I:\Nzero^2\longrightarrow\mathcal A_{\{2,3\}}^2,
 \qquad I(x,y)=(2^x,3^y).
\]
The valuation map
\[
 V(u,v)=(\nu_2(u),\nu_3(v))
\]
satisfies $V\circ I=\operatorname{id}_{\Nzero^2}$, so $I$ is injective.  For
$z\neq0$,
the solution set is exactly
\[
 I^{-1}\!\left(\Delta_{\{2,3\}}^{-1}(z)\right).
\]
Corollary~\ref{cor:Polya-fixed-difference-fibres} makes the inner fibre
finite, and injectivity makes its preimage finite.  For $z=0$, unique
factorization gives only
$2^x=3^y=1$, hence $(x,y)=(0,0)$.  This proves Crandall's literal integer-
exponent statement; it is not merely the positive-$z$ statement printed in
Pillai's Corollary~1.
\end{proof}

\subsubsection{Herschfeld's eventual-uniqueness theorem}

Crandall's reference~10 spells the author's surname ``Herschefeld.''  The
primary paper's title and byline give \emph{Aaron Herschfeld}
\cite[pp.~231--234]{Herschfeld1936}.  Crandall reports at most one integer
solution for sufficiently large $z$ \cite[p.~1288]{Crandall1978}.  The
primary theorem is stronger in the sign parameter but narrower in the
exponent domain: it treats sufficiently large $|d|$ and positive exponents.
The exact domain bridge is proved below rather than hidden in the citation.

\begin{sourceissue}[Herschfeld source boundaries]
\label{sourceissue:Herschfeld-boundaries}
On printed p.~232, a congruence argument for the specialized equation (2)
ends by referring to equation (1); the objects and displayed congruence make
the intended label unambiguous, but the defect is retained.  Printed p.~233
reports a further $|d|\leq100$ calculation from an unpublished paper; it is
not used here.  The final paragraph's general at-most-nine result for
$a^x-b^y=d$ has a separate Siegel-1929 dependency and is not a premise of the
special $(2,3)$ theorem.  That general result remains outside the admitted
scope until its exact cited theorem and coefficient hypotheses are read.
\end{sourceissue}

\begin{lemma}[The two order coordinates]
\label{lem:Herschfeld-orders}
For $r\geq1$ and $s\geq3$,
\begin{equation}
\label{eq:Herschfeld-orders}
 \operatorname{ord}_{3^r}(2)=2\cdot3^{r-1},
 \qquad
 \operatorname{ord}_{2^s}(3)=2^{s-2}.
\end{equation}
\end{lemma}

\begin{proof}
If $2^n\equiv1\pmod{3^r}$, reduction modulo three shows that $n=2m$ is
even.  The lifting-the-exponent identity gives
\[
 v_3(2^{2m}-1)=v_3(4^m-1)=1+v_3(m).
\]
Thus $3^r\mid2^n-1$ exactly when $2\cdot3^{r-1}\mid n$.  If
$3^n\equiv1\pmod{2^s}$ with $s\geq3$, reduction modulo eight shows that
$n$ is even, and
\[
 v_2(3^n-1)=v_2(3-1)+v_2(3+1)+v_2(n)-1=2+v_2(n).
\]
Therefore $2^s\mid3^n-1$ exactly when $2^{s-2}\mid n$.  This also makes
precise Herschfeld's older phrase that $3$ ``belongs to'' $2^{s-2}$ modulo
$2^s$ \cite[pp.~232--234]{Herschfeld1936}.
\end{proof}

\begin{theorem}[Herschfeld's positive-exponent theorem, with dependency repaired]
\label{thm:Herschfeld-positive-eventual-uniqueness}
There is a real number $D>0$ such that, for every $d\in\Z$ with $|d|>D$,
\begin{equation}
\label{eq:Herschfeld-positive-fibre}
 \#\{(x,y)\in\Npos^2:2^x-3^y=d\}\leq1.
\end{equation}
The proof supplies no effective numerical value of $D$.
\end{theorem}

\begin{proof}
Fix $0<\delta<1/2$.  Apply Theorem~
\ref{thm:Pillai-repaired-quantitative} first to bases $(2,3)$ and then to
$(3,2)$.  Let $x_1,y_1\in\Nzero$ be thresholds such that
\begin{align}
 0<2^u-3^v,\quad u>x_1
 &\Longrightarrow 2^u-3^v>2^{u(1-\delta)},
 \label{eq:Herschfeld-Pillai-positive}\\
 0<3^v-2^u,\quad v>y_1
 &\Longrightarrow 3^v-2^u>3^{v(1-\delta)}.
 \label{eq:Herschfeld-Pillai-negative}
\end{align}
Set
\begin{equation}
\label{eq:Herschfeld-threshold}
 D=\max\{2^{x_1+5},3^{y_1+2}\}.
\end{equation}

Suppose two distinct solutions have the same $d$.  They cannot agree in one
coordinate without agreeing in both.  Relabel them so that $X>x$; then
\[
 2^X-2^x=3^Y-3^y>0,
\]
so $Y>y$.  Factoring and using coprimality gives
\begin{equation}
\label{eq:Herschfeld-two-solution-factorization}
 2^x(2^{X-x}-1)=3^y(3^{Y-y}-1),
\end{equation}
and hence
\begin{equation}
\label{eq:Herschfeld-two-orders}
 2^{X-x}\equiv1\pmod{3^y},
 \qquad
 3^{Y-y}\equiv1\pmod{2^x}.
\end{equation}
Lemma~\ref{lem:Herschfeld-orders} yields
\begin{equation}
\label{eq:Herschfeld-exponent-separation}
 X-x\geq2\cdot3^{y-1},
 \qquad
 Y-y\geq2^{x-2}\quad(x>2).
\end{equation}

If $d>0$, then $d>2^{x_1+5}$ and
$d=2^x-3^y<2^x$, so $x>x_1+5>5$.  Therefore
\[
 2^X=3^Y+d>3^Y
 \geq3^{2^{x-2}}>2^{2^{x-2}},
 \qquad X>2^{x-2}>2x.
\]
Applying \eqref{eq:Herschfeld-Pillai-positive} at the larger solution gives
\[
 d=2^X-3^Y>2^{X(1-\delta)}>2^{X/2}>2^x>d,
\]
a contradiction.

If $d<0$, then $|d|>3^{y_1+2}$ and
$|d|=3^y-2^x<3^y$, so $y>y_1+2>2$.  Now
\[
 3^Y=2^X-d>2^X
 \geq2^{2\cdot3^{y-1}}>3^{3^{y-1}},
 \qquad Y>3^{y-1}>2y.
\]
Applying \eqref{eq:Herschfeld-Pillai-negative} at the larger solution gives
\[
 |d|=3^Y-2^X>3^{Y(1-\delta)}>3^{Y/2}>3^y>|d|,
\]
again a contradiction.  The elementary inequalities used above are strict
for exactly the displayed ranges: $2^{x-2}>2x$ for $x>5$ and
$3^{y-1}>2y$ for $y>2$.  This reconstructs Herschfeld's printed proof while
replacing only its reliance on Pillai's defective literal derivation by the
separately attributed repair already proved in this edition
\cite[pp.~233--234]{Herschfeld1936}.
\end{proof}

The next statement is an edition consequence, not part of Herschfeld's
printed exponent domain.

\begin{proposition}[Exact all-integer extension]
\label{prop:Herschfeld-integer-eventual-uniqueness}
For the same $D$ as in \eqref{eq:Herschfeld-threshold} and every
$d\in\Z$ with $|d|>D$,
\begin{equation}
\label{eq:Herschfeld-integer-fibre}
 \#\{(x,y)\in\Z^2:2^x-3^y=d\}\leq1.
\end{equation}
\end{proposition}

\begin{proof}
The exponent-domain argument in Proposition~
\ref{prop:Crandall-fixed-z-finiteness} proves that any integer-valued
difference has $x,y\in\Nzero$.  It remains to include a possible zero
coordinate without importing it into Herschfeld's statement.

Suppose two nonnegative solutions are ordered as in the preceding proof.  If
$d>0$, then $d<2^x$ still forces $x>x_1+5>5$, even if $y=0$.  The second
congruence in \eqref{eq:Herschfeld-two-orders} and the second bound in
\eqref{eq:Herschfeld-exponent-separation} remain valid, and the larger pair
$(X,Y)$ has positive coordinates.  The positive-$d$ contradiction is
unchanged.  If $d<0$, then $|d|<3^y$ still forces $y>y_1+2>2$, even if
$x=0$.  The first congruence and first order bound remain valid, the larger
pair is positive, and the negative-$d$ contradiction is unchanged.  Thus
the nonnegative fibre has size at most one.  Composing this conclusion with
the proved integer-to-nonnegative exponent reduction gives
\eqref{eq:Herschfeld-integer-fibre}.  This proves the exact domain used by
Crandall and strengthens his one-sided large-$z$ wording to Herschfeld's
two-sided large-$|d|$ parameter without changing either source statement.
\end{proof}

\begin{nonclaim}
Herschfeld's eventual uniqueness is not inferred from P\'olya's fixed-$d$
finiteness, and fixed-$d$ finiteness is not inferred from eventual
uniqueness.  The threshold in \eqref{eq:Herschfeld-threshold} is ineffective;
the theorem supplies neither an effective search bound nor a Collatz
endpoint or cycle exclusion.
\end{nonclaim}

\begin{proposition}[Exact primitive-root solution-to-return construction]
\label{prop:Crandall-primitive-root-return}
Put $t=\log_2 3$.  Let $p$ be an odd prime such that
\(\operatorname{ord}_p(2)=p-1\), and suppose
\begin{equation}
\label{eq:Crandall-primitive-root-input}
 (x,y)\in\Z^2,
 \qquad 2^x-3^y=p,
 \qquad y\geq\frac{p}{t-1}.
\end{equation}
Then one can construct integers
\[
 0=A_0<A_1<\cdots<A_y=x
\]
and a positive odd integer $m>1$ such that
\begin{equation}
\label{eq:Crandall-primitive-root-return}
 m(2^{A_y}-3^y)
 =\sum_{j=0}^{y-1}2^{A_j}3^{y-1-j}.
\end{equation}
Consequently $m$ occurs in its own Crandall trajectory, at accelerated
return time $y$; its least period is only asserted to divide $y$.
\end{proposition}

\begin{proof}
The integer-exponent denominator argument in the proof of Proposition~
\ref{prop:Crandall-fixed-z-finiteness} gives $x,y\in\Nzero$.  Under the last
hypothesis in \eqref{eq:Crandall-primitive-root-input}, $y>0$.  The equality
then excludes $p=3$ after reduction modulo $3$, so $p\geq5$.  Since
$0<t-1<1$, the last hypothesis also gives $y>p\geq5$; in particular all
index ranges below are nonempty.

Because $2^x>3^y$, one has $x>ty$.  Hence the integer
\begin{equation}
\label{eq:Crandall-primitive-root-slack}
 L=x-y>(t-1)y\geq p
 \qquad\Longrightarrow\qquad L\geq p+1.
\end{equation}
Equivalently, the assumed lower bound on $y$ is the exact integer inequality
$3^y\geq2^{p+y}$; no decimal approximation to $t$ is being used.

First put $A_j=j$ for $0\leq j\leq y-2$, and set
\[
 R_0=\sum_{j=0}^{y-2}2^j3^{y-1-j}.
\]
If $R_0\not\equiv0\pmod p$, choose $A_{y-1}$ to be the least integer in
$[y-1,x-1]$ satisfying
\begin{equation}
\label{eq:Crandall-primitive-root-selector-first}
 2^{A_{y-1}}\equiv-R_0\pmod p.
\end{equation}
Such an exponent exists: the interval contains at least $p-1$ consecutive
integers by \eqref{eq:Crandall-primitive-root-slack}, and
\(\operatorname{ord}_p(2)=p-1\) says exactly that the powers of $2$ over any
complete exponent-residue system modulo $p-1$ biject onto
\((\Z/p\Z)^\times\).

If instead $R_0\equiv0\pmod p$, keep $A_j=j$ for $0\leq j\leq y-3$, replace
$A_{y-2}$ by $y-1$, and write
\[
 R_1=\sum_{j=0}^{y-3}2^j3^{y-1-j}+2^{y-1}3.
\]
The replacement changes the partial sum by
\[
 R_1-R_0=3\cdot2^{y-2}\not\equiv0\pmod p,
\]
so $R_1\not\equiv0\pmod p$.  Choose $A_{y-1}$ to be the least integer in
$[y,x-1]$ satisfying
\begin{equation}
\label{eq:Crandall-primitive-root-selector-second}
 2^{A_{y-1}}\equiv-R_1\pmod p.
\end{equation}
This interval again contains at least $p-1$ consecutive integers, since its
cardinality is $L\geq p+1$.  Thus the selector is defined in both cases, and
the resulting sequence is strictly increasing with $A_0=0$ and $A_y=x$.

Let
\[
 S=\sum_{j=0}^{y-1}2^{A_j}3^{y-1-j}.
\]
The selected congruence gives $p\mid S$.  The $j=0$ summand is odd and every
other summand is even, so $S$ and $m=S/p$ are odd.  Moreover $0<t-1<1$
gives $y>p$, whence
\[
 S\geq3^{y-1}\geq3^p>p;
\]
The last strict inequality follows by induction from $3^n>n$ for $n\geq1$;
therefore $m$ is a positive odd integer greater than $1$.  Finally,
$2^{A_y}-3^y=2^x-3^y=p$, so $mp=S$ is exactly
\eqref{eq:Crandall-primitive-root-return}.  Corollary~7.1, in its exact form
proved above, now gives the asserted self-return.
\end{proof}

\begin{corollary}[Crandall's conditional primitive-root bound]
\label{cor:Crandall-primitive-root-bound}
Assume Conjecture~2.1.  If $p$ is a prime for which $2$ is a primitive root,
then every integer solution of $2^x-3^y=p$ satisfies
\[
 y<\frac{p}{\log_2 3-1}.
\]
\end{corollary}

\begin{proof}
A solution violating the strict bound satisfies the hypotheses of
Proposition~\ref{prop:Crandall-primitive-root-return} and therefore produces
a positive odd $m>1$ of infinite height, contrary to Conjecture~2.1.
\end{proof}

\begin{nonclaim}
\label{nonclaim:Crandall-primitive-root}
The source prints the conditional bound but omits this construction.  The
proposition does not assert that any solution violating the bound exists, and
it is not the adjacent sufficient obstruction with $A_j=j$.  The selected
return time $y$ need not be the least period, and the corollary remains
conditional on the Collatz conjecture.
\end{nonclaim}

\subsubsection{Continued fractions and the repaired cycle bound}

Put $t=\log_2 3$, and let $p_n/q_n$ be its continued-fraction convergents,
with
\begin{equation}
\label{eq:Crandall-convergent-recurrence}
 p_{-1}=1,\ q_{-1}=0,\ p_0=a_0,\ q_0=1,
 \qquad
 p_n=a_np_{n-1}+p_{n-2},\quad
 q_n=a_nq_{n-1}+q_{n-2}.
\end{equation}

\begin{lemma}[Exact continued-fraction transfer for the Crandall variables]
\label{lem:Crandall-CF-dependency}
The number $t=\log_2 3$ is irrational.  Let $p_n/q_n$ be its simple
continued-fraction convergents, with the recurrence
\eqref{eq:Crandall-convergent-recurrence}.  Then, for $n\geq2$,
\begin{equation}
\label{eq:Crandall-CF-error-bounds}
 \frac{1}{q_n+q_{n+1}}
 < |p_n-q_nt|
 < \frac{1}{q_{n+1}}.
\end{equation}
If $x\in\Z$ and $0<y<q_n$, then
\begin{equation}
\label{eq:Crandall-CF-best-approximation}
 |p_n-q_nt|<|x-yt|.
\end{equation}
For $n\geq4$, if $y=q_n$ and $x\in\Z$ with $x\ne p_n$, then the same
strict inequality holds with $y=q_n$.
\end{lemma}

\begin{proof}
Since $t>1$, a rational representation $t=a/b$ with $b>0$ would have
$a>0$ and $2^a=3^b$, contradicting unique factorization.  The continued
fraction is therefore infinite.  Steuding's Theorems~3.1--3.2
\cite[printed p.~38]{Steuding2005} give the
displayed recurrence and
\[
 p_nq_{n-1}-p_{n-1}q_n=(-1)^{n-1},
\]
so each $p_n/q_n$ is reduced.  Steuding's Theorem~3.6
\cite[printed p.~42]{Steuding2005}, applied to the
complete tail $\alpha_{n+1}$ of $t$, gives
\[
 t-\frac{p_n}{q_n}
 =\frac{(-1)^n}{q_n(\alpha_{n+1}q_n+q_{n-1})}.
\]
Because $a_{n+1}<\alpha_{n+1}<a_{n+1}+1$ and
$q_{n+1}=a_{n+1}q_n+q_{n-1}$, multiplying by $q_n$ yields
\eqref{eq:Crandall-CF-error-bounds}.  For $x>0$, Steuding's
Theorem~3.8 \cite[printed pp.~44--45]{Steuding2005}, with $(p,q)=(x,y)$,
gives \eqref{eq:Crandall-CF-best-approximation};
the equality case $x/y=p_n/q_n$ cannot occur because $p_n/q_n$ is reduced
and $0<y<q_n$.  If $x\leq0$, then $|x-yt|\geq yt>1$, since
$t=\log_2 3>1$, while the upper bound
in \eqref{eq:Crandall-CF-error-bounds} is $<1$.

For the final assertion, \eqref{eq:Crandall-CF-error-bounds} gives
$|p_n-q_nt|<1/q_{n+1}<1/2$ for $n\geq4$.  If $x\ne p_n$, integrality gives
$|x-p_n|\geq1$, and the triangle inequality gives
$|x-q_nt|\geq1-|p_n-q_nt|>|p_n-q_nt|$.
\end{proof}

\begin{sourceissue}[four repairs in the continued-fraction chain]
\label{sourceissue:Crandall-continued-fraction-chain}
The proof chain on printed pp.~1289--1290 needs four distinct repairs
\cite[pp.~1289--1290]{Crandall1978}.
\begin{enumerate}[label=(\roman*)]
\item Lemma~7.1 omits $0<y$.  As printed, $(x,y)=(0,0)$ contradicts its
strict best-approximation inequality.  The used form has $0<y<q_n$ and is
proved from Lemma~\ref{lem:Crandall-CF-dependency}.
\item Lemma~7.3 introduces $p_n/q_n$ but then writes $p_m,q_m$, leaves $x$
unquantified, and uses $|e^u-1|>|u|$.  This analytic inequality holds for
$u>0$, not for $u<0$.  The exact form needed for cycles is
\begin{equation}
\label{eq:Crandall-repaired-exponential-bound}
 \left.
 \begin{gathered}
  x,y\in\Z,\quad 0<y<q_n,\\
  x-y\log_2 3>0
 \end{gathered}
 \right\}
 \Longrightarrow
 |2^x-3^y|>3^y\log2\,|p_n-q_nt|.
\end{equation}
Indeed, for $u=(x-yt)\log2>0$ the positive-term power series gives
$e^u-1-u=\sum_{j\geq2}u^j/j!>0$, and best approximation supplies
$|x-yt|>|p_n-q_nt|$.  In the cycle application, the missing sign is exactly
\eqref{eq:Crandall-positive-cycle-denominator}.
\item Theorem~7.2 explicitly reduces to $k\leq q_n$, but then invokes
Lemma~7.3, whose repaired denominator hypothesis is strict: $y<q_n$.
The final assertion of Lemma~\ref{lem:Crandall-CF-dependency} handles
$k=q_n$ when $A_k\ne p_n$.  If $A_k=p_n$, the positive exponential
inequality is itself strict.  Thus the same denominator estimate holds at
the equality boundary that the printed proof does not justify.
\item Corollary~7.2 claims that its minimum becomes unbounded as $n$
increases.  For a fixed cycle its second entry instead tends to zero.  The
correct proof first fixes $n$ and only then chooses a sufficiently large
cycle minimum; this quantifier order is proved below.
\end{enumerate}
The later references to Eq.~(7.5) when computing $q_n$ are correct only as
references to the printed continued-fraction list; the recurrence itself is
Eq.~(7.6).
\end{sourceissue}

\begin{theorem}[Crandall's cycle bound, with proof gaps repaired]
\label{thm:Crandall-cycle-bound}
Let $m>1$ be the least member of a $C_{\mathrm{Cr}}$-cycle of least
accelerated period $k$.  For every integer $n\geq4$,
\begin{equation}
\label{eq:Crandall-cycle-bound}
 k>\min\left\{q_n,\frac{2m}{q_n+q_{n+1}}\right\}.
\end{equation}
Consequently, only finitely many cycles have any prescribed accelerated
period \cite[p.~1290]{Crandall1978}.
\end{theorem}

\begin{proof}
Minimality of $m$ and the branch equation give, for the coordinates
\eqref{eq:Crandall-cumulative-time},
\begin{equation}
\label{eq:Crandall-cycle-growth-bound}
 2^{A_j}\leq(3+1/m)^j.
\end{equation}
The return equation and this bound imply
\begin{equation}
\label{eq:Crandall-cycle-minimum-bound}
 m<\frac{k(3+1/m)^{k-1}}{2^{A_k}-3^k}.
\end{equation}
If $k>q_n$, \eqref{eq:Crandall-cycle-bound} is immediate.  Otherwise the
repaired best-approximation argument in
Source issue~\ref{sourceissue:Crandall-continued-fraction-chain}, including
its equality case, gives
\[
 2^{A_k}-3^k>3^k\log2\,|p_n-q_nt|.
\]
Combining this with \eqref{eq:Crandall-cycle-minimum-bound}, and writing
$v=1/(3m)>0$, first yields
\[
 k>m(3+1/m)(\log2)|p_n-q_nt|(1+v)^{-k}.
\]
The function $f(v)=(1+v)^{-k}$ has
$f''(v)=k(k+1)(1+v)^{-k-2}>0$; strict convexity above its tangent at zero
therefore gives $f(v)>f(0)+f'(0)v=1-kv$.  Hence
\begin{equation}
\label{eq:Crandall-cycle-bound-intermediate}
 k>m(\log2)(3+1/m)|p_n-q_nt|\left(1-\frac{k}{3m}\right).
\end{equation}
For $n\geq4$, the recurrence gives
$q_n+q_{n+1}\geq q_4+q_5=12+41=53>20$.  If $k\geq m/10$, the second member
of the minimum in \eqref{eq:Crandall-cycle-bound} is strictly below
$m/10\leq k$.  Otherwise $k<m/10$, so
$1-k/(3m)>29/30$, and the continued-fraction lower bound
$|p_n-q_nt|>(q_n+q_{n+1})^{-1}$ turns
\eqref{eq:Crandall-cycle-bound-intermediate} into
\[
 k>\frac{(3m+1)\log2}{q_n+q_{n+1}}\frac{29}{30}
   >\frac{2m}{q_n+q_{n+1}}.
\]
The final strict inequality is exact: the positive-term expansion
$\log2=2(1/3+1/(3^3\cdot3)+\cdots)$ gives
$\log2>56/81>20/29$, while $(3m+1)/m>3$.
This proves \eqref{eq:Crandall-cycle-bound}.

For finiteness, fix a proposed period $L$ and choose an integer $n\geq4$ with
$q_n>L$; the convergent denominators are unbounded.  An infinite family of
distinct period-$L$ cycles would have unbounded distinct
least members, so choose one with
$m>L(q_n+q_{n+1})/2$.  Both entries of the minimum would then exceed $L$,
contradicting \eqref{eq:Crandall-cycle-bound} with $k=L$.
\end{proof}

\phantomsection\label{historical-10^9-premise}
Printed p.~1282 says that Conjecture~(2.1) had been verified for
$m<10^9$, citing references~[1] and~[5] \cite[p.~1282]{Crandall1978}.
Those references are not two independent witnesses.  Reference~[1] is
Conway's \emph{Unpredictable iterations}; its correct proceedings pages are
49--52, not Crandall's 45--52.  Reference~[5], printed as
``\emph{Zentralblatt f\"ur Mathematik}, Band~233, p.~10041,'' is a defective
rendering of \emph{Zbl}~0337.10041, the Zentralblatt record for that same
Conway paper: $10041$ is part of the accession identifier, not a printed
page.  Thus the two routes collapse to one work, rather than corroborating
one another \cite[p.~1292]{Crandall1978}.

Conway defines the ordinary one-step map $g$ by $g(n)=n/2$ for even $n$ and
$g(n)=3n+1$ for odd $n$, literally the map $U$ in
\eqref{eq:unshortened-U}.  On the first page of his article he reports the
inclusive range $n\leq10^9$ and attributes its verification to D.~H. Lehmer,
Emma Lehmer, and J.~L. Selfridge \cite[p.~219]{Conway1972}.  He supplies no
citation for that sentence and no program, machine, range partition,
stopping criterion, independent check, checksum, log, or preserved output.
His proved undecidability theorem is expressly said to say nothing about the
particular Collatz game; the verification sentence is historical background,
not a theorem of the article.  The published contemporaneous report and its
attribution are therefore content-verified, while the underlying computation
remains unauthenticated and nonreproducible at this checkpoint.

\begin{proposition}[One-step verification to accelerated exclusion]
\label{prop:Conway-Crandall-historical-transfer}
Let $N\in\Npos$.  Suppose that for every $n\in\Npos$ with $n\leq N$ there
exists $r\in\Nzero$ such that $U^r(n)=1$.  Then for every
$m\in\Opos$ with $m\leq N$ there exists $j\in\Nzero$ such that
$C_{\mathrm{Cr}}^j(m)=1$.  Consequently no $m\in\Opos$ with
$1<m\leq N$ can satisfy $C_{\mathrm{Cr}}^k(m)=m$ for any $k\in\Npos$.
\end{proposition}

\begin{proof}
Starting from an odd $m$, the successive odd states of the $U$-orbit are
exactly the $C_{\mathrm{Cr}}$-orbit, with the $j$th odd return occurring at
the unshortened time $A_j+j$ by \eqref{eq:Crandall-three-clocks}.  Since $1$
is odd, any occurrence $U^r(m)=1$ is one of those successive odd returns;
hence $C_{\mathrm{Cr}}^j(m)=1$ for a unique first such index $j$.  Finally
$C_{\mathrm{Cr}}(1)=1$.  If $m>1$ also satisfied
$C_{\mathrm{Cr}}^k(m)=m$, its accelerated orbit would be periodic from $m$
and could not enter the fixed point $1$, a contradiction.
\end{proof}

Printed p.~1290 uses the ``known result'' that no $1<m<10^9$ appears in
its own trajectory, without a new inline citation
\cite[p.~1290]{Crandall1978}.  Proposition
\ref{prop:Conway-Crandall-historical-transfer} proves the exact clock and
domain transfer from Conway's stronger inclusive ordinary-map report to that
accelerated statement; it does not authenticate the computation.

For a hypothetical nontrivial cycle, its least member $m_0$ lies in
$\Opos$.  Conway's inclusive report, after the proved transfer, directly
excludes $m_0\leq10^9$.  Independently, Crandall's weaker strict-range form
first gives $m_0\geq10^9$, and oddness excludes equality with the even
endpoint.  Both routes therefore give, exactly as Crandall writes,
\[
 m_0>10^9,
 \qquad\text{equivalently}\qquad m_0\geq1{,}000{,}000{,}001.
\]
With
\begin{equation}
\label{eq:Crandall-historical-convergents}
 q_{10}=31867,
 \qquad q_{11}=79335,
\end{equation}
Theorem~\ref{thm:Crandall-cycle-bound} gives the least accelerated period
$\ell$ the bound
\[
 \ell>\min\left\{31867,\frac{2\cdot10^9}{111202}\right\}>17985.
\]
Every positive return exponent $k$ is a multiple of $\ell$, so this recovers
Crandall's Theorem~7.3: $C_{\mathrm{Cr}}^k(m)=m$ with $m>1$ implies
$k>17985$, conditional on the published 1972 report reused by Crandall in
1978.  The source variable is an accelerated return time, not a current or
reproducible verification record.
For the same
return, \eqref{eq:Crandall-three-clocks} gives
$A_k\geq k>17985$ shortened steps and $A_k+k>k>17985$ unshortened steps;
these are transferred inequalities, not an identification of the three
clocks.

\phantomsection\label{sourceissue:Crandall-bounded-supremum}
The opening of Section~7 also says every member of a bounded trajectory is
strictly below its supremum.  The correct relation is $\leq$.  The intended
finite pigeonhole argument survives: boundedness plus absence of nontrivial
self-returns would force the trajectory to contain $1$.

\subsubsection{The generalized \texorpdfstring{$qx+r$}{qx+r} programme}

For positive odd $q,r$, $q>1$, Crandall defines
\begin{equation}
\label{eq:Crandall-qxr}
 C_{q,r}(m)=\frac{qm+r}{2^{e_{q,r}(m)}},
 \qquad
 e_{q,r}(m)=\vtwo(qm+r),
 \qquad m\in\Opos.
\end{equation}
Conjecture~8.1 predicts that, except for $(q,r)=(3,1)$, some positive odd
starting value fails to reach $1$.  The exceptional pair is not an assertion
that the ordinary conjecture is true.

\begin{proposition}[Crandall's proved generalized failures]
\label{prop:Crandall-generalized-failures}
The following statements hold.
\begin{enumerate}[label=(\alph*)]
\item If $r>1$, every $m\in\Opos$ divisible by $r$ fails to reach $1$.
\item For $(q,r)=(5,1)$,
\[
 13\longmapsto33\longmapsto83\longmapsto13,
 \qquad(e_0,e_1,e_2)=(1,1,5).
\]
\item For $(q,r)=(181,1)$,
\[
 27\longmapsto611\longmapsto27,
 \qquad(e_0,e_1)=(3,12).
\]
\end{enumerate}
\end{proposition}

\begin{proof}
If $r\mid m$, then $qm+r\equiv0\pmod r$; division by a power of $2$, a unit
modulo odd $r$, preserves divisibility by $r$.  Induction excludes the value
$1$.  The two cycles follow by exact substitution:
\[
 5(13)+1=2(33),\quad5(33)+1=2(83),\quad5(83)+1=2^5(13),
\]
and
\[
 181(27)+1=2^3(611),\qquad181(611)+1=2^{12}(27).
\]
The associated differences are $2^7-5^3=3$ and
$2^{15}-181^2=7$.  The printed starting value in the second cycle is $27$;
$21$ is an OCR error \cite[p.~1291]{Crandall1978}.
\end{proof}

\begin{proposition}[the omitted finite certificate for $q=1093$]
\label{prop:Crandall-q1093}
For $C_{1093,1}$, a positive odd integer has height one exactly when
\begin{equation}
\label{eq:Crandall-q1093-height-one}
 m=\frac{2^{364p}-1}{1093}
 \qquad(p\in\Z_{>0}).
\end{equation}
No positive odd integer has height two.  Every positive odd integer not in
\eqref{eq:Crandall-q1093-height-one} therefore has infinite height, and this
infinite-height set has natural density one in $\Opos$.
\end{proposition}

\begin{proof}
Exact modular exponentiation gives
\begin{equation}
\label{eq:Crandall-q1093-certificate}
 \operatorname{ord}_{1093}(2)=364,
 \qquad
 2^{364}\equiv1\pmod{1093^2}.
\end{equation}
For the first equality it suffices to check the proper divisors of
$364=2^2\cdot7\cdot13$; convenient witnesses include
\[
 2^{182}\equiv1092,
 \qquad2^{52}\equiv27,
 \qquad2^{28}\equiv121\pmod{1093}.
\]
A height-one value satisfies $1093m+1=2^a$, so the order calculation is
equivalent to $364\mid a$ and proves
\eqref{eq:Crandall-q1093-height-one}.  The congruence modulo $1093^2$ shows
that every such $m$ is itself divisible by $1093$.  If it had a predecessor
$u$, then $1093u+1=2^b m$ would be simultaneously congruent to $1$ and $0$
modulo $1093$, impossible.  Thus height two does not occur.  Any finite
trajectory of height greater than one would contain a height-two tail, so no
other finite height occurs.  Finally, the number of parameters $p$ producing
$m\leq x$ is $O(\log x)$; division by the number of odd integers at most $x$
gives density zero for the height-one family.

The source states the classification and calls the height-two impossibility
easy, but omits the modular certificate
\cite[p.~1291]{Crandall1978}.  The calculation is reproduced by
\path{certificates/crandall_1978_checks.py}.
\end{proof}

\begin{nonclaim}
\label{nonclaim:Crandall-qxr-unboundedness}
Infinite height need not mean an unbounded orbit; it may arise from a
nontrivial cycle.  Crandall explicitly states that no unbounded trajectory
was then known for any $qx+1$ problem.  The $7x+1$ trajectory from $3$
exceeding $10^{2000}$ is a finite empirical excursion only.  The heuristic
drift $\log(q/4)>0$ for $q>3$ is not an unboundedness theorem
\cite[pp.~1291--1292]{Crandall1978}.
\end{nonclaim}

\subsection{Everett's finite parity coordinates and density-one descent}
\label{subsec:Everett-1977}

Everett works on \(\Nzero\) with the map
\begin{equation}
\label{eq:Everett-f}
 f_{\mathrm{Ev}}(m)=
 \begin{cases}
  m/2,&m\text{ even},\\
  (3m+1)/2,&m\text{ odd}.
 \end{cases}
\end{equation}
The subscript is editorial: the source writes \(f\).  Thus
\(f_{\mathrm{Ev}}=T|_{\Nzero}\), exactly; it is neither the unshortened map
\(U\) nor Crandall's odd first-return map.  Put
\begin{equation}
\label{eq:Everett-parity-coordinates}
 m_n=f_{\mathrm{Ev}}^n(m),
 \qquad
 x_n=m_n\bmod2\in\{0,1\}.
\end{equation}
The source word is chronological.  When encoded as an integer below, \(x_0\)
is the least significant digit \cite[pp.~42--43]{Everett1977}.

\begin{theorem}[Everett, Theorem 1: finite parity fibres]
\label{thm:Everett-finite-fibres}
Let \(N\ge1\), let
\(w=(x_0,\ldots,x_{N-1})\in\{0,1\}^N\), and put
\[
 X=X(w)=\sum_{n=0}^{N-1}x_n.
\]
There are unique integers \(a(w),b(w)\) satisfying
\[
 0\le a(w)<2^N,
 \qquad
 0\le b(w)<3^X,
\]
such that the complete nonnegative fibre of \(w\) and its \(N\)-th iterate
are
\begin{equation}
\label{eq:Everett-affine-fibre}
 m=a(w)+2^NQ,
 \qquad
 m_N=b(w)+3^XQ,
 \qquad Q\in\Nzero.
\end{equation}
Consequently every length-\(N\) binary word occurs at one and only one
representative \(0\le m<2^N\).
\end{theorem}

\begin{proof}
This is the source induction with its implicit quotient domain made explicit.
For \(N=1\), \(x_0=0\) gives
\[
 m=2Q,\qquad m_1=Q,\qquad (a,b,X)=(0,0,0),
\]
whereas \(x_0=1\) gives
\[
 m=1+2Q,\qquad m_1=2+3Q,\qquad (a,b,X)=(1,2,1).
\]

Assume \eqref{eq:Everett-affine-fibre} at length \(N\), and abbreviate its
coordinates to \(a,b,X,Q_N\).  Because \(3^X\) is odd, the parity of
\(m_N=b+3^XQ_N\) fixes the parity of \(Q_N\) once \(b\) and the next required
bit \(x_N\) are fixed.  There are exactly four cases.

If \(b\) is even and \(x_N=0\), write \(Q_N=2Q_{N+1}\).  Then
\[
 a_N=a,\qquad b_N=b/2,\qquad X_N=X.
\]
If \(b\) is even and \(x_N=1\), write \(Q_N=1+2Q_{N+1}\).  Then
\[
 a_N=a+2^N,\qquad
 b_N=\frac{3b+3^{X+1}+1}{2},\qquad X_N=X+1.
\]
Since \(b\le3^X-1\),
\[
 b_N\le
 \frac{3(3^X-1)+3^{X+1}+1}{2}
 =3^{X+1}-1<3^{X+1}.
\]

If \(b\) is odd and \(x_N=0\), write \(Q_N=1+2Q_{N+1}\).  Then
\[
 a_N=a+2^N,\qquad
 b_N=\frac{b+3^X}{2},\qquad X_N=X,
\]
and \(b<3^X\) gives \(b_N<3^X\).  Finally, if \(b\) is odd and
\(x_N=1\), write \(Q_N=2Q_{N+1}\).  Then
\[
 a_N=a,\qquad
 b_N=\frac{3b+1}{2},\qquad X_N=X+1,
\]
and
\[
 b_N<\frac{3^{X+1}+1}{2}<3^{X+1}.
\]
In each case \(0\le a_N<2^{N+1}\), and the forced parity decomposition of
\(Q_N\) is unique.  This proves existence, uniqueness, both coefficient
bounds, and the retained quotient coordinate at stage \(N+1\)
\cite[pp.~43--44, Eq.~(4) and proof]{Everett1977}.
\end{proof}

\begin{proposition}[the exact finite parity morphism and inverse residue]
\label{prop:Everett-finite-parity-morphism}
For \(N\ge1\), define
\begin{equation}
\label{eq:Everett-PN}
 P_N:\Z/2^N\Z\longrightarrow\{0,1\}^N,
 \qquad
 P_N([m])=(T^n(m)\bmod2)_{0\le n<N}.
\end{equation}
Then \(P_N\) is a bijection.  If
\[
 c(w)=\sum_{j=0}^{N-1}x_j2^j
       3^{\sum_{i=j+1}^{N-1}x_i},
\]
its inverse is the unique residue
\begin{equation}
\label{eq:Everett-inverse-residue}
 P_N^{-1}(w)\equiv-3^{-X(w)}c(w)\pmod{2^N}.
\end{equation}
Reduction \(\Z/2^{N+1}\Z\to\Z/2^N\Z\) corresponds to deletion of the last
word digit.  On \(\Nzero\), the fibre over \(w\) is exactly
\(a(w)+2^N\Nzero\), and \eqref{eq:Everett-affine-fibre} records the quotient
that the word forgets.
\end{proposition}

\begin{proof}
Theorem~\ref{thm:Everett-finite-fibres} proves well-definedness on residue
classes, singleton fibres at residue level, surjectivity, and the stated
nonnegative fibres.  Iterating the two branch equations gives
\begin{equation}
\label{eq:Everett-word-affine-composition}
 2^NT^N(m)=3^{X(w)}m+c(w)
\end{equation}
whenever the first \(N\) parities are \(w\).  This follows by induction: the
next zero branch multiplies the existing numerator by no new factor of \(3\),
while the next one branch multiplies it by \(3\) and adds \(2^N\).
Divisibility of the right side by \(2^N\) gives
\eqref{eq:Everett-inverse-residue}; \(3^{X(w)}\) is odd and therefore a unit
modulo \(2^N\).  The same prefix read at levels \(N+1\) and \(N\) proves the
reduction identity.
\end{proof}

The source corollary uses the \(2^N\) positive representatives
\(1,\ldots,2^N\), replacing the zero residue by \(2^N\).  This implicit
representative change is valid, but it must be stated.  Numerically encoding
the target word identifies \(P_N\) exactly with the finite chronological
parity map \(\overline Q_{3,N}\) in
\eqref{eq:LDW-finite-parity-map}.  Everett proves the bijection and its
integer affine fibres.  Later sources add permutation-order, lifting,
inverse-limit, and graph-relation structure.  Everett's quotient \(Q_N\) in
the proof of Theorem~1 is not the later parity encoder \(Q_3\).

If two nonnegative integers have the same infinite parity sequence, the
finite theorem makes their difference divisible by \(2^N\) for every \(N\),
so they are equal.  This is the source's infinite injectivity conclusion.  It
does not assert that every infinite binary sequence is realized by an
ordinary nonnegative integer.

\begin{theorem}[Everett, Theorem 2: natural-density-one descent]
\label{thm:Everett-density-one-descent}
Let
\begin{equation}
\label{eq:Everett-AM}
 A(M)=\#\{1\le m\le M:\exists k\ge1,\ f_{\mathrm{Ev}}^k(m)<m\}.
\end{equation}
Then
\begin{equation}
\label{eq:Everett-density-limit}
 \lim_{M\to\infty}\frac{A(M)}M=1.
\end{equation}
Thus the source phrase ``almost every'' means ordinary natural density one
in the intervals \([1,M]\).
\end{theorem}

\begin{proof}
First take \(M=2^N\).  For a length-\(N\) word put
\begin{equation}
\label{eq:Everett-L-epsilon}
 X=\sum_{n=0}^{N-1}x_n,
 \qquad
 L=\frac{\log2}{\log(10/3)},
 \qquad
 \varepsilon=L-\frac12.
\end{equation}
Here \(\varepsilon>0\), since \(L>1/2\) is equivalent to \(4>10/3\).
Let \(H_N\) be the set of words satisfying
\begin{equation}
\label{eq:Everett-HN}
 \frac12-\varepsilon<\frac XN<\frac12+\varepsilon=L.
\end{equation}
Under uniform counting on \(\{0,1\}^N\), each coordinate is Bernoulli with
mean \(1/2\), and every prescribed pair of coordinate values has
\(2^{N-2}\) extensions.  Hence
\[
 \mathbb E X=N/2,
 \qquad
 \operatorname{Var}(X)=N/4.
\]
Chebyshev's inequality gives the exact estimate cited, but not proved, in the
source:
\begin{equation}
\label{eq:Everett-Chebyshev}
 \frac{\#H_N}{2^N}
 \ge1-\frac1{4\varepsilon^2N}.
\end{equation}
This is a finite uniform count transported by
Proposition~\ref{prop:Everett-finite-parity-morphism}; it does not assume that
one ordinary orbit has independent random steps.

Let \(D_N\) contain the words with \(X/N<L\).  Then \(H_N\subseteq D_N\).
If \(x_n=0\), then \(m_{n+1}/m_n=1/2\).  If \(x_n=1\) and \(m_n>1\), write
\(m_n=2Q+1\), where \(Q\ge1\); then
\begin{equation}
\label{eq:Everett-odd-ratio}
 \frac{m_{n+1}}{m_n}
 =\frac{3Q+2}{2Q+1}\le\frac53.
\end{equation}
Moreover,
\begin{equation}
\label{eq:Everett-word-contraction}
 \frac XN<L
 \quad\Longleftrightarrow\quad
 (1/2)^{N-X}(5/3)^X<1.
\end{equation}

Use the source corollary's representatives \(1\le m\le2^N\).  For a word in
\(D_N\), either some \(m_n=1<m\) with \(n<N\), or no such value occurs.  In
the latter case every odd factor satisfies \eqref{eq:Everett-odd-ratio}, so
\[
 m_N=m_0\prod_{n=0}^{N-1}\frac{m_{n+1}}{m_n}
 \le(1/2)^{N-X}(5/3)^Xm_0<m_0.
\]
The one excluded starting value is \(m=1\).  Therefore
\begin{equation}
\label{eq:Everett-dyadic-count}
 A(2^N)\ge\#D_N-1\ge\#H_N-1,
\end{equation}
and \eqref{eq:Everett-Chebyshev} gives
\begin{equation}
\label{eq:Everett-dyadic-density}
 \frac{A(2^N)}{2^N}
 \ge1-\frac1{4\varepsilon^2N}-2^{-N}\longrightarrow1.
\end{equation}

It remains to retain the source's interpolation, rather than claim that a
dyadic subsequence proves the full limit.  Put \(A_N=A(2^N)\) and
\begin{equation}
\label{eq:Everett-nN}
 n_N
 =2^N+\bigl((2^{N+1}-2^N)-(A_{N+1}-A_N)\bigr)
 =A_N+(2^{N+1}-A_{N+1}),
 \qquad 2^N\le n_N\le2^{N+1}.
\end{equation}
The difference \(n_N-2^N\) is exactly the number of failures in the dyadic
shell \((2^N,2^{N+1}]\).  Its bounds follow because that shell contains
\(2^N\) integers.
For \(2^N<M\le n_N\), monotonicity gives
\[
 \frac{A(M)}M\ge\frac{A_N}{n_N}.
\]
For \(M=n_N+k\), the dyadic interval has exactly
\(n_N-2^N\) failures in total, so its first
\(n_N-2^N+k\) members contain at least \(k\) successes.  Consequently
\[
 \frac{A(M)}M
 \ge\frac{A_N+k}{n_N+k}
 \ge\frac{A_N}{n_N},
\]
where the last inequality is equivalent to \(k(n_N-A_N)\ge0\).  Finally,
\begin{equation}
\label{eq:Everett-interpolation-limit}
 \frac{A_N}{n_N}
 =\frac{A_N/2^N}
 {A_N/2^N+2(1-A_{N+1}/2^{N+1})}\longrightarrow1
\end{equation}
by \eqref{eq:Everett-dyadic-density}.  The two interval estimates and
\eqref{eq:Everett-interpolation-limit} prove
\eqref{eq:Everett-density-limit}
\cite[pp.~44--45, Eqs.~(5)--(12) and Thm.~2]{Everett1977}.
\end{proof}

\begin{calculation}[an exact rational concentration witness]
\label{calc:Everett-rational-bound}
The strict inequality \(L>23/40\) is equivalent to
\[
 2^{40}3^{23}>10^{23},
\]
and direct integer subtraction gives
\[
 2^{40}3^{23}-10^{23}=3511519796081828298752>0.
\]
Thus \(\varepsilon>3/40\).  Applying the same Chebyshev proof with the
smaller rational window \(3/40\) gives the completely explicit edition
bound
\begin{equation}
\label{eq:Everett-rational-density-bound}
 \frac{A(2^N)}{2^N}
 \ge1-\frac{400}{9N}-2^{-N}.
\end{equation}
This is a proof-checking witness for the source limit, not an attribution of
the displayed constant to Everett.
\end{calculation}

\begin{proposition}[Everett--Crandall first-descent crosswalk]
\label{prop:Everett-Crandall-crosswalk}
Define
\begin{align*}
 D_T&=\{m\in\Opos:\exists k\ge1,\ T^k(m)<m\},\\
 D_C&=\{m\in\Opos:\exists j\ge1,\ C_{\mathrm{Cr}}^j(m)<m\}.
\end{align*}
Then \(D_T=D_C\) pointwise.  Moreover,
\begin{equation}
\label{eq:Everett-Crandall-relative-density}
 \lim_{M\to\infty}
 \frac{\#(D_C\cap[1,M])}{\#(\Opos\cap[1,M])}=1.
\end{equation}
This is the exact primary verification of the theorem reported by Crandall.
\end{proposition}

\begin{proof}
If \(C_{\mathrm{Cr}}^j(m)<m\), then
\(C_{\mathrm{Cr}}^j(m)=T^{A_j}(m)\) by
\eqref{eq:Crandall-three-clocks}, so \(m\in D_T\).  Conversely, suppose
\(T^k(m)<m\).  If \(T^k(m)\) is odd, it is one of the successive odd-return
states.  If it is even, put \(r=\vtwo(T^k(m))\ge1\).  The next \(r\)
shortened steps are divisions by two, and
\[
 T^{k+r}(m)=\frac{T^k(m)}{2^r}<m
\]
is odd.  It is therefore a later \(C_{\mathrm{Cr}}\)-iterate.  This proves
the pointwise equality without identifying shortened time \(k\) with
accelerated time \(j\).

Every positive even \(m\) belongs to the full shortened descent set because
\(T(m)=m/2<m\).  Put
\[
 B(M)=\#(D_C\cap[1,M]),
 \qquad
 O(M)=\#(\Opos\cap[1,M])=\lceil M/2\rceil.
\]
The exact finite-cutoff identity is therefore
\begin{equation}
\label{eq:Everett-Crandall-finite-count}
 A(M)=\lfloor M/2\rfloor+B(M).
\end{equation}
Consequently
\[
 1-\frac{B(M)}{O(M)}
 =\frac{M}{O(M)}\left(1-\frac{A(M)}M\right).
\]
The multiplier tends to \(2\).  Hence Everett's ordinary-density statement
and Crandall's relative-density statement are equivalent, and
\eqref{eq:Everett-Crandall-relative-density} follows.
\end{proof}

\begin{nonclaim}
Everett proves first descent below the starting value, not convergence to
\(1\).  The exceptional set may be infinite of density zero.  The proof gives
no uniform descent-time or total-stopping-time bound and no accelerated-time
bound.  Its finite parity surjectivity is not ordinary-integer surjectivity
onto infinite binary sequences.  It proves no statement about \(U\),
generalized \(qx+r\) maps, nontrivial cycles, or unbounded trajectories.  The
introductory verification ``up to the millions'' is historical context
without a supplied computational certificate.
\end{nonclaim}

\begin{sourceissue}[manifestation and presentation boundaries]
\label{sourceissue:Everett-manifestation}
The four raster pages are clear, but their OCR layer corrupts formulas and is
not used as mathematical evidence.  The publisher identity is fixed by DOI
and PII; the inspected content witness is a matching VOR facsimile, while a
publisher-hosted PDF hash remains unavailable.  In the article itself, the
abstract's phrase ``every integer'' must be read as every positive integer:
the fixed point \(0\) has no iterate equal to \(1\).  The body correctly uses
\(m\ge1\) for that conjecture.  The positive-representative swap and
\(Q_N\in\Nzero\) are implicit, and the probability estimate is cited rather
than proved.  The preceding reconstruction makes each step explicit.  No
theorem-level mathematical error was found.
\end{sourceissue}

\subsection{Steuding's orbit continued fraction: an exact endpoint recoding}
\label{subsec:Steuding-522}

Steuding's starred Exercise~5.22 uses the unshortened positive-integer
recursion \eqref{eq:unshortened-U}, asks for an equivalence-class
reformulation of the Collatz conjecture, and asks for a Markoff constant
\cite[physical PDF p.~100; printed p.~87, Ex.~5.22]{Steuding2005}.
It is an exercise prompt, not a theorem accompanied by a proof.  The
propositions and calculation below supply an edition solution without a
novelty claim.

The target printed in the exercise is the purely periodic irrational
\[
 \beta=[\overline{4,2,1}],
\]
with a visible overbar.  Ordinary PDF text extraction loses that bar.  The
unbarred finite fraction is $[4,2,1]=13/3$ and cannot be substituted:
Steuding's rational and irrational equivalence classes are disjoint
\cite[printed p.~79]{Steuding2005}.

\begin{proposition}[whole orbit to continued-fraction morphism]
\label{prop:Steuding-orbit-CF-morphism}
For $m\in\Npos$, define
\begin{equation}
\label{eq:Steuding-Theta}
 x_j(m)=U^j(m),
 \qquad
 \Theta(m)=[x_0(m),x_1(m),x_2(m),\ldots].
\end{equation}
Every $\Theta(m)$ is an irrational real number.  If
\[
 \mathcal G(y)=\frac{1}{y-\lfloor y\rfloor}
\]
is the Gauss tail map on nonintegral reals, then
\begin{equation}
\label{eq:Steuding-Gauss-conjugacy}
 \lfloor\Theta(m)\rfloor=m,
 \qquad
 \mathcal G(\Theta(m))=\Theta(U(m)).
\end{equation}
Consequently $\Theta$ is a bijection from $\Npos$ onto its image, its inverse
there is the floor map, and it conjugates $U$ to $\mathcal G$ restricted to
that image.
\end{proposition}

\begin{proof}
Every $x_j(m)$ is a positive integer.  The finite convergents of
\eqref{eq:Steuding-Theta} have denominator continuants tending to infinity,
and consecutive convergents differ by the reciprocal of the product of
their denominators.  Their alternating nested intervals therefore have a
unique common limit.  That limit cannot be rational: if it were the reduced
fraction $p/q$, choose $n$ so large that $q_n>q$ and $q_{n+1}>q$.  Since each
convergent $p_n/q_n$ is reduced, it is then distinct from $p/q$.  The source
bound
\[
 \left|\frac{p}{q}-\frac{p_n}{q_n}\right|<
 \frac{1}{q_nq_{n+1}}
\]
would make the left-hand side strictly smaller than $1/(q q_n)$.  But two
distinct reduced rationals with
denominators $q$ and $q_n$ are separated by at least $1/(q q_n)$.  This is
the constructive
content of Steuding's infinite-fraction discussion
\cite[Sec.~3.4, printed pp.~42--43]{Steuding2005}.

The tail $[x_1(m),x_2(m),\ldots]$ is greater than one, so
\[
 \Theta(m)=m+\frac{1}{[x_1(m),x_2(m),\ldots]}
\]
has floor $m$.  Applying $\mathcal G$ deletes the first digit, while the
recursion gives $x_{j+1}(m)=x_j(U(m))$.  This proves both identities in
\eqref{eq:Steuding-Gauss-conjugacy}.  The floor identity proves injectivity
and supplies the displayed inverse; surjectivity is only onto the stated
image.
\end{proof}

Steuding defines $\alpha\sim\gamma$ by
\begin{equation}
\label{eq:Steuding-unimodular-equivalence}
 \alpha=\frac{a\gamma+b}{c\gamma+d},
 \qquad
 a,b,c,d\in\Z,
 \qquad
 ad-bc=\pm1.
\end{equation}
Serret's Theorem~5.7 says that two irrational simple continued fractions are
equivalent exactly when finite prefixes can be removed from both to leave the
same infinite tail \cite[printed pp.~79--80]{Steuding2005}.

\begin{sourceissue}[the matrix step in the printed Serret proof]
\label{sourceissue:Steuding-Serret-proof}
Printed p.~78 calls all integer matrices of determinant $\pm1$ the special
linear group $\mathrm{SL}_2(\Z)$; the group is $\mathrm{GL}_2(\Z)$, since
$\mathrm{SL}_2(\Z)$ has determinant $+1$.  Lemma~5.6 also suppresses the
choice of simultaneous matrix sign and leaves its constructive sufficiency
direction to an induction.

The converse proof on printed p.~80 then (i) invokes $c>d>0$ for the original
matrix before that reduction has been proved, (ii) duplicates the numerator
entry $R$ by printing
$S=ap_{m-1}+bq_{m-1}$ instead of
$S=cp_{m-1}+dq_{m-1}$, and (iii) writes $\alpha$ where the convergents and
error terms belong to $\gamma$ (called $\beta$ in the source).  The factors
$c\alpha+d$ on the same page must likewise be $c\gamma+d$.

Here is the exact repair.  Start with an arbitrary matrix $M$ in
\eqref{eq:Steuding-unimodular-equivalence} and replace it by $-M$, if
necessary, so that $D=c\gamma+d>0$.  For convergents $p_m/q_m$ to $\gamma$,
put $\varepsilon_m=p_m-\gamma q_m$.  The lower row of the composite matrix
is
\begin{align*}
 Q_m&=cp_m+dq_m=Dq_m+c\varepsilon_m,\\
 S_m&=cp_{m-1}+dq_{m-1}=Dq_{m-1}+c\varepsilon_{m-1}.
\end{align*}
Since $\varepsilon_m\to0$ and
\[
 q_m-q_{m-1}=(a_m-1)q_{m-1}+q_{m-2}\longrightarrow\infty,
\]
one has $Q_m>S_m>0$ for every sufficiently large $m$.  The Euclidean
algorithm on the coprime pair $(Q_m,S_m)$ constructs its positive
continued-fraction denominator word.  The two finite presentations obtained
by splitting the last partial quotient choose the required determinant sign.
After that choice, any two numerator rows with the same lower row and the
same determinant differ by an integral multiple of $(Q_m,S_m)$; changing
the first partial quotient absorbs precisely that multiple.  Thus the
composite matrix is a finite continued-fraction prefix, and the two numbers
have a common tail.  This proves the needed Serret direction without the
printed circular step.  The complete locator audit is in
\path{qa/STEUDING-2005-F011-exercise-522-audit.md}.
\end{sourceissue}

\begin{theorem}[edition solution of Steuding's Exercise 5.22]
\label{thm:Steuding-522-equivalence}
Let $b_{3r}=4$, $b_{3r+1}=2$, and $b_{3r+2}=1$.  The following statements
are equivalent:
\begin{align*}
 &\text{for every $m\in\Npos$, the $U$-orbit of $m$ eventually enters the
 cycle $4,2,1$};\\
 &\text{for every $m\in\Npos$, }\Theta(m)\sim[\overline{4,2,1}].
\end{align*}
More precisely, for each fixed $m$,
\begin{equation}
\label{eq:Steuding-pointwise-equivalence}
 \Theta(m)\sim[\overline{4,2,1}]
 \quad\Longleftrightarrow\quad
 \exists r,s\in\Nzero\ \forall j\in\Nzero:
 U^{r+j}(m)=b_{s+j}.
\end{equation}
\end{theorem}

\begin{proof}
If the orbit eventually enters the cycle, its digit sequence in
\eqref{eq:Steuding-Theta} is eventually a cyclic shift of
$(4,2,1,4,2,1,\ldots)$.  Serret's theorem therefore gives the forward
implication of \eqref{eq:Steuding-pointwise-equivalence}.  Conversely,
Serret's theorem turns the equivalence of the two irrational continued
fractions into the displayed equality of tails for some $r,s$.  Hence
$U^r(m)\in\{4,2,1\}$.  The literal identities
\[
 U(4)=2,
 \qquad U(2)=1,
 \qquad U(1)=4
\]
then force every subsequent value.  Universal quantification over $m$ gives
the two global statements.  The phase variable $s$ is retained; no preferred
entry point into the cycle is assumed.
\end{proof}

\begin{calculation}[the requested Markoff constant]
\label{calc:Steuding-522-Markoff}
Let $d=\sqrt{229}$.  The exact period matrix is
\[
 \begin{pmatrix}4&1\\1&0\end{pmatrix}
 \begin{pmatrix}2&1\\1&0\end{pmatrix}
 \begin{pmatrix}1&1\\1&0\end{pmatrix}
 =\begin{pmatrix}13&9\\3&2\end{pmatrix}.
\]
Therefore
\begin{equation}
\label{eq:Steuding-beta-quadratic}
 3\beta^2-11\beta-9=0,
 \qquad
 \beta=\frac{11+d}{6}.
\end{equation}
Index its partial quotients by $a_{3r}=4$, $a_{3r+1}=2$,
$a_{3r+2}=1$.  Formula~(5.11) of the source is
\begin{equation}
\label{eq:Steuding-Markoff-Perron}
 \lambda(\beta)=\liminf_{n\to\infty}
 \left(
  [a_{n+1},a_{n+2},\ldots]
  +[0,a_n,a_{n-1},\ldots,a_1]
 \right)^{-1}.
\end{equation}
The three phases, with no cyclic normalization, are
\[
\begin{array}{c|c|c|c|c}
n\bmod3 & [a_{n+1},a_{n+2},\ldots]
 & \displaystyle\lim[0,a_n,\ldots,a_1]
 & \text{sum} & \text{reciprocal}\\ \hline
0 & (13+d)/10 & (d-13)/10 & d/5 & 5/d\\
1 & (7+d)/18  & (d-7)/18  & d/9 & 9/d\\
2 & (11+d)/6  & (d-11)/6  & d/3 & 3/d
\end{array}
\]
The reversed finite fractions converge on each residue subsequence to the
displayed reversed periodic tail; equivalently, Steuding's Galois formula
identifies each with the negative conjugate of the corresponding cyclic
complete quotient \cite[Thm.~5.3 and Eq.~(5.11)]{Steuding2005}.  Taking the
liminf in \eqref{eq:Steuding-Markoff-Perron} gives
\begin{equation}
\label{eq:Steuding-Markoff-value}
 \boxed{\lambda([\overline{4,2,1}])
 =\min\left\{\frac5{\sqrt{229}},\frac9{\sqrt{229}},
                    \frac3{\sqrt{229}}\right\}
 =\frac3{\sqrt{229}}.}
\end{equation}
The exact quadratic-field identities are reproduced by
\path{certificates/steuding_522_checks.py}.
\end{calculation}

\begin{proposition}[exact Crandall clock crosswalk]
\label{prop:Steuding-Crandall-crosswalk}
For an odd start $y_0$, put $y_i=C_{\mathrm{Cr}}^i(y_0)$,
$e_i=\vtwo(3y_i+1)$, and $A_j=\sum_{i<j}e_i$.  Then the Steuding digit
sequence and the accelerated orbit satisfy
\begin{equation}
\label{eq:Steuding-Crandall-clock}
 y_j=U^{A_j+j}(y_0),
\end{equation}
and, for $1\leq r\leq e_j+1$,
\begin{equation}
\label{eq:Steuding-inserted-even-states}
 U^{A_j+j+r}(y_0)=\frac{3y_j+1}{2^{r-1}}.
\end{equation}
For an arbitrary $m=2^s y_0$ with $y_0$ odd, the first $s$ Steuding digits
are the exact halving prefix before this odd-return expansion begins.
\end{proposition}

\begin{proof}
Equation \eqref{eq:Steuding-Crandall-clock} is
\eqref{eq:Crandall-three-clocks}.  Starting from the odd state $y_j$, the
first $U$ step forms $3y_j+1$ and the next $e_j$ steps divide successively by
$2$, proving \eqref{eq:Steuding-inserted-even-states}; at $r=e_j+1$ the
value is $y_{j+1}$.  The even-start statement is the literal iteration
$U^r(2^s y_0)=2^{s-r}y_0$ for $0\leq r\leq s$.
\end{proof}

\begin{nonclaim}
\label{nonclaim:Steuding-Markoff-boundary}
Theorem~\ref{thm:Steuding-522-equivalence} is a lossless endpoint recoding,
not a proof of either side: the full orbit is stored as the full digit
sequence and the initial integer is recovered by the floor map.  Steuding's
Theorem~5.8 gives only
\[
 \Theta(m)\sim\beta
 \ \Longrightarrow\ 
 \lambda(\Theta(m))=\frac3{\sqrt{229}}.
\]
Equality of Markoff constants is not a converse and cannot certify Collatz
termination.  The orbit-digit continued fraction $\Theta(m)$ is also not
Crandall's Section~7 continued fraction for the fixed number $\log_2 3$;
no map between those two constructions is evidenced here.  Finally,
\eqref{eq:Steuding-Crandall-clock} is a variable-time expansion, not a
one-step identification of $U$ and $C_{\mathrm{Cr}}$.
\end{nonclaim}


\subsection{The \texorpdfstring{$2$}{2}-adic parity conjugacy and notation
orientation}

Let $S:\Ztwo\to\Ztwo$ be the digit shift
\[
 S(z)=\frac{z-(z\bmod2)}2.
\]

\begin{proposition}[Bernstein--Lagarias, with maps oriented]
\label{prop:BL-Q}
Define
\begin{equation}
\label{eq:BL-Q}
 Q_3(x)=\sum_{i=0}^{\infty}(T^i(x)\bmod2)2^i
 \qquad(x\in\Ztwo).
\end{equation}
If $z\ne0$ has one positions
$0\leq d_1<d_2<\cdots$ (a finite or infinite list), put
\begin{equation}
\label{eq:BL-Phi}
 \Phi_3(z)=-\sum_{r}3^{-r}2^{d_r},
 \qquad \Phi_3(0)=0.
\end{equation}
Then $Q_3$ and $\Phi_3$ are mutually inverse $2$-adic isometries and
\begin{equation}
\label{eq:BL-conjugacy}
 Q_3\circ T=S\circ Q_3,
 \qquad
 T=\Phi_3\circ S\circ\Phi_3^{-1}.
\end{equation}
\end{proposition}

\begin{proof}
If $x\equiv y\pmod{2^N}$, induction through their common parity branches
shows that the first $N$ digits of $Q_3(x)$ and $Q_3(y)$ agree.  Conversely,
if those digits agree and contain $r$ ones, the common affine word gives
\[
 T^N(x)-T^N(y)=\frac{3^r}{2^N}(x-y).
\]
The left side is in $\Ztwo$ and $3^r$ is a $2$-adic unit, so
$x\equiv y\pmod{2^N}$.  Hence $Q_3$ preserves every $2$-adic congruence
distance and is injective.

The series \eqref{eq:BL-Phi} converges because its term valuations tend to
infinity.  Splitting off the first digit gives
\[
 \Phi_3(2w)=2\Phi_3(w),
 \qquad
 \Phi_3(1+2w)=\frac{2\Phi_3(w)-1}{3}.
\]
The first value is even and the second odd, and applying the corresponding
branch of $T$ gives $T(\Phi_3(z))=\Phi_3(Sz)$.  Thus the chronological
parities of $\Phi_3(z)$ are the digits of $z$, so
$Q_3\Phi_3=\operatorname{id}_{\Ztwo}$.  The injective $Q_3$ is therefore
bijective with inverse $\Phi_3$.  Removing the first parity digit proves the
first identity in \eqref{eq:BL-conjugacy}; conjugating by the inverse proves
the second.  This reproduces Bernstein--Lagarias Eqs.~(1.3), (1.5), and
(1.6) with the inverse orientation explicit
\cite[author-copy pp.~1--2]{BernsteinLagarias1996}.
\end{proof}

\begin{corollary}[periodic parity word to fixed-return fraction]
\label{cor:periodic-word-resummation}
With the word data of
Proposition~\ref{prop:formal-chronological-composition}, put
\[
 q_{\boldsymbol\varepsilon}=\sum_{i=0}^{L-1}\varepsilon_i2^i,
 \qquad
 z_{\boldsymbol\varepsilon}
 =\frac{q_{\boldsymbol\varepsilon}}{1-2^L}\in\Ztwo.
\]
Then
\begin{equation}
\label{eq:periodic-word-resummation}
 \Phi_3(z_{\boldsymbol\varepsilon})
 =R(\boldsymbol\varepsilon)
 =\frac{\displaystyle\sum_{r=0}^{m-1}3^{m-r-1}2^{v_r}}
        {2^L-3^m}.
\end{equation}
Thus the B\"ohm--Sontacchi return candidate is the inverse parity image of
the purely periodic chronological word.  Its exact $T$-period is the least
shift period of that infinite word.  In particular, a length-$L$ return has
exact period $L$ precisely when the repeated block has least shift period
$L$.
\end{corollary}

\begin{proof}
The one positions of $z_{\boldsymbol\varepsilon}$ are
$v_r+sL$ for $0\leq r<m$ and $s\geq0$.  Formula \eqref{eq:BL-Phi} gives
\begin{align*}
 \Phi_3(z_{\boldsymbol\varepsilon})
 &=-\sum_{s\geq0}\sum_{r=0}^{m-1}
       3^{-(sm+r+1)}2^{v_r+sL}\\
 &=-\left(\sum_{r=0}^{m-1}3^{-(r+1)}2^{v_r}\right)
       \sum_{s\geq0}\left(\frac{2^L}{3^m}\right)^s\\
 &=\frac{\displaystyle\sum_{r=0}^{m-1}3^{m-r-1}2^{v_r}}
        {2^L-3^m},
\end{align*}
where the geometric series converges in $\Ztwo$ because
$|2^L/3^m|_2=2^{-L}<1$.  The conjugacy
$Q_3T=SQ_3$ and bijectivity of $Q_3$ identify the least positive return time
of the point with the least positive shift period of its parity sequence.
\end{proof}

Laarhoven--de Weger use the symbol $\Phi$ for the chronological series
\eqref{eq:BL-Q}, first modulo $2^k$ and then on $\Ztwo$.  Their identities are
$T_k=\Phi_k^{-1}\sigma_{2,k}\Phi_k$ and
$T=\Phi^{-1}S\Phi$ \cite[Secs.~2--3]{LaarhovenDeWeger2012}.
Consequently their $\Phi$ is $Q_3$ in the notation of
Proposition~\ref{prop:BL-Q}, not Bernstein--Lagarias's inverse $\Phi_3$.
This is an exact notation crosswalk, not a mathematical disagreement.

\subsection{Finite De Bruijn relations and the projective parity map}
\label{subsec:finite-debruijn-relations}

Laarhoven--de Weger's finite objects must first be typed as relations.  For
$k\geq1$ put
\[
 V_k=\Z/2^k\Z,
 \qquad
 \mathcal P(V_k)=\{A:A\subseteq V_k\}.
\]
Their modular Collatz graph has an edge $[a]_k\to[b]_k$ precisely when
integer lifts $a_1,b_1$ exist with $T(a_1)=b_1$.  Its forward relation is
therefore
\begin{equation}
\label{eq:LDW-finite-Collatz-relation}
 \mathcal T_k([a]_k)
 =\bigl\{[T(a)]_k,[T(a)+2^{k-1}]_k\bigr\}
 \quad:V_k\longrightarrow\mathcal P(V_k).
\end{equation}
Here $a$ may be any integer representative.  All representatives have the
same parity, and changing $a$ by $2^k$ changes its branch value by an odd
multiple of $2^{k-1}$.  Thus the two displayed residues are exactly all
possible outputs and are distinct.

Write a binary word least-significant digit first and identify
$b_0\cdots b_{k-1}$ with $\sum_{i<k}b_i2^i\in V_k$.  The binary De Bruijn
forward relation is
\begin{equation}
\label{eq:LDW-finite-shift-relation}
 \sigma_{2,k}(b_0\cdots b_{k-1})
 =\{b_1\cdots b_{k-1}c:c\in\{0,1\}\}.
\end{equation}
Numerically, this is the right shift together with that value plus
$2^{k-1}$.  Neither \eqref{eq:LDW-finite-Collatz-relation} nor
\eqref{eq:LDW-finite-shift-relation} is a deterministic self-map.

The source calls its finite encoder $\Phi_k$.  To prevent its dimension
subscript from colliding with Bernstein--Lagarias's multiplier subscript in
$\Phi_3=Q_3^{-1}$, this edition writes the same map as
\begin{equation}
\label{eq:LDW-finite-parity-map}
 \overline Q_{3,k}:V_k\longrightarrow V_k,
 \qquad
 \overline Q_{3,k}([a]_k)
 =\sum_{i=0}^{k-1}(T^i(a)\bmod2)2^i.
\end{equation}
Thus $\overline Q_{3,k}$ is source $\Phi_k$ without a normalization or a
change of orientation.  The symbol $3$ in $Q_3$ and $\Phi_3$ denotes the
$3x+1$ multiplier; it is not the graph dimension.

Everett's earlier Theorem~1 is exactly the bijectivity of this same finite
chronological map on the representatives $0\leq a<2^k$, together with the
stronger integer-fibre coordinates in
Theorem~\ref{thm:Everett-finite-fibres}.  Laarhoven--de Weger's later theorem
adds the two-successor graph relation; that relation is not retroactively
attributed to Everett.

For a point map $f:V_k\to V_k$ and $A\subseteq V_k$, write
$f[A]=\{f(a):a\in A\}$.  This elementwise extension is required in the next
statement.

\begin{theorem}[Laarhoven--de Weger finite graph theorem, reconstructed]
\label{thm:LDW-finite-debruijn}
For every $k\geq1$, $\overline Q_{3,k}$ is a bijection and
\begin{equation}
\label{eq:LDW-finite-relation-conjugacy}
 \overline Q_{3,k}[\mathcal T_k(a)]
 =\sigma_{2,k}(\overline Q_{3,k}(a))
 \qquad(a\in V_k).
\end{equation}
Consequently it is a directed-graph isomorphism from the binary modular
Collatz graph to $B(2,k)$.  Its inverse is the reduction modulo $2^k$ of
Bernstein--Lagarias's $\Phi_3=Q_3^{-1}$.
\end{theorem}

\begin{proof}
Proposition~\ref{prop:BL-Q} gives the exact congruence equivalence
\[
 x\equiv y\pmod{2^k}
 \quad\Longleftrightarrow\quad
 Q_3(x)\equiv Q_3(y)\pmod{2^k}.
\]
It proves that \eqref{eq:LDW-finite-parity-map} is well-defined and injective;
finiteness makes it bijective.  It also identifies it pointwise as the
reduction of $Q_3$, so reduction of $\Phi_3$ supplies the inverse.

Let $a_0$ and $a_0+2^k$ represent the two lift classes modulo $2^{k+1}$.
Their first $k$ parity digits agree, while their next parity digits differ,
again by the exact isometry.  Applying $T$ deletes the common first digit.
The two resulting length-$k$ itinerary words therefore have common first
$k-1$ digits and both possible terminal digits.  They are precisely
$\sigma_{2,k}(\overline Q_{3,k}([a_0]_k))$.  On the state side their residues
are exactly the two values in
\eqref{eq:LDW-finite-Collatz-relation}.  This proves
\eqref{eq:LDW-finite-relation-conjugacy} as equality of two-element subsets.
\end{proof}

The active source theorem has no proof environment.  Its preceding prose
depends on an earlier Lagarias congruence, and source line 131 names the next
digit as $x_{k+1}$ although the definitions force $x_k$.  A proof at source
lines 160--166 is commented out; its first sentence reverses the image edge,
although its following equivalence has the orientation used above.  These are
source defects, not alternative graph conventions
\cite[Sec.~2]{LaarhovenDeWeger2012}.

\begin{proposition}[exact path multiplicities]
\label{prop:LDW-path-counts}
Let $A_k$ be the adjacency matrix of either graph in
Theorem~\ref{thm:LDW-finite-debruijn}, and let $J$ be the $2^k$-by-$2^k$
all-ones matrix.  Then
\begin{equation}
\label{eq:LDW-path-counts}
 A_k^k=J,
 \qquad
 A_k^\ell=2^{\ell-k}J\quad(\ell\geq k).
\end{equation}
\end{proposition}

\begin{proof}
After $k$ De Bruijn steps all $k$ initial digits have been discarded.  For a
specified terminal word, exactly one ordered choice of the $k$ appended
digits reaches it.  After $\ell\geq k$ steps the first $\ell-k$ appended
digits are discarded and arbitrary, while the last $k$ are forced by the
terminal word.  This gives respectively $1$ and $2^{\ell-k}$ directed paths.
The graph isomorphism transports the count to the modular Collatz graph.
\end{proof}

\begin{nonclaim}
Equation \eqref{eq:LDW-path-counts} yields a uniform endpoint modulo $2^k$
only after a probability model is specified---for example, independent
uniform choices of the two outgoing finite-graph edges, equivalently uniform
unresolved higher input bits.  It does not make the deterministic integer
$T^k(n)$ random, and it does not prove that the known low bits are
statistically independent of an integer iterate.  The source's phrase that
they tell us ``absolutely nothing'' is not used without this qualification.
\end{nonclaim}

\subsection{Edge projections, the inverse limit, and literal integer edges}

Let $\rho_{k+1,k}:V_{k+1}\to V_k$ be reduction.  Directly from itinerary
prefixes,
\begin{equation}
\label{eq:LDW-projective-encoder}
 \rho_{k+1,k}\circ\overline Q_{3,k+1}
 =\overline Q_{3,k}\circ\rho_{k+1,k}.
\end{equation}
Let $E_k\subseteq V_k\times V_k$ be the edge relation of $\mathcal T_k$.
Division by two loses exactly one congruence digit, so
\begin{equation}
\label{eq:LDW-edge-lower-precision}
 ([a]_k,[b]_k)\in E_k
 \quad\Longleftrightarrow\quad
 b\equiv T(a)\pmod{2^{k-1}}.
\end{equation}
The right side is independent of the representative $a$ because its parity
and branch value modulo $2^{k-1}$ are fixed by $[a]_k$.

\begin{proposition}[exact graph inverse limit and natural-edge recovery]
\label{prop:LDW-edge-inverse-limit}
Reduction of both endpoints maps $E_{k+1}$ onto $E_k$.  The inverse limit of
the vertex systems is $\Ztwo$, and its compatible edge relation is exactly
the graph of the deterministic map $T:\Ztwo\to\Ztwo$.  Under
\eqref{eq:LDW-projective-encoder}, the inverse-limit graph isomorphism is
$Q_3$, and the target relation is the graph of the deterministic digit shift
$S$.

Moreover, for fixed literal $a,b\in\Nzero$, once both lie in
$\{0,\ldots,2^k-1\}$,
\begin{equation}
\label{eq:LDW-eventual-natural-edge}
 \begin{aligned}
 b=T(a)
 &\quad\Longleftrightarrow\quad
 ([a]_k,[b]_k)\in E_k\text{ for all sufficiently large }k,\\
 &\quad\Longleftrightarrow\quad
 ([a]_k,[b]_k)\in E_k\text{ for infinitely many }k.
 \end{aligned}
\end{equation}
\end{proposition}

\begin{proof}
Equation \eqref{eq:LDW-edge-lower-precision} shows that every higher edge
reduces to a lower edge.  Conversely, either lift of the initial vertex fixes
one more input bit and hence selects one of the two lower successors; choosing
that lift and either next higher bit supplies a higher edge.  Thus edge
projection is onto.  Notice that for one fixed higher vertex its two
successors reduce to a singleton, not to the entire lower successor set.

Suppose compatible $x,y\in\Ztwo$ form an edge at every level.  Then
\eqref{eq:LDW-edge-lower-precision} gives
$y\equiv T(x)\pmod{2^{k-1}}$ for every $k$, so $y=T(x)$.  The converse follows
by reduction of an actual edge.  The same argument for suffixes gives the
graph of $S$, while \eqref{eq:LDW-projective-encoder} and
Proposition~\ref{prop:BL-Q} identify the limiting map as $Q_3$.

If $b=T(a)$, every sufficiently large literal pair is an edge.  Conversely,
an occurrence in $E_k$ gives $b\equiv T(a)\pmod{2^{k-1}}$.  Infinitely many
occurrences include one with $2^{k-1}>|b-T(a)|$, when the congruence forces
equality in $\Z$.  This proves \eqref{eq:LDW-eventual-natural-edge}.
\end{proof}

This is the typed repair of the source phrase ``edges that appear in
infinitely many'' changing finite graphs.  It distinguishes the common
ambient pair $(a,b)\in\Nzero^2$ from compatible residue vertices and does not
call the set-valued $\mathcal T_k$ a deterministic quotient map.

\subsection{Periodic components, predecessor fibres, and source corrections}

\begin{proposition}[exact binary cycle count]
\label{prop:LDW-cycle-count}
Let $M_k$ be the number of components of $C(\Ztwo)$ containing a cycle of
least length $k$.  Then
\begin{equation}
\label{eq:LDW-cycle-count}
 2^k=\sum_{d\mid k}dM_d,
 \qquad
 M_k=\frac1k\sum_{d\mid k}\mu(d)2^{k/d}
 =\frac{2^k}{k}+O(2^{k/2}).
\end{equation}
In particular,
\[
 (M_k)_{k\geq1}=(2,1,2,3,6,9,\ldots).
\]
\end{proposition}

\begin{proof}
The $Q_3$ conjugacy identifies $T$-cycles with shift orbits of purely
periodic binary sequences.  The infinite periodic repetition of every
length-$k$ block has a unique least shift period $d\mid k$.  A shift orbit of
least period $d$ contains exactly $d$ distinct phases, so counting all
length-$k$ blocks gives the first identity.
M\"obius inversion gives the second.  In the error sum every proper divisor
contributes at most $2^{k/2}$, and the number of divisors is at most $k$, which
gives the displayed asymptotic.  Direct substitution gives the listed values.
\end{proof}

The source prints the correct divisor identity and M\"obius formula, but its
asymptotic changes the bound variable from $k$ to an unbound $n$.  Its printed
sequence $(1,2,1,2,3,6,\ldots)$ contradicts both the formula and its own two
fixed-point components; Proposition~\ref{prop:LDW-cycle-count} records the
forced correction.  Likewise, a cyclic De Bruijn sequence gives a Hamiltonian
cycle.  The source calls it a path while its parenthesized wrap-around digits
supply the closing edge \cite[Secs.~2 and 4]{LaarhovenDeWeger2012}.

\begin{proposition}[exact predecessor fibres]
\label{prop:LDW-predecessor-fibres}
For every $y\in\Ztwo$,
\begin{equation}
\label{eq:LDW-predecessor-fibres}
 T^{-1}(y)=\left\{2y,\frac{2y-1}{3}\right\},
 \qquad
 S^{-1}(y)=\{2y,2y+1\}.
\end{equation}
Both fibres have two distinct elements.  In particular, the forward Collatz
map is not a function automorphism on any weak component.
\end{proposition}

\begin{proof}
The point $2y$ is even and maps to $y$.  Since $3$ is a unit in $\Ztwo$,
$(2y-1)/3$ is a $2$-adic integer; it is odd and its odd branch also maps to
$y$.  The two branch equations give no other preimages.  Equality of the two
would imply $4y=-1$, impossible in $\Ztwo$.  The shift formula is immediate
from its deleted first digit.  Both predecessors of a point lie in its weak
component, proving the final assertion.
\end{proof}

Source line 241 says the forward mapping on an aperiodic component is an
automorphism.  Equation \eqref{eq:LDW-predecessor-fibres} contradicts that
sentence under the ordinary meaning of function automorphism.  The separate
source assertion that all aperiodic components are pairwise isomorphic is not
used downstream.

\subsection{Generalized maps: exact unit gate and scaling quotient}

For odd $a,b$, Laarhoven--de Weger define
\[
 T^{(a,b)}(n)=
 \begin{cases}
  n/2,&n\equiv0\pmod2,\\
  (an+b)/2,&n\equiv1\pmod2.
 \end{cases}
\]
At retained modulus $2$, this is precisely the Matthews presentation
\[
 (m_0,r_0)=(1,0),
 \qquad
 (m_1,r_1)=(a,-b).
\]
Oddness of $a,b$ supplies the branch congruence and makes both multipliers
units.  Hence the finite and infinite parity encoders in source Section~5.1
are the $d=2$ specialization of
Theorem~\ref{thm:general-d-itinerary-conjugacy}; their orientation is
$\operatorname{It}_{T^{(a,b)}}$, not its inverse.  Only at
$(a,b)=(3,1)$ does this encoder specialize to $Q_3$.

For a retained base $d\geq2$, consider instead
\begin{equation}
\label{eq:LDW-general-p-branch}
 F(n)=\frac{a_i n+b_i}{d}
 \qquad(n\equiv i\pmod d).
\end{equation}
The exact hypotheses for the full digit-shift conclusion are
\begin{equation}
\label{eq:LDW-general-p-gate}
 a_i i+b_i\equiv0\pmod d,
 \qquad
 \gcd(a_i,d)=1
 \quad(0\leq i<d).
\end{equation}
The first is branch integrality.  For $n=i+dt$, the next digit is
\begin{equation}
\label{eq:LDW-next-digit-affine}
 F(i+dt)\equiv \frac{a_i i+b_i}{d}+a_it\pmod d.
\end{equation}
The second condition makes this an affine permutation of all $d$ next digits
with singleton fibres.  With $m_i=a_i$ and $r_i=-b_i$, conditions
\eqref{eq:LDW-general-p-gate} are exactly the retained Matthews congruence and
unit gate.  Theorem~\ref{thm:general-d-itinerary-conjugacy} therefore proves
the full finite and $d$-adic De Bruijn conclusions.  If a multiplier is not a
unit, Proposition~\ref{prop:nonunit-itinerary-obstruction} gives the missing
digits and non-singleton fibres explicitly.

The source's Section~5.2 says only ``appropriately chosen'' coefficients and
does not print \eqref{eq:LDW-general-p-gate}.  If that phrase means merely
integer-valued branches, choose $a_i=d,b_i=0$ for every $i$.  Then
$F=\operatorname{id}_{\Z}$, whose finite graph is not a De Bruijn graph.  The
source also leaves $x_i^{(f)}$ undefined and does not say that $p$ is prime
despite writing $\Z_p$.  The edition admits the claim only at the exact scope
above; for prime $d=p$ it is the usual $p$-adic statement.  The source's
ternary example has multipliers $2,4,4$, all units modulo $3$, and lies inside
the repaired scope.

\begin{proposition}[rational cycles and primitive $3n+b$ data]
\label{prop:LDW-rational-cycle-scaling}
Let $(x_0,\ldots,x_{L-1})$ be a rational $T$-cycle in $\Ztwo$, with cyclic
indices, and let $b>0$ be a common odd denominator.  Put $n_j=bx_j$.  Then
$(n_j)$ is an integer cycle of $T^{(3,b)}$, and
\begin{equation}
\label{eq:LDW-rational-scaling}
 T(x_j)=\frac1bT^{(3,b)}(n_j).
\end{equation}
Conversely every such integer cycle scales to a rational $T$-cycle.  The
resulting bijection is between rational cycles and the equivalence classes
generated by the scalings
\begin{equation}
\label{eq:LDW-scaling-equivalence}
 (b,(n_j))\sim(cb,(cn_j))
 \qquad(c\in\Npos\text{ odd}),
\end{equation}
or, equivalently, their unique representatives with $b>0$ and
$\gcd(b,n_0,\ldots,n_{L-1})=1$.  For a rational $T$-cycle this primitive
$b$ is automatically coprime to $3$.
\end{proposition}

\begin{proof}
Multiplication by odd $b$ preserves parity in $\Ztwo$.  On an even phase,
$bT(x_j)=n_j/2$.  On an odd phase,
$bT(x_j)=(3n_j+b)/2$.  This proves
\eqref{eq:LDW-rational-scaling} and both cycle directions.  Scaling numerator
and denominator by a common odd positive integer leaves every $x_j$ fixed.
Dividing by the common gcd terminates at a primitive representative, and any
two primitive positive-denominator representatives of the same rational
tuple are equal.  If the cycle has return length $L$ and contains $m$ odd
steps, the fixed-return equation shows that the reduced denominator of every
phase divides $2^L-3^m$.  Hence their primitive common denominator also
divides $2^L-3^m$.  If $m\geq1$, this integer is coprime to $3$.  If $m=0$,
every phase uses $x\mapsto x/2$, so periodicity forces the zero cycle and the
primitive common denominator is $1$.  Thus it is coprime to $3$ in both
cases.
\end{proof}

The source's first, unqualified ``one-to-one'' sentence requires
\eqref{eq:LDW-scaling-equivalence}: for every odd $b$, the integer cycle
$\{b,2b\}$ of $T^{(3,b)}$ scales to the same rational cycle $\{1,2\}$.
Its subsequent phrase ``having denominator $b$'' can instead be read as the
exact reduced common denominator; on that reading it already selects the
primitive representative and is correct.  Lagarias's primary treatment
explicitly retains the raw gcd strata and then defines this primitive layer
\cite[pp.~39--40]{Lagarias1990}.  The proposition makes both scopes explicit
without weakening the scaling identity or suppressing the raw scaling
fibres.

\begin{sourceissue}[remaining prose and hypothesis boundaries]
\label{issue:LDW-source-boundaries}
The source introduction describes absence of divergent paths by
$\lim_kT^k(n)<\infty$, although the $1\leftrightarrow2$ orbit has no limit;
the boundedness condition is $\sup_kT^k(n)<\infty$.  Source line 79 says
``for all $n$'' where its fixed-input argument requires ``for all $k$''.
The rational-periodicity statement in Sections~4 is explicitly a conjecture.
The average-multiplier discussion and the predicted $5n+1$ components are
heuristics; no theorem is attributed to either discussion.  Finally, the source's positive-integer
reformulation through $\frac13\Z\setminus\Z$ depends on earlier cited
literature and is not enlarged here before that exact dependency is checked.
\end{sourceissue}

\begin{nonclaim}
The finite graph theorem, exact path counts, inverse limit, cycle enumeration,
and generalized unit-gated conjugacy do not characterize
$Q_3(\Npos)\subset\Ztwo$.  They therefore do not prove that a positive
integer reaches $1$.  The arithmetic difficulty lies in the embedded subset
and its labels, not in the existence of the ambient De Bruijn isomorphism.
\end{nonclaim}

\subsection{Primary parity and rational-cycle dependencies}
\label{subsec:primary-parity-dependencies}

The De Bruijn source invokes three earlier results whose exact scopes now
have controlling primary witnesses.  Everett already proves the finite
parity bijection and its nonnegative affine fibres \cite{Everett1977};
Lagarias adds the finite permutation structure and proves the infinite parity
map \cite{Lagarias1985}; Bernstein proves the inverse-series inclusion
\cite{Bernstein1994}, and Lagarias's rational-cycle coordinates
\cite{Lagarias1990}.  This section keeps all three map orientations and both
the raw and reduced denominator layers visible.

Put
\[
 \mathcal R=\Q\cap\Ztwo=\Z_{(2)}.
\]
Lagarias 1985 denotes this ring by $\mathbf Q_2$; it is the set of rationals
with odd reduced denominator, not the full $2$-adic field.  Its elements are
exactly the finite or eventually periodic binary expansions in $\Ztwo$.

\begin{proposition}[Everett's finite bijection; Lagarias's permutation structure]
\label{prop:Lagarias-finite-parity}
For $k\geq1$, define
\[
 \overline Q_k([n]_k)
 =\left[\sum_{i=0}^{k-1}(T^i(n)\bmod2)2^i\right]_k.
\]
Then $\overline Q_k$ is a permutation of $\Z/2^k\Z$.  Its order is a power
of two, in fact it divides $2^k$.  Every permutation cycle of length $r$ at
level $k$ lifts at level $k+1$ either to one cycle of length $2r$ or to two
cycles of length $r$.
\end{proposition}

\begin{proof}
Everett's Theorem~1 proves the bijection by the complete four-case affine
induction in Theorem~\ref{thm:Everett-finite-fibres}.  Lagarias's Theorem B
states the permutation and power-of-two order
\cite[p.~7, Thm.~B]{Lagarias1985}; its proof is explicitly a sketch and omits
the induction details.  Bernstein independently states the lift alternative
\cite[p.~407]{Bernstein1994}.  For the later structure at the scope used here,
Proposition~\ref{prop:BL-Q} gives
\[
 a\equiv b\pmod{2^k}
 \quad\Longleftrightarrow\quad
 Q_3(a)\equiv Q_3(b)\pmod{2^k}.
\]
Thus $\overline Q_k$ is the reduction of the same isometric bijection $Q_3$
and is a permutation.  Each residue has two lifts.  A lifted permutation
cycle either returns each chosen lift after $r$ steps, yielding two
$r$-cycles, or exchanges the two lifts after $r$ steps, yielding one
$2r$-cycle.  Induction from level one proves that every cycle length and the
permutation order divide $2^k$.
\end{proof}

This is the finite labeling called $\Phi_k$ by Laarhoven--de Weger.  Its
cycles are cycles of a permutation of labels.  They are not cycles of the
two-successor modular Collatz relation $\mathcal T_k$, and
Proposition~\ref{prop:Lagarias-finite-parity} alone is not the graph
isomorphism of Theorem~\ref{thm:LDW-finite-debruijn}.

\begin{sourceissue}[Lagarias 1985 finite table and examples]
\label{sourceissue:Lagarias1985-finite-table}
The sentence immediately before Table~2 on printed p.~7 says that one-cycles
are omitted; the table itself is on printed p.~8.  Its $k=6$ row also omits
the nontrivial transposition $(14,46)$.  Direct evaluation gives
\[
 \overline Q_6(14)=46,
 \qquad
 \overline Q_6(46)=14.
\]
The displayed permutation order $4$ remains correct.  The calculation
certificate \path{certificates/lagarias_dependency_checks.py} enumerates
all residues for $k\leq6$ and verifies that this is the only nontrivial cycle
missing from the printed $k=6$ row.

On printed p.~18 the series for $Q_\infty(1)$ has summation index $i$ but
exponent $2n$; the intended exponent is $2i$.  More substantially, the
printed value $Q_\infty(3)=-20/3$ is inconsistent with the paper's own map:
the trajectory $3,5,8,4,2,1,2,\ldots$ has parity word
$11000(10)^\infty$, and hence
\[
 Q_\infty(3)=1+2+\frac{2^5}{1-2^2}=-\frac{23}{3}.
\]
The printed value is not used.
\end{sourceissue}

Lagarias's Eq.~(2.33) defines the infinite map in the forward orientation
\[
 Q_\infty(x)=\sum_{i\geq0}(T^i(x)\bmod2)2^i.
\]

\begin{proposition}[Lagarias infinite parity theorem, proof completed]
\label{prop:Lagarias-infinite-parity}
The map $Q_\infty$ is a continuous, bijective, measure-preserving self-map of
$\Ztwo$.  It is $Q_3$ and Laarhoven--de Weger $\Phi$; it is not the inverse
map $\Phi_3$.
\end{proposition}

\begin{proof}
Theorem L states the result \cite[pp.~17--18]{Lagarias1985}.  Its printed
argument proves prefix preservation, injectivity, surjectivity through finite
inverse residues, and inverse continuity.  Compatibility of those inverse
residues follows because every $\overline Q_k$ is a reduction of the same
prefix map.  Every residue cylinder modulo $2^k$ is therefore permuted onto
one cylinder of the same Haar measure, which supplies the measure-preservation
step not printed at the end of the source proof.
\end{proof}

\begin{proposition}[one proved half of Periodicity]
\label{prop:Bernstein-periodicity-half}
Let $\Phi_3=Q_3^{-1}$.  Then
\begin{equation}
\label{eq:Bernstein-rational-inclusion}
 \Phi_3(\mathcal R)\subseteq\mathcal R,
\end{equation}
or equivalently
\begin{equation}
\label{eq:Bernstein-rational-inclusion-forward}
 \mathcal R\subseteq Q_3(\mathcal R).
\end{equation}
The Periodicity Conjecture is the equality
\begin{equation}
\label{eq:Lagarias-periodicity-conjecture}
 Q_3(\mathcal R)=\mathcal R.
\end{equation}
\end{proposition}

\begin{proof}
Bernstein writes a rational parity address as a finite or eventually
periodic sequence of one-positions and substitutes it into his explicit
series for $\Phi_3$.  The resulting geometric relation proves that the image
is rational \cite[p.~407, Cor.~1]{Bernstein1994}.  Applying the bijection
$Q_3$ gives the equivalent forward inclusion.  Lagarias prints
\eqref{eq:Lagarias-periodicity-conjecture} as an unproved conjecture on
printed p.~18 \cite{Lagarias1985}; no reverse inclusion is added.
\end{proof}

The positive-integer ``third form'' in Lagarias 1985 is likewise a conjecture:
\begin{equation}
\label{eq:Lagarias-third-form}
 Q_3(\Npos)\subseteq\frac13\Z\setminus\Z.
\end{equation}
Bernstein instead proves that the inverse-oriented assertion
$\Z^+\subseteq\Phi_3((1/3)\Z)$ is equivalent to the ordinary conjecture
\cite[pp.~406--407, Thms.~2--3]{Bernstein1994}.  An equivalence to an endpoint
conjecture is not a proof of either formulation.

Lagarias's last exploratory value on printed p.~18 is left as
$Q_3(-82/7)=?/15$.  Exact iteration gives a prefix of length $13$ followed by
the parity block $1100$, hence the independent finite calculation
\begin{equation}
\label{calc:Lagarias1985-minus82}
 Q_3(-82/7)=\frac{18098}{5}=\frac{54294}{15}.
\end{equation}
The same certificate reproduces this value.  It is not attributed to the
source and does not prove Periodicity.

\begin{theorem}[Lagarias rational fixed-return coordinates]
\label{thm:Lagarias-rational-return}
Let $n\geq1$ and $v=(v_0,\ldots,v_{n-1})\in\{0,1\}^n$, and put
\[
 m(v)=\sum_{i=0}^{n-1}v_i,
 \qquad
 D(v)=2^n-3^{m(v)}\neq0,
\]
\[
 N(v)=\sum_{j=0}^{n-1}
 v_j2^j3^{v_{j+1}+\cdots+v_{n-1}}.
\]
There is exactly one $x(v)\in\mathcal R$ with
$T^n(x(v))=x(v)$ and initial parity word $v$, namely
\begin{equation}
\label{eq:Lagarias-rational-return}
 x(v)=\frac{N(v)}{D(v)}.
\end{equation}
Its least $T$-period is $n$ exactly when $v$ is not a proper repeated block.
\end{theorem}

\begin{proof}
The branch composition determined by $v$ is
\[
 U_{v_{n-1}}\circ\cdots\circ U_{v_0}(x)
 =\frac{3^{m(v)}x+N(v)}{2^n},
\]
and $D(v)\neq0$ because $2^n$ is even while $3^{m(v)}$ is odd.  Thus its
unique formal fixed point is
\eqref{eq:Lagarias-rational-return}.  Lagarias proves that this formal point
has the prescribed actual branches: for each $y\in\mathcal R$, exactly one of
the two forward branch expressions $U_0(y)=y/2$ and
$U_1(y)=(3y+1)/2$ lies in $\mathcal R$, with the choice determined by the
parity of $y$; if $y\notin\mathcal R$, both expressions remain outside
$\mathcal R$ \cite[pp.~36--37, Thm.~2.1]{Lagarias1990}.  The denominator
$D(v)$ is odd, so the formal fixed point lies in $\mathcal R$.  If its
prescribed composition ever selected the other forward expression, it would
leave $\mathcal R$ and could never return, contradicting that the composition
ends again at $x(v)\in\mathcal R$.  Thus every prescribed branch is actual.
For every $d\geq1$, the conjugacy in
Proposition~\ref{prop:BL-Q} gives
$Q_3(T^d(x(v)))=S^d(Q_3(x(v)))$.  Hence a $T$-return after $d$ steps implies
a shift return after $d$ steps.  Conversely, a shift return and the
injectivity of $Q_3$ from Proposition~\ref{prop:Lagarias-infinite-parity}
imply $T^d(x(v))=x(v)$.  The two sets of positive return times are equal, so
their least elements agree; this least period is $n$ exactly when $v$ is not
a proper repeated block.
\end{proof}

Write the positive reduced common denominator as
\begin{equation}
\label{eq:Lagarias-reduced-denominator}
 \delta(v)=\frac{|D(v)|}{\gcd(|N(v)|,|D(v)|)}.
\end{equation}
It is always odd and prime to $3$.  The proof cannot omit the all-even word.
If $m(v)\geq1$, then $3\nmid D(v)$.  If $m(v)=0$, the return equation is
$x=x/2^n$, so $x=0$ and $\delta(v)=1$.  Thus
\begin{equation}
\label{eq:Lagarias-denominator-mod-six}
 \delta(v)\equiv\pm1\pmod6.
\end{equation}

\begin{proposition}[raw scaling fibres and the primitive layer]
\label{prop:Lagarias-primitive-scaling}
For odd $b$, let
\[
 T^{(3,b)}(n)=
 \begin{cases}
 (3n+b)/2,&n\text{ odd},\\
 n/2,&n\text{ even}.
 \end{cases}
\]
Multiplication by a common positive odd denominator gives
\[
 bT(x_j)=T^{(3,b)}(bx_j).
\]
Consequently rational $T$-cycles correspond to common-odd-scaling classes
of integer $T^{(3,b)}$-cycles.  Without identifying the raw representatives,
the unique member with $b$ equal to the positive reduced common denominator
is characterized by
\[
 \gcd(b,n_0,\ldots,n_{L-1})=1.
\]
It is exactly a primitive integer $T^{(3,b)}$-cycle, and division by $b$ is
the two-sided inverse.  Its parameter satisfies $b\equiv\pm1\pmod6$.
\end{proposition}

\begin{proof}
The branch identity is the calculation in
Proposition~\ref{prop:LDW-rational-cycle-scaling}.  Lagarias explicitly
defines the primitive condition and the inverse correspondence on printed
pp.~39--40 \cite{Lagarias1990}.  For the raw strata, if $d\mid b$ and
$n=dm$, $b=dc$, then
\[
 T^{(3,b)}(dm)=dT^{(3,c)}(m).
\]
This proves every scaling fibre without suppressing it.  Choosing the reduced
common denominator gives the displayed gcd condition and a literal inverse.
Equation~\eqref{eq:Lagarias-denominator-mod-six} supplies the parameter
restriction.
\end{proof}

\begin{sourceissue}[Lagarias 1990 local coordinate defects]
\label{sourceissue:Lagarias1990-local-defects}
Six printed statements require exact local repair.

First, the itinerary sentence at the top of printed p.~34 defines
$v=(v_0,\ldots,v_{n-1})$ but quantifies
$T^i(x(v))\equiv v_i\pmod2$ for $1\leq i\leq n$, which skips $v_0$ and invokes
the undefined $v_n$.  The correct range is $0\leq i\leq n-1$.  Second,
printed p.~39 says that $T$ never changes the denominator of an
arbitrary element of $\mathcal R$, but $T(1/3)=1$.  The correct preserved
scope is a reduced odd denominator prime to $3$; every periodic rational has
that property by the case split above.  Third, the printed formula for
$\delta(v)$ needs $|D(v)|$ for a positive denominator unless a signed-gcd
convention is supplied.  Fourth, printed p.~40 writes
$T^{(3,b)}(n)=T^{(3,-b)}(-n)$, whereas direct branch evaluation gives
\[
 T^{(3,b)}(n)=-T^{(3,-b)}(-n).
\]
Fifth, printed p.~34 ends one correspondence sentence with
$b\equiv1\pmod6$; the preceding sentence, printed p.~40, and the examples
require $b\equiv\pm1\pmod6$.  Sixth, the invariant gcd stratum on p.~40 uses
the standing $\gcd(b,3)=1$ hypothesis and is not a statement for arbitrary
odd $b$.
\end{sourceissue}

\begin{nonclaim}
The Periodicity Conjecture does not follow from the one-sided Bernstein
inclusion.  Lagarias's primitive-cycle existence and finiteness statements
are conjectures.  A periodic parity address gives a rational periodic state,
but this does not classify the positive integers inside $\Ztwo$, exclude
nontrivial integer cycles, or prove convergence to $1$.
\end{nonclaim}


\subsection{Relatively-prime generalized maps and the full
\texorpdfstring{$d$}{d}-adic shift}

Matthews's author form first permits a retained presentation integer
$d\geq2$, nonzero integers
$m_0,\ldots,m_{d-1}$, and integers $r_i\equiv i m_i\pmod d$, and writes
\begin{equation}
\label{eq:Matthews-general-map}
 \mathcal T(x)=\frac{m_i x-r_i}{d}
 \qquad(x\equiv i\pmod d).
\end{equation}
The source symbol is $T$; the calligraphic letter here prevents collision with
the already fixed shortened $3x+1$ map.  The source separately reserves
$d(\mathcal T)$ for the least modulus on whose residue classes the map is
affine.  On an arbitrary retained presentation impose the unit condition
\begin{equation}
\label{eq:Matthews-unit-gate}
 \gcd(m_0m_1\cdots m_{d-1},d)=1,
\end{equation}
equivalently when every $m_i$ is a unit modulo $d$.  This condition forces the
retained $d$ to equal $d(\mathcal T)$: if affine data at a smaller modulus $q$
described the same map, equality of rational slopes on each nonempty
intersection of a $q$-class and a $d$-class would give
$q m_i=d m'_j$; coprimality gives $d\mid q$, a contradiction.  At the least
modulus this is exactly Matthews's \emph{relatively-prime type}
\cite[author-form pp.~2--3]{Matthews2010AuthorForm}.

The 1984 Matthews--Watts paper uses a nonidentical presentation: positive
unit multipliers $m_1,\ldots,m_d$, a complete residue set
$R_d=\{x_1,\ldots,x_d\}$, and branch data indexed by those representatives.
The exact crosswalk sends the later digit $i$ to the unique $x_j\equiv i\pmod
d$ and relabels the corresponding $m_j,r_j$.  This reindexing does not erase
the earlier positivity hypothesis.

\begin{proposition}[source extension and cylinder theorems]
\label{prop:matthews-extension-cylinders}
Every map \eqref{eq:Matthews-general-map} at the retained $d$ has a unique
continuous extension
$\widehat{\mathcal T}:\mathbb Z_d\to\mathbb Z_d$ with the same branch formula.
Relative primeness is not needed for this assertion.

Assume \eqref{eq:Matthews-unit-gate}.  For a word
$\mathbf i=(i_0,\ldots,i_{\alpha-1})\in\{0,\ldots,d-1\}^{\alpha}$, put
\[
 b_{\alpha}(\mathbf i)=
 \sum_{j=0}^{\alpha-1}
 \frac{r_{i_j}d^j}{m_{i_0}m_{i_1}\cdots m_{i_j}}
 \in\mathbb Z_d.
\]
Then its itinerary cylinder is exactly
\begin{equation}
\label{eq:Matthews-cylinder-coordinate}
 \left\{x\in\mathbb Z_d:
 \widehat{\mathcal T}^{j}(x)\equiv i_j\pmod d
 \text{ for }0\leq j<\alpha\right\}
 =b_{\alpha}(\mathbf i)+d^{\alpha}\mathbb Z_d.
\end{equation}
Thus the $d^{\alpha}$ length-$\alpha$ cylinders are exactly the
$d^{\alpha}$ residue classes modulo $d^{\alpha}$.
\end{proposition}

\begin{proof}
On the clopen cylinder $x\equiv i\pmod d$, the congruence
$r_i\equiv i m_i\pmod d$ puts $m_ix-r_i$ in $d\mathbb Z_d$; division by $d$
therefore gives a continuous branch into $\mathbb Z_d$.  The branches assemble
continuously, and uniqueness follows from density of $\mathbb Z$ in
$\mathbb Z_d$.  When $d=d(\mathcal T)$ this is Matthews's Theorem~6.1
\cite[author-form p.~12]{Matthews2010AuthorForm}; the displayed clopen-branch
proof also establishes the stated retained-presentation version.

For the second assertion, every denominator in $b_{\alpha}$ is a unit.  Put
\[
 X_k=\sum_{j=k}^{\alpha-1}
 \frac{r_{i_j}d^{j-k}}{m_{i_k}\cdots m_{i_j}}
 \qquad(0\leq k<\alpha).
\]
Then $X_k\equiv r_{i_k}/m_{i_k}\equiv i_k\pmod d$ and
$(m_{i_k}X_k-r_{i_k})/d=X_{k+1}$ for $k<\alpha-1$.  Adding any element of
$d^{\alpha}\mathbb Z_d$ to $X_0=b_{\alpha}$ preserves these first $\alpha$
digits, because the difference after $j$ common branches remains divisible by
$d^{\alpha-j}$.  Conversely, repeated rearrangement of the branch equation for
an $x$ with the displayed itinerary gives
\[
 x=b_{\alpha}(\mathbf i)
 +\frac{d^{\alpha}\widehat{\mathcal T}^{\alpha}(x)}
 {m_{i_0}\cdots m_{i_{\alpha-1}}},
\]
whose second term lies in $d^{\alpha}\mathbb Z_d$.  This proves
\eqref{eq:Matthews-cylinder-coordinate}.  It reconstructs author-form
Theorem~2.2(ii) and is the later indexing of Matthews--Watts Eq.~(2.4) and
Section~3, Lemma~4
\cite[author-form p.~4]{Matthews2010AuthorForm}
\cite[pp.~169, 172--173]{MatthewsWatts1984}.
\end{proof}

Let $S_d:\mathbb Z_d\to\mathbb Z_d$ be the digit shift
\[
 S_d\!\left(\sum_{j\geq0}a_jd^j\right)=
 \sum_{j\geq0}a_{j+1}d^j.
\]
For $x\in\mathbb Z_d$, let $i_j(x)\in\{0,\ldots,d-1\}$ be defined by
$\widehat{\mathcal T}^{j}(x)\equiv i_j(x)\pmod d$.

\begin{theorem}[constructive generalized itinerary conjugacy]
\label{thm:general-d-itinerary-conjugacy}
Under \eqref{eq:Matthews-unit-gate}, the map
\begin{equation}
\label{eq:general-itinerary-encoder}
 \operatorname{It}_{\mathcal T}:\mathbb Z_d\longrightarrow\mathbb Z_d,
 \qquad
 \operatorname{It}_{\mathcal T}(x)=\sum_{j\geq0}i_j(x)d^j,
\end{equation}
is an isometric homeomorphism for the digit-prefix metric.  Its inverse is the
coordinate series
\begin{equation}
\label{eq:general-itinerary-inverse}
 \operatorname{It}_{\mathcal T}^{-1}
 \!\left(\sum_{j\geq0}i_jd^j\right)
 =\sum_{j\geq0}
 \frac{r_{i_j}d^j}{m_{i_0}m_{i_1}\cdots m_{i_j}}.
\end{equation}
It conjugates the extended map to the full shift:
\begin{equation}
\label{eq:general-itinerary-conjugacy}
 \operatorname{It}_{\mathcal T}\circ\widehat{\mathcal T}
 =S_d\circ\operatorname{It}_{\mathcal T},
 \qquad
 \widehat{\mathcal T}
 =\operatorname{It}_{\mathcal T}^{-1}
  \circ S_d\circ\operatorname{It}_{\mathcal T}.
\end{equation}
Every fibre is a singleton.  No additive or ring-homomorphism property is
asserted, so no group-theoretic kernel is attached to this map.
\end{theorem}

\begin{proof}
For a prescribed digit sequence $(i_j)_{j\geq0}$, define the convergent suffix
series
\[
 X_k=\sum_{j=k}^{\infty}
 \frac{r_{i_j}d^{j-k}}{m_{i_k}m_{i_{k+1}}\cdots m_{i_j}}.
\]
The unit condition gives
\[
 X_k\equiv\frac{r_{i_k}}{m_{i_k}}\equiv i_k\pmod d,
 \qquad
 \frac{m_{i_k}X_k-r_{i_k}}d=X_{k+1}.
\]
Hence $X_0$ has exactly the prescribed itinerary, proving that
\eqref{eq:general-itinerary-inverse} is a right inverse.

For an arbitrary $x$, repeated rearrangement of its actual branch equations
gives, for every $\alpha\geq1$,
\[
 x=
 \sum_{j=0}^{\alpha-1}
 \frac{r_{i_j(x)}d^j}{m_{i_0(x)}\cdots m_{i_j(x)}}
 +\frac{d^{\alpha}\widehat{\mathcal T}^{\alpha}(x)}
 {m_{i_0(x)}\cdots m_{i_{\alpha-1}(x)}}.
\]
The final term tends to zero in the digit-prefix metric.  Thus
\eqref{eq:general-itinerary-inverse} is also a left inverse.

If $x\equiv y\pmod{d^{\alpha}}$, common-branch induction gives
$\widehat{\mathcal T}^{j}(x)\equiv
\widehat{\mathcal T}^{j}(y)\pmod{d^{\alpha-j}}$ for
$0\leq j\leq\alpha$; hence their first $\alpha$ itinerary digits agree.
The converse follows from Proposition~\ref{prop:matthews-extension-cylinders}.
Therefore
\[
 x\equiv y\pmod{d^{\alpha}}
 \quad\Longleftrightarrow\quad
 \operatorname{It}_{\mathcal T}(x)\equiv
 \operatorname{It}_{\mathcal T}(y)\pmod{d^{\alpha}},
\]
which proves the isometry and continuity in both directions.  Finally,
$i_j(\widehat{\mathcal T}x)=i_{j+1}(x)$, so deleting the first itinerary digit
proves \eqref{eq:general-itinerary-conjugacy}.
\end{proof}

This is a constructive proof, with no novelty claim, of the general-$d$
statement that the author form presents only as something that ``should hold''
\cite[author-form p.~14]{Matthews2010AuthorForm}.  At $d=2$ for the shortened
$3x+1$ map, \eqref{eq:general-itinerary-encoder} is pointwise the
Bernstein--Lagarias $Q_3$ and Laarhoven--de Weger's $\Phi$ orientation.  It is
an explicit choice having the orientation of Matthews's $C_T$.  Matthews's
Theorem~6.3 alone does not state a coordinate normalization that would identify
its named $C_T$ pointwise merely from the conjugacy equation.

\begin{corollary}[Haar cylinder independence and strong mixing]
\label{cor:Matthews-strong-mixing}
The map $\operatorname{It}_{\mathcal T}$ preserves normalized Haar measure
$\mu_d$.  If
$A=a+d^{\alpha}\mathbb Z_d$ and $B=b+d^{\beta}\mathbb Z_d$, then for
$K\geq\beta$,
\begin{equation}
\label{eq:Matthews-cylinder-mixing}
 \mu_d(\widehat{\mathcal T}^{-K}(A)\cap B)
 =\mu_d(A)\mu_d(B)=d^{-\alpha-\beta}.
\end{equation}
Consequently $\widehat{\mathcal T}$ is strongly mixing and ergodic on
$(\mathbb Z_d,\mu_d)$.
\end{corollary}

\begin{proof}
The isometry permutes the $d^{\alpha}$ residue balls at every level, so it
preserves their measures and hence Haar measure on the Borel sigma-algebra.
Under the conjugacy, $B$ fixes the first $\beta$ digits and $S_d^{-K}(A)$ fixes
digits $K$ through $K+\alpha-1$.  For $K\geq\beta$ these blocks are disjoint,
giving exactly $d^{-\alpha-\beta}$.  Finite unions of residue balls approximate
every Borel set in Haar measure.  Replacing two Borel sets by such unions makes
the error in the intersection at most the sum of the two symmetric-difference
measures, uniformly in $K$ because the maps preserve measure.  Letting that
error tend to zero proves strong mixing; ergodicity follows by applying the
mixing identity to an invariant set.  This reproduces Matthews--Watts Lemma~5
and Corollary~1 and Matthews's later Theorem~6.2
\cite[pp.~173--174]{MatthewsWatts1984}
\cite[author-form pp.~12--13]{Matthews2010AuthorForm}.
\end{proof}

\begin{proposition}[exact nonunit obstruction]
\label{prop:nonunit-itinerary-obstruction}
Suppose instead that $g_i=\gcd(m_i,d)>1$ for some branch $i$.  Among the
$d$ lifts $i+dt$ modulo $d^2$, the next itinerary digit assumes exactly
$d/g_i$ values, each with $g_i$ lifts.  Hence some two-digit words beginning
with $i$ are absent and the canonical itinerary encoder is not a bijection
onto the full $d$-shift.
\end{proposition}

\begin{proof}
The branch equation gives
\[
 \mathcal T(i+dt)=\mathcal T(i)+m_it.
\]
Modulo $d$, the image of multiplication by $m_i$ has cardinality $d/g_i$ and
every image element has $g_i$ preimages.  These are exactly the possible
second digits and their lift multiplicities.
\end{proof}

For the unshortened Collatz map in Matthews's presentation,
$d=2,m_1=6,r_1=-2$, so every odd input maps to an even value and the word
$11$ is absent.  Proposition~\ref{prop:nonunit-itinerary-obstruction} obstructs
the canonical itinerary map from being a full-shift conjugacy; it does not
assert categorical nonexistence of every differently defined conjugacy.

\begin{remark}[formal certificate boundary]
The pinned Lean certificate checks the literal lift identity, the explicit
unit-affine inverse, the equality between the residue gcd and the ordinary gcd
of the retained integer multiplier, the
$d/g_i$ attainable-digit count, and the $g_i$-element fibre over every
attainable digit.  In particular, the range and fibre assertions in
Proposition~\ref{prop:nonunit-itinerary-obstruction} are kernel-checked.  The
artifact deliberately does not certify the inverse limit, infinite series,
homeomorphism, or full shift conjugacy. Those assertions therefore remain
outside the formal certificate; extending the formalization is project work,
not an open mathematical problem.
\end{remark}

\begin{nonclaim}
The Haar-almost-every conclusions live in $\mathbb Z_d$.  The embedded ordinary
integers are countable and Haar-null.  Neither the conjugacy nor strong mixing
proves an endpoint statement for positive integer Collatz orbits; it transports
the unresolved embedding of that special subset into shift coordinates.
\end{nonclaim}

\subsection{Integer rcwa presentations: exact crosswalk and a signed-slope defect}

Kohl calls a map $f:\mathbb Z\to\mathbb Z$ \emph{residue-class-wise affine}
(rcwa) when some positive integer $q$ admits coefficients
$a_i,b_i,c_i\in\mathbb Z$ for which
\begin{equation}
\label{eq:Kohl-rcwa-presentation}
 f(n)=\frac{a_i n+b_i}{c_i}
 \qquad(n\equiv i\pmod q).
\end{equation}
The least such $q$ is $\operatorname{Mod}(f)$.  At the least modulus the source
stores each coefficient triple in primitive form and takes $c_i>0$
\cite[Defs.~1.1 and 2.1]{Kohl2008}.  Algorithm~2.2 separately permits a
nonminimal modulus and unreduced input before reduction.  The following bridge
therefore retains its presentation modulus; it does not normalize either source.

\begin{proposition}[fixed-modulus Matthews--rcwa crosswalk]
\label{prop:Matthews-Kohl-crosswalk}
Fix $d\geq2$.  The presentation-level assignment
\begin{equation}
\label{eq:Matthews-to-Kohl}
 (m_i,r_i)_{0\leq i<d}
 \longmapsto
 (a_i,b_i,c_i)_{0\leq i<d}
 =(m_i,-r_i,d)_{0\leq i<d}
\end{equation}
sends the Matthews presentation \eqref{eq:Matthews-general-map} to an rcwa
presentation of the same function $\mathbb Z\to\mathbb Z$.  Its image consists
of the fixed-modulus, common-denominator presentations with nonzero branch
slopes, and its inverse on that presentation class is
$m_i=a_i$, $r_i=-b_i$.

Conversely, let \eqref{eq:Kohl-rcwa-presentation} be a reduced rcwa coefficient
list at the retained modulus $d$, and assume $a_i\neq0$ for every $i$.  Then
\begin{equation}
\label{eq:Kohl-to-Matthews}
 m_i=\frac{a_i d}{c_i},
 \qquad
 r_i=-\frac{b_i d}{c_i}
\end{equation}
are integers, satisfy $r_i\equiv i m_i\pmod d$, and reproduce the same branch
map in Matthews form.  Returning through \eqref{eq:Matthews-to-Kohl} scales the
$i$th reduced triple by $d/c_i$.  Reducing that scaled triple recovers the
starting reduced triple exactly; the literal common-denominator list and the
reduced list are distinct but reversibly related while $d$ is retained.
\end{proposition}

\begin{proof}
The Matthews congruence puts $m_i n-r_i$ in $d\mathbb Z$ throughout
$i+d\mathbb Z$, which proves the first assertion and its displayed inverse.
For the converse, Kohl's reduced-list divisibility gives
$c_i\mid d$ and $c_i\mid a_i i+b_i$.  Hence the integers in
\eqref{eq:Kohl-to-Matthews} obey
\[
 i m_i-r_i=\frac{d(a_i i+b_i)}{c_i}\equiv0\pmod d,
\]
and direct substitution gives
$(m_i n-r_i)/d=(a_i n+b_i)/c_i$.

For Matthews data, put $g_i=\gcd(m_i,d)$.  The congruence
$r_i\equiv i m_i\pmod d$ gives
\[
 \gcd(m_i,r_i,d)=g_i.
\]
Therefore branchwise reduction sends
$(m_i,-r_i,d)$ exactly to
$(m_i/g_i,-r_i/g_i,d/g_i)$.  Applying
\eqref{eq:Kohl-to-Matthews} with the retained $d$ recovers $(m_i,r_i)$.
Conversely, primitivity and $c_i\mid a_i i+b_i$ imply
$\gcd(a_i,c_i)=1$; hence for the reconstructed data
$\gcd(a_i d/c_i,d)=d/c_i$, so reduction returns precisely
$(a_i,b_i,c_i)$.  This proves exact reversibility at fixed $d$.
\end{proof}

This bridge separates three conditions which are not interchangeable.  In the
common-denominator presentation, the unit condition is
$\gcd(m_i,d)=1$ for every branch.  As proved above it forces $d$ to be the
least modulus, makes every triple $(m_i,-r_i,d)$ primitive, and leaves every
reduced denominator equal to $d$.  Hence Kohl's divisor is exactly $d$.
Kohl calls a reduced rcwa map \emph{integral} when this divisor is one, so for
$d\geq2$ Matthews relatively-prime type excludes Kohl integrality.  Neither
condition says that the integer map is bijective.  For example, the shortened
$3x+1$ map has reduced triples
$(1,0,2)$ and $(3,1,2)$: it is Matthews-relatively-prime, has Kohl divisor two,
and is surjective but not injective on $\mathbb Z$.

\begin{sourceissue}[signed slopes in Kohl Remark 1.2]
\label{issue:Kohl-signed-density}
For a reduced coefficient list at modulus $q$, the printed remark uses
\begin{equation}
\label{eq:Kohl-printed-density}
 \frac1q\sum_{i=0}^{q-1}\frac{c_i}{a_i}
\end{equation}
and asserts that it is at most one for an injection, at least one for a
surjection, and equal to one for a bijection
\cite[Rem.~1.2]{Kohl2008}.  As printed, the expression is undefined on a
zero-slope branch and false for negative slopes.  The same article lists
$f(n)=-n$ as an rcwa permutation; at modulus one its reduced triple is
$(-1,0,1)$, so \eqref{eq:Kohl-printed-density} equals $-1$, not $1$
\cite[Def.~2.4 and following example]{Kohl2008}.
\end{sourceissue}

\begin{proposition}[corrected rcwa image-density statistic]
\label{prop:Kohl-density-repair}
For a reduced rcwa list at modulus $q$, put
\begin{equation}
\label{eq:Kohl-corrected-density}
 D(f)=\frac1q\sum_{\substack{0\leq i<q\\a_i\neq0}}
 \frac{c_i}{|a_i|}.
\end{equation}
If $f$ is injective, then $D(f)\leq1$; if $f$ is surjective, then
$D(f)\geq1$; and if $f$ is bijective, then $D(f)=1$.
\end{proposition}

\begin{proof}
On $i+q\mathbb Z$ with $a_i\neq0$, the image of the branch is one residue
class of modulus $|a_i|q/c_i$ and therefore has natural density
$c_i/(|a_i|q)$.  A zero-slope branch has a singleton image and density zero.
If $f$ is injective, the branch images are disjoint, so their density sum is at
most one.  If $f$ is surjective, their finite union, together with any singleton
images, is all of $\mathbb Z$, so the sum of their densities counted with
overlap is at least one.  A bijection satisfies both bounds.
\end{proof}

The absolute value and omission of zero slopes are forced, while the positive-
slope Collatz calculation in the source remains unchanged.  This repaired
integer-image statistic neither proves the $d$-adic itinerary conjugacy nor
transfers a Haar-almost-every statement to ordinary integer orbits.

\subsection{Siegel's Hydra and numen formalism: typed source boundary}
\label{sec:Siegel-Hydra-numen}

The controlling witness in this subsection is Siegel's version-pinned source
for arXiv:2601.17030v3 \cite{Siegel2026v3}.  Its Hydra and numen sections are
unchanged from v2 to v3; the v3 change is confined to the later Fourier
calculation audited at the end of this subsection.  The source permits both
number fields and finite extensions of $\mathbb F_q(T)$.  Accordingly,
$\mathcal O_K$ is a finite free module over the integer ring
$\mathcal O_{\mathbb F}$ of the base field; it has a finite
$\mathbb Z$-basis only in the number-field case.

Let $K/\mathbb F$ be such a global field, let
$\Lambda\subsetneq\mathcal O_K$ be a nonzero ideal, and retain an enumeration
of the $p=\lvert\mathcal O_K/\Lambda\rvert$ cosets by
$\{0,\ldots,p-1\}$, with $0$ denoting $\Lambda$.  If $\ell$ is a place, an
$\mathbb F$-linear map used on $K_\ell$ will always be assumed bounded at
$\ell$.  This assumption is necessary: an arbitrary
$\mathbb F$-linear endomorphism of $K$ need not extend continuously to a fixed
completion.  The operator norm is
\[
 \lVert r\rVert_\ell
 =\sup_{z\in K\setminus\{0\}}\frac{|r(z)|_\ell}{|z|_\ell},
\]
when this supremum is finite.  The source display includes $z=0$ in a quotient
by $|z|_\ell$ and its following remark asserts extendibility without this
boundedness condition.

\begin{definition}[Siegel's printed branch data]
\label{def:Siegel-Hydra-v3}
A pre-Hydra has, for every $j\in\mathcal O_K/\Lambda$, an invertible affine
$\mathbb F$-linear branch
\[
 H_j:K\longrightarrow K,\qquad H_j(z)=r_j(z)+c_j,
 \qquad
 r_j\in\operatorname{End}_{\mathbb F}^{\times}(K),\quad c_j\in K.
\]
The selected piecewise evaluator is
\[
 H(z)=H_{[z]_\Lambda}(z)\qquad(z\in\mathcal O_K).
\]
The source's literal Definition~5.3 additionally requires
$H_j(z)\in\mathcal O_K$ for every $z\in\mathcal O_K$ and every $j$.
\end{definition}

The enumeration in the source should read
$i=\Lambda_{\beta(i)}$ for a coset $i\in\mathcal O_K/\Lambda$, not
$i\in\Lambda_{\beta(i)}$.  After the cosets have been replaced by digit
labels, $\operatorname{DigSum}_p$ uses those labels directly; applying
$\beta$ to them a second time is ill-typed.  We retain the enumeration rather
than normalize it away.

\begin{sourceissue}[the two printed integrality predicates]
\label{issue:Siegel-Hydra-integrality}
Definition~5.3 requires
\[
 H_j(\mathcal O_K)\subseteq\mathcal O_K
 \qquad\text{for every }j,
\tag{G}
\]
whereas Definition~5.4 calls the same Hydra \emph{integral} when
\[
 H_j(z)\in\mathcal O_K
 \quad\Longleftrightarrow\quad z\in j+\Lambda
 \qquad(z\in\mathcal O_K).
\tag{E}
\]
These conditions cannot coexist.  Since $\Lambda$ is proper, choose a coset
different from $j+\Lambda$ and an integer $z$ in it.  Condition~(G) says
$H_j(z)$ is integral and condition~(E) says it is not.  The source's own
shortened Collatz example displays the intended distinction:
$(3z+1)/2$ is integral exactly on odd $z$, so it satisfies~(E) but not~(G).
\end{sourceissue}

We therefore use three explicitly different predicates:
\[
\begin{array}{ll}
\text{\emph{selected-branch admissible}}&
 H_j(j+\Lambda)\subseteq\mathcal O_K,\\[2mm]
\text{\emph{exact-integral on global integers}}&
 H_j^{-1}(\mathcal O_K)\cap\mathcal O_K=j+\Lambda,\\[2mm]
\text{\emph{inverse-integral at a completion $v$}}&
 H_j^{-1}(\mathcal O_{K_v})=j+\Lambda\mathcal O_{K_v}.
\end{array}
\]
The first makes the piecewise evaluator a map
$\mathcal O_K\to\mathcal O_K$.  The second detects the assigned coset among
global integers.  The third is the stronger local inverse-image statement
needed by a backward branch-correctness argument.  It is not implied merely
by writing the second predicate.

\begin{proposition}[scalar Hydra--Matthews bridge at retained modulus]
\label{prop:Siegel-scalar-Hydra-Matthews}
Let $K=\mathbb Q$, $\mathcal O_K=\mathbb Z$, and
$\Lambda=D\mathbb Z$, with the residue labels $0,\ldots,D-1$ retained.
Scalar selected-branch-admissible data
\[
 H_i(n)=a_i n+b_i
\]
are in exact two-sided correspondence with Matthews data
\[
 H_i(n)=\frac{m_i n-s_i}{D},
 \qquad m_i,s_i\in\mathbb Z,\qquad
 s_i\equiv i m_i\pmod D,
\]
by
\[
 m_i=Da_i,\qquad s_i=-Db_i.
\]
The complete integrality locus of the $i$th extension is
\[
 i+\frac{D}{\gcd(m_i,D)}\mathbb Z.
\]
Consequently, exact integrality on global integers is equivalent to
$\gcd(m_i,D)=1$ for every $i$.
\end{proposition}

\begin{proof}
Admissibility gives $H_i(i),H_i(i+D)\in\mathbb Z$.  Subtracting gives
$Da_i\in\mathbb Z$, and then
\[
 Db_i=D H_i(i)-(Da_i)i\in\mathbb Z.
\]
Thus $m_i$ and $s_i$ are integral, while
\[
 m_i i-s_i=D(a_i i+b_i)=D H_i(i)
\]
proves the congruence.  Conversely, that congruence makes
$(m_i n-s_i)/D$ integral on $i+D\mathbb Z$, and the two displayed
assignments are literal inverses.

For arbitrary $n\in\mathbb Z$, branch integrality is equivalent to
$m_i(n-i)\equiv0\pmod D$.  Writing $g_i=\gcd(m_i,D)$, this is equivalent to
$n-i\in(D/g_i)\mathbb Z$, which proves the locus and the final assertion.
\end{proof}

This proposition applies only to scalar branches.  A genuinely non-scalar
$\mathbb F$-linear Hydra is not silently placed in the Matthews or rcwa
classes.  For the scalar subclass, the further reduced rcwa crosswalk is
Proposition~\ref{prop:Matthews-Kohl-crosswalk}; neither step minimizes the
retained modulus.

\begin{sourceissue}[uniform separation is not bounded finiteness]
\label{issue:Siegel-discrete-orbit-classification}
Source Proposition~4.2 claims that if $Y$ is uniformly separated in an
arbitrary valued field, every orbit of every map $T:Y\to Y$ is preperiodic or
norm-divergent.  Its proof uses the stronger assertion that every bounded
intersection with $Y$ is finite.

For a nontrivially valued counterexample, take $K=\mathbb Q((t))$ with the
$t$-adic absolute value, $Y=\mathbb Z\subset K$, and $T(n)=n+1$.  Distinct
integers have distance one, but the orbit of zero is injective and has norm
one after the first term.  It is neither preperiodic nor divergent.

The proof becomes valid under the exact replacement
\[
 \#\{y\in Y:|y|\leq R\}<\infty
 \qquad\text{for every }R>0.
\tag{BF}
\]
For a number field, $\mathcal O_K$ is boundedly finite in its full Minkowski
embedding.  It need not be discrete at one archimedean embedding:
$\mathbb Z[\sqrt2]$ has nonzero elements arbitrarily close to zero there.
Nor does escape in the full Minkowski norm imply escape at every individual
archimedean absolute value.
\end{sourceissue}

For the numen, let $\Sigma_p^*$ be the finite words in
$\{0,\ldots,p-1\}$, including the empty word.  The source fixes
least-significant digit first:
\[
 \operatorname{DigSum}_p(j_1,\ldots,j_N)
 =\sum_{a=1}^{N}j_a p^{a-1}.
\]

\begin{proposition}[finite affine word and numen cocycle]
\label{prop:Siegel-word-numen}
For $\mathbf j=(j_1,\ldots,j_N)$, define
\[
 H_{\mathbf j}=H_{j_1}\circ\cdots\circ H_{j_N},
 \qquad
 H_{\mathbf j}(z)=M_{\mathbf j}(z)+X_{\mathbf j}.
\]
Then
\[
 M_{\mathbf j}=r_{j_1}\circ\cdots\circ r_{j_N},
 \qquad
 X_{\mathbf j}
 =\sum_{m=0}^{N-1}
   (r_{j_1}\circ\cdots\circ r_{j_m})(c_{j_{m+1}}),
\tag{W}
\label{eq:Siegel-word-formula}
\]
with the empty operator product equal to $I$.  For concatenation,
\[
 M_{\mathbf i\wedge\mathbf j}
 =M_{\mathbf i}\circ M_{\mathbf j},
 \qquad
 X_{\mathbf i\wedge\mathbf j}
 =M_{\mathbf i}(X_{\mathbf j})+X_{\mathbf i}.
\tag{C}
\label{eq:Siegel-word-cocycle}
\]
\end{proposition}

\begin{proof}
The rightmost branch in $H_{j_1}\circ\cdots\circ H_{j_N}$ acts first.
Composing two affine maps gives
\[
 (M_i,X_i)\circ(M_j,X_j)
 =(M_i\circ M_j,M_i(X_j)+X_i).
\]
Induction proves both~\eqref{eq:Siegel-word-formula} and
\eqref{eq:Siegel-word-cocycle}; the empty word gives $(I,0)$.
\end{proof}

Here and below a word is an address.  It is not automatically an actual
piecewise-Hydra itinerary.

\begin{proposition}[centered numen descent and multiplier non-descent]
\label{prop:Siegel-centered-numen-descent}
Let $\mathbf i\sim\mathbf j$ when their DigSum values agree.  If
$c_0=H_0(0)=0$, then $X_{\mathbf i}=X_{\mathbf j}$ whenever
$\mathbf i\sim\mathbf j$, so the finite numen factors as
\[
 \Sigma_p^*/{\sim}\ \xrightarrow[\cong]{\operatorname{DigSum}_p}\mathbb N_0
 \xrightarrow{X_H}K.
\]
The multiplier does not generally factor through this quotient:
\[
 M_{\mathbf j\wedge(0)}=M_{\mathbf j}\circ r_0
\]
need not equal $M_{\mathbf j}$.
\end{proposition}

\begin{proof}
Every finite DigSum fibre consists of the canonical base-$p$ word and its
finitely many terminal-zero extensions.  Since the appended zero is the
rightmost map in the composition,
\[
 X_{\mathbf j\wedge(0)}
 =H_{\mathbf j}(H_0(0))=H_{\mathbf j}(0)=X_{\mathbf j}.
\]
Iteration proves descent.  The displayed multiplier identity follows directly
from Proposition~\ref{prop:Siegel-word-numen} and supplies the obstruction.
\end{proof}

Thus every occurrence of $M_H(n)$ requires a canonical digit word and its
retained length.  The source proves descent only for $X_H$.

\begin{proposition}[finite digit functional equations]
\label{prop:Siegel-numen-functional-equations}
Functions $f:\mathbb N_0\to K$ satisfying
\begin{equation}
\label{eq:Siegel-finite-recursion}
 f(pn+j)=r_j(f(n))+c_j
 \qquad(n\in\mathbb N_0,\ 0\leq j<p)
\end{equation}
are in bijection with the affine fibre
\[
 \{a\in K:(I-r_0)a=c_0\},
\]
by $a=f(0)$.  If $I-r_0$ is invertible, there is one solution, with
$a=(I-r_0)^{-1}c_0$.  If
$n=\sum_{m=0}^{N-1}d_m p^m$ is its canonical digit expansion and
\[
 P_0=I,\qquad P_m=r_{d_0}\circ\cdots\circ r_{d_{m-1}},
\]
then
\begin{equation}
\label{eq:Siegel-corrected-finite-truncation}
 f(n)=P_N(a)+\sum_{m=0}^{N-1}P_m(c_{d_m}).
\end{equation}
\end{proposition}

\begin{proof}
Putting $n=j=0$ in~\eqref{eq:Siegel-finite-recursion} gives
$(I-r_0)f(0)=c_0$.  Conversely, once an admissible $a$ is chosen, repeated
Euclidean division by $p$ reaches zero and recursively defines one value at
every nonnegative integer.  Applying the recursion along the canonical digits
proves~\eqref{eq:Siegel-corrected-finite-truncation}.
\end{proof}

The finite theorem is sound.  It also isolates the term omitted in the
source's proof of its continuation lemma.

\begin{sourceissue}[Lemma 6.1 as printed]
\label{issue:Siegel-numen-extension}
Lemma~6.1 assumes only that $I-r_0$ is invertible, but its finite formula uses
the centered word sum and omits $P_N(a)$ from
\eqref{eq:Siegel-corrected-finite-truncation}.  Its equation~(6.29) prints
equality between the norm of a sum and the sum of termwise norm bounds; only
$\leq$ follows.  Its pointwise density definition uses a limsup frequency
when $\log\lVert r_j\rVert_\ell\leq0$ and a liminf when that logarithm is
positive.  These are the reverse of the choices needed for an upper Lyapunov
bound.  Almost-everywhere convergence leaves a null nonconvergence set on
which no value is assigned.  The proof establishes uniqueness on
$\mathbb N_0$, not uniqueness of arbitrary functions on $\mathbb Z_p$
satisfying the same recursion.  Finally, in the nonarchimedean case the
displayed bound omits $|a|_\ell$.

The last omission has a concrete instance.  Over $\mathbb Q_2$, take
$r_0=3$ and $c_0=1$.  Then
$a=(1-3)^{-1}=-1/2$, so $|a|_2=2>|c_0|_2=1$.
\end{sourceissue}

\begin{theorem}[corrected numen continuation]
\label{thm:Siegel-numen-extension-repair}
Fix a completion $K_\ell$ at which all $r_j$ are bounded, assume
$I-r_0$ is invertible, and put $a=(I-r_0)^{-1}c_0$.  For
\[
 \mathfrak z=\sum_{m\geq0}d_m p^m\in\mathbb Z_p,
\]
define
\begin{equation}
\label{eq:Siegel-corrected-continuation}
 f_N(\mathfrak z)
 =P_N(a)+\sum_{m=0}^{N-1}P_m(c_{d_m}),
 \qquad
 P_m=r_{d_0}\circ\cdots\circ r_{d_{m-1}}.
\end{equation}
For $q_{j,m}=m^{-1}\#\{0\leq k<m:d_k=j\}$ and
$\alpha_j=\log\lVert r_j\rVert_\ell$, set
\begin{equation}
\label{eq:Siegel-upper-Lyapunov}
 L^+(\mathfrak z)=
 \sum_{\alpha_j<0}\alpha_j\liminf_{m\to\infty}q_{j,m}
 +\sum_{\alpha_j>0}\alpha_j\limsup_{m\to\infty}q_{j,m}.
\end{equation}
Then:
\begin{enumerate}[label=\textnormal{(\roman*)}]
\item if $L^+(\mathfrak z)<0$, the sequence
  \eqref{eq:Siegel-corrected-continuation} converges;
\item if $\prod_{j=0}^{p-1}\lVert r_j\rVert_\ell<1$, it converges for
  Haar-almost every $\mathfrak z$ and defines a measurable function on its
  conull convergence domain;
\item if $\rho=\max_j\lVert r_j\rVert_\ell<1$, it converges uniformly on
  $\mathbb Z_p$ to a continuous function;
\item whenever $\mathfrak z$ lies in the convergence domain,
  $p\mathfrak z+j$ also lies there and
  \[
   f(p\mathfrak z+j)=r_j(f(\mathfrak z))+c_j;
  \]
\item if $\ell$ is nonarchimedean and
  $\max_j\lVert r_j\rVert_\ell\leq1$, every finite truncation satisfies
  \[
   |f_N(\mathfrak z)|_\ell
   \leq\max\{|a|_\ell,\max_j|c_j|_\ell\}.
  \]
\end{enumerate}
In case~(iii), the continuous extension is the unique continuous function
with the digit recursion and initial value $a$.
\end{theorem}

\begin{proof}
Submultiplicativity gives
\[
 \lVert P_m\rVert_\ell
 \leq\prod_{j=0}^{p-1}
      \lVert r_j\rVert_\ell^{m q_{j,m}}.
\]
For a negative $\alpha_j$, the limsup of
$\alpha_jq_{j,m}$ is $\alpha_j\liminf q_{j,m}$; for a positive
$\alpha_j$, it is $\alpha_j\limsup q_{j,m}$.  Since there are finitely many
digits,
\[
 \limsup_{m\to\infty}\frac1m\log\lVert P_m\rVert_\ell
 \leq L^+(\mathfrak z).
\]
Under~(i), there are $\eta<1$ and $m_0$ such that
$\lVert P_m\rVert_\ell\leq\eta^m$ for $m\geq m_0$.  The translations form a
finite set, so the series in
\eqref{eq:Siegel-corrected-continuation} is absolutely convergent and
$P_N(a)\to0$.

Haar-almost every base-$p$ digit sequence has each digit frequency $1/p$.
Its Lyapunov exponent is bounded above by
\[
 \frac1p\sum_j\log\lVert r_j\rVert_\ell
 =\frac1p\log\prod_j\lVert r_j\rVert_\ell,
\]
which proves~(ii); a pointwise limit of measurable truncations is measurable
on the convergence domain.  Under~(iii),
$\lVert P_m\rVert_\ell\leq\rho^m$ for every digit sequence, and the geometric
tail estimate is uniform.  The truncations are locally constant, hence their
uniform limit is continuous.

The finite identity
\[
 f_N(p\mathfrak z+j)=r_j(f_{N-1}(\mathfrak z))+c_j
\]
proves~(iv) on passage to the limit.  In the nonarchimedean case, the
ultrametric inequality and $\lVert P_m\rVert_\ell\leq1$ give~(v).
Finally, $\mathbb N_0$ is dense in $\mathbb Z_p$, and
Proposition~\ref{prop:Siegel-numen-functional-equations} fixes all its values;
two continuous extensions therefore agree everywhere.
\end{proof}

Functional equations alone do not give literal uniqueness among arbitrary
functions on $\mathbb Z_p$.  For example, with $p=2$,
$r_0=r_1=1/2$, and $c_0=c_1=0$, choose a non-eventually-periodic digit
sequence.  Its tail-equivalence class is countable and contains no shift
cycle.  A nonzero seed value propagates consistently over that class by the
invertible branch factors, while the function may be set to zero off the
class.  The result is Borel and satisfies the recursion, but is not the zero
function.

\begin{proposition}[operator-typed periodic word]
\label{prop:Siegel-periodic-word-fixed-point}
For a nonempty finite word $\mathbf w$, put
$M=M_{\mathbf w}$ and $X=X_{\mathbf w}$.  Then
\begin{equation}
\label{eq:Siegel-repeated-word}
 X_{\mathbf w^{\wedge m}}
 =\sum_{k=0}^{m-1}M^k(X).
\end{equation}
If $\lVert M\rVert_\ell<1$, then $I-M$ is invertible on $K_\ell$ and
\[
 z_{\mathbf w}=(I-M)^{-1}X
\]
is the unique fixed point of the affine extension $H_{\mathbf w}$.
It is a point of an actual $N$-step piecewise-Hydra cycle only if, for
$\mathbf w=(w_1,\ldots,w_N)$,
\begin{equation}
\label{eq:Siegel-branch-correctness}
\begin{split}
 z_{\mathbf w}&\in w_N+\Lambda,\\
 H_{w_N}(z_{\mathbf w})&\in w_{N-1}+\Lambda,\\
 &\ \vdots\\
 H_{w_2}\circ\cdots\circ H_{w_N}(z_{\mathbf w})
   &\in w_1+\Lambda.
\end{split}
\end{equation}
\end{proposition}

\begin{proof}
The cocycle gives
$X_{\mathbf w^{\wedge m}}=M(X_{\mathbf w^{\wedge(m-1)}})+X$;
induction proves~\eqref{eq:Siegel-repeated-word}.  The Neumann series
$\sum_{k\geq0}M^k$ converges in operator norm and is the inverse of $I-M$.
Thus $M(z_{\mathbf w})+X=z_{\mathbf w}$, and subtracting two fixed-point
equations proves uniqueness.  Since the rightmost branch acts first,
conditions~\eqref{eq:Siegel-branch-correctness} are exactly the successive
branch-selection requirements.  Without them, the calculation concerns only
the affine extension.
\end{proof}

The source's fractions $(1-M^m)/(1-M)$ and $X/(1-M)$ must be replaced by
operator composition with $(I-M)^{-1}$.  Its equality
$\lVert M^m\rVert=\lVert M\rVert^m$ is false for a general operator norm;
the sufficient relation is
$\lVert M^m\rVert\leq\lVert M\rVert^m$.

\begin{sourceissue}[Siegel Theorem 6.1 as stated]
\label{thm:Siegel-correspondence-registered}
The source assumes a proper, centered, integral $\Lambda$-Hydra and the two
conditions
\begin{enumerate}
\item for every $n\geq1$, some place $\ell_n$ satisfies
  $\lVert M_H(n)\rVert_{\ell_n}<1$;
\item some nontrivial nonarchimedean place $v$ satisfies
  $\lVert r_j\rVert_v>1$ for every digit $j$ and $v(\Lambda)\geq1$.
\end{enumerate}
It then states:
\begin{enumerate}
\item an element $z\in\mathcal O_K$ is periodic for $H$ if and only if
  $z=X_H(\mathfrak z)$ for some
  $\mathfrak z\in(\mathbb Z_p\setminus\mathbb N_0)\cap\mathbb Q$;
\item if $K$ is a number field and
  $X_H(\mathfrak z)\in\mathcal O_K$ for an irrational
  $\mathfrak z\in\mathbb Z_p\setminus\mathbb Q$, then the source calls
  $X_H(\mathfrak z)$ a trajectory divergent with respect to every
  nontrivial archimedean absolute value.
\end{enumerate}
This is a source-statement record, not an admitted theorem of this edition.
\end{sourceissue}

The proof sketch does not certify either conclusion.  The exact blocking
points are as follows.
\begin{enumerate}
\item The literal Hydra and integral hypotheses have no simultaneous models,
  by Source Issue~\ref{issue:Siegel-Hydra-integrality}.
\item $M_H$ does not descend through finite DigSum equivalence.  Thus
  $M_H(n)$ is not defined until a canonical word and its retained length are
  specified.  The displayed $\lambda_p(n)$ is likewise not defined.
\item Two formulas use an undefined $M_q$.  Scalar fractions such as
  $(1-M^m)/(1-M)$ are not typed for general endomorphisms, and the equality
  $\lVert M^m\rVert=\lVert M\rVert^m$ is false.  The operator-typed,
  submultiplicative replacement is
  Proposition~\ref{prop:Siegel-periodic-word-fixed-point}.
\item Passing from an arbitrary finite word to the canonical digit word of
  its DigSum may remove terminal zeroes.  Repeating those two words need not
  give the same multiplier or affine map.  The centered all-zero word and
  the fixed point zero also lie outside the displayed $n\geq1$ construction.
\item The notation $v(\Lambda)$, the \emph{correctness} predicate, and the
  claimed extension of the branches to $\mathcal O_{K_v}$ are not defined or
  proved in this source.  A formal affine fixed point is not an actual
  piecewise orbit until every condition in
  \eqref{eq:Siegel-branch-correctness} is established.
\item The converse uses integrality backwards: from an integral later value
  it infers that a possibly nonintegral predecessor is integral.  The
  printed global predicate does not imply the completion-level inverse-image
  identity needed for this step, nor is a local-to-global arrow supplied.
\item The forward direction gives, at most, an existential rational address
  for a periodic point.  It does not prove that the same numen fibre contains
  no irrational address.  Consequently it cannot support the later inference
  that an irrational address has a nonpreperiodic value.  The zero fibre is
  also unclassified.
\item Part II begins by invoking a nonexistent hypothesis~(iii), then imports
  a zero-fibre assertion and a correctness result without rebuilding their
  hypotheses.  It also depends on the false uniformly-discrete orbit
  classification repaired in
  Source Issue~\ref{issue:Siegel-discrete-orbit-classification}.
\item A numen value is a point, not a trajectory.  Even after replacing it
  by its forward orbit, escape in the full Minkowski embedding would not by
  itself prove divergence at every individual archimedean place.
\end{enumerate}

Accordingly, the periodic-word fixed-point proposition is retained at its
proved scope, while the source's rational-fibre correspondence and
irrational-address divergence conclusions are not proved and are not used.
The legacy audit identifier \texttt{PO-COL-000006} records that gap; it is not
a mathematical problem or a commitment attributed to the source or reader.

\begin{sourceissue}[the v3 nonarchimedean Fourier inversion]
\label{issue:Siegel-v3-prop76}
Let $L=K_\ell$ be a finite extension of $\mathbb Q_q$, let
$\mathcal O_L$ be its integer ring, let $\pi$ be a uniformizer, and let
$\mathfrak D=\mathfrak D_{L/\mathbb Q_q}$ be its different.  The source uses
the trace character
\[
 e_\ell(x)=
 \exp\!\left(2\pi i
 \{\operatorname{Tr}_{L/\mathbb Q_q}(x)\}_q\right).
\]
The annihilator of $\mathcal O_L$ under the pairing
$(t,x)\mapsto e_\ell(tx)$ is
\[
 \mathfrak D^{-1}
 =\{t\in L:\operatorname{Tr}_{L/\mathbb Q_q}(t\mathcal O_L)
                    \subseteq\mathbb Z_q\},
\]
not $\mathcal O_L$ in general.  Thus the source's asserted dual
$L/\mathcal O_L$ is correct for this character only after a
conductor normalization, for example when the relevant different is
trivial.  With the printed character the dual is
$L/\mathfrak D^{-1}$.

For comparison, if an $L$-valued random variable $X$ has first been proved
to lie in $\mathcal O_L$ almost surely, define
\[
 \widehat\mu(t)=\mathbb E[e_\ell(-tX)]
 \qquad(t\in L/\mathfrak D^{-1}).
\]
Writing $e$ and $f=[L:\mathbb Q_q]/e$ for the ramification index and residue
degree, finite character orthogonality gives the correctly typed identity
\begin{equation}
\label{eq:Siegel-corrected-local-Fourier}
\Pr(X\in w+\pi^n\mathcal O_L)
=q^{-nf}
 \sum_{t\in\pi^{-n}\mathfrak D^{-1}/\mathfrak D^{-1}}
   \widehat\mu(t)e_\ell(tw)
\qquad(w\in\mathcal O_L,\ n\geq0).
\end{equation}
Indeed, the displayed frequency quotient is the complete character group of
$\mathcal O_L/\pi^n\mathcal O_L$, which has $q^{nf}$ elements.  Averaging
those characters is one on the zero residue class and zero on every other
class, proving~\eqref{eq:Siegel-corrected-local-Fourier}.

This correction does not certify source Proposition~7.6.  Its surrounding
Fourier chain has further independent defects:
\begin{enumerate}
\item for a nonscalar bounded local branch $r_j$, the frequency map is the
  trace adjoint $r_j^\dagger$, defined by
  \[
   \operatorname{Tr}_{L/\mathbb Q_q}(t\,r_jx)
   =\operatorname{Tr}_{L/\mathbb Q_q}(r_j^\dagger t\,x),
  \]
  not $t\mapsto r_jt$ unless the required self-adjoint or scalar identity is
  proved;
\item the characteristic function is naturally defined first on $L$; its
  descent to a quotient requires an almost-sure support statement matched to
  the annihilator;
\item source Proposition~7.5 multiplies values in an object assumed only to
  be an abelian group, shadows its summation variable, and leaves another
  variable unbound, so it cannot justify the later rearrangement;
\item $B_{H,\ell}=e\max_j\log_q|c_j|_\ell$ is undefined when every
  translation is zero and may be negative, although Proposition~7.5 is
  invoked only for a nonnegative integer parameter;
\item the correct size relation is
  $q^{B/e}=|\pi^{-B}|_\ell$, not $|\pi^B|_\ell$; a real absolute value cannot
  be compared to the field element $\pi^B$; and division by $\pi^B$ cannot
  be replaced by division by the real number $q^{B/e}$.
\end{enumerate}
The visible v3 changes repair neither the dual lattice nor these typing and
scaling defects.  No proposition depending on this Fourier chain is used
downstream.  The legacy audit identifier
\texttt{PO-COL-}\allowbreak\texttt{000007} records the resulting dependency
gap; it does not convert the defective chain into an open mathematical problem.
\end{sourceissue}

\begin{nonclaim}
The finite numen is an address-to-state construction.  It is not identified
with a state-to-itinerary conjugacy, and its fibres are not classified.
Neither the corrected continuation theorem nor the formal periodic-word
fixed point proves an ordinary-integer Collatz endpoint.  The source's
Correspondence Principle and Fourier consequences are recorded as open
audits, not imported conclusions.
\end{nonclaim}


\subsection{Current almost-bounded-orbit theorem: registered boundary}

For $A\subset\Npos$, Tao's logarithmic density is the limit, when it exists,
of
\[
 \frac{\displaystyle\sum_{n\in A\cap[1,x]}1/n}
      {\displaystyle\sum_{n\in\Npos\cap[1,x]}1/n}.
\]
The current source is the version-pinned arXiv source $1909.03562v7$, not the
older routed local witness.

\begin{theorem}[Tao, source theorem]
\label{thm:Tao-main-registered}
Let $f:\Npos\to\mathbb R$ satisfy $f(N)\to+\infty$.  With $U$ as in
\eqref{eq:unshortened-U} and
\[
 U_{\min}(N)=\inf_{j\in\Nzero}U^j(N),
\]
one has $U_{\min}(N)<f(N)$ on a subset of $\Npos$ of logarithmic density one.
\end{theorem}

This is Theorem~1.3 of the source
\cite[Sec.~1.1, Thm.~1.3]{Tao2026v7}.  Its quantifier over every divergent
$f$ and its logarithmic-density conclusion are retained exactly; it is not a
natural-density theorem and not a bounded-orbit theorem with a universal
constant.  At the checkpoint at which this theorem was first registered, its
proof-dependency audit was still open.  The exact repaired source chain
through Proposition~1.11, Theorem~3.1, Theorem~1.6, and this theorem is now
closed in Subsection~\ref{subsec:Tao-main-reduction}; no local claim may use
more than the displayed conclusion.

\begin{proposition}[exact Siegel--Tao Syracuse coordinate]
\label{prop:Siegel-Tao-Syracuse-coordinate}
Let
\[
 \Omega=(\Npos)^{\Npos},\qquad
 \mathbb P_{\rm geom}(a_k=r)=2^{-r},
\]
with independent coordinates, and put
\(A_k=a_1+\cdots+a_k\).  The map
\begin{equation}
\label{eq:Siegel-Tao-run-coordinate}
 \mathcal R:\Omega\longrightarrow\Ztwo,
 \qquad
 \mathcal R((a_k)_{k\geq1})=
       \sum_{k\geq1}2^{A_k-1},
\end{equation}
is a bijection onto the set \(\Ztwo^{\infty1}\) of $2$-adic digit strings
having infinitely many one digits.  Its inverse sends the increasing one
positions \(0\leq n_1<n_2<\cdots\) to
\begin{equation}
\label{eq:Siegel-Tao-run-inverse}
 a_1=n_1+1,\qquad a_k=n_k-n_{k-1}\quad(k\geq2).
\end{equation}
It pushes \(\mathbb P_{\rm geom}\) to normalized Haar measure on \(\Ztwo\);
the complement of \(\Ztwo^{\infty1}\) is Haar-null.

For Siegel's numen \(\chi_3:\Ztwo\to\mathbb Z_3\), one has the pointwise
identity
\begin{equation}
\label{eq:Siegel-Tao-pointwise-Syracuse}
 \chi_3(\mathcal R((a_k)))
   =\sum_{k\geq1}3^{k-1}2^{-A_k}
   \quad\hbox{in }\mathbb Z_3.
\end{equation}
For \(n\geq1\), define the pointwise realization
\[
 \operatorname{Syrac}_n(a_1,\ldots,a_n)
 :=F_n(a_n,\ldots,a_1)\pmod{3^n}.
\]
Consequently,
\begin{equation}
\label{eq:Siegel-Tao-finite-Syracuse}
 \chi_3(\mathcal R((a_k)))\pmod{3^n}
 =\operatorname{Syrac}_n(a_1,\ldots,a_n)
 =\sum_{k=1}^{n}3^{k-1}2^{-A_k}\pmod{3^n},
\end{equation}
and \(\operatorname{Syrac}_n(a_1,\ldots,a_n)\), under
\(\mathbb P_{\rm geom}\), has exactly Tao's law
\(\operatorname{Syrac}(\mathbb Z/3^n\mathbb Z)\).  Thus the equality is
pointwise for this explicitly declared coupling, rather than an unnamed
equality in distribution.
\end{proposition}

\begin{proof}
The partial sums in \eqref{eq:Siegel-Tao-run-coordinate} converge in
\(\Ztwo\), because \(A_k\to\infty\).  Their one positions are exactly
\(A_k-1\), which proves \eqref{eq:Siegel-Tao-run-inverse}, both inverse
identities, and singleton fibres.

Fix \((a_1,\ldots,a_k)\).  Its run-length cylinder has product mass
\[
 \prod_{j=1}^k2^{-a_j}=2^{-A_k}.
\]
Its image fixes every binary digit from position zero through position
\(A_k-1\), with ones exactly at \(A_j-1\); that Haar cylinder also has mass
\(2^{-A_k}\).  These cylinders generate the product sigma-algebra.  Hence
\(\mathcal R_*\mathbb P_{\rm geom}\) is Haar measure.  A $2$-adic string
with finitely many ones is a nonnegative integer, so the omitted set is
countable and Haar-null.

Siegel's historical v1 one-position formula is
\[
 \chi_p\!\left(\sum_{k\geq1}2^{n_k}\right)
 =\sum_{k\geq1}\frac{p^{k-1}}{2^{n_k+1}}.
\]
The current v3 source obtains the same formula from its finite word cocycle
and $q$-adic rising limit
\cite[historical v1 source lines 405--568]{Siegel2007v1}
\cite[current v3 source lines 1851--2000]{Siegel2007v3}.
Substitution of \(p=3\) and \(n_k=A_k-1\) proves
\eqref{eq:Siegel-Tao-pointwise-Syracuse}.  Tao's source formula is the same
series with \(A_k=a_{[1,k]}\)
\cite[source lines 384--393, 436--440]{Tao2026v7}.  Modulo \(3^n\), every
term with \(k>n\) vanishes.  Expanding the source definition of \(F_n\)
after the displayed reversal \((a_1,\ldots,a_n)\mapsto(a_n,\ldots,a_1)\)
gives the middle expression in
\eqref{eq:Siegel-Tao-finite-Syracuse}.  Reversal preserves the iid product
law, which proves the final law statement without identifying two
deterministic functions that the source defines only in distribution.
\end{proof}

This proves, with sample space, measure, coordinate order, inverse, fibres,
and projections exposed, the equality of laws stated in Siegel v3 source
lines~3180--3184.  It does not identify the shortened and accelerated maps or
their clocks.  The same indexed Siegel source contains no first-passage
interval, Proposition~5.2 event chain, or Section~5 function \(c_n\); it is a
coordinate source for the later Fourier layer, not a repair of that argument.
The exact version and source audit is
\path{qa/SIEGEL-2007-V1-V3-F023-syracuse-crosswalk-audit.md}; the finite exact
certificate is
\path{certificates/siegel_v1_tao_syracuse_crosswalk_checks.py}.

\begin{sourceissue}[current Siegel v3 Fourier presentation]
\label{sourceissue:Siegel-2007-v3-Fourier}
The targeted v3 passage has three forced local defects that remain quarantined
from later use.  Source line~3146 writes \(\widehat{\mathbb Z}_2\) where the
denominators and the declared characteristic-function domain require
\(\widehat{\mathbb Z}_3\).  With the source convention
\(\varphi_3(t)=\mathbb E e^{-2\pi i\{tX\}_3}\), line~3144's inversion
exponential requires the opposite sign; its printed sign reconstructs the
mass at \(-m\).  Already modulo three it swaps the masses $1/3$ at $1$ and
$2/3$ at $2$.  Finally, line~3164's last characteristic function omits its
subscript~$3$.  None affects Proposition~\ref{prop:Siegel-Tao-Syracuse-coordinate},
which was proved directly from the coordinate series.
\end{sourceissue}

\begin{sourceissue}[current v7 proof-text audit]
\label{sourceissue:Tao-v7-audit}
A line-level audit of the pinned v7 source records several forced local
corrections: a function/value omission and the printed $N-1$ in place of the
induction index $n-1$ in
Lemma~2.1; a scalar/vector mismatch and an $n^{-1/2}$ factor where a
$d$-dimensional Gaussian integral gives $n^{-d/2}$ in Lemma~2.2; undefined
or inconsistent indices in the proofs of Proposition~1.9 and
Proposition~5.2; a first-passage location written where a first-passage time
is required; an unsigned $\exp(O(\log^{0.6}x))$ used as a lower bound; the
malformed expression $q^{-1}\log O(M_1/M_0)$ where the preceding calculation
gives $q^{-1}(1+\log(M_1/M_0))$; and malformed formulas in the black-set and
renewal arguments.
The v7 text does correct an older time/location interchange at the end of
Proposition~5.2 and an older complemented-event mismatch in Section~6.
The latter is forced: the factorized Fourier coefficient contains
\(1_{E_k\cap B_k\cap C_{k,l}}\), so v5's replacement by
\(1_{\overline E_k\cap B_k\cap C_{k,l}}\) has no equality or domination
behind it.

The complete v5--v7 comparison has 28 non-whitespace hunks.  The journal
agrees with v7 on 15 distinguishing hunks, with v5 on 12, and is hybrid on
the remaining Lemma~7.9 propagation hunk; it is not a uniform intermediate
version.  In particular, its equations~(6.7)--(6.8) retain the v5 reserve,
while equation~(6.11) contains the v7 event correction.

The reserve change is substantive.  With
\(a=\log(4/3)\) and \(b=\log2\), the optimized fluctuation cost is
\[
 \frac{b^2}{4a}C_A^2\log n
 =0.4175208154\ldots C_A^2\log n.
\]
The v5/journal half-window retains only
\(\frac12bC_A^2\log n=0.3465735903\ldots C_A^2\log n\), and its printed
conversion also drops the factor \(b\); an explicit admissible family makes
one claimed sub-\(3^n\) summand exceed \(3^n\).  V7's
\(0.99bC_A^2\log n\) reserve closes with exact remaining coefficient
\(0.2686948933\ldots\).  The stopping-time derivation and Archimedean
separation are proved in
Source Issue~\ref{sourceissue:Tao-v5-v7-section-six-reserve} and
Proposition~\ref{prop:Tao-v7-Fourier-to-mixing}.  This identifies a failure
of the printed v5 argument, not a proof that its stated corollary is false
by every method.

Two issues are not merely typographical.  Lines~758--763 try to infer
$T_x(\mathbf N_y)\in I_y$ from an additive endpoint restriction
$\mathbf N_y\in[y+2\log^{0.8}x,y^\alpha-2\log^{0.8}x]$.  Since the preceding
time estimate depends on $\log(\mathbf N_y/x)$, this additive displacement
does not supply a margin of order $\log^{0.8}x$.  The following proposition
proves a replacement and fixes its exact downstream boundary.  In Lemma~5.3,
lines~909--914 replace summable terms by
$2^{-(a_1+\cdots+a_{m_0})/2}$ and then call the resulting product uniformly
bounded.  Its exact sum is $(1+\sqrt2)^{m_0}$, which grows with $x$; a second
proposition below retains the stronger pre-weakening estimates and proves the
uniform bound.  Separately,
the Lemma~7.9
Markov event must be intersected with $\{R\leq\mathbf r\}$; the next line
works under exactly that restriction.  The complete v7 renewal dependency is
reconstructed and certified at its declared scope in
Corollary~\ref{cor:Tao-v7-fine-scale-branch}, while the stronger v5/journal
form remains outside that proof.  These findings neither enlarge the theorem
nor constitute a claim that its conclusion is false.

At the final proof step, v7 line~1851 and journal p.~54 name
\(E_{p,4^A(1+p^3)}\), whereas \(E_*\) was defined using
\(E_{p,4^A(1+p)^3}\) and the ensuing inequality again uses
\(4^A(1+p)^3\).  The edition retains the definitionally supported latter
event and does not assert an identity between the two parameters.  The
terminal boundary, all 28 source hunks, the journal matrix, and the active
citation census are recorded in
\texttt{qa/TAO-V5-V7-JOURNAL-F028-whole-version-citation-audit.md}.
\end{sourceissue}

\begin{proposition}[Multiplicative endpoint repair for v7 Proposition~5.2]
\label{prop:Tao-v7-endpoint-buffer}
Retain the source constants
\[
 \alpha=1.001,\qquad \lambda=\log(4/3),\qquad L=\log x,
\]
let $y\in\{x^\alpha,x^{\alpha^2}\}$, and let $\mathbf N_y$ have the source's
logarithmic distribution on the odd integers in $[y,y^\alpha]$.  Let
\[
 \mathcal G_x=
 \left\{\vec a^{(n_0)}(\mathbf N_y)\in
                  \mathcal A^{(n_0)}\right\}
\]
be the good valuation event in source lines~733--740, and retain the source
interval
\[
 I_y=\left[
 \frac{\log(y/x)}{\lambda}+L^{0.8},
 \frac{\log(y^\alpha/x)}{\lambda}-L^{0.8}
 \right].
\]
For all sufficiently large $x$, the interval
\[
 R_y=\left[y\exp(L^{0.8}),
            y^\alpha\exp(-L^{0.8})\right]
\]
is nonempty and
\[
 \mathcal G_x\cap\{\mathbf N_y\in R_y\}
 \subseteq\{T_x(\mathbf N_y)\in I_y\}.
\]
Uniformly for the two displayed values of $y$,
\[
 \mathbb P(\mathbf N_y\notin R_y)=O(L^{-1/5})
 \quad\hbox{and hence}\quad
 \mathbb P(T_x(\mathbf N_y)\in I_y)=1-O(L^{-1/5}).
\]
\end{proposition}

\begin{proof}
On $\mathcal G_x$, source lines~741--752 give, with one absolute constant
$K_0$ uniform in the two choices of $y$,
\begin{equation}
 \label{eq:Tao-good-orbit-bound}
 \left|
  \log\frac{\operatorname{Syr}^{j}(\mathbf N_y)}{\mathbf N_y}
  +j\lambda
 \right|\leq K_0L^{0.6}
 \qquad(0\leq j\leq n_0).
\end{equation}
This estimate reaches the threshold before $n_0$.  Indeed,
$\mathbf N_y\leq x^{\alpha^3}$ and
\[
 \log\frac{\operatorname{Syr}^{n_0}(\mathbf N_y)}x
 \leq (\alpha^3-1)L-\lambda n_0+K_0L^{0.6}.
\]
Here $\alpha^3-1<1/300$, whereas
\[
 \frac{\lambda}{10\log2}>\frac1{40}:
 \qquad
 \log(4/3)>\frac{1/3}{1+1/3}=\frac14,
 \quad \log2<1.
\]
Since $n_0=\lfloor L/(10\log2)\rfloor$, the preceding upper bound is
negative for all sufficiently large $x$.  Thus $1\leq T_x(\mathbf N_y)\leq
n_0$ on $\mathcal G_x$.  Apply
\eqref{eq:Tao-good-orbit-bound} at $j=T_x(\mathbf N_y)$ and at
$j=T_x(\mathbf N_y)-1$.  The defining inequalities
\[
 \operatorname{Syr}^{T_x(\mathbf N_y)}(\mathbf N_y)\leq x
 <\operatorname{Syr}^{T_x(\mathbf N_y)-1}(\mathbf N_y)
\]
then give an absolute $K\geq K_0$ such that
\begin{equation}
 \label{eq:Tao-first-time-two-sided}
 \left|T_x(\mathbf N_y)-
       \frac{\log(\mathbf N_y/x)}{\lambda}\right|
 \leq KL^{0.6}.
\end{equation}

Since $\lambda<1$, the constant $\delta=1/\lambda-1$ is positive.  If
$\mathbf N_y\in R_y$, then
\begin{align*}
 \frac{\log(\mathbf N_y/x)}\lambda
 -\left(\frac{\log(y/x)}\lambda+L^{0.8}\right)
 &\geq \delta L^{0.8},\\
 \left(\frac{\log(y^\alpha/x)}\lambda-L^{0.8}\right)
 -\frac{\log(\mathbf N_y/x)}\lambda
 &\geq \delta L^{0.8}.
\end{align*}
Because $\delta L^{0.8}/L^{0.6}=\delta L^{0.2}\to\infty$, both margins
dominate the error in \eqref{eq:Tao-first-time-two-sided}; this proves the
event inclusion.  The logarithmic width of $R_y$ is
$(\alpha-1)\log y-2L^{0.8}>0$ for large $x$.

For $1\leq a<b$, subtracting the even terms from the ordinary harmonic sum
gives the uniform integral estimate
\[
 \sum_{\substack{a\leq m\leq b\\m\equiv1\pmod2}}\frac1m
 =\frac12\log\frac ba+O(1/a).
\]
The denominator of the law of $\mathbf N_y$ is consequently
$\frac12(\alpha-1)\log y+O(1/y)$, while the two discarded intervals have
combined mass $L^{0.8}+O(1/y)$.  Their probability is $O(L^{-0.2})=
O(L^{-1/5})$.  Source line~737 gives
$\mathbb P(\mathcal G_x^c)=O(L^{-10})$, and the last assertion follows.
\end{proof}

The multiplicative restriction is used only to replace source lines~759--762.
Every subsequent line through the end of the section at line~925 uses the
unchanged interval $I_y$ and the resulting probability estimate, not either
endpoint restriction.  In particular, lines~816--820 have an $L^{0.8}$
interval margin against an $O(L^{0.7})$ reconstruction error, and
$\#(I_y\cap\mathbb Z)$ in lines~841--845 is unchanged.  One other local
inequality in this path must be written with its sign exposed: on
$\mathcal G_x$, for $0\leq j\leq n-m_0$, there is an absolute $K_1$ such that
\[
 \operatorname{Syr}^{j}(\mathbf N_y)
 \geq e^{-2K_1L^{0.6}}(4/3)^{n-m_0-j}
       \operatorname{Syr}^{n-m_0}(\mathbf N_y).
\]
If the last factor lies in $E'$, its source lower bound is
$e^{-L^{0.7}}(4/3)^{m_0}x$; since
$m_0=\lfloor(\alpha-1)L/100\rfloor$, the resulting lower bound is strictly
larger than $x$ for sufficiently large $x$.  This is the typed replacement
for the unsigned $\exp(O(L^{0.6}))$ comparison at lines~790--793.  Likewise,
the calculation immediately preceding source line~889 gives, when $q\leq
M_0$,
\[
 \sum_{\substack{M_0\leq M\leq M_1\\M\equiv a\pmod q}}\frac1M
 \leq\frac1q\left(1+\log\frac{M_1}{M_0}\right),
\]
which is the exact expression used in the following estimate.

\begin{proposition}[completion of the v7 Section~5 internal argument]
\label{prop:Tao-v7-Prop52-cn-repair}
Retain the notation of the pinned source and of
Proposition~\ref{prop:Tao-v7-endpoint-buffer}.  Put
\[
 r=n-m_0,
 \qquad
 G_x=\mathcal G_x
 =\left\{\vec a^{(n_0)}(\mathbf N_y)
                 \in\mathcal A^{(n_0)}\right\}.
\]
For the chosen \(E\subseteq2\mathbb N+1\), write
\[
\begin{split}
 E'=\{M\in2\mathbb N+1:\;&T_x(M)=m_0,\ 
       \operatorname{Pass}_x(M)\in E,\\
 &e^{-L^{0.7}}(4/3)^{m_0}x
       \leq M\leq e^{L^{0.7}}(4/3)^{m_0}x\}.
\end{split}
\]
For \(n\in I_y\cap\mathbb Z\), define
\[
\begin{split}
 B_{n,y}={}&
 \{T_x(\operatorname{Syr}^{r}(\mathbf N_y))=m_0\}\\
 &{}\cap
 \{\operatorname{Pass}_x(\operatorname{Syr}^{r}(\mathbf N_y))\in E\}
 \cap G_x.
\end{split}
\]
Then the source event \(B_{n,y}\) satisfies
the exact equality
\begin{equation}
\label{eq:Tao-v7-Bny-repaired}
 B_{n,y}
 =\left\{\operatorname{Syr}^{r}(\mathbf N_y)\in E'\right\}\cap G_x.
\end{equation}
For a fixed positive tuple
\(\vec a\in\mathcal A^{(r)}\) and fixed \(M\in E'\), the final equivalence
needed in the approximate formula is
\begin{equation}
\label{eq:Tao-v7-fixed-M-repaired}
 \left\{\operatorname{Syr}^{r}(\mathbf N_y)=M,\,
        \vec a^{(r)}(\mathbf N_y)=\vec a\right\}
 =\left\{\operatorname{Aff}_{\vec a}(\mathbf N_y)=M\right\}.
\end{equation}
In particular, the displayed approximate formula of source
Proposition~5.2 follows with \(n-m_0\) in place of the undefined \(n-m\).

For \(X\in\mathbb Z/3^r\mathbb Z\), define exactly as in source line~854
\begin{equation}
\label{eq:Tao-v7-cn}
 c_n(X)=3^r
   \sum_{\substack{M\in E'\\M\equiv X\pmod{3^r}}}\frac1M.
\end{equation}
There is an absolute constant \(C\) such that
\begin{equation}
\label{eq:Tao-v7-cn-uniform}
 0\leq c_n(X)\leq C
\end{equation}
for every sufficiently large \(x\), both choices of \(y\), every
\(n\in I_y\cap\mathbb Z\), and every \(X\in\mathbb Z/3^r\mathbb Z\).

To type the remaining dependency exactly, put
\[
 p_r(Y)=\mathbb P\!\left(
  \operatorname{Syrac}(\mathbb Z/3^r\mathbb Z)=Y\right).
\]
For \(n\in I_y\), the source parameter inequalities give
\(r=n-m_0\geq m_0\geq1\) for all sufficiently large \(x\), so the two
levels below lie in the declared domain of the mixing proposition.
The fine-scale-mixing input used by the source is, for every fixed \(B>0\),
\begin{equation}
\label{eq:Tao-v7-fine-mixing-input}
 \sum_{Y\bmod3^r}\left|
 p_r(Y)-3^{m_0-r}
 \sum_{\substack{Y'\bmod3^r\\Y'\equiv Y\pmod{3^{m_0}}}}p_r(Y')
 \right|\ll_B m_0^{-B}.
\end{equation}
After the indicator of \(\mathcal A^{(r)}\) is removed by a correctly typed
Chernoff estimate, the exact implication from
\eqref{eq:Tao-v7-fine-mixing-input} changes the expectation by
\(O_B(m_0^{-B})\).  Its resulting main term collapses, without a further
estimate, to the source quantity
\begin{equation}
\label{eq:Tao-v7-Z-repaired}
 Z=\sum_{M\in E'}\frac{3^{m_0}}M
       \mathbb P\!\left(
       \operatorname{Syrac}(\mathbb Z/3^{m_0}\mathbb Z)
          \equiv M\pmod{3^{m_0}}\right).
\end{equation}
\end{proposition}

\begin{proof}
Assume first that \(B_{n,y}\) holds.  Thus
\(T_x(\operatorname{Syr}^{r}(\mathbf N_y))=m_0\),
\(\operatorname{Pass}_x(\operatorname{Syr}^{r}(\mathbf N_y))\in E\), and
\(G_x\) holds.  Time \(n=r+m_0\) is a crossing:
\[
 \operatorname{Syr}^{n}(\mathbf N_y)\leq x
 <\operatorname{Syr}^{n-1}(\mathbf N_y).
\]
The good-orbit estimate used in
\eqref{eq:Tao-first-time-two-sided} places both this crossing time and the
actual first-passage time within \(O(L^{0.6})\) of
\(\log(\mathbf N_y/x)/\log(4/3)\).  Hence their difference is
\(O(L^{0.6})\).  Since
\(m_0=\lfloor L/100000\rfloor\), this gives
\(T_x(\mathbf N_y)\geq n-m_0=r\) for all sufficiently large \(x\).  The
semigroup identity for first passage may therefore be used at time \(r\),
and gives
\[
 T_x(\mathbf N_y)
 =r+T_x(\operatorname{Syr}^{r}(\mathbf N_y))=n.
\]
Writing \(M=\operatorname{Syr}^{r}(\mathbf N_y)\), the same signed
good-orbit comparison between times \(r\) and \(n\) gives a fixed absolute
\(K_2\).  Moreover, the crossing at time \(n\) and the two adjacent
good-prefix bounds imply
\(e^{-K_2L^{0.6}}x\leq\operatorname{Syr}^n(\mathbf N_y)
\leq x\).  Consequently,
\[
 e^{-K_2L^{0.6}}(4/3)^{m_0}x
 \leq M\leq
 e^{K_2L^{0.6}}(4/3)^{m_0}x.
\]
Since \(K_2L^{0.6}\leq L^{0.7}\) for all sufficiently large \(x\), this
lies in the defining window for \(E'\).  The two other defining conditions
for \(E'\), namely
\(T_x(M)=m_0\) and \(\operatorname{Pass}_x(M)\in E\), are already part of
\(B_{n,y}\).  Hence the right side of
\eqref{eq:Tao-v7-Bny-repaired} holds.

Conversely, suppose that
\(M=\operatorname{Syr}^{r}(\mathbf N_y)\in E'\) and \(G_x\) holds.  The
signed lower comparison displayed immediately before this proposition gives,
for every \(0\leq j\leq r\),
\[
 \operatorname{Syr}^{j}(\mathbf N_y)
 \geq e^{-2K_1L^{0.6}}(4/3)^{r-j}M
 \geq e^{-2K_1L^{0.6}-L^{0.7}}
       (4/3)^{n-j}x.
\]
The exponent at \(j=r\) is
\[
 m_0\log(4/3)-2K_1L^{0.6}-L^{0.7},
\]
which tends to \(+\infty\) because
\(m_0=\lfloor L/100000\rfloor\); for \(j<r\) the exponent is larger.
Thus the displayed quantity is \(>x\), and no passage has occurred by
time \(r\).  The defining properties \(T_x(M)=m_0\) and
\(\operatorname{Pass}_x(M)\in E\) now give passage time \(n\) and passage
location in \(E\), proving the reverse inclusion in
\eqref{eq:Tao-v7-Bny-repaired}.  Replacing the full good event by its prefix
of length \(r\) changes a fixed-\(n\) probability by at most
\(O(L^{-10})\); summing over the \(O(L)\) integers in \(I_y\) remains
\(O(L^{-9})\).  Finally, the prefix-valuation uniqueness lemma says that
\(\operatorname{Aff}_{\vec a}(\mathbf N_y)=M\), with \(M\) an odd positive
integer, is equivalent to the conjunction on the left of
\eqref{eq:Tao-v7-fixed-M-repaired}.  This proves the repaired approximate
formula.

It remains to prove \eqref{eq:Tao-v7-cn-uniform}.  Write
\[
 m=m_0,\qquad A=a_1+\cdots+a_m,\qquad b=a_m,
\]
where every \(a_i\geq1\); the source's \(\mathbb N^{m_0}\) at lines~862
and~865 must therefore be \((\mathbb N+1)^{m_0}\).  For
\(\vec a=(a_1,\ldots,a_m)\) define
\[
 c_{n,\vec a}(X)
 =3^r\!\!\sum_{\substack{M\in E'\\M\equiv X\pmod{3^r}\\
                  \vec a^{(m)}(M)=\vec a}}\frac1M.
\]
Uniqueness of the valuation prefix gives the exact decomposition
\[
 c_n(X)=\sum_{\vec a\in(\mathbb N+1)^m}c_{n,\vec a}(X).
\]
For an \(M\) in the summand, first passage at time \(m\) gives
\[
 3^m2^{-A}M+F_m(\vec a)\leq x
 <3^{m-1}2^{-(A-b)}M+F_{m-1}(a_1,\ldots,a_{m-1}).
\]
The coefficient in the second term is \(3^{m-1}\), not the printed \(3^m\).
Since \(0\leq F_{m-1}\leq3^{m-1}=o(x)\), the second inequality gives
\(M>3^{-m}2^{A-b}x\) for all sufficiently large \(x\); the first gives the
upper bound below.  Thus every such \(M\) lies in the definite interval
\[
 M_0\leq M\leq M_1,\qquad
 M_0=3^{-m}2^{A-b}x,\qquad
 M_1=3^{-m}2^A x.
\]

Multiplying the \(m\)-step affine identity by \(2^A\) shows that \(M\) lies
in one residue class modulo \(2^{A+1}\).  Together with
\(M\equiv X\pmod{3^r}\), the Chinese remainder theorem gives one residue
class modulo
\[
 q=2^{A+1}3^r.
\]
If \(2^A\leq x^{1/2}\), then \(b\leq A\), \(n\leq n_0\), and
\(3^n\leq x^{4/25}\) give
\[
 \frac q{M_0}=\frac{2^{b+1}3^n}{x}
 \leq2x^{-17/50}=o(1),
 \qquad \frac{M_1}{M_0}=2^b.
\]
Thus \(q\leq M_0\) for all sufficiently large \(x\).  The exact progression
estimate already displayed above gives
\begin{equation}
\label{eq:Tao-v7-cn-low-A}
 c_{n,\vec a}(X)\ll b2^{-A}.
\end{equation}
If \(2^A>x^{1/2}\), the exact last-step valuation identity is
\[
 b=\nu_2\!\left(
 3\left(3^{m-1}2^{-(A-b)}M
          +F_{m-1}(a_1,\ldots,a_{m-1})\right)+1\right).
\]
The inner coefficient \(3^{m-1}\) corrects the second occurrence of the
printed \(3^m\) defect at source line~899.  Put
\(H=3^{m-1}2^{-(A-b)}M\).  The preceding lower bound gives
\(H>x-F_{m-1}\), hence \(F_{m-1}=o(H)\).  Because \(2^b\) divides the
positive integer \(3(H+F_{m-1})+1\), one has
\[
 2^b\leq3(H+F_{m-1})+1\leq5H
\]
for all sufficiently large \(x\).  Therefore
\(M\geq(1/5)3^{1-m}2^A\gg3^{-m}2^A\).  Applying the unsimplified
progression estimate with the larger of the two lower bounds yields
\begin{equation}
\label{eq:Tao-v7-cn-high-A}
 c_{n,\vec a}(X)\ll3^n2^{-A}+b2^{-A}.
\end{equation}

The \(b\)-terms in both branches have the exact total mass
\[
 \sum_{a_1,\ldots,a_m\geq1}b2^{-A}
 =\left(\sum_{a\geq1}2^{-a}\right)^{m-1}
   \sum_{b\geq1}b2^{-b}=2.
\]
On the high-\(A\) branch,
\[
 \sum_{2^A>x^{1/2}}2^{-A}
 \leq x^{-1/4}
       \sum_{a_1,\ldots,a_m\geq1}2^{-A/2}
 =x^{-1/4}(1+\sqrt2)^m.
\]
Since \(3^5<2^8\), \(1+\sqrt2<e\), \(n\leq n_0\), and
\(m\leq L/100000\),
\[
 3^n x^{-1/4}(1+\sqrt2)^m
 \leq x^{4/25-1/4+1/100000}
 =x^{-8999/100000}.
\]
Equations~\eqref{eq:Tao-v7-cn-low-A} and
\eqref{eq:Tao-v7-cn-high-A} now prove the absolute uniform bound.  This is
the step missing from the printed summation: its weakened majorant has total
mass \((1+\sqrt2)^m\), which is not uniformly bounded.

For the iid geometric tuple in source line~852, the same union--Chernoff
calculation used earlier gives
\[
 \mathbb P\!\left((\mathsf a_1,\ldots,\mathsf a_r)
          \notin\mathcal A^{(r)}\right)=O(L^{-10}).
\]
This is the typed replacement for citing source line~737, which concerns the
actual valuation tuple.  Boundedness of \(c_n\) permits removal of the
indicator with the same error.  Applying the displayed sourced
\(\ell^1\) input \eqref{eq:Tao-v7-fine-mixing-input} and multiplying by
\(c_n\) changes the expectation by \(O_B(m_0^{-B})\).  Projective
consistency of the source variables gives, for each residue \(X\bmod3^{m_0}\),
\[
 \sum_{\substack{Y\bmod3^r\\Y\equiv X\pmod{3^{m_0}}}}p_r(Y)
 =\mathbb P\!\left(
   \operatorname{Syrac}(\mathbb Z/3^{m_0}\mathbb Z)=X\right);
\]
this is source line~355, obtained by reducing the defining finite series.
Therefore the main term is
\[
 \sum_{X\bmod3^r}c_n(X)3^{m_0-r}
 \mathbb P\!\left(
   \operatorname{Syrac}(\mathbb Z/3^{m_0}\mathbb Z)
       \equiv X\pmod{3^{m_0}}\right).
\]
Substitution of \eqref{eq:Tao-v7-cn} and interchange of the finite sums
multiplies \(3^r\) by \(3^{m_0-r}\), leaves one \(X=M\bmod3^r\) for every
\(M\in E'\), and gives \eqref{eq:Tao-v7-Z-repaired} exactly.  Taking, for
example, \(B=12\) makes the uniform mixing-output error more than sufficient
after the \(O(L)\) passage-time sum.  This last paragraph certifies the exact
dependency use of the stated mixing estimate; it does not supply a proof of
that estimate.
\end{proof}

Thus Proposition~5.2 and the uniform \(c_n\) lemma are certified at their
stated scope, and the Section~5 deduction is checked through source line~925.
At this dependency boundary the fine-scale-mixing, Fourier, and renewal
branch was still open; the next subsection reconstructs and discharges that
branch without treating the durable whole-paper release gate as complete.
In the source's nonperiodicity remark, line~699 additionally
requires an exponent representing \(-\mathbf m\) modulo the least positive
period, or a period multiple at least \(\mathbf m\); the counting conclusion
then remains unchanged.  The earlier descent estimate also requires
\(1.9n_0\), not the undefined \(1.9n\), at source line~689.

\subsection{The fine-scale Fourier and renewal branch}

This subsection reconstructs the branch left open after
Proposition~\ref{prop:Tao-v7-Prop52-cn-repair}.  The controlling text is the
version-pinned arXiv v7 source, lines~926--1860
\cite{Tao2026v7}.  ArXiv v5 and the published paper are comparison
manifestations, not silent replacements for v7.  Exact source hashes,
line-level defects, and the proof calculations abbreviated here are recorded
in the F024 Fourier--renewal audit and the F028 whole-version audit in the
external QA register.

\begin{sourceissue}[version boundary in the many-triangles estimate]
\label{sourceissue:Tao-v5-v7-many-triangles}
ArXiv v5, source lines~1654--1659, and the journal's Lemma~7.9 state the
unconditional estimate
\begin{equation}
\label{eq:Tao-v5-rip}
 \mathbb E\exp\left(
 -\sum_{p=1}^{t_{\min(r,R)}}1_W(X_p)
 +\varepsilon\min(r,R)\right)\le e^\varepsilon.
\end{equation}
ArXiv v7, source lines~1660--1665, states only
\begin{equation}
\label{eq:Tao-v7-rip}
 \mathbb E1_{\{r\ge R\}}
 \exp\left(-\sum_{p=1}^{t_R}1_W(X_p)+\varepsilon R\right)
 \le e^\varepsilon.
\end{equation}
On \(r\ge R\) the integrands agree, but the v7 formula deletes every path
with \(r<R\); the two statements are not identical.  Moreover, the displayed
v5/journal induction crosses the first triangle's top after \(t_1\).  On the
event \(r=1<R\), the left side of \eqref{eq:Tao-v5-rip} stops at \(t_1\), so
the additional factors used by that induction are absent.  We therefore
record the stronger statement as printed but do not use it.  The weaker v7
statement is proved below.
\end{sourceissue}

\begin{sourceissue}[Section 6 reserve is a substantive version repair]
\label{sourceissue:Tao-v5-v7-section-six-reserve}
Write \(a=\log(4/3)\), \(b=\log2\), and
\[
 q=\frac{b^2}{4a}=0.4175208154\ldots.
\]
The optimized square-root fluctuation cost in Corollary~6.3 is
\(qC_A^2\log N\).  ArXiv v5 and journal equations~(6.7)--(6.8)
retain only \(\frac12bC_A^2\log N=0.3465735903\ldots C_A^2\log N\);
the v5 calculation also drops the factor \(b\) when converting \(2^l\)
to an exponential.  Its printed size estimate therefore has the wrong
leading sign.  An explicit family satisfying the broad v5 concentration
conditions makes one of its natural-number summands exceed \(3^N\); the
construction and its exact scope are recorded in F028.

V7 uses \(0.99bC_A^2\log N=0.6862157088\ldots C_A^2\log N\), and the
stopping event supplies precisely that tightened window after enlarging
\(C_A\).  More generally, a retained coefficient \(\rho\) closes this
calculation exactly when
\[
 \rho>\frac{\log2}{4\log(4/3)}
 =0.6023552099\ldots.
\]
The v5 failure is a defect in the printed proof.  It neither produces two
colliding coordinate tuples nor proves the separation corollary false by a
different argument.
\end{sourceissue}

\begin{proposition}[exact Section 6 reserve boundary]
\label{prop:Tao-Section6-reserve-repair}
Put
\[
 a=\log(4/3),\qquad b=\log2,\qquad
 \rho_*=\frac{b}{4a}.
\]
Let \(n\ge2\), \(L=\log n\), \(C_A>0\), and let
\(a_1,\ldots,a_K\) be positive integers satisfying
\begin{equation}
\label{eq:Tao-Section6-prefix-lower}
 a_{[1,j]}\ge2j-C_A(\sqrt{jL}+L)
 \qquad(1\le j\le K).
\end{equation}
If
\[
 l=a_{[1,K]}
 \le\frac{n\log3}{b}-\rho C_A^2L
\]
for a fixed \(\rho>\rho_*\), then, after first choosing \(C_A\)
sufficiently large in terms of \(\rho-\rho_*\) and then \(n\) sufficiently
large in terms of \(C_A,\rho\),
\begin{equation}
\label{eq:Tao-Section6-Archimedean-separation}
 \sum_{j=1}^{K}3^{j-1}2^{\,l-a_{[1,j]}}<3^n.
\end{equation}
Here
\[
 \rho_*=0.6023552099\ldots,\qquad 0.99>\rho_*.
\]

For the v5 value \(\rho=\frac12\), the universal conclusion
\eqref{eq:Tao-Section6-Archimedean-separation} is false even on sequences
satisfying the source's all-interval concentration inequalities.  By
contrast, on the v7 stopping event \(E_k\cap B_k\), \(C_A\ge600\) gives
\[
 a_{[1,k+1]}
 \le\frac{n\log3}{\log2}-0.99C_A^2\log n.
\]
\end{proposition}

\begin{proof}
From \eqref{eq:Tao-Section6-prefix-lower},
\[
\begin{split}
 S
 :=\sum_{j=1}^{K}3^{j-1}2^{\,l-a_{[1,j]}}
 &\ll
 3^n e^{-\rho bC_A^2L+bC_AL}
 \sum_{j\ge1}e^{-aj+bC_A\sqrt{jL}}.
\end{split}
\]
The exact completion of the square is
\[
 -aj+bC_A\sqrt{jL}
 =-a\left(\sqrt j-\frac{bC_A\sqrt L}{2a}\right)^2
 +\frac{b^2}{4a}C_A^2L.
\]
Splitting the sum into unit intervals around its maximum, or comparing it
with the corresponding Gaussian integral after \(u=\sqrt j\), gives
\[
 \sum_{j\ge1}e^{-aj+bC_A\sqrt{jL}}
 \ll(1+C_A^2L)e^{b^2C_A^2L/(4a)}.
\]
The coefficient of \(C_A^2L\) in the resulting exponent is
\[
 -\rho b+\frac{b^2}{4a}=-b(\rho-\rho_*).
\]
It is strictly negative under the stated hypothesis.  It absorbs first the
linear \(bC_AL\) term after \(C_A\) is chosen, and then the polynomial factor
after \(n\) is chosen.  This proves \(S<3^n\).

For the v5 counterexample, set
\[
 \lambda=\frac{2a}{b},\qquad \gamma=\lambda^{-2},\qquad
 J=\lfloor\gamma C_A^2L\rfloor,
\]
and, for \(1\le i\le J\), put
\[
 d_i=\lfloor\lambda i\rfloor-\lfloor\lambda(i-1)\rfloor,\qquad
 a_i=2-d_i;
\]
put \(d_i=0\) and \(a_i=2\) afterwards.  Thus \(a_i\in\{1,2\}\), and the
deficit on every interval of length \(r\) is at most \(\lambda r+1\).
For \(r\le J\), \(\lambda r\le C_A\sqrt{rL}\); for \(r>J\), the entire
deficit is at most \(C_A\sqrt{rL}+O(1)\).  Hence the source's additional
\(C_AL\) tolerance makes every printed interval inequality valid.

Let \(D=\lfloor\lambda J\rfloor\).  Choose \(l\), with \(l+D\) even, within
two of
\[
 \frac{n\log3}{b}-\frac12C_A^2L,
\]
and put \(K=(l+D)/2\).  For sufficiently large \(n\), \(K\ge J\) and
\(a_{[1,K]}=2K-D=l\).  The \(j=J\) summand satisfies
\[
 \log\frac{3^{J-1}2^{\,l-a_{[1,J]}}}{3^n}
 =
 \left(\frac{b^2}{4a}-\frac12b\right)C_A^2L+O(1).
\]
The coefficient is
\[
 0.4175208154\ldots-0.3465735903\ldots
 =0.0709472252\ldots>0,
\]
so that single summand eventually exceeds \(3^n\).

Finally put \(T=n\log3/b-C_A^2L\).  On \(B_k\),
\(a_{[1,k]}\le T<a_{[1,k+1]}\).  The \(i=k,j=k+1\) inequality on \(E_k\)
and \(a_k\ge1\) yield
\[
 a_{[1,k+1]}\le T+1+C_A(\sqrt L+L).
\]
For \(n\ge2\), the last overshoot is at most \(6C_AL\), and for
\(C_A\ge600\) this is at most \(0.01C_A^2L\).  This gives the v7
\(0.99\)-window.
\end{proof}

\begin{nonclaim}
Proposition~\ref{prop:Tao-Section6-reserve-repair} refutes the universal
Archimedean size assertion in the printed v5 proof, not the separation
corollary by every possible proof.  The construction supplies no pair of
colliding tuples.  The coefficient \(0.99\) is convenient rather than
canonical, and the threshold \(\rho_*\) belongs to this exact
optimization-and-absorption argument.  No Fourier-decay theorem, main
theorem, or Collatz endpoint follows from this finite margin calculation.
\end{nonclaim}

\begin{definition}[uniform lifting and fibre averaging]
\label{def:Tao-uniform-lift}
For \(0\le r\le N\), let
\[
 L_{r\to N}:\mathbb R^{\mathbb Z/3^r\mathbb Z}
 \longrightarrow\mathbb R^{\mathbb Z/3^N\mathbb Z},
 \qquad
 (L_{r\to N}\mu)(y)=3^{r-N}\mu(y\bmod3^r).
\]
For the reduction map
\(\pi_{N,r}:\mathbb Z/3^N\mathbb Z\to\mathbb Z/3^r\mathbb Z\), its fibre
above \(x\) is \(x+3^r\mathbb Z/3^N\mathbb Z\).  On signed measures on the
larger group define the fibre-constant projection
\[
 P_{r,N}:\mathbb R^{\mathbb Z/3^N\mathbb Z}
 \longrightarrow\mathbb R^{\mathbb Z/3^N\mathbb Z},\qquad
 (P_{r,N}\nu)(y)=3^{r-N}
 \sum_{\substack{z\bmod3^N\\z\equiv y\ (3^r)}}\nu(z).
\]
\end{definition}

Both maps preserve total mass, \(P_{r,N}\) is an \(\ell^1\)-contraction,
and \(L_{r\to N}\) is an \(\ell^1\)-isometry.  If \(\mu_N\) is the law of
Tao's \(N\)-Syracuse random variable, projective consistency says
\((\pi_{N,r})_*\mu_N=\mu_r\).  Substitution in the two displayed maps gives
the typed identity
\begin{equation}
\label{eq:Tao-lift-projective}
 P_{r,N}\mu_N=L_{r\to N}\mu_r.
\end{equation}

\begin{proposition}[exact reduction from characteristic decay to fine-scale mixing]
\label{prop:Tao-v7-Fourier-to-mixing}
For every \(q\in\Npos\), let \(\mu_q\) be the law on
\(\mathbb Z/3^q\mathbb Z\) just defined.  Assume the v7
characteristic-function estimate, uniformly for
\(\xi\in\mathbb Z/3^q\mathbb Z\) with \(3\nmid\xi\),
\begin{equation}
\label{eq:Tao-v7-characteristic-input}
 \left|\mathbb E
 e^{-2\pi i\xi\operatorname{Syrac}(\mathbb Z/3^q\mathbb Z)/3^q}
 \right|\ll_Bq^{-B}
\end{equation}
for every fixed \(B>0\), with an implied constant depending only on \(B\).
After correcting source line~1051 from \(3^q\)
to \(3^{q-1}\), and line~1093 from \(m\) to \(k+1\), the source argument
proves
\[
 \|\mu_N-L_{r\to N}\mu_r\|_1\ll_A r^{-A}
 \qquad(N,r\in\Npos,\ 1\le r\le N).
\]
This is exactly Proposition~1.14, including its domain and uniformity.
\end{proposition}

\begin{proof}
First suppose \(0.9N\le r\le N\).  Let \(E\) be the source's all-interval
concentration event.  For a fixed \(k\), let \(E_k\supseteq E\) retain only
the concentration inequalities involving \(a_1,\ldots,a_{k+1}\).  The source's
concentration estimate gives
\(\mathbb P(E_k\setminus E)=O_{A,C_A}(N^{-A-1})\), so the triangle
inequality permits \(E\) to be replaced by \(E_k\) at the stated error
scale.  On \(E_k\), stop at the unique \(k\) satisfying
\[
 a_{[1,k]}\le \frac{N\log3}{\log2}-C_A^2\log N
 <a_{[1,k+1]},
\]
and condition also on \(a_{[1,k+1]}=l\).  Write \(B_k\) for this stopping
event and \(C_{k,l}=\{a_{[1,k+1]}=l\}\).  Put
\[
 T_N=\frac{N\log3}{\log2}-C_A^2\log N.
\]
The \(i=k,j=k+1\) inequality in \(E_k\), together with \(a_k\ge1\), gives
\[
 a_{k+1}\le1+C_A(\sqrt{\log N}+\log N),
\]
and hence
\begin{equation}
\label{eq:Tao-v7-point99-window}
 l=a_{[1,k+1]}
 \le T_N+1+C_A(\sqrt{\log N}+\log N)
 \le\frac{N\log3}{\log2}-0.99C_A^2\log N
\end{equation}
after choosing \(C_A\ge600\) and then \(N\ge2\).  The last displayed
sufficient choice follows from
\(\sqrt{\log N}\le2\log N\), \(2\le3\log N\), and
\(6C_A\le0.01C_A^2\).  The prefix and suffix offset maps then have the
exact typed decomposition
\[
 X_N=F_{k+1}(a_{k+1},\ldots,a_1)
 +3^{k+1}2^{-l}F_{N-k-1}(a_N,\ldots,a_{k+2})
 \pmod{3^N}.
\]
The events \(E_k,B_k,C_{k,l}\) and the prefix offset depend only on
\(a_1,\ldots,a_{k+1}\), and hence are independent of the suffix variables.
Let
\(\eta\in\mathbb Z/3^N\mathbb Z\) be a frequency surviving fibre averaging,
 so \(\eta\notin3^{N-r}\mathbb Z/3^N\mathbb Z\).  There are unique
\[
 0\le v<N-r,
 \qquad \bar\eta'\in(\mathbb Z/3^{N-v}\mathbb Z)^\times,
\]
where \(v=\nu_3(\eta)\) and
\[
 \eta/3^v=2^l\bar\eta'\pmod{3^{N-v}}.
\]
Thus the quotient is asserted unique only in the correct quotient ring, not
as a residue modulo \(3^N\).  Put
\[
 q=N-k-v-1.
\]
The concentration range for \(k\), together with \(v<N-r\le0.1N\), gives
\[
 q\ge0.9N-\frac{N\log3}{2\log2}
       -O(C_A\sqrt{N\log N})-1\gg N
\]
for sufficiently large \(N\).  Let
\(\eta'_q\in(\mathbb Z/3^q\mathbb Z)^\times\) be the reduction of
\(\bar\eta'\).  Cancelling the exact powers of three and two in the suffix
character gives the typed frequency morphism
\begin{align*}
 &\mathbb E\exp\left(
 -\frac{2\pi i\eta\,3^{k+1}2^{-l}
 F_{N-k-1}(a_N,\ldots,a_{k+2})}{3^N}\right)\\
 &\qquad=
 \mathbb E\exp\left(-\frac{2\pi i\eta'_q
 F_{N-k-1}(a_N,\ldots,a_{k+2})}{3^q}\right)\\
 &\qquad=
 \mathbb E\exp\left(-\frac{2\pi i\eta'_q
 \operatorname{Syrac}(\mathbb Z/3^q\mathbb Z)}{3^q}\right).
\end{align*}
The final equality is projective consistency of the offset law from
\(N-k-1\) to \(q\), and
\eqref{eq:Tao-v7-characteristic-input} bounds this factor by
\(O_B(N^{-B})\), uniformly in every surviving frequency and admissible
\((k,l)\).

Write \(\mathcal E_{r,N;k,l}^2\) for the squared oscillation at this fixed
pair \((k,l)\).  Fourier inversion, removal of the frequencies divisible by
\(3^{N-r}\), the preceding suffix bound, and Plancherel give
\begin{equation}
\label{eq:Tao-prefix-collision}
 \mathcal E_{r,N;k,l}^2\ll_B N^{-2B}3^N\sum_y p_y^2,
\end{equation}
where, with the source events \(E_k,B_k,C_{k,l}\),
\[
 p_y=\mathbb P\bigl(
 F_{k+1}(a_{k+1},\ldots,a_1)\equiv y\pmod{3^N},
 E_k\cap B_k\cap C_{k,l}\bigr).
\]
Thus \((p_y)_y\) is the unnormalised restricted joint mass, not a
conditional probability distribution.

For positive coordinates,
\[
 F_q(a_1,\ldots,a_q)
 =\sum_{u=1}^q3^{q-u}2^{-a_{[u,q]}}.
\]
The unique term of least 2-adic valuation is
\(3^{q-1}2^{-a_{[1,q]}}\).  Hence equality of two rational offsets forces
equality of their total valuation sums.  Subtraction gives equality of the
two \((q-1)\)-coordinate tails, and induction proves that \(F_q\) has
singleton fibres on \(\Npos^q\).

Under the fixed-total and concentration conditions, multiplying a congruence
of two prefix offsets by \(2^l\) gives a congruence modulo \(3^N\) between
two positive integers.  Their prefix inequalities and
\eqref{eq:Tao-v7-point99-window} meet the hypotheses of
Proposition~\ref{prop:Tao-Section6-reserve-repair} with \(\rho=0.99\).
Hence both positive integers are strictly below \(3^N\).  They are therefore
equal as integers.  Dividing by \(2^l\) and applying rational injectivity
proves residue-level separation.
Thus
\[
 0\le p_y\le2^{-l},\qquad \sum_y p_y\le1,
\]
and consequently
\[
 3^N\sum_yp_y^2\le3^N(\max_yp_y)\sum_yp_y
 \le3^N2^{-l}=N^{O(C_A^2)}.
\]
Taking \(B\) sufficiently large in
\eqref{eq:Tao-v7-characteristic-input} absorbs this precise polynomial and
closes \eqref{eq:Tao-prefix-collision} with an arbitrary prescribed
\(N^{-2A-4}\) reserve.  There are \(O_A(N)\) admissible stopping indices
\(k\), and, from the displayed fixed-total window, \(O_A(\log N)\)
admissible totals \(l\).  Cauchy--Schwarz converts the squared bound back to
the fixed-pair \(\ell^1\)-oscillation, and the union bound over these
\((k,l)\) uses less than the two reserved powers of \(N\).  The complement
of the all-interval concentration event is \(O_A(N^{-A-1})\) after enlarging
\(C_A\).  Adding these disjoint fixed-pair contributions proves the
restricted \(\ell^1\) estimate, rather than silently identifying a
conditioned estimate with the unconditional one.

For general \(10\le r\le N\), construct
\(N=n_0>n_1>\cdots>n_t=r\) by setting
\(n_{i+1}=\max(r,\lceil0.9n_i\rceil)\).  Because every \(n_i\ge10\)
before termination, this sequence decreases strictly, and every adjacent
pair is in the proved range.  By \eqref{eq:Tao-lift-projective}, functoriality of the
lifts, and preservation of the \(\ell^1\)-norm,
\[
\begin{aligned}
 \|\mu_N-L_{r\to N}\mu_r\|_1
 &\le\sum_{i=0}^{t-1}
 \|L_{n_i\to N}\mu_{n_i}
       -L_{n_{i+1}\to N}\mu_{n_{i+1}}\|_1\\
 &=\sum_{i=0}^{t-1}
 \|\mu_{n_i}-L_{n_{i+1}\to n_i}\mu_{n_{i+1}}\|_1.
\end{aligned}
\]
Applying the restricted estimate with exponent \(A+1\), and summing the
geometrically separated levels, gives \(O_A(r^{-A})\).  This supplies the
lifting sequence implicit in source lines~932--936.  When \(1\le r<10\),
the triangle inequality gives a bound of two, which is
\(O_A(r^{-A})\) after enlarging the constant by at most \(2\,9^A\).
The finitely many \(N\) excluded while proving the restricted large-\(N\)
estimate are absorbed in the same \(A\)-dependent constant.  These two
finite-scale clauses are separate from the strictly decreasing chain.
\end{proof}

For \(d\in\Npos\), \(t\in\mathbb R_{\ge0}\), and \(y\in\mathbb R^d\),
we use Tao's Gaussian--exponential kernel
\begin{equation}
\label{eq:Tao-G-kernel}
 G_t(y)=
 \begin{cases}
  \exp(-\|y\|^2/t)+\exp(-\|y\|),&t>0,\\
  \exp(-\|y\|),&t=0.
 \end{cases}
\end{equation}

\begin{lemma}[the corrected two-dimensional local factor]
\label{lem:Tao-v7-local-factor}
In the proof of Tao's Lemma~2.2, the source-line-516 factor is
\(n^{-d/2}\), not \(n^{-1/2}\).  In particular, let
\(\mathbf H_1,\mathbf H_2,\ldots\) be iid copies of the holding vector
\(\mathbf H\in\mathbb Z^2\) of
Proposition~\ref{prop:Tao-v7-pairing-holding}, put
\(\mathbf H_{[1,q]}=\sum_{i=1}^q\mathbf H_i\) with the empty sum at
\(q=0\), and recall \(\mathbb E\mathbf H=(4,16)\).  There exist constants
\(c,C>0\), depending only on the law of \(\mathbf H\), such that for every
\(q\in\Nzero\) and \((x,t)\in\mathbb Z^2\),
\[
 \mathbb P(\mathbf H_{[1,q]}=(x,t))
 \le C(q+1)^{-1}G_q(c((x,t)-q(4,16))).
\]
\end{lemma}

\begin{proof}
After the source's contour shift, the Fourier integral contains
\(e^{-cn|u|^2}\) on a bounded subset of \(\mathbb R^d\).  The change of
variables \(v=\sqrt n\,u\) gives
\[
 \int e^{-cn|u|^2}\,du
 \le n^{-d/2}\int_{\mathbb R^d}e^{-c|v|^2}\,dv.
\]
All remaining source steps preserve this dimension.  Setting \(d=2\)
gives the displayed factor.
\end{proof}

\begin{proposition}[pairing, phase, and exact holding vector]
\label{prop:Tao-v7-pairing-holding}
Let \(n\in\Npos\), let
\(\xi\in\mathbb Z/3^n\mathbb Z\) satisfy \(3\nmid\xi\), and define
\[
 \chi:\mathbb Z[1/2]\longrightarrow\mathbb C,\qquad
 \chi(x)=\exp\left(-\frac{2\pi i\xi(x\bmod3^n)}{3^n}\right).
\]
Fix \(0<\varepsilon\le e^{-10}\).  For
\((r,l)\in\Npos\times\mathbb Z\), define the
signed fractional phase
\[
 \theta(r,l)=\left\{
 \frac{\xi3^{2r-2}(2^{-l+1}\bmod3^n)}{3^n}
 \right\}\in(-1/2,1/2],
\]
and define the typed partition of the strip
\[
 B=\{(r,l):1\le r\le\lfloor n/2\rfloor,
                  |\theta(r,l)|\le\varepsilon\},\qquad
 W=\bigl(\{1,\ldots,\lfloor n/2\rfloor\}\times\mathbb Z\bigr)\setminus B.
\]
Let \((a_i)_{i\ge1}\) be iid with
\(\mathbb P(a_i=u)=2^{-u}\) for \(u\in\Npos\), and put
\[
 S_\chi(m)=\mathbb E\chi\left(
 \sum_{v=1}^{m}3^{v-1}2^{-\sum_{w=1}^{v}a_w}
 \right)\qquad(m\in\Nzero),
\]
with the empty-sum convention \(S_\chi(0)=1\).
For independent copies \(a_1,a_2\) of this geometric variable, define
\[
 \begin{aligned}
 f&:\mathbb Z[1/2]\times\{2,3,\ldots\}\longrightarrow\mathbb C,\\
 f(x,b)&=\mathbb E\!\left[
 \chi\bigl(x(2^{a_2}+3)\bigr)\mid a_1+a_2=b\right],\\
 g&:\mathbb Z[1/2]\longrightarrow\mathbb C,\\
 g(x)&=\mathbb E\chi\bigl(x2^{-a_1}\bigr).
 \end{aligned}
\]
Let \(q=\lfloor n/2\rfloor\),
\(b_r=a_{2r-1}+a_{2r}\), and
\(B_r=\sum_{u=1}^rb_u\), with \(B_0=0\).  Then the \(b_r\) are iid with
\(\mathbb P(b_r=b)=(b-1)2^{-b}\) for \(b\ge2\).  Expectation in the next
two displays is over this joint Pascal law, and conditioning on the
\(b_r\) gives the exact formulas
\[
 S_\chi(2q)=\mathbb E
 \prod_{r=1}^qf(3^{2r-2}2^{-B_r},b_r)
\]
and
\[
 S_\chi(2q+1)=\mathbb E\left[
 \prod_{r=1}^qf(3^{2r-2}2^{-B_r},b_r)\,
 g(3^{2q}2^{-B_q})\right].
\]
At every white point \((r,l)\),
\[
 |f(3^{2r-2}2^{-l},3)|\le e^{-\varepsilon^3}.
\]
The occurrences \(b_r=3\) form a renewal process with iid holding vector
\[
 \mathbf H=(1,3)+\sum_{i=1}^{R}(1,Q_i),
 \qquad
 \mathbb P(R=s)=\frac14\left(\frac34\right)^s
 \qquad(s\in\Nzero),
\]
where \(R\) is independent of the iid \(Q_i\), and the \(Q_i\) have the
Pascal law conditioned on \(Q_i\ne3\).  This vector
has exponential tail, full lattice span, and mean \((4,16)\).
\end{proposition}

\begin{proof}
The pairing identities follow by conditioning while retaining the random
odd-\(n\) terminal factor inside the expectation.  Conditional on \(b=3\),
the fibres \((a_1,a_2)=(1,2),(2,1)\) have equal mass.  Therefore
\[
 f(x,3)=\frac{\chi(5x)+\chi(7x)}2
 =\frac{\chi(5x)}2(1+\chi(2x)).
\]
For \(x=3^{2r-2}2^{-l}\), the corrected character identity is
\(\chi(2x)=e^{-2\pi i\theta(r,l)}\).  Since
\(\theta(r,l)\in(-1/2,1/2]\),
\[
 |f(x,3)|=\cos(\pi\theta(r,l)).
\]
The elementary bound
\(\cos(\pi u)\le e^{-2u^2}\) for \(|u|\le1/2\) gives, at a white point,
\[
 |f(x,3)|\le e^{-2\varepsilon^2}\le e^{-\varepsilon^3},
\]
where the last inequality uses \(0<\varepsilon\le e^{-10}<2\).

The success probability is \(\mathbb P(b=3)=1/4\), giving the displayed
failure-before-success representation.  For
\(k=(k_1,k_2)\in\mathbb R^2\), put
\(M_Q(k)=\mathbb E e^{k_1+k_2Q}\).  On the open neighbourhood of zero on
which \(M_Q(k)<\infty\) and \(3M_Q(k)<4\), the normalized moment generating
function is
\begin{equation}
\label{eq:Tao-hold-mgf-corrected}
 \mathbb E e^{\mathbf H\cdot k}
 =\sum_{s=0}^\infty\frac14\left(\frac34\right)^s
 e^{(1,3)\cdot k}M_Q(k)^s
 =\frac{e^{(1,3)\cdot k}}{4-3M_Q(k)}.
\end{equation}
It equals one at \(k=0\) and is finite near zero.  This corrects the
off-by-one series at source line~1408 and proves the exponential tail.
The support contains \((1,3),(2,5),(2,7),(2,8)\); its differences generate
\((0,1)\) and \((1,0)\), so it is contained in no coset of a proper subgroup
of \(\mathbb Z^2\).  Finally,
\(\mathbb E Q=13/3\), \(\mathbb ER=3\), and direct substitution gives
\(\mathbb E\mathbf H=(4,16)\).
\end{proof}

\begin{proposition}[corrected black-triangle decomposition]
\label{prop:Tao-v7-black-triangles}
Retain \(n,\xi,\varepsilon,\theta,B,W\) from
Proposition~\ref{prop:Tao-v7-pairing-holding}.
After the forced corrections at v7 lines~1272, 1285--1287, and
1311--1313, \(B\) is a disjoint union of lattice triangles
\[
 \Delta=\{(j,l):j\ge j_\Delta,\ l\le l_\Delta,\
 (j-j_\Delta)\log9+(l_\Delta-l)\log2\le s_\Delta\}.
\]
Here \(j_\Delta\in\Npos\), \(l_\Delta\in\mathbb Z\),
\(s_\Delta\in\mathbb R_{\ge0}\), and the displayed set is intersected with
the lattice \(\Npos\times\mathbb Z\).
With \(\delta=\frac1{10}\log(1/\varepsilon)\), every triangle satisfies
\[
 \Delta\subseteq
 \{(j,l)\in\Npos\times\mathbb Z:j\le n/2-\delta\},
\]
and
distinct triangles have Euclidean distance at least \(\delta\).
\end{proposition}

\begin{proof}
The exact recurrences are
\[
 \theta(j+1,l)\equiv9\theta(j,l),\qquad
 \theta(j,l-1)\equiv2\theta(j,l)\pmod{\mathbb Z}.
\]
Call a point weakly black when \(|\theta|\le1/100\).  The recurrences give
three exact propagation rules.  First, if \((j,l)\) is weakly black and
\((j+1,l)\) is black, then \(|9\theta(j,l)|\le9/100<1/2\), so there is no
wrap and \(|\theta(j,l)|=|\theta(j+1,l)|/9\le\varepsilon\).  The same
argument with \(|2\theta(j,l)|\le2/100<1/2\) applies when \((j,l-1)\) is
black.  Second, if \((j+1,l)\) and \((j,l-1)\) are weakly black, then
\[
 |\theta(j,l)|\le |\theta(j+1,l)|+4|\theta(j,l-1)|\le5/100;
\]
now \(|2\theta(j,l)|\le1/10<1/2\), so the doubling recurrence has no wrap
and \(|\theta(j,l)|=|\theta(j,l-1)|/2\le1/200\).  Third, if
\((j-1,l)\) and \((j,l-1)\) are weakly black, the factor-nine recurrence
first gives \(|\theta(j,l)|\le9/100\); hence
\(|2\theta(j,l)|\le18/100<1/2\), and doubling again gives
\(|\theta(j,l)|\le1/200\).

If \((j,l)\) is black, multiplying its phase to the primitive residue at
level \(3^n\) gives
\[
 3^{n+1-2j}\varepsilon\ge\frac13.
\]
Taking logarithms gives
\[
 j\le \frac n2+1-\frac{\log(1/\varepsilon)}{2\log3}
 \le \frac n2-\frac1{10}\log(1/\varepsilon)
\]
for \(\varepsilon\le e^{-10}\).  There is no infinite upward black run:
if \((j,l')\) were black for every \(l'\ge l\), repeated nonwrapping
doubling would give
\(|\theta(j,l')|\le2^{l-l'}\varepsilon\), contradicting the same primitive
residue lower bound as \(l'\to\infty\).

Starting from a black \((j,l)\), let \(l_*\ge l\) be the unique top of its
consecutive black vertical run, and let \(j_*\le j\) be the unique left end
of the consecutive black run on row \(l_*\).  Existence and uniqueness follow
from the preceding no-infinite-run statement and the lower boundary
\(j=1\).  Because \(2j_*-2<n\), while neither \(\xi\) nor \(2\) is
divisible by three, the defining rational phase is not integral; hence
\(\theta(j_*,l_*)\ne0\).  Write
\(|\theta(j_*,l_*)|=\varepsilon e^{-s_*}\).  On the corrected domain
\(j'\ge j_*\), \(l'\le l_*\), multiplication by
\(9^{j'-j_*}2^{l_*-l'}\) gives
\[
 |\theta(j',l')|\le
 \varepsilon e^{-s_*+(j'-j_*)\log9+(l_*-l')\log2},
\]
with equality before the right side reaches \(1/2\).  Thus all points of the
anchored triangle are black.  If an exterior Case~1 point is within
\(\delta\), and
\(X=-s_*+(j'-j_*)\log9+(l_*-l')\log2\), then
\[
 \varepsilon<\varepsilon e^X
 \le\varepsilon^{1-(\log9+\log2)/10}<\tfrac12,
\]
so \(|\theta(j',l')|=\varepsilon e^X>\varepsilon\).  The corrected Case~2
argument is as follows.  If \(l'>l_*\) and a nearby exterior point
\((j',l')\) were black, Euclidean distance at most \(\delta\) gives
\[
 0<(l'-l_*)\log2\le\delta\log2,
 \qquad
 (j'-j_*)\log9\le s_*+\delta\log9.
\]
Exact downward phase transport then gives
\[
 |\theta(j',l_*+1)|
 \le\varepsilon^{1-\log2/10}<1/100.
\]
When \(j'>j\), the anchored phase bound also gives
\[
 |\theta(j'-1,l_*)|
 \le\varepsilon^{1-\log9/10}<1/100.
\]
The second propagation rule iterates left from \(j'-1\) to \(j\).  If
\(j'=j\), \((j,l_*+1)\) is already the weakly black point supplied by the
first display.  Thus in either case \((j,l_*+1)\) is weakly black, and
blackness of \((j,l_*)\) makes it black by the first rule, contradicting the
definition of \(l_*\).  When \(j'<j\), the black row from \(j_*\) through
\(j\) and the third rule instead propagate weak blackness right from
\((j',l_*+1)\) to \((j,l_*+1)\), giving the same contradiction.

In Case~3, \(j'<j_*\).  The same distance condition gives
\[
 |l'-l_*|\log2\le s_*+\delta\log2,
 \qquad 0<(j_*-j')\log9\le\delta\log9.
\]
Assuming \((j',l')\) black, exact horizontal phase transport gives
\[
 |\theta(j_*-1,l')|
 \le\varepsilon^{1-\log9/10}<1/100.
\]
If \(l'\ge l_*\), vertical transport gives the explicit bound
\[
 |\theta(j_*-1,l_*)|
 \le\varepsilon^{1-(\log9+\log2)/10}<1/100.
\]
If \(l'<l_*\), the anchored bound gives
\(|\theta(j_*,l'+1)|\le\varepsilon^{1-\log2/10}<1/100\), and the second
rule propagates upward until \((j_*-1,l_*)\) is weakly black.  In either
subcase, the first propagation rule and blackness of \((j_*,l_*)\) make
\((j_*-1,l_*)\) black, contradicting the defining minimality of \(j_*\).

Hence every exterior point within \(\delta\) is white.  Since
\(\delta\ge1\), the black horizontal and vertical paths defining the anchor
cannot cross the boundary: the first black path point outside the triangle
would be an exterior black point within distance one.  Thus every original
black point lies in its anchored triangle.

Conversely, every point of this triangle reconstructs the same anchor.
For each column occurring in the triangle, the entire vertical segment up to
row \(l_*\) is black and the point at height \(l_*+1\) is an exterior collar
point, hence white.  On row \(l_*\), the entire horizontal segment back to
\(j_*\) is black and, unless \(j_*=1\), the point at \(j_*-1\) is an
exterior collar point, hence white.  Therefore its reconstructed pair is
exactly \((j_*,l_*)\).  The anchored triangles partition \(B\), and the
white collar makes distinct members \(\delta\)-separated.

Confinement of the entire abstract triangle follows directly from the
anchor denominator, without presupposing confinement of its lattice points.
The nonzero signed phase satisfies
\[
 |\theta(j_*,l_*)|\ge3^{2j_*-2-n}.
\]
Since its other expression is \(\varepsilon e^{-s_*}\),
\[
 s_*\le(n-2j_*+2)\log3-\log(1/\varepsilon).
\]
Consequently its real right tip satisfies
\[
 j_*+\frac{s_*}{\log9}
 \le\frac n2+1-\frac{\log(1/\varepsilon)}{2\log3}
 \le\frac n2-\delta.
\]
Every point of the triangle lies weakly to the left of that tip, proving the
asserted confinement and, in particular, the real-tip bound needed later.
\end{proof}

\begin{lemma}[first-passage Green kernel]
\label{lem:Tao-v7-first-passage-repaired}
Let \((\mathbf H_i)_{i\ge1}\) be iid copies of \(\mathbf H\), let
\(\mathbf V_q=\sum_{i=1}^q\mathbf H_i\), and set
\(\mathbf V_0=(0,0)\).  For \(s\in\Nzero\), put
\(\kappa_s=\inf\{q\in\Npos:(\mathbf V_q)_2>s\}\).  For every
\(j,l\in\mathbb Z\) with \(l>s\),
\[
 \mathbb P(\mathbf V_{\kappa_s}=(j,l))
 \ll\frac{e^{-c(l-s)}}{\sqrt{1+s}}
 G_{1+s}(c(j-s/4)).
\]
Moreover, there is an absolute \(C\ge1\) such that the following typed
post-passage estimate holds.  Let \(m\in\Npos\) be sufficiently large,
let \(s,p\in\Nzero\), let \(r_0\in[1,\infty)\), and let \(h>0\).  Suppose
\[
 s>\frac{m}{\log^2m},\qquad 0\le p\le m^{0.1},\qquad
 4r_0\le h\le m^{0.4},
\]
and let \(\Sigma\subset\mathbb Z\) have pairwise spacing at least \(h\).
Then, for all sufficiently large \(m\),
\[
 \mathbb P\!\left(
 \operatorname{dist}((\mathbf V_{\kappa_s+p})_1,\Sigma)\le r_0
 \right)\le C\frac {r_0}h.
\]
\end{lemma}

\begin{proof}
Set
\(U(x,t)=\sum_{q\ge0}\mathbb P(\mathbf V_q=(x,t))\) for
\((x,t)\in\mathbb Z\times\Nzero\).  By
Lemma~\ref{lem:Tao-v7-local-factor},
\[
 \mathbb P(\mathbf V_q=(x,t))
 \ll(q+1)^{-1}G_q(c((x,t)-q(4,16))).
\]
Writing \(z=x-t/4\), the invertible coordinate identity
\[
 (x-4q,t-16q)=\left(z+\frac{t-16q}{4},t-16q\right)
\]
preserves norms up to two absolute constants.  Put
\(\langle t\rangle=1+t\), and, after harmlessly weakening the denominator
\(q\) in the Gaussian to \(q+1\), set
\[
 A_q(z,t)=\frac1{q+1}
 \exp\left(-c\frac{z^2+(t-16q)^2}{q+1}\right).
\]
In the central range the complete Gaussian calculation is
\[
 \sum_{t/2\le16q\le2t}A_q(z,t)
 \ll \langle t\rangle^{-1/2}
 e^{-c'z^2/\langle t\rangle}.
\]
Indeed \(q+1\asymp\langle t\rangle\) there and the discrete Gaussian in
\(t-16q\) has mass \(O(\sqrt{\langle t\rangle})\).  In the lower range,
\(|t-16q|\gg t\) and \(q+1\ll\langle t\rangle\); in the upper range use
\[
 (q+1)+\frac{z^2}{q+1}
 \gg \langle t\rangle+\min\left(
 \frac{z^2}{\langle t\rangle},|z|\right).
\]
Geometric summation therefore gives, respectively,
\[
 \sum_{16q<t/2}A_q(z,t)
 \ll\langle t\rangle^{-1/2}e^{-c'z^2/\langle t\rangle},
 \qquad
 \sum_{16q>2t}A_q(z,t)
 \ll\langle t\rangle^{-1/2}
 \left(e^{-c'z^2/\langle t\rangle}+e^{-c'|z|}\right).
\]
The exponential half of \(G_q\) is even smaller, since
\[
 \sum_{q\ge0}\frac{e^{-c(|z|+|t-16q|)}}{q+1}
 \ll\langle t\rangle^{-1/2}e^{-c'|z|}.
\]
The cases \(t=0,1\) are absorbed by changing the absolute constants.  Thus
\[
 U(x,t)\ll(1+t)^{-1/2}G_{1+t}(c'(x-t/4)).
\]

At first passage the final vertical increment is at least \(l-s\).  Using
its exponential tail and independence from the earlier increments gives
\[
 \mathbb P(\mathbf V_{\kappa_s}=(j,l))
 \ll e^{-c(l-s)}\sum_{r=0}^s\sum_{u\ge1}
 e^{-c(r+u)}U(j-u,s-r).
\]
For
\[
 K_t(y)=(1+t)^{-1/2}
 \left(e^{-c y^2/(1+t)}+e^{-c|y|}\right),
\]
the two needed convolution inequalities, with absolute
\(a,c_1,c_2>0\), are
\begin{align}
 \sum_{r=0}^{s}e^{-ar}K_{s-r}(x+r/4)
 &\ll K_s(c_1x),\label{eq:Tao-vertical-convolution}\\
 \sum_{u\ge1}e^{-au}K_s(x-u)
 &\ll K_s(c_2x).\label{eq:Tao-horizontal-convolution}
\end{align}
For \eqref{eq:Tao-vertical-convolution}, split at \(r=s/2\).  In the first
part the two variances are comparable and
\[
 ar+c\frac{(x+r/4)^2}{1+s-r}
 \ge c'\frac{x^2}{1+s}+c'r;
\]
in the second part the factor \(e^{-ar}\), together with the retained
\(x+r/4\) decay, obeys the explicit bound
\[
 ar+c\frac{(x+r/4)^2}{1+s-r}
 \ge c'\left(s+\min\left(\frac{x^2}{1+s},|x|\right)\right)
 \qquad(r>s/2).
\]
Since
\(\sum_{r>s/2}(1+s-r)^{-1/2}e^{-c's}\ll
(1+s)^{-1/2}\), this absorbs the lost square-root prefactor.  For the
exponential-kernel half, the corresponding inequalities are
\[
 ar+c|x+r/4|\ge c'(|x|+r)\quad(r\le s/2),
 \qquad
 ar+c|x+r/4|\ge c'(|x|+s)\quad(r>s/2).
\]
For \eqref{eq:Tao-horizontal-convolution}, split at \(u=|x|/2\).  In the
first range \(|x-u|\ge|x|/2\) whenever \(x\ge0\), while for \(x<0\) it
holds for every \(u\ge1\).  In the complementary range the geometric sum
of \(e^{-au}\) is \(O(e^{-a|x|/2})\).  These statements apply separately
to the Gaussian and exponential halves of \(K_s\) and prove the second
displayed inequality.  In the actual summand \(U(j-u,s-r)\), the kernel
coordinate is \(j-s/4-u+r/4\).  Thus first substitute
\(x=j-s/4-u\) in \eqref{eq:Tao-vertical-convolution}; after that summation,
substitute \(x=j-s/4\) in
\eqref{eq:Tao-horizontal-convolution}.  This proves the first-passage display and
fills the unexpanded calculation at source lines~1457--1463.

For the post-passage assertion, first sum the displayed first-passage bound
over its vertical coordinate.  The horizontal mass function is bounded by
\(C K_s(x-s/4)\), where
\[
 K_s(y)=(1+s)^{-1/2}
 \left(e^{-c y^2/(1+s)}+e^{-c|y|}\right).
\]
If \(\Sigma'\) is \(h\)-separated and \(r_0\le h/4\), the intervals of
radius \(r_0\) about its points are disjoint.  Since \(r_0\ge1\), direct
summation of the Gaussian and exponential pieces gives
\[
 \sum_{\operatorname{dist}(x,\Sigma')\le r_0}K_s(x-s/4)
 \le C\left(\frac {r_0}{\sqrt{1+s}}+\frac {r_0}h\right).
\]
Here \(h\le m^{0.4}=o(\sqrt{s})\), because
\(s>m/\log^2m\), so the right side is at most \(Cr_0/h\).  Conditional on the
sum of the \(p\) later horizontal increments, the target set is merely a
translate of \(\Sigma\) and has the same spacing.  Averaging the uniform
bound proves the assertion.  The condition \(h\le m^{0.4}\) is essential;
without it a singleton at the kernel peak would contradict a bare
\(O(r_0/h)\) claim.
\end{proof}

\begin{proposition}[repaired renewal monotonicity]
\label{prop:Tao-v7-renewal-monotonicity}
Retain \(n,\xi\) and the black--white partition from
Proposition~\ref{prop:Tao-v7-black-triangles}.  There is an absolute
\(\varepsilon_0>0\), fixed in the exit-rectangle calculation below, and we
assume \(0<\varepsilon\le\varepsilon_0\).  Extend \(1_W\) by zero off
\([n/2]\times\mathbb Z\).  For \(x\in\mathbb Z^2\), let
\[
 Q(x)=\mathbb E_x\prod_{q\ge0}e^{-\varepsilon^31_W(X_q)},
 \qquad X_q=x+\mathbf V_q,
\]
and define
\[
 Q_m=\sup_{\substack{j\in\Npos,\ j\ge\lfloor n/2\rfloor-m\\l\in\mathbb Z}}
 \max(\lfloor n/2\rfloor-j,1)^A Q(j,l).
\]
For every fixed sufficiently large \(A\), there is an explicit finite
threshold \(C_{A,\varepsilon}\) such that
\[
 Q_m\le Q_{m-1}\qquad
 (m\in\Npos,\ C_{A,\varepsilon}\le m\le\lfloor n/2\rfloor).
\]
The proof uses the weaker v7 Lemma~7.9 and does not assume
\eqref{eq:Tao-v5-rip}.
\end{proposition}

\begin{proof}
If \(m=\lfloor n/2\rfloor\), the restrictions defining \(Q_m\) and
\(Q_{m-1}\) both reduce to \(j\in\Npos\), so the asserted inequality is an
equality.  Hence suppose \(m\le\lfloor n/2\rfloor-1\).  It then suffices to
treat the genuine boundary row
\(j=\lfloor n/2\rfloor-m\in\Npos\).
If \((j,l)\) is white, the recursion
\[
 Q(j,l)=e^{-\varepsilon^3}\mathbb E Q((j,l)+\mathbf H)
\]
and
\[
 \max(m-r,1)^{-1}\le m^{-1}
 \exp\left(Cr\frac{\log m}{m}\right)
\]
show that the immediate white factor absorbs the horizontal exponential
moment for sufficiently large \(m\).  This is source Case~1.

Suppose next that \((j,l)\) lies at depth
\(0\le s\le m/\log^2m\) in a black triangle.  Stop at \(\kappa_s\) and
apply the recursion once more than the source does:
\[
 Q(j,l)=\mathbb E\left[
 e^{-\varepsilon^3\sum_{q=0}^{\kappa_s}1_W(X_q)}
 Q(X_{\kappa_s+1})\right].
\]
The definition of \(Q_{m-1}\) now applies at the terminal point.  With
\(\lambda=CA\log m/m\),
\[
 Q(j,l)\le m^{-A}Q_{m-1}
 \mathbb E\left[e^{-\varepsilon^31_W(X_{\kappa_s})}
 e^{\lambda(\mathbf V_{\kappa_s+1})_1}\right].
\]
Lemma~\ref{lem:Tao-v7-first-passage-repaired} gives
\(\mathbb E e^{\lambda(\mathbf V_{\kappa_s})_1}
=1+O_A(1/\log m)\).  We now prove, rather than assume, the uniform exit
probability.  Summing that lemma first over
\(l-s>M_v\), and then over
\(|j-s/4|>M_h\sqrt{1+s}\), gives
\[
 \mathbb P\left(
  1\le(\mathbf V_{\kappa_s})_2-s\le M_v,
  \ |(\mathbf V_{\kappa_s})_1-s/4|
       \le M_h\sqrt{1+s}
 \right)\ge\frac12
\]
once the two absolute integers \(M_v,M_h\) are chosen sufficiently large.
Write \(J=(\mathbf V_{\kappa_s})_1\) and
\(L=(\mathbf V_{\kappa_s})_2\).  In this rectangle
\(1\le L-s\le M_v\), so the exit point is above the top row by at most
\(M_v\).  Horizontally,
\[
 j+J-j_\Delta\ge \frac s4-M_h\sqrt{1+s}\ge-C_-(M_h),
\]
while the triangle inequality and \(J\le s/4+M_h\sqrt{1+s}\) give
\begin{align*}
 (j+J-j_\Delta)\log9
 &\le s_\Delta-s\log2+\frac{s\log9}{4}
       +M_h\log9\sqrt{1+s}\\
 &\le s_\Delta+C_+(M_h).
\end{align*}
The last constants are finite because
\(\log2-\frac14\log9>0\), so a negative linear term absorbs the square-root
term.  Thus the horizontal coordinate lies within an absolute distance of
the top-row segment
\([j_\Delta,j_\Delta+s_\Delta/\log9]\).  Combining the horizontal and
vertical bounds gives a Euclidean distance
\(C_{\rm exit}=C(M_v,M_h)\) from the triangle.  Choose the global
\(\varepsilon_0\le e^{-10}\) so that
\[
 \frac1{10}\log(1/\varepsilon_0)>C_{\rm exit}+1.
\]
The white collar then gives
\(\mathbb P(X_{\kappa_s}\in W)\ge c_*\) with \(c_*=1/2\), uniformly in
the triangle, its depth, \(n\), and \(\xi\).

Strong Markov factors the final holding increment.  Its exponential moment
is \(1+O_A(\log m/m)\).  Since the exponential weight accumulated before
that increment is at least one, the remaining expectation is at most
\[
 1+O_A(1/\log m)-c_*(1-e^{-\varepsilon^3}),
\]
before multiplication by the final \(1+o(1)\) factor, and is at most one for
\(m\ge C_{A,\varepsilon}\).  This repairs the unsupported source line~1556.

It remains to treat a starting point at depth
\(s>m/\log^2m\).  We first prove the two exact auxiliary bounds used in this
case.  Membership in the initial triangle and the real right-tip bound give
the exact upper range
\begin{equation}
\label{eq:Tao-deep-depth-upper}
 \frac m{\log^2m}<s\le\frac{\log9}{\log2}\,m.
\end{equation}
Put \(x=(j,l)\), retain the first-passage time \(\kappa_s\), and reset the
post-exit process by
\[
 \widetilde X_p=X_{\kappa_s+p}=x+\mathbf V_{\kappa_s+p},\qquad
 \widetilde{\mathbf V}_p
   =\mathbf V_{\kappa_s+p}-\mathbf V_{\kappa_s}
 \quad(p\in\Nzero).
\]
Conditionally on \(\mathcal F_{\kappa_s}\), the increments of
\(\widetilde{\mathbf V}\) are iid copies of \(\mathbf H\) and are
independent of \(\widetilde X_0\).

We first prove the v7 many-triangles estimate uniformly for a renewal process
\(Z_p=z+\mathbf V_p\) started at an arbitrary \(z\in\mathbb Z^2\).  Define
\(t_1\) as its first black time, allowing \(t_1=0\).  After entry into the
\(i\)-th triangle \(\Delta_i\), let \(\sigma_i>t_i\) be the first time whose
vertical coordinate exceeds \(l_{\Delta_i}\), and let \(t_{i+1}\ge\sigma_i\)
be the first subsequent black time.  Let \(r\in\Nzero\cup\{\infty\}\) be
the number of entries.  The exit-rectangle argument gives, on \(r\ge i\),
\[
 \mathbb P_z(Z_{\sigma_i}\in W\mid\mathcal F_{t_i})\ge c_*.
\]
Set \(\rho=1-(1-e^{-1})c_*\in(0,1)\), and decrease the fixed
\(\varepsilon_0\) once more so that \(e^{\varepsilon_0}\rho\le1\).  For
\(q\in\Nzero\), define the uniformly typed contraction
\[
 H_q(z)=\mathbb E_z\left[
  1_{\{r\ge q+1\}}\prod_{i=1}^{q}e^{-1_W(Z_{\sigma_i})}
 \right].
\]
Since \(H_0(z)\le1\), strong Markov at the first exit gives inductively
\begin{align*}
 H_q(z)
 &=\mathbb E_z\left[
  1_{\{r\ge1\}}e^{-1_W(Z_{\sigma_1})}
  H_{q-1}(Z_{\sigma_1})\right]\\
 &\le\rho^{q-1}\mathbb E_z\left[
  1_{\{r\ge1\}}e^{-1_W(Z_{\sigma_1})}\right]
 \le\rho^q.
\end{align*}
On \(r\ge R\), the distinct exits
\(\sigma_1,\ldots,\sigma_{R-1}\) occur no later than \(t_R\); hence the
complete white count dominates their indicators and
\begin{equation}
\label{eq:Tao-v7-many-triangles-proved}
 \mathbb E_z1_{\{r\ge R\}}
 e^{-\sum_{p=1}^{t_R}1_W(Z_p)+\varepsilon R}
 \le e^{\varepsilon R}\rho^{R-1}\le e^\varepsilon.
\end{equation}
This proves the weaker v7 Lemma~7.9 with a defined time-zero atom and the
required uniform induction.

We next reconstruct the parameterised large-triangle estimate.  For
\(p\in\Nzero\) and \(s'>0\), define the event
\[
 E_{p,s'}=\{\widetilde X_p\in\Delta'
   \text{ for some }\Delta'\in\mathcal T
   \text{ with }s_{\Delta'}\ge s'\}.
\]
Let \(D\ge D_0\), where \(D_0\) is a sufficiently large absolute constant.
The first-passage overshoot bound and the exponential tail of the next
\(p\) increments give
\[
 \mathbb P\bigl((\widetilde X_p)_2
     \ge l_\Delta+2D(1+p)\bigr)
 \le C e^{-cD(1+p)}.
\]
The horizontal first-passage kernel and the same increment tail, using
\(p\le m^{0.1}\), \(s>m/\log^2m\), and \(s'\le m^{0.4}\), give
\[
 \mathbb P\left(
 \left|(\mathbf V_{\kappa_s+p})_1-\frac s4\right|>2s^{0.6}
 \right)
 \le C\frac{D(1+p)}{s'}.
\]
Outside these two exceptional events, membership of \(\widetilde X_p\) in
an eligible \(\Delta'\) forces its lower tip to satisfy
\begin{equation}
\label{eq:Tao-lower-tip-localisation}
 l_\Delta-10<l_{\Delta'}-\frac{s_{\Delta'}}{\log2}
 \le l_\Delta+2D(1+p).
\end{equation}
For if the left inequality failed, moving horizontally by
\(O(D(1+p))\) on the row \(l_\Delta\) would produce a lattice point of
\(\Delta'\); the horizontal bound and
\(\frac14\log9<\log2\) put the same point in the initial \(\Delta\),
contradicting disjointness.  The right inequality follows directly from
the triangle inequality for \(\widetilde X_p\in\Delta'\).  Substituting
\eqref{eq:Tao-lower-tip-localisation} back into that triangle inequality
gives
\[
 |(\widetilde X_p)_1-j_{\Delta'}|\le C D(1+p).
\]
If \(\Delta',\Delta''\) are two eligible triangles, intersecting each with
the row \(l_\Delta+\lfloor s'/2\rfloor\) produces disjoint horizontal
integer intervals of length at least
\(c s'-CD(1+p)\).  Thus, provided
\(s'\ge C_1D(1+p)\), their top-left coordinates are at least \(c_1s'\)
apart.  Decrease the separation constant once, if necessary, so that
\(0<c_1\le1\).  Apply the post-passage part of
Lemma~\ref{lem:Tao-v7-first-passage-repaired} with neighbourhood radius
\(r_0=CD(1+p)\) and spacing \(h=c_1s'\).  Enlarging \(C_1\) enforces
\(4r_0\le h\), and \(s'\le m^{0.4}\) enforces the smoothing-scale bound.
Consequently there are absolute \(C_0,c_0,C_1>0\) such that
\begin{equation}
\label{eq:Tao-large-triangle-parametric}
 \mathbb P(E_{p,s'})\le C_0\left(
 \frac{D(1+p)}{s'}+e^{-c_0D(1+p)}\right)
\end{equation}
whenever
\[
 0\le p\le m^{0.1},\qquad
 C_1D(1+p)\le s'\le m^{0.4}.
\]
This exposes the lower-tip and separated-top steps preceding the
post-\(p\) convolution at source lines~1805--1812.

Now put \(\delta_A=10^{-A-2}\), and choose
\[
 D_A\ge\max\left(D_0,
 \frac1{c_0}\log\frac{6C_0}{\delta_A(1-e^{-c_0})}\right),
 \qquad
 K_A\ge\max\left(C_1D_A,
 \frac{6C_0D_A\zeta(2)}{\delta_A}\right).
\]
Set \(H=10A/\varepsilon^3\) and
\[
 R=\left\lceil
 \frac{H+\log3+(A+2)\log10+\varepsilon+1}{\varepsilon}
 \right\rceil.
\]
Let \(h=\lceil H\rceil\), put \(u_1=h\), and define, for
\(1\le i<R\),
\[
 u_{i+1}=u_i+\left\lceil10K_A(1+u_i)^3\right\rceil+h+1,
 \qquad P=u_R+2.
\]
All of \(D_A,K_A,H,R,h,u_i,P\) depend only on \(A,\varepsilon\).  Enlarge
\(C_{A,\varepsilon}\) so that
\begin{equation}
\label{eq:Tao-case3-applicability}
 P\le m^{0.1},\qquad K_A(1+P)^3\le m^{0.4}.
\end{equation}

Only now define the finite exceptional event
\[
 E_*^{(P)}=\bigcup_{p=0}^{P-1}E_{p,K_A(1+p)^3}.
\]
Every term lies in the domain of
\eqref{eq:Tao-large-triangle-parametric}, and the union bound gives
\begin{align}
 \mathbb P(E_*^{(P)})
 &\le \frac{C_0D_A}{K_A}
       \sum_{p=0}^{P-1}\frac1{(1+p)^2}
   +C_0\sum_{p=0}^{P-1}e^{-c_0D_A(1+p)}\notag\\
 &\le\delta_A/6+\delta_A/6=\delta_A/3.
 \label{eq:Tao-E-star-budget}
\end{align}
This replaces the source's insufficient
\(O(A^2 4^{-A})\): its ratio to \(10^{-A-2}\) is
\(100A^2(10/4)^A\), which is unbounded.

Condition on \(\widetilde X_0\) and let \(r,t_i,\sigma_i\) be the preceding
triangle statistics for the shifted process \((\widetilde X_p)_{p\ge0}\).
Put
\[
 S_R=\sum_{p=1}^{t_R}1_W(\widetilde X_p),\qquad
 F_*=\left\{r\ge R,
 e^{-S_R+\varepsilon R}>3\delta_A^{-1}e^\varepsilon\right\}.
\]
The uniform estimate \eqref{eq:Tao-v7-many-triangles-proved}, integrated in
\(\widetilde X_0\), and Markov's inequality give
\(\mathbb P(F_*)\le\delta_A/3\).  Outside \(F_*\), the event \(r\ge R\)
forces
\[
 S_R\ge\varepsilon R-\log3-\log(\delta_A^{-1})-\varepsilon>H.
\]

Define the high-horizontal event
\[
 G=\{(\mathbf V_{\kappa_s+P})_1\ge0.9m\}.
\]
By \eqref{eq:Tao-deep-depth-upper} and the first-passage kernel,
\[
 \mathbb P((\mathbf V_{\kappa_s})_1\ge0.8m)\le Ce^{-cm},
\]
because \(0.8>\frac14\frac{\log9}{\log2}\).  Strong Markov and the
exponential tail of the fixed sum
\((\widetilde{\mathbf V}_P)_1\) give
\[
 \mathbb P(G)\le C_Pe^{-cm}.
\]
Enlarge \(C_{A,\varepsilon}\) again so that
\begin{equation}
\label{eq:Tao-G-budget}
 \mathbb P(G)\le\delta_A/3,
 \qquad m^A\mathbb P(G)\le\frac12.
\end{equation}
Outside \(G\), positivity of the horizontal increments and
\(j=\lfloor n/2\rfloor-m\) keep every
\(\widetilde X_p\), \(0\le p<P\), inside the strip where
\(B\uplus W\) is the full partition.

Define the shifted white count
\[
 N_W=\sum_{p=0}^{P-1}1_W(\widetilde X_p).
\]
Suppose \(N_W\le H\) and none of \(E_*^{(P)},F_*,G\) occurs.  Among the
\(h+1\) indices \(0,\ldots,h\) at least one is black.  Therefore the actual
first entry time defined above satisfies \(t_1\le h=u_1\).  Induct on the
recursively defined entry times, not on an arbitrary selection of black
points.  If \(t_i\le u_i\), then, outside \(E_*^{(P)}\), its triangle
\(\Delta_i\) satisfies
\[
 s_{\Delta_i}<K_A(1+t_i)^3\le K_A(1+u_i)^3.
\]
Every holding increment has vertical component at least three, so by time
\[
 b_i=t_i+\left\lceil10K_A(1+u_i)^3\right\rceil+1
\]
the process is strictly above \(l_{\Delta_i}\), and in particular
\(b_i\ge\sigma_i\).  The block \(b_i,\ldots,b_i+h\) lies in
\(\{0,\ldots,P-1\}\) and contains a black point; otherwise that block alone
would contribute \(h+1>H\) white points.  By the defining minimality of the
next entry time,
\[
 t_{i+1}\le b_i+h
 \le u_i+\left\lceil10K_A(1+u_i)^3\right\rceil+h+1
 =u_{i+1}.
\]
Thus \(t_{i+1}\) exists and is the entry into the next distinct triangle.
Iteration gives \(R\) entries and \(t_R\le u_R\le P-2\).  Hence
\(S_R\le N_W\le H\), contradicting the preceding consequence outside
\(F_*\).  Therefore
\begin{equation}
\label{eq:Tao-low-white-budget}
 \mathbb P(N_W\le H)
 \le\mathbb P(E_*^{(P)})+\mathbb P(F_*)+\mathbb P(G)
 \le\delta_A.
\end{equation}

It remains to connect this probability estimate to the recursion.  Dropping
the factors before \(\kappa_s\) and then iterating exactly \(P\) times gives
\begin{align*}
 Q(j,l)
 &\le\mathbb E\left[e^{-\varepsilon^3N_W}Q(\widetilde X_P)\right]\\
 &\le m^{-A}Q_{m-1}\,
 \mathbb E\left[e^{-\varepsilon^3N_W}
 \max\left(1-\frac{(\mathbf V_{\kappa_s+P})_1}{m},
                 \frac1m\right)^{-A}\right].
\end{align*}
The second line uses the fact that the terminal horizontal coordinate has
advanced by at least one, so the defining range of \(Q_{m-1}\) applies.
On \(G\) the weight is at most \(m^A\); on \(G^c\) it is at most \(10^A\).
Moreover,
\[
 \mathbb E e^{-\varepsilon^3N_W}
 \le\mathbb P(N_W\le H)+e^{-\varepsilon^3H}
 \le\delta_A+e^{-10A}\le10^{-A-1}
\]
for sufficiently large \(A\).  By \eqref{eq:Tao-G-budget}, the complete
weighted expectation is at most
\[
 m^A\mathbb P(G)+10^A(\delta_A+e^{-10A})\le1
\]
after the stated enlargement of \(C_{A,\varepsilon}\).  This proves
\(Q_m\le Q_{m-1}\) in the final case.
\end{proof}

\begin{corollary}[closure of the v7 fine-scale branch]
\label{cor:Tao-v7-fine-scale-branch}
With the repairs proved above,
\[
 Q(j,l)\ll_A\max(\lfloor n/2\rfloor-j,1)^{-A}
 \quad(j\in\Npos,\ l\in\mathbb Z),
 \qquad \mathbb E Q(\mathbf H)\ll_A n^{-A}.
\]
Consequently the v7 implications
\[
 \text{Proposition 7.8}\Rightarrow
 \text{Proposition 7.3}\Rightarrow
 \text{Proposition 7.1}\Rightarrow
 \text{Proposition 1.17}\Rightarrow
 \text{Proposition 1.14}
\]
hold at their printed uniform scope.  The exact downstream dependency in the
source is
\[
 \text{Proposition 1.14}+\text{Proposition 1.9}
 \Rightarrow\text{Proposition 1.11}
 \Rightarrow\text{Theorem 1.6}
 \Rightarrow\text{Theorem 1.3}.
\]
\end{corollary}

\begin{proof}
Monotonicity and the finite base range bound all \(Q_m\) independently of
\(m\).  The first coordinate of \(\mathbf H\) has exponential tail, so
splitting at \(\mathbf H_1=n/4\) gives
\(\mathbb E Q(\mathbf H)=O_A(n^{-A})\).  This is Proposition~7.3.
The white-point pairing in
Proposition~\ref{prop:Tao-v7-pairing-holding} gives Proposition~7.1, whose
reversed-variable form is Proposition~1.17.  Finally
Proposition~\ref{prop:Tao-v7-Fourier-to-mixing} gives Proposition~1.14.
The second display records every consequence edge leaving this branch.
Proposition~\ref{prop:Tao-v7-Prop19-valuation-law} now supplies the separately
typed Proposition~1.9 input.  The exact step combining it with
Proposition~1.14 to prove Proposition~1.11, and the two later theorem edges,
remain behind their named proof-audit gate.  This is a dependency record, not
a handwaving theorem claim.
\end{proof}

\begin{nonclaim}
Corollary~\ref{cor:Tao-v7-fine-scale-branch} certifies this versioned proof
branch; it does not replace the durable whole-paper release audit.  By itself
it does not certify Proposition~1.9, Proposition~1.11, Theorem~1.6, or
Theorem~1.3; Proposition~\ref{prop:Tao-v7-Prop19-valuation-law} supplies the
first of those results by an independent argument.  This branch does
not prove the stronger v5/journal Lemma~7.9, use Baker's theorem, establish
exponential Fourier decay, prove natural-density mixing of the full affine
map, give a fixed absolute orbit bound, or prove convergence to one.
\end{nonclaim}


\subsection{The valuation-law gate: Tao's Proposition~1.9}

The proposition numbered 1.9 in the controlling v7 text is the distribution
theorem for the finite Syracuse valuation tuple, not the later first-passage
stabilisation theorem.  Its statement is at v7 source lines~269--278 and its
proof occupies lines~597--647 \cite{Tao2026v7}.  Proposition~1.11 is the
first-passage result.  This numbering matters in the dependency graph:
Proposition~1.14 and Proposition~1.9 are the two inputs used to obtain
Proposition~1.11.

The frozen routes identify the v5 and journal manifestations of the same
work.  Direct comparison uses v5 source lines~269--278 and~591--641, v7
source lines~269--278 and~597--647, and journal physical pages~6 and~17--18.
Routing metadata is not used as mathematical evidence.

\begin{sourceissue}[three-manifestation defects in Proposition~1.9]
\label{sourceissue:Tao-Prop19-three-manifestations}
All three inspected manifestations print the same substantive residue-fibre
argument and preserve three operative presentation defects.  In the
statement, the ambient ring is printed as \(\mathbb Z/2^\ell\mathbb Z\),
although \(\ell\) is undefined; both the displayed law and every subsequent
use require \(\mathbb Z/2^{n'}\mathbb Z\).  In the tail proof, one intermediate
line changes the first-crossing event
\(A_k<n'\leq A_{k+1}\) into \(A_k\leq n'<A_{k+1}\); the former is the event
used before and after that line.  Finally, the two central truncations are
printed as \(|\vec a|<m\), although no \(m\) is defined in that section; the
proof requires \(|\vec a|<n'\).

The source also invokes the empty prefix sum \(A_0\) without extending its
earlier nonempty-sum convention, and it does not dispatch \(n=0\).  In the
v5/v7 TeX proof of the prerequisite valuation-uniqueness lemma, a function is
once written without its argument and \(\operatorname{Syr}^{N-1}(N)\) is
printed where the induction requires \(\operatorname{Syr}^{n-1}(N)\).
The reconstruction below supplies all of these data explicitly.
\end{sourceissue}

For \(s\geq1\), put
\[
 \mathcal O_s=(2\mathbb Z+1)/2^s\mathbb Z.
\]
Also set \(\mathcal O_0=\mathbb Z/\mathbb Z\), the singleton residue space,
so that the \(n=0,n'=0\) endpoint in the source statement is typed.  For
\(s\geq1\), \(\mathcal O_s\) has \(2^{s-1}\) elements.  For
\(1\leq t\leq s\), let
\[
 \pi_{s,t}:\mathcal O_s\longrightarrow\mathcal O_t,
 \qquad [u]_{2^s}\longmapsto[u]_{2^t}
\]
be reduction.  Every fibre has exactly \(2^{s-t}\) elements, so
\((\pi_{s,t})_*\operatorname{Unif}(\mathcal O_s)\) is
\(\operatorname{Unif}(\mathcal O_t)\).  For signed measures \(\mu,\nu\) on
\(\mathcal O_s\), the source's unhalved total-variation convention gives the
exact contraction
\begin{equation}
\label{eq:Tao-Prop19-TV-contraction}
\begin{split}
 \| (\pi_{s,t})_*\mu-(\pi_{s,t})_*\nu\|_1
 & =\sum_{x\in\mathcal O_t}
 \left|\sum_{u\in\pi_{s,t}^{-1}(x)}(\mu(u)-\nu(u))\right|\\
 &\leq\sum_{u\in\mathcal O_s}|\mu(u)-\nu(u)|.
\end{split}
\end{equation}

\begin{lemma}[valuation prefix and its exact residue fibre]
\label{lem:Tao-Prop19-residue-fibre}
Let \(n\geq1\), let \(N\in2\mathbb N+1\), and let
\(\vec a=(a_1,\ldots,a_n)\in(\mathbb N+1)^n\).  Define
\[
 A_0=0,\qquad A_j=a_1+\cdots+a_j\quad(1\leq j\leq n),
 \qquad A=A_n,
\]
and
\[
 C_n(\vec a)=\sum_{i=1}^n3^{n-i}2^{A_{i-1}}.
\]
Then
\begin{equation}
\label{eq:Tao-Prop19-affine-numerator}
 \operatorname{Aff}_{\vec a}(N)
 =\frac{3^nN+C_n(\vec a)}{2^A},
\end{equation}
and
\begin{equation}
\label{eq:Tao-Prop19-prefix-equivalence}
 \vec a^{(n)}(N)=\vec a
 \quad\Longleftrightarrow\quad
 \operatorname{Aff}_{\vec a}(N)\in2\mathbb N+1.
\end{equation}
Consequently the event in \eqref{eq:Tao-Prop19-prefix-equivalence} is the
single odd residue cylinder
\begin{equation}
\label{eq:Tao-Prop19-cylinder}
 N\equiv b_{\vec a}\pmod{2^{A+1}},
 \qquad
 b_{\vec a}\equiv
 3^{-n}\bigl(2^A-C_n(\vec a)\bigr)\pmod{2^{A+1}}.
\end{equation}
If \(A<n'\), the inverse image of this cylinder in \(\mathcal O_{n'}\) has
exactly \(2^{n'-A-1}\) elements and hence uniform mass \(2^{-A}\).
\end{lemma}

\begin{proof}
Multiplying the source's affine-offset formula by \(2^A\) gives
\eqref{eq:Tao-Prop19-affine-numerator}.  The forward implication in
\eqref{eq:Tao-Prop19-prefix-equivalence} is the defining Syracuse iteration.
For the converse, induct on \(n\).  When \(n=1\), the equality
\(
  2^{a_1}\operatorname{Aff}_{a_1}(N)=3N+1
\)
and the oddness of \(\operatorname{Aff}_{a_1}(N)\) give
\(a_1=\nu_2(3N+1)\) directly.  Now let \(n\geq2\) and write
\[
 X=\operatorname{Aff}_{a_1,\ldots,a_{n-1}}(N),
 \qquad
 Y=\operatorname{Aff}_{\vec a}(N)\in2\mathbb N+1.
\]
Both affine maps preserve positivity and \(X\in\mathbb Z[1/2]\).  The last
step gives
\[
 2^{a_n}Y=3X+1.
\]
Thus \(3X\) is an odd integer.  A rational number in
\(\mathbb Z[1/2]\) whose product by \(3\) is integral is itself integral;
hence \(X\) is a positive odd integer.  The induction hypothesis identifies
the first \(n-1\) valuation coordinates.  Since
\((3X+1)/2^{a_n}=Y\) is odd, \(a_n=\nu_2(3X+1)\), proving the last coordinate.

An integer quotient in \eqref{eq:Tao-Prop19-affine-numerator} is odd exactly
when its numerator is congruent to \(2^A\) modulo \(2^{A+1}\).  Because
\(3^n\) is a unit modulo \(2^{A+1}\), this gives the unique residue in
\eqref{eq:Tao-Prop19-cylinder}.  It is odd: the first term of
\(C_n(\vec a)\) is \(3^{n-1}\), every other term is even, and \(2^A\) is
even.  Finally, reduction
\(\mathcal O_{n'}\to\mathcal O_{A+1}\) has fibres of size
\(2^{n'-(A+1)}\), proving the last assertion.
\end{proof}

The exponential estimate used in both tails and in the accumulated cylinder
error can be written without an unnamed appeal to Stirling's formula.  Let
\[
 H(p)=-p\log p-(1-p)\log(1-p),
 \qquad
 p_0=\frac1{2+c_0/2}<\frac12,
 \qquad
 \eta(c_0)=\log2-H(p_0)>0.
\]
If \(M=n'-1\), \(n'\geq(2+c_0)n\), and \(n\geq2/c_0\), then
\(n/M\leq p_0\).  Set \(z_0=p_0/(1-p_0)\in(0,1)\).  Since
\(z_0^k\geq z_0^n\) for \(0\leq k\leq n\),
\begin{equation}
\label{eq:Tao-Prop19-binomial-bound}
 2^{-M}\sum_{k=0}^n\binom Mk
 \leq z_0^{-n}\left(\frac{1+z_0}{2}\right)^M
 \leq \exp\bigl(-\eta(c_0)M\bigr).
\end{equation}
The finitely many positive \(n<2/c_0\) are absorbed into a constant depending
on \(c_0\).  This proves a bound \(C_{c_0}2^{-c_1n}\), for example with any
fixed \(0<c_1<\eta(c_0)/\log2\).

\begin{proposition}[Tao Proposition~1.9, residue-fibre reconstruction]
\label{prop:Tao-v7-Prop19-valuation-law}
Let \(n\in\mathbb N\), let \(\mathbf N\) take values in
\(2\mathbb N+1\), let \(c_0>0\), and let
\(n'\in\mathbb N\) satisfy \(n'\geq(2+c_0)n\).  Suppose, with the unhalved
total variation used by the source, that
\begin{equation}
\label{eq:Tao-Prop19-hypothesis}
 d_{\mathrm{TV}}\left(
 \mathbf N\bmod2^{n'},\operatorname{Unif}(\mathcal O_{n'})
 \right)\leq K2^{-n'}
\end{equation}
for a hypothesis constant \(K=K(c_0)\geq0\), independent of
\(n,n'\), and the law of \(\mathbf N\).  Then there are
\(c_1=c_1(c_0)>0\) and \(C=C(c_0)>0\) such that
\begin{equation}
\label{eq:Tao-Prop19-conclusion}
 d_{\mathrm{TV}}\left(
 \vec a^{(n)}(\mathbf N),\operatorname{Geom}(2)^n
 \right)\leq C(1+K)2^{-c_1n}.
\end{equation}
Thus the source statement follows with exactly the source-permitted
\(c_0\)-dependence of the implicit constants.
\end{proposition}

\begin{proof}
For \(n=0\), both laws are the point mass at the empty tuple, so the distance
is zero.  Assume \(n\geq1\), and write
\[
 \vec a^{(n)}(\mathbf N)=(\mathsf a_1,\ldots,\mathsf a_n),
 \qquad
 S_0=0,\qquad S_j=\mathsf a_1+\cdots+\mathsf a_j.
\]
Because every \(\mathsf a_j\) is positive, the actual-valuation tail is the
disjoint first-crossing union
\begin{equation}
\label{eq:Tao-Prop19-first-crossing}
 \{S_n\geq n'\}
 =\bigsqcup_{k=0}^{n-1}\{S_k<n'\leq S_{k+1}\}.
\end{equation}
 Fix \(k\) and a positive prefix \((a_1,\ldots,a_k)\), put
 \(A_0=0\) and \(A_j=\sum_{r=1}^{j}a_r\) for \(1\leq j\leq k\), and assume
 \(A_k<n'\).
On the corresponding crossing event, the exact affine numerator at time
\(k+1\) is divisible by \(2^{S_{k+1}}\), hence
\begin{equation}
\label{eq:Tao-Prop19-crossing-residue}
 3^{k+1}\mathbf N+
 \sum_{i=1}^{k+1}3^{k+1-i}2^{A_{i-1}}
 \equiv0\pmod{2^{n'}}.
\end{equation}
The displayed sum depends only on the fixed \(k\)-prefix, including
\(A_0=0\).  Its first term is odd and all later terms are even, so
\eqref{eq:Tao-Prop19-crossing-residue} fixes one odd residue of
\(\mathbf N\) modulo \(2^{n'}\).  Its probability is at most
\((2+K)2^{-n'}\): the uniform mass of one element of \(\mathcal O_{n'}\) is
\(2^{1-n'}\), and \eqref{eq:Tao-Prop19-hypothesis} controls the error.

Positive \(k\)-tuples with \(A_k<n'\) correspond bijectively to the
\(k\)-element subsets
\[
 \{A_1<\cdots<A_k\}\subseteq\{1,\ldots,n'-1\},
\]
and therefore number \(\binom{n'-1}{k}\), including the one empty prefix at
\(k=0\).  Summing \eqref{eq:Tao-Prop19-first-crossing} and applying
\eqref{eq:Tao-Prop19-binomial-bound} gives
\begin{equation}
\label{eq:Tao-Prop19-actual-tail}
 \mathbb P(S_n\geq n')\ll_{c_0}(1+K)2^{-c_1n}.
\end{equation}

For the model law, realize independent \(\operatorname{Geom}(2)\)
coordinates as waiting times between successive successes in fair Bernoulli
trials.  Their sum is at least \(n'\) exactly when the first \(n'-1\) trials
contain at most \(n-1\) successes.  Hence
\begin{equation}
\label{eq:Tao-Prop19-model-tail}
 \mathbb P\left(|\operatorname{Geom}(2)^n|\geq n'\right)
 =2^{-(n'-1)}\sum_{k=0}^{n-1}\binom{n'-1}{k}
 \ll_{c_0}2^{-c_1n}.
\end{equation}

It remains to compare the two laws on positive tuples
\(\vec a=(a_1,\ldots,a_n)\) with \(A=|\vec a|<n'\).  Lemma
\ref{lem:Tao-Prop19-residue-fibre}, the exact pushforward
\eqref{eq:Tao-Prop19-TV-contraction}, and
\eqref{eq:Tao-Prop19-hypothesis} give
\begin{equation}
\label{eq:Tao-Prop19-point-mass}
 \left|
 \mathbb P(\vec a^{(n)}(\mathbf N)=\vec a)-2^{-A}
 \right|\leq K2^{-n'}.
\end{equation}
The term \(2^{-A}\) is exactly the mass of \(\vec a\) under
\(\operatorname{Geom}(2)^n\).  Positive \(n\)-tuples with \(A<n'\) are in
bijection with \(n\)-element subsets of \(\{1,\ldots,n'-1\}\), so their
number is \(\binom{n'-1}{n}\).  Thus their accumulated error is bounded by
\[
 K2^{-n'}\binom{n'-1}{n}\ll_{c_0}K2^{-c_1n}
\]
using \eqref{eq:Tao-Prop19-binomial-bound}.  Adding the two tail bounds
\eqref{eq:Tao-Prop19-actual-tail} and
\eqref{eq:Tao-Prop19-model-tail} is justified by the exact decomposition
\[
\begin{split}
 d_{\mathrm{TV}}\!\left(
 \vec a^{(n)}(\mathbf N),\operatorname{Geom}(2)^n\right)
 &\leq
 \sum_{\substack{\vec a\in(\mathbb N+1)^n\\|\vec a|<n'}}
 \left|\mathbb P(\vec a^{(n)}(\mathbf N)=\vec a)-2^{-|\vec a|}\right|\\
 &\quad+
 \mathbb P(S_n\geq n')
 +\mathbb P\!\left(|\operatorname{Geom}(2)^n|\geq n'\right).
\end{split}
\]
This proves \eqref{eq:Tao-Prop19-conclusion}.
\end{proof}

The proof exposes the exact morphism suppressed in the printed argument.
Define
\[
 \mathcal U_{n,n'}=
 \left\{u\in\mathcal O_{n'}:
 \begin{array}{l}
 \text{some positive odd representative \(N\) of \(u\) satisfies}\\[-2pt]
 |\vec a^{(n)}(N)|<n'
 \end{array}\right\}.
\]
Lemma~\ref{lem:Tao-Prop19-residue-fibre} shows that the displayed condition
and the resulting tuple are independent of the chosen positive
representative: a tuple of total \(A<n'\) is determined already modulo
\(2^{A+1}\).  Hence
\[
 V_{n,n'}:\mathcal U_{n,n'}\longrightarrow\mathcal D_{n,n'},
 \qquad
 V_{n,n'}([N]_{2^{n'}})=\vec a^{(n)}(N),
\]
is a well-defined map, where
\[
 \mathcal D_{n,n'}=
 \{\vec a\in(\mathbb N+1)^n:|\vec a|<n'\}.
\]
It is surjective, and its fibre over \(\vec a\) is precisely the cylinder
\(b_{\vec a}\pmod{2^{A+1}}\), containing
\(2^{n'-A-1}\) odd residues modulo \(2^{n'}\).  Its set-theoretic kernel pair
is
\[
 \mathcal U_{n,n'}\mathbin{\times_{\mathcal D_{n,n'}}}\mathcal U_{n,n'}
 =\{(u,v):V_{n,n'}(u)=V_{n,n'}(v)\},
\]
namely equality of the first \(n\) valuation coordinates.  There is an
explicit section: for each \(\vec a\), take the least positive odd
representative \(\widetilde b_{\vec a}\) of the cylinder
\(b_{\vec a}\pmod{2^{A+1}}\) and put
\[
 s_{n,n'}(\vec a)=[\widetilde b_{\vec a}]_{2^{n'}},
 \qquad V_{n,n'}\circ s_{n,n'}=\operatorname{id}_{\mathcal D_{n,n'}}.
\]
This section chooses all unrecovered higher binary digits to be zero.  The map
has no left inverse, and therefore no two-sided inverse: the tuple
\((1,\ldots,1)\) has \(A=n\) and fibre size
\(2^{n'-n-1}\geq2\), because
\(n'\geq(2+c_0)n\) with \(n\geq1\).

\begin{nonclaim}
Proposition~\ref{prop:Tao-v7-Prop19-valuation-law} certifies only the
valuation-law theorem numbered Proposition~1.9.  It neither assumes nor
proves first-passage stabilisation.  The separate edge combining this result
with Proposition~1.14 to obtain Proposition~1.11 must still be checked before
Theorem~1.6 or Theorem~1.3 is admitted as proved in this edition.  No
natural-density statement, convergence-to-one statement, or fixed universal
orbit bound follows from this proposition.
\end{nonclaim}


\subsection{The first-passage gate: Tao's Proposition~1.11}

The controlling statement is Proposition~1.11 at v7 source
lines~316--335, and the deduction from Propositions~1.9 and~1.14 occupies
v7 lines~649--925 \cite{Tao2026v7}.  The corresponding comparison loci are
v5 lines~643--919 and journal physical pages~18--27.  The three
manifestations have the same substantive modular, logarithmic-mass, and
integer-window calculations.  ArXiv v7 corrects the v5/journal interchange
of passage time and passage location in the converse implication (v5
lines~789--791; journal physical page~23; v7 lines~795--797).  No conclusion
below is transferred silently between versions.

\begin{sourceissue}[local defects at the Proposition~1.11 edge]
\label{sourceissue:Tao-Prop111-edge}
The opening heuristic at v7 lines~651--652 writes
\(\operatorname{Pass}_x\) where the quantity being approximated is the time
\(T_x\), and its numerator requires \(\log(N_y/x)\).  Line~689 has
\(1.9n\) where the fixed descent time is \(1.9n_0\).  The periodic-orbit
aside at line~699 requires a representative of \(-\mathbf m\) modulo the
least positive period, or a period multiple at least \(\mathbf m\); that
aside is not used here.

The additive endpoint deletion printed at lines~759--762 does not produce a
time margin of order \((\log x)^{0.8}\).  It has already been replaced by the
multiplicative deletion in
Proposition~\ref{prop:Tao-v7-endpoint-buffer}.  The unsigned lower comparison,
the event equality, the undefined \(m\) in the tuple length, the fixed-
\(M\) affine equality, the two previous-iterate coefficients, the malformed
progression estimate, the nonuniform weakened \(c_n\)-majorant, and the
iid-versus-actual Chernoff citation have already been repaired in
Proposition~\ref{prop:Tao-v7-Prop52-cn-repair}.  This subsection supplies the
remaining modular law, descent exponent, affine inverse, integer count,
normalization cancellation, and arbitrary-event-to-total-variation arrows.
\end{sourceissue}

Throughout this subsection retain the source constants and floors
\begin{equation}
\label{eq:Tao-Prop111-parameters}
 \alpha=\frac{1001}{1000},\qquad
 L=\log x,\qquad \lambda=\log(4/3),\qquad
 n_0=\left\lfloor\frac{L}{10\log2}\right\rfloor,
 \qquad m_0=\left\lfloor\frac{\alpha-1}{100}L\right\rfloor.
\end{equation}
For \(j\in\{1,2\}\), put \(y_j=x^{\alpha^j}\) and
\begin{equation}
\label{eq:Tao-Prop111-support-normalizer}
 \mathcal R_{y_j}=(2\mathbb N+1)\cap[y_j,y_j^\alpha],
 \qquad
 H_{y_j}=\sum_{N\in\mathcal R_{y_j}}\frac1N,
 \qquad
 \mathbb P(\mathbf N_{y_j}=N)=\frac{1}{H_{y_j}N}.
\end{equation}
Real endpoints cause no ambiguity: the intersection is with the integer set
\(2\mathbb N+1\).

\begin{lemma}[progression mass with real endpoints]
\label{lem:Tao-Prop111-progression-mass}
Let \(1\le A<B\), let \(q\ge2\) be even, put
\(J=\log(B/A)\), and fix an odd residue \(r\pmod q\).  Then
\begin{equation}
\label{eq:Tao-Prop111-one-progression}
 \left|
  \sum_{\substack{A\le N\le B\\N\equiv r\ (\mathrm{mod}\ q)}}\frac1N
  -\frac Jq
 \right|\le\frac1A.
\end{equation}
If \(W_r\) denotes the sum in
\eqref{eq:Tao-Prop111-one-progression} and
\(H=\sum_{r\ {\rm odd}\ (q)}W_r\), then
\begin{equation}
\label{eq:Tao-Prop111-total-progression-error}
 \sum_{r\ {\rm odd}\ (q)}\left|W_r-\frac Jq\right|
 \le\frac q{2A},
 \qquad
 \left|H-\frac J2\right|\le\frac q{2A}.
\end{equation}
\end{lemma}

\begin{proof}
If the progression contributes points
\(t_0<t_1<\cdots<t_s\), consecutive points differ by \(q\).  Monotonicity
of \(1/t\) gives
\[
 \sum_{i=0}^s\frac1{t_i}
 \le \frac1{t_0}+\frac1q\int_{t_0}^{t_s}\frac{dt}{t}
 \le \frac1A+\frac Jq.
\]
It also gives
\[
 \sum_{i=0}^s\frac1{t_i}
 \ge \frac1q\int_{t_0}^{t_s+q}\frac{dt}{t}
 \ge \frac Jq-\frac1q\log\frac{t_0}{A}
 \ge \frac Jq-\frac1A,
\]
because \(0\le t_0-A<q\) and
\(\log(t_0/A)\le(t_0-A)/A\).  If the progression contributes no point,
then \(B-A<q\), whence \(J/q<(B-A)/(qA)<1/A\).  This proves
\eqref{eq:Tao-Prop111-one-progression}.  There are exactly \(q/2\) odd
residue classes.  Summing and using the triangle inequality proves
\eqref{eq:Tao-Prop111-total-progression-error}.
\end{proof}

\begin{proposition}[exact modular input to Proposition~1.9]
\label{prop:Tao-Prop111-modular-input}
Let \(q=2^{3n_0}\).  Uniformly for \(j=1,2\),
\begin{equation}
\label{eq:Tao-Prop111-modular-law}
 d_{\mathrm{TV}}\!\left(
  \mathbf N_{y_j}\bmod q,
  \operatorname{Unif}((2\mathbb Z+1)/q\mathbb Z)
 \right)\ll q^{-1}=2^{-3n_0},
\end{equation}
where total variation is the source's unhalved \(\ell^1\) distance.  Hence
Proposition~\ref{prop:Tao-v7-Prop19-valuation-law}, with
\[
 n=n_0,\qquad n'=3n_0,\qquad c_0=1,
\]
gives
\begin{equation}
\label{eq:Tao-Prop111-valuations}
 d_{\mathrm{TV}}\!\left(
  \vec a^{(n_0)}(\mathbf N_{y_j}),
  \operatorname{Geom}(2)^{n_0}
 \right)\ll x^{-c_2}
\end{equation}
for one absolute \(c_2>0\).
\end{proposition}

\begin{proof}
Set \(A=y_j\), \(B=y_j^\alpha\), and
\(J=(\alpha-1)\log y_j\).  For every odd \(r\pmod q\), let \(W_r\) be
the sum in Lemma~\ref{lem:Tao-Prop111-progression-mass}.  For sufficiently
large \(x\), \(AJ\ge2q\), and therefore
\(H_{y_j}\ge J/4\).  Since uniform measure assigns mass \(2/q\), not
\(1/q\), to each of the \(q/2\) odd classes,
\begin{align}
 d_{\mathrm{TV}}
 &=\sum_{r\ {\rm odd}\ (q)}
   \left|\frac{W_r}{H_{y_j}}-\frac2q\right| \notag\\
 &\le \frac1{H_{y_j}}
 \left(
  \sum_{r\ {\rm odd}\ (q)}\left|W_r-\frac Jq\right|
  +\left|H_{y_j}-\frac J2\right|
 \right)
 \le\frac{2q}{AJ-q}\le\frac{4q}{AJ}.
\label{eq:Tao-Prop111-modular-L1-bound}
\end{align}
The floor in \(n_0\) is retained exactly:
\[
 \frac{x^{3/10}}8\le q\le x^{3/10}.
\]
Consequently
\[
 \frac{4q}{AJ}
 \le q^{-1}
 \frac{4x^{0.6-\alpha}}
      {(\alpha-1)\alpha L}
 =o(q^{-1}),
\]
uniformly for \(j=1,2\).  This proves
\eqref{eq:Tao-Prop111-modular-law} with room to spare.  The size hypothesis
in Proposition~\ref{prop:Tao-v7-Prop19-valuation-law} is the exact equality
\(3n_0=(2+c_0)n_0\).  Finally,
\[
 2^{-c_1n_0}\le2^{c_1}x^{-c_1/10},
\]
so its conclusion has the form \eqref{eq:Tao-Prop111-valuations} despite the
floor.
\end{proof}

\begin{proposition}[the finite-time descent estimate]
\label{prop:Tao-Prop111-finite-descent}
There is an absolute \(c_3>0\) such that, uniformly for \(j=1,2\),
\begin{equation}
\label{eq:Tao-Prop111-failure}
 \mathbb P\bigl(T_x(\mathbf N_{y_j})=+\infty\bigr)\ll x^{-c_3}.
\end{equation}
More precisely, outside an event of probability \(O(x^{-c_3})\),
\(T_x(\mathbf N_{y_j})\le n_0\).
\end{proposition}

\begin{proof}
For iid \(G_1,\ldots,G_{n_0}\sim\operatorname{Geom}(2)\), the lower-tail
Chernoff bound applied to
\[
 \mathbb E e^{-tG_1}=\frac{e^{-t}}{2-e^{-t}}
\]
gives
\(\mathbb P(G_1+\cdots+G_{n_0}\le1.9n_0)\le e^{-\gamma n_0}\)
for an absolute \(\gamma>0\).  For example, take
\(e^{-t}=18/19\); the base is
\(0.9(19/18)^{1.9}<1\).  Combining this with
\eqref{eq:Tao-Prop111-valuations} gives
\begin{equation}
\label{eq:Tao-Prop111-low-sum}
 \mathbb P\!\left(
  |\vec a^{(n_0)}(\mathbf N_{y_j})|\le1.9n_0
 \right)\ll x^{-c_3}
\end{equation}
after decreasing \(c_3\) if necessary.

On the complementary event, the exact affine formula and
\(0\le F_{n_0}\le3^{n_0}\) give
\begin{equation}
\label{eq:Tao-Prop111-descent-bound}
 \operatorname{Syr}^{n_0}(\mathbf N_{y_j})
 \le 3^{n_0}2^{-1.9n_0}x^{\alpha^3}+3^{n_0}.
\end{equation}
The numerical exponent check is strict, not heuristic.  Put
\[
 \beta_1=\alpha^3+
  \frac{\log3-1.9\log2}{10\log2},
 \qquad
 \beta_2=\frac{\log3}{10\log2}.
\]
Rational interval arithmetic gives
\begin{equation}
\label{eq:Tao-Prop111-exponents}
 \beta_1<0.972<1,
 \qquad \beta_2<0.159<1.
\end{equation}
One elementary certificate for these decimal inequalities uses
\[
 \log u=2\sum_{k=0}^{M}\frac{z^{2k+1}}{2k+1}+R_M,
 \quad z=\frac{u-1}{u+1},
 \quad 0<R_M<\frac{2z^{2M+3}}{(2M+3)(1-z^2)},
\]
with \(u=2,3\), \(z=1/3,1/2\), and \(M=20\); every endpoint is rational.
Since
\(n_0\ge L/(10\log2)-1\) and
\(\log3-1.9\log2<0\), the first term on the right side of
\eqref{eq:Tao-Prop111-descent-bound} is at most
\[
 e^{1.9\log2-\log3}x^{\beta_1}.
\]
The positive coefficient in the second term instead uses
\(n_0\le L/(10\log2)\), and that term is at most \(x^{\beta_2}\).
Thus
\(e^{1.9\log2-\log3}x^{\beta_1}+x^{\beta_2}<x\) for all sufficiently
large \(x\).
Thus passage occurs by time \(n_0\), and
\eqref{eq:Tao-Prop111-low-sum} proves the claim.
\end{proof}

For the passage-location part, let \(E\subseteq(2\mathbb N+1)\cap[1,x]\).
For either choice of \(y=y_j\), retain the source interval
\begin{equation}
\label{eq:Tao-Prop111-Iy}
 I_y=\left[
  \frac{\log(y/x)}\lambda+L^{4/5},
  \frac{\log(y^\alpha/x)}\lambda-L^{4/5}
 \right]
\end{equation}
and the source set, independent of \(y\),
\begin{equation}
\label{eq:Tao-Prop111-Eprime}
\begin{split}
 E'=\{M\in2\mathbb N+1:\;&T_x(M)=m_0,
 \ \operatorname{Pass}_x(M)\in E,\\
 &e^{-L^{7/10}}(4/3)^{m_0}x
 \le M\le
 e^{L^{7/10}}(4/3)^{m_0}x\}.
\end{split}
\end{equation}
For an integer \(n\in I_y\), put \(r=n-m_0\).  For sufficiently large
\(x\), every such \(n\) satisfies
\begin{equation}
\label{eq:Tao-Prop111-range}
 m_0\le r\le n_0-m_0,
 \qquad 1\le m_0\le r.
\end{equation}

\begin{proposition}[the singleton affine inverse]
\label{prop:Tao-Prop111-affine-inverse}
Let \(n\in I_y\cap\mathbb Z\), \(r=n-m_0\), and
\(\vec a=(a_1,\ldots,a_r)\in\mathcal A^{(r)}\).  Put
\[
 A_0=0,
 \qquad A_k=a_1+\cdots+a_k,
 \qquad A=A_r,
 \qquad
 C_r(\vec a)=\sum_{i=1}^r3^{r-i}2^{A_{i-1}},
 \qquad F_r(\vec a)=2^{-A}C_r(\vec a).
\]
For \(M\in E'\), the equation
\(\operatorname{Aff}_{\vec a}(N)=M\) has at most one solution, namely
\begin{equation}
\label{eq:Tao-Prop111-Phi}
 \Phi_{r}(\vec a,M)
 =\frac{2^AM-C_r(\vec a)}{3^r}
 =2^A\frac{M-F_r(\vec a)}{3^r}.
\end{equation}
It has a solution in \(\mathcal R_y\) if and only if
\begin{equation}
\label{eq:Tao-Prop111-congruence}
 M\equiv F_r(\vec a)\pmod{3^r}
\end{equation}
in \(\mathbb Z[1/2]/3^r\mathbb Z[1/2]\).  When this congruence holds,
\(\Phi_r(\vec a,M)\) is a positive odd integer in \([y,y^\alpha]\),
\[
 \vec a^{(r)}(\Phi_r(\vec a,M))=\vec a,
 \qquad
 \operatorname{Syr}^r(\Phi_r(\vec a,M))=M,
\]
and its logarithmic point mass is
\begin{equation}
\label{eq:Tao-Prop111-point-mass}
 \mathbb P\!\left(
  \operatorname{Aff}_{\vec a}(\mathbf N_y)=M
 \right)
 =\frac{1+O(x^{-c_4})}{\frac{\alpha-1}{2}\log y}
  \frac{2^{-A}3^r}{M}
\end{equation}
for one absolute \(c_4>0\), uniformly in all displayed data.
\end{proposition}

\begin{proof}
The affine identity
\[
 \operatorname{Aff}_{\vec a}(N)=\frac{3^rN+C_r(\vec a)}{2^A}
\]
proves uniqueness and forces \eqref{eq:Tao-Prop111-Phi}.  Congruence
\eqref{eq:Tao-Prop111-congruence} is necessary.  Conversely, it makes
\eqref{eq:Tao-Prop111-Phi} an integer: \(M-F_r(\vec a)\) belongs to
\(3^r\mathbb Z[1/2]\), and its denominator divides \(2^A\).  The integer is
odd because \(C_r(\vec a)\) is odd, \(2^AM\) is even, and \(3^r\) is odd.

It remains to prove positivity and the interval bounds, rather than infer
them from presentation.  Since \(r\le n_0\),
\[
 0\le F_r(\vec a)\le3^r\le x^{0.159}
\]
by \eqref{eq:Tao-Prop111-exponents}.  From \eqref{eq:Tao-Prop111-Eprime},
\(M\ge e^{-L^{7/10}}x\).  Since
\(L^{7/10}\le(141/1000)L\) for all sufficiently large \(L\), the preceding
two bounds give
\begin{equation}
\label{eq:Tao-Prop111-offset-ratio}
 0\le\frac{F_r(\vec a)}M\le x^{-7/10}<\frac12.
\end{equation}
Write \(A=2r+\delta\) and
\(M=x(4/3)^{m_0}e^\eta\).  Membership in
\(\mathcal A^{(r)}\) and \(E'\) gives
\(|\delta|<L^{3/5}\) and \(|\eta|\le L^{7/10}\).  Hence
\begin{equation}
\label{eq:Tao-Prop111-Phi-log}
 \log\Phi_r(\vec a,M)
 =L+n\lambda+\delta\log2+\eta+
   \log\left(1-\frac{F_r(\vec a)}M\right).
\end{equation}
The final three terms have absolute value at most \(2L^{7/10}\) for all
sufficiently large \(x\).  The endpoints of \(I_y\) give margins
\(\lambda L^{4/5}\), and
\(2L^{7/10}<\lambda L^{4/5}\) eventually.  Thus
\(\log y\le\log\Phi_r\le\log y^\alpha\), proving both positivity and the
claimed range.  Substitution proves
\(\operatorname{Aff}_{\vec a}(\Phi_r)=M\); the exact valuation-prefix lemma,
Lemma~\ref{lem:Tao-Prop19-residue-fibre}, then proves the valuation and
iterate identities.

Finally, Lemma~\ref{lem:Tao-Prop111-progression-mass} with modulus \(2\)
gives
\begin{equation}
\label{eq:Tao-Prop111-harmonic-normalizer}
 \left|H_y-\frac{\alpha-1}{2}\log y\right|\le\frac1y,
 \qquad
 H_y=\frac{\alpha-1}{2}\log y\,(1+O(x^{-1})).
\end{equation}
Using \eqref{eq:Tao-Prop111-Phi} exactly,
\[
 \frac1{\Phi_r(\vec a,M)}
 =\frac{2^{-A}3^r}{M-F_r(\vec a)}.
\]
Equations~\eqref{eq:Tao-Prop111-offset-ratio} and
\eqref{eq:Tao-Prop111-harmonic-normalizer} now give
\eqref{eq:Tao-Prop111-point-mass}.
\end{proof}

The last proposition is an exact morphism, not merely a pointwise formula.
Define
\[
 \mathcal D_{r,y}=
 \{(\vec a,M)\in\mathcal A^{(r)}\times E':
       M\equiv F_r(\vec a)\pmod{3^r}\}
\]
and
\[
 \mathcal X_{r,y}=
 \{N\in\mathcal R_y:
   \vec a^{(r)}(N)\in\mathcal A^{(r)},\
   \operatorname{Syr}^r(N)\in E'\}.
\]
Then
\begin{equation}
\label{eq:Tao-Prop111-bijection}
 \Theta_{r,y}:\mathcal X_{r,y}\longrightarrow\mathcal D_{r,y},
 \qquad
 N\longmapsto
 \bigl(\vec a^{(r)}(N),\operatorname{Syr}^r(N)\bigr)
\end{equation}
is a bijection with inverse \(\Phi_r\).  Both compositions are identities;
every fibre is a singleton; its set-theoretic kernel pair is the diagonal;
and no coordinate or higher binary digit is discarded.  This is the exact
map underlying source lines~808--827.

\begin{lemma}[the exact integer time-window count]
\label{lem:Tao-Prop111-window-count}
For \(y_j=x^{\alpha^j}\), let the endpoints of \(I_{y_j}\) be
\[
 a_j=\frac{\alpha^j-1}{\lambda}L+L^{4/5},
 \qquad
 b_j=\frac{\alpha^{j+1}-1}{\lambda}L-L^{4/5}.
\]
For sufficiently large \(x\), the interval is nonempty and
\begin{equation}
\label{eq:Tao-Prop111-exact-count}
\begin{split}
 \#(I_{y_j}\cap\mathbb Z)
 &=\lfloor b_j\rfloor-\lceil a_j\rceil+1\\
 &=\frac{\alpha-1}{\lambda}\log y_j-2L^{4/5}+\varepsilon_j,
 \qquad -1<\varepsilon_j\le1.
\end{split}
\end{equation}
Consequently
\begin{equation}
\label{eq:Tao-Prop111-normalization-cancellation}
 \frac{\#(I_{y_j}\cap\mathbb Z)}
      {\frac{\alpha-1}{2}\log y_j}
 =\frac2\lambda
  -\frac4{(\alpha-1)\alpha^j}L^{-1/5}
  +O(L^{-1})
 =\frac2\lambda+O(L^{-1/5}).
\end{equation}
The leading term is independent of \(j\); the displayed finite-
\(x\) buffer correction is not.
\end{lemma}

\begin{proof}
The first line is the exact count of integer points in a closed real
interval.  Subtracting the endpoints gives the second line with
\[
 \varepsilon_j=(\lfloor b_j\rfloor-b_j)
 +(a_j-\lceil a_j\rceil)+1,
 \qquad -1<\varepsilon_j\le1.
\]
Since \(\log y_j=\alpha^jL\), division proves
\eqref{eq:Tao-Prop111-normalization-cancellation}.  It also shows why the
source's variable exponent \(c\) must be decreased to at most \(1/5\) at
this step.  The lower endpoint satisfies \(a_j\ge2m_0\), and
\[
 \frac{\alpha^3-1}{\lambda}<0.010439
 <0.144269<\frac1{10\log2};
\]
therefore \(I_y\subseteq[2m_0,n_0]\) for all sufficiently large \(x\),
which also proves \eqref{eq:Tao-Prop111-range} with every floor retained.
\end{proof}

\begin{theorem}[Tao Proposition~1.11, exact reconstructed edge]
\label{thm:Tao-v7-Prop111-first-passage}
Let \(\alpha=1.001\), and let \(\mathbf N_y\) have the logarithmic law on
the odd integers in \([y,y^\alpha]\).  There is an absolute \(c>0\) such
that, for all sufficiently large \(x\),
\begin{equation}
\label{eq:Tao-Prop111-theorem-failure}
 \mathbb P\bigl(T_x(\mathbf N_y)=+\infty\bigr)\ll x^{-c}
 \qquad (y=x^\alpha,x^{\alpha^2})
\end{equation}
and
\begin{equation}
\label{eq:Tao-Prop111-theorem-TV}
 d_{\mathrm{TV}}\!\left(
  \operatorname{Pass}_x(\mathbf N_{x^\alpha}),
  \operatorname{Pass}_x(\mathbf N_{x^{\alpha^2}})
 \right)\ll L^{-c}
\end{equation}
in the source's unhalved convention.
\end{theorem}

\begin{proof}
Equation~\eqref{eq:Tao-Prop111-theorem-failure} is
Proposition~\ref{prop:Tao-Prop111-finite-descent}.  Fix an arbitrary
\(E\subseteq(2\mathbb N+1)\cap[1,x]\), form \(E'\) by
\eqref{eq:Tao-Prop111-Eprime}, and define
\begin{equation}
\label{eq:Tao-Prop111-ZE}
 Z_E=\sum_{M\in E'}\frac{3^{m_0}}M
 \mathbb P\!\left(
  \operatorname{Syrac}(\mathbb Z/3^{m_0}\mathbb Z)
  \equiv M\pmod{3^{m_0}}
 \right).
\end{equation}
This quantity depends on \(x\) and \(E\), but not on whether
\(y=x^\alpha\) or \(y=x^{\alpha^2}\).

Proposition~\ref{prop:Tao-v7-endpoint-buffer} places the first-passage time
in \(I_y\) outside an \(O(L^{-1/5})\) event.  The exact event equality and
signed pre-passage estimate in
Proposition~\ref{prop:Tao-v7-Prop52-cn-repair}, together with the bijection
\eqref{eq:Tao-Prop111-bijection} and point mass
\eqref{eq:Tao-Prop111-point-mass}, give
\begin{equation}
\label{eq:Tao-Prop111-pre-mixing}
\begin{split}
 \mathbb P(\operatorname{Pass}_x(\mathbf N_y)\in E)
 ={}&\frac{1+O(x^{-c_4})}
          {\frac{\alpha-1}{2}\log y}
 \sum_{n\in I_y\cap\mathbb Z} A_{n,E}
 +O(L^{-c_5}),
\end{split}
\end{equation}
where, with \(r=n-m_0\),
\[
 A_{n,E}=3^r
 \sum_{\vec a\in\mathcal A^{(r)}}2^{-|\vec a|}
 \sum_{\substack{M\in E'\\M\equiv F_r(\vec a)\ (3^r)}}\frac1M.
\]
All sums are over their displayed discrete sets; no conditional law or
coupling is being inserted.

The uniform \(c_n\)-bound in
Proposition~\ref{prop:Tao-v7-Prop52-cn-repair} permits deletion of the bad
iid path tube.  Proposition~\ref{prop:Tao-v7-Fourier-to-mixing}, whose
dependency branch is closed in
Corollary~\ref{cor:Tao-v7-fine-scale-branch}, then replaces the law at level
\(r\) by the exact uniform lift of its level-\(m_0\) projection.  The finite
pairing is
\begin{equation}
\label{eq:Tao-Prop111-exact-collapse}
\begin{split}
 &\sum_{X\bmod3^r}
 \left(3^r\!\sum_{\substack{M\in E'\\M\equiv X\ (3^r)}}\frac1M\right)
 3^{m_0-r}
 \mathbb P\!\left(
  \operatorname{Syrac}(\mathbb Z/3^{m_0}\mathbb Z)
  \equiv X\pmod{3^{m_0}}
 \right)\\
 &\hspace{35mm}=Z_E.
\end{split}
\end{equation}
Indeed, for each \(M\) there is exactly one residue
\(X=M\pmod{3^r}\), and \(3^r3^{m_0-r}=3^{m_0}\).  Thus, choosing the
arbitrary mixing exponent large enough before summing over the
\(O(L)\) possible times,
\begin{equation}
\label{eq:Tao-Prop111-uniform-An}
 A_{n,E}=Z_E+O(L^{-c_6})
\end{equation}
uniformly for both \(y\)'s and all \(n\in I_y\cap\mathbb Z\).  Moreover
\(A_{n,E}=O(1)\) by the uniform \(c_n\)-bound, so
\eqref{eq:Tao-Prop111-uniform-An} gives \(Z_E=O(1)\) without circular use
of the probability conclusion.

Substitute \eqref{eq:Tao-Prop111-uniform-An} into
\eqref{eq:Tao-Prop111-pre-mixing}.  The \(O(L)\) summation factor is
cancelled by the denominator of size \(L\); after decreasing the common
positive exponent,
\[
 \mathbb P(\operatorname{Pass}_x(\mathbf N_y)\in E)
 =\frac{\#(I_y\cap\mathbb Z)}
        {\frac{\alpha-1}{2}\log y}Z_E+O(L^{-c}).
\]
Lemma~\ref{lem:Tao-Prop111-window-count} therefore yields the uniform,
\(y\)-independent main term
\begin{equation}
\label{eq:Tao-Prop111-arbitrary-event}
 \mathbb P(\operatorname{Pass}_x(\mathbf N_y)\in E)
 =\frac{2}{\lambda}Z_E+O(L^{-c}).
\end{equation}
Taking the difference of \eqref{eq:Tao-Prop111-arbitrary-event} for the two
values of \(y\), and then the supremum over all subsets \(E\) of the common
state space \((2\mathbb N+1)\cap[1,x]\), gives an \(O(L^{-c})\) event
distance.  The convention \(\operatorname{Pass}_x(N)=1\) when
\(T_x(N)=+\infty\) keeps both laws in this same state space.  Finally the
source inequality
\[
 d_{\mathrm{TV}}(P,Q)
 \le2\sup_E|P(E)-Q(E)|
\]
proves \eqref{eq:Tao-Prop111-theorem-TV} in the unhalved convention.
\end{proof}

\begin{nonclaim}
Theorem~\ref{thm:Tao-v7-Prop111-first-passage} proves exactly the
first-passage stabilisation proposition and closes the edge
\[
 \text{Proposition~1.9}+\text{Proposition~1.14}
 \Longrightarrow\text{Proposition~1.11}.
\]
It does not assert exact residue uniformity, independence of actual valuation
coordinates, arbitrary-\(y\) uniformity, natural-density transport, a coupling
of the two passage locations, convergence to \(1\), exclusion of divergent or
periodic orbits, or a fixed universal orbit bound.  The finite-
\(x\) time-window correction depends on \(y\), even though its leading
normalization does not.  The exact later implications to Theorem~3.1,
Theorem~1.6, and Theorem~1.3, including their source defects and repaired
morphisms, are proved in the next subsection; they had not yet been proved at
the checkpoint represented by this subsection.
\end{nonclaim}


\subsection{The exact reduction from first passage to the main theorem}
\label{subsec:Tao-main-reduction}

The source reduction from Proposition~1.11 to Theorems~3.1, 1.6, and~1.3
is v7 lines~535--593, with the definitions and statements at lines~159--206
\cite{Tao2026v7}.  The corresponding v5 loci are lines~529--587 and
159--206; after deletion of whitespace, the two reduction blocks are
identical.  The journal manifestation reproduces the same argument on
physical pages~15--17, with its logarithmic-density and Syracuse setup on
physical pages~3--4 and Proposition~1.11 on physical page~8.  The comparison
below is therefore version-pinned but not a claim that the complete
manifestations are text-identical.

\begin{sourceissue}[defects and suppressed arrows in the printed reduction]
\label{sourceissue:Tao-main-reduction}
Eight points must be separated.

First, v7 lines~591--593 define
\(\widetilde f(x)=\inf_{N\geq x}f(N)\) and then use that lower envelope on
the range \(N\leq x\), where it supplies no comparison.  The displayed
bound is false: on the odd positive integers, take \(f(1)=0\) and
\(f(N)=e^N\) for \(N\geq3\).  The term \(N=1\) contributes one to the
printed left side, while its asserted right side is
\(O(\log x/x^c)\).  Second, the complement of
\(\operatorname{Syr}_{\min}(N)<f(N)\) is
\(\operatorname{Syr}_{\min}(N)\geq f(N)\), not the printed strict
inequality in the other direction; equality cannot be discarded.  Third,
the statement permits real-valued \(f\), while line~591 silently restricts
its codomain to \([0,+\infty)\).

Fourth, line~206 assigns logarithmic density \(2^{-a}\) to the stratum
\(\nu_2(N)=a\) and \(1-2^{-a_0}\) to the union
\(0\leq a\leq a_0\).  The exact values are respectively
\(2^{-(a+1)}\) and \(1-2^{-(a_0+1)}\); the function on the odd coordinate
must moreover be \(M\mapsto f(2^aM)\), not \(f(M)\).

Fifth, the iteration at lines~569--580 omits the possible value \(J=0\).
Sixth, its probabilistic blocks need not be defined at small scales because
the narrow interval can contain no odd integer.  Seventh, the parenthetical
at line~552 says that finite passage is redundant when \(x\geq N_0\); this
is false under the source convention \(\operatorname{Pass}_x=1\) when
\(T_x=+\infty\), since \(1\in E_{N_0}\).  The displayed proof does retain
the finite-passage condition.  Eighth, the substitution at lines~569--574
produces
\(O((\alpha^{j-2}\log y)^{-c})\), not the printed
\(O((\alpha^j\log y)^{-c})\).  Their ratio is the fixed factor
\(\alpha^{2c}\), so the iteration survives only after that factor is
explicitly absorbed.

The same passage also suppresses the indexed first-passage morphism, the
actual multiplicative block cover, its harmonic normalisation, and the
dyadic weight transport.  These are proved below.  None of the repairs is
silently attributed to the source.
\end{sourceissue}

For a real cutoff \(X\geq1\), set
\[
 H_{\rm odd}(X)=
 \sum_{\substack{1\leq M\leq X\\M\ {\rm odd}}}\frac1M,
 \qquad
 H(X)=\sum_{1\leq N\leq X}\frac1N.
\]
The integral test, with the integer cutoffs retained, gives
\begin{equation}
\label{eq:Tao-main-harmonic-asymptotics}
 H_{\rm odd}(X)=\frac12\log X+O(1),
 \qquad H(X)=\log X+O(1).
\end{equation}
For a real threshold \(N_0\geq2\), write
\begin{equation}
\label{eq:Tao-main-EN0}
 E_{N_0}=\{M\in2\mathbb N+1:
        \operatorname{Syr}_{\min}(M)\leq N_0\}.
\end{equation}

\begin{lemma}[the indexed nested-threshold morphism]
\label{lem:Tao-nested-passage-tail}
Let \(1\leq u\leq v\), and put
\[
 \mathcal D_u=\{N\in2\mathbb N+1:T_u(N)<+\infty\}.
\]
For \(N\in\mathcal D_u\), define
\[
 d_{u,v}(N)=T_u(N)-T_v(N)\in\mathbb N,
 \qquad \sigma_d(k)=k+d.
\]
Then
\begin{equation}
\label{eq:Tao-nested-pass-point}
 \operatorname{Pass}_u(N)
 =\operatorname{Syr}^{d_{u,v}(N)}
    (\operatorname{Pass}_v(N)).
\end{equation}
More precisely, for the indexed tails
\[
 \mathcal O_w^N(k)=
 \operatorname{Syr}^k(\operatorname{Pass}_w(N)),
 \qquad k\in\mathbb N,
\]
one has the exact identity
\begin{equation}
\label{eq:Tao-nested-tail-shift}
 \mathcal O_u^N=\mathcal O_v^N\circ\sigma_{d_{u,v}(N)}.
\end{equation}
The shift \(\sigma_d\) is injective, has image \(d+\mathbb N\), has
singleton fibres over that image, and has diagonal equality kernel.  Hence
the value-set of the \(u\)-tail is contained in the value-set of the
\(v\)-tail.  In particular,
\begin{equation}
\label{eq:Tao-nested-EN0-transport}
 T_u(N)<+\infty,\quad \operatorname{Pass}_u(N)\in E_{N_0}
 \quad\Longrightarrow\quad
 \operatorname{Pass}_v(N)\in E_{N_0}.
\end{equation}
\end{lemma}

\begin{proof}
Every iterate at most \(u\) is also at most \(v\), so
\(T_v(N)\leq T_u(N)<+\infty\).  The semigroup law for the iterates gives
\[
 \operatorname{Syr}^{T_u(N)}(N)
 =\operatorname{Syr}^{T_u(N)-T_v(N)}
   (\operatorname{Syr}^{T_v(N)}(N)),
\]
which is \eqref{eq:Tao-nested-pass-point}.  Applying a further \(k\)
iterates gives \eqref{eq:Tao-nested-tail-shift}.  If the later tail from
\(\operatorname{Pass}_u(N)\) reaches \(N_0\), the containing tail from
\(\operatorname{Pass}_v(N)\) reaches the same value, proving
\eqref{eq:Tao-nested-EN0-transport}.  No injectivity of
\(\mathcal O_v^N\) is used or asserted: a periodic orbit can repeat values
at different indices.
\end{proof}

For every scale \(s\) for which the odd logarithmic block is nonempty,
retain
\[
 \mathcal R_s=(2\mathbb N+1)\cap[s,s^\alpha],
 \qquad H_s=\sum_{N\in\mathcal R_s}\frac1N,
 \qquad
 \mathbb P(\mathbf N_s=N)=\frac1{H_sN}.
\]
Define the numerical event probability
\begin{equation}
\label{eq:Tao-main-bs}
 b_{N_0}(s)=\mathbb P\left(
  T_s(\mathbf N_{s^\alpha})<+\infty,
  \ \operatorname{Pass}_s(\mathbf N_{s^\alpha})\in E_{N_0}
 \right).
\end{equation}

\begin{proposition}[the exact block recurrence]
\label{prop:Tao-main-block-recurrence}
There are absolute constants \(X_*,C,c>0\) such that, whenever
\(s\geq X_*\),
\begin{equation}
\label{eq:Tao-main-block-recurrence}
 b_{N_0}(s^\alpha)
 \geq b_{N_0}(s)-C(\log s)^{-c}.
\end{equation}
The cutoff \(X_*\) is chosen so that Proposition~1.11 applies and every
block occurring here contains an odd integer.
\end{proposition}

\begin{proof}
On the probability space supporting \(\mathbf N_{s^{\alpha^2}}\), let
\[
 A_s=\{T_s(\mathbf N_{s^{\alpha^2}})<+\infty,
 \ \operatorname{Pass}_s(\mathbf N_{s^{\alpha^2}})\in E_{N_0}\}.
\]
Lemma~\ref{lem:Tao-nested-passage-tail}, with \(u=s\) and
\(v=s^\alpha\), gives
\[
 A_s\subseteq
 \{T_{s^\alpha}(\mathbf N_{s^{\alpha^2}})<+\infty,
 \ \operatorname{Pass}_{s^\alpha}
       (\mathbf N_{s^{\alpha^2}})\in E_{N_0}\}.
\]
The probability of the event on the right is
\(b_{N_0}(s^\alpha)\).  Proposition~1.11 and the event inequality for its
unhalved total-variation convention therefore give
\begin{align*}
 b_{N_0}(s^\alpha)
 &\geq\mathbb P(A_s)\\
 &\geq
 \mathbb P(\operatorname{Pass}_s(\mathbf N_{s^{\alpha^2}})
             \in E_{N_0})-Cs^{-c}\\
 &\geq
 \mathbb P(\operatorname{Pass}_s(\mathbf N_{s^\alpha})
             \in E_{N_0})
       -C(\log s)^{-c}-Cs^{-c}\\
 &\geq b_{N_0}(s)-C'(\log s)^{-c}.
\end{align*}
Increasing \(C\) proves \eqref{eq:Tao-main-block-recurrence}.  This uses
the two marginal passage-location laws and failure control; it does not
construct a coupling.
\end{proof}

\begin{theorem}[repaired exact form of Tao's Theorem~3.1]
\label{thm:Tao-main-alt-repaired}
There are absolute constants \(C,c>0\) such that, for all real
\(N_0\geq2\) and all real \(X\geq2\),
\begin{equation}
\label{eq:Tao-main-alt-odd}
 \sum_{\substack{M\in2\mathbb N+1,\ M\leq X\\
       \operatorname{Syr}_{\min}(M)>N_0}}
       \frac1M
 \leq C\frac{\log X}{(\log N_0)^c},
\end{equation}
and
\begin{equation}
\label{eq:Tao-main-alt-collatz}
 \sum_{\substack{N\in\mathbb N+1,\ N\leq X\\
       \operatorname{Col}_{\min}(N)>N_0}}
       \frac1N
 \leq C\frac{\log X}{(\log N_0)^c}.
\end{equation}
Consequently the corresponding good sets have, uniformly in \(X\), odd
relative and full logarithmic probabilities at least
\(1-O((\log N_0)^{-c})\).
\end{theorem}

\begin{proof}
Choose \(X_*\) as in
Proposition~\ref{prop:Tao-main-block-recurrence}.  We first prove the block
bound
\begin{equation}
\label{eq:Tao-main-block-bad-probability}
 \mathbb P(\operatorname{Syr}_{\min}(\mathbf N_s)>N_0)
 \ll(\log N_0)^{-c}.
\end{equation}
For bounded \(N_0\), this follows after increasing the implied constant.
Thus assume \(N_0\) is large enough that
\(N_0^{1/\alpha^3}\geq X_*\).

If \(s\leq N_0^{1/\alpha}\), every integer in
\([s,s^\alpha]\) is at most \(N_0\), so the probability in
\eqref{eq:Tao-main-block-bad-probability} is exactly zero.  This is the
omitted \(J=0\) case.

Now suppose \(s>N_0^{1/\alpha}\).  Let \(J\geq1\) be least such that
\[
 y=s^{\alpha^{-J}}<N_0^{1/\alpha}.
\]
Minimality gives
\begin{equation}
\label{eq:Tao-main-y-bounds}
 y^\alpha\geq N_0^{1/\alpha},
 \qquad y\geq N_0^{1/\alpha^2},
 \qquad y^\alpha<N_0.
\end{equation}
Put \(r_0=y^{1/\alpha}\).  Then
\(r_0\geq N_0^{1/\alpha^3}\geq X_*\), and the input in
\(b_{N_0}(r_0)\) is \(\mathbf N_y\).  On finite passage below \(r_0\),
the passage location is less than \(N_0\), so Proposition~1.11 gives
\begin{equation}
\label{eq:Tao-main-base-event}
 b_{N_0}(r_0)\geq1-Cr_0^{-c}.
\end{equation}
For \(j=1,\ldots,J\), apply
\eqref{eq:Tao-main-block-recurrence} at the exact scale
\(y^{\alpha^{j-2}}\).  The accumulated logarithmic error is
\begin{align}
 \sum_{j=1}^J
  (\alpha^{j-2}\log y)^{-c}
 &\leq (\log y)^{-c}
     \sum_{j=1}^{\infty}\alpha^{-c(j-2)}\notag\\
 &=\frac{\alpha^c}{1-\alpha^{-c}}(\log y)^{-c}
 \ll(\log N_0)^{-c}.
\label{eq:Tao-main-geometric-error}
\end{align}
Here \eqref{eq:Tao-main-y-bounds} was used only after the exact factor in
the first line was retained.  Equations
\eqref{eq:Tao-main-base-event} and
\eqref{eq:Tao-main-geometric-error} yield
\[
 b_{N_0}(s^{1/\alpha})\geq1-C(\log N_0)^{-c}.
\]
Its input variable is exactly \(\mathbf N_s\).  Whenever its event holds,
the orbit of \(\mathbf N_s\) reaches \(E_{N_0}\), proving
\eqref{eq:Tao-main-block-bad-probability}.

For \(s\geq2\), the integral test gives \(H_s\ll\log s\), with an
absolute constant because \(\alpha=1001/1000\) is fixed.  An applied cover
scale with \(1<s<2\) has no odd integer in \([s,s^\alpha]\), so its local
sum is zero; every nonempty cover block to which the estimate is applied
satisfies \(s\geq2\).  Hence
\begin{equation}
\label{eq:Tao-main-local-harmonic-bound}
 \sum_{\substack{M\in(2\mathbb N+1)\cap[s,s^\alpha]\\
          \operatorname{Syr}_{\min}(M)>N_0}}
          \frac1M
 \ll\frac{\log s}{(\log N_0)^c}.
\end{equation}
For a target \(X>N_0\), set
\(z_k=X^{\alpha^{-k}}\), and let \(L\geq1\) be least such that
\(z_L\leq N_0\).  Then
\[
 z_{L-1}>N_0,
 \qquad z_L>N_0^{1/\alpha},
 \qquad
 [z_L,X]=\bigcup_{k=1}^L[z_k,z_k^\alpha].
\]
Every bad starting value is greater than \(N_0\), since its orbit contains
the starting value.  Thus the leftover interval \([1,z_L]\) is bad-free,
and endpoint duplication in the displayed cover can only enlarge an upper
bound.  Finally,
\begin{equation}
\label{eq:Tao-main-cover-log-sum}
 \sum_{k=1}^L\log z_k
 =\log X\sum_{k=1}^L\alpha^{-k}
 \leq\frac{\log X}{\alpha-1}.
\end{equation}
Summing \eqref{eq:Tao-main-local-harmonic-bound} proves
\eqref{eq:Tao-main-alt-odd}.  If \(X\leq N_0\), its left side is zero.

For the all-integer assertion, write every \(N\) uniquely as
\(N=2^aM\), with \(a\in\mathbb N\) and \(M\) odd.  The exact minimum
identity proved again below gives
\begin{equation}
\label{eq:Tao-main-fixed-dyadic-sum}
 \sum_{\substack{N\leq X\\
       \operatorname{Col}_{\min}(N)>N_0}}\frac1N
 =\sum_{a\geq0}2^{-a}
  \sum_{\substack{M\leq X/2^a,\ M\ {\rm odd}\\
       \operatorname{Syr}_{\min}(M)>N_0}}\frac1M.
\end{equation}
When \(X/2^a<2\), the inner bad sum vanishes because \(N_0\geq2\).
For the remaining terms, \eqref{eq:Tao-main-alt-odd} and the truncated,
termwise nonnegative estimate
\[
 \sum_{\substack{a\geq0\\X/2^a\geq2}}
     2^{-a}\log(X/2^a)
 \leq \log X\sum_{a\geq0}2^{-a}
 \leq2\log X
\]
prove \eqref{eq:Tao-main-alt-collatz}.  The probability statements follow
from \eqref{eq:Tao-main-harmonic-asymptotics}; no identity of unnormalised
and normalised harmonic mass is being asserted.
\end{proof}

\begin{theorem}[repaired exact form of Tao's Theorem~1.6]
\label{thm:Tao-main-syr-repaired}
Let \(f:2\mathbb N+1\to\mathbb R\) satisfy
\(f(M)\to+\infty\).  Then
\begin{equation}
\label{eq:Tao-main-syr-strict}
 \operatorname{Syr}_{\min}(M)<f(M)
\end{equation}
for a set of odd positive integers of relative odd logarithmic density one.
Equivalently, the set of odd \(M\) on which
\(\operatorname{Syr}_{\min}(M)\geq f(M)\) has relative odd logarithmic
density zero.
\end{theorem}

\begin{proof}
Fix an integer \(K\geq2\).  Since \(f(M)\to+\infty\), there is a finite
cutoff \(Y_K\) such that
\[
 f(M)>K\qquad(M\geq Y_K,\ M\ {\rm odd}).
\]
Let
\(D_f=\{M\ {\rm odd}:\operatorname{Syr}_{\min}(M)\geq f(M)\}\).
For every \(X\), the exact tail inclusion gives
\begin{align}
 \sum_{\substack{M\leq X\\M\in D_f}}\frac1M
 &\leq H_{\rm odd}(Y_K)
 +\sum_{\substack{M\leq X,\ M\ {\rm odd}\\
          \operatorname{Syr}_{\min}(M)>K}}\frac1M\notag\\
 &\leq H_{\rm odd}(Y_K)
 +C\frac{\log X}{(\log K)^c}.
\label{eq:Tao-main-fixed-K-tail}
\end{align}
Divide first by \(H_{\rm odd}(X)\), let \(X\to+\infty\) with \(K\)
fixed, and only then let \(K\to+\infty\).  Equation
\eqref{eq:Tao-main-harmonic-asymptotics} gives
\[
 \limsup_{X\to\infty}
 \frac{\sum_{M\leq X,\ M\in D_f}1/M}{H_{\rm odd}(X)}
 \leq\frac{2C}{(\log K)^c}\longrightarrow0.
\]
This proves the strict inequality in \eqref{eq:Tao-main-syr-strict}, includes
the equality part of its complement, permits arbitrary nonmonotone real
\(f\), and confines every negative or otherwise exceptional value of \(f\)
to the finite prefix.  No tail lower envelope is used.
\end{proof}

\begin{proposition}[the exact dyadic stratum morphism]
\label{prop:Tao-main-dyadic-morphism}
Let
\[
 \mathcal O=2\mathbb N+1,
 \qquad
 \mathcal S_a=\{N\in\mathbb N+1:\nu_2(N)=a\}.
\]
For each \(a\in\mathbb N\), the maps
\[
 \iota_a:\mathcal O\longrightarrow\mathcal S_a,
 \quad \iota_a(M)=2^aM,
 \qquad
 \pi_a:\mathcal S_a\longrightarrow\mathcal O,
 \quad \pi_a(N)=N/2^a
\]
are two-sided inverses.  Their fibres are singletons, their equality kernels
are diagonal, and they lose no starting-number coordinate information.
Globally,
\begin{equation}
\label{eq:Tao-main-global-dyadic-bijection}
 \mathbb N+1\longleftrightarrow
 \coprod_{a\in\mathbb N}\{a\}\times\mathcal O,
 \qquad
 N\longmapsto(\nu_2(N),N/2^{\nu_2(N)})
\end{equation}
is a bijection with inverse \((a,M)\mapsto2^aM\).

For every \(A\subseteq\mathcal O\) and real \(X\geq1\), the cutoff and
harmonic weight transport exactly as
\begin{equation}
\label{eq:Tao-main-dyadic-weight}
 \sum_{\substack{N\in\iota_a(A)\\N\leq X}}\frac1N
 =2^{-a}\sum_{\substack{M\in A\\M\leq X/2^a}}\frac1M.
\end{equation}
In particular,
\begin{equation}
\label{eq:Tao-main-dyadic-density}
 \sum_{\substack{N\leq X\\\nu_2(N)=a}}\frac1N
 =2^{-a}H_{\rm odd}(X/2^a)
 =2^{-(a+1)}\log X+O_a(1).
\end{equation}
Finally,
\begin{equation}
\label{eq:Tao-main-minimum-identity}
 \operatorname{Col}_{\min}(2^aM)
 =\operatorname{Syr}_{\min}(M).
\end{equation}
\end{proposition}

\begin{proof}
The defining condition \(\nu_2(N)=a\) says uniquely that \(N=2^aM\)
with \(M\) odd, which proves every bijection and fibre assertion.  Replacing
\(N\) by \(2^aM\) in the finite sum proves
\eqref{eq:Tao-main-dyadic-weight}; then
\eqref{eq:Tao-main-harmonic-asymptotics} proves
\eqref{eq:Tao-main-dyadic-density}.

It remains to prove the minimum identity rather than infer it from notation.
The first \(a\) Collatz steps starting at \(2^aM\) are divisions by two and
terminate at \(M\); every value on that initial segment is at least \(M\).
If \(M_j\) is an odd Syracuse state and
\(q_j=\nu_2(3M_j+1)\), then the corresponding Collatz segment is
\[
 M_j,\ 2^{q_j}M_{j+1},\ 2^{q_j-1}M_{j+1},\ldots,
 2M_{j+1},\ M_{j+1},
 \qquad M_{j+1}=\operatorname{Syr}(M_j).
\]
Every suppressed intermediate value is at least its terminal odd value.
Thus the minimum over the complete Collatz orbit is exactly the minimum over
the odd Syracuse states, proving \eqref{eq:Tao-main-minimum-identity}.
The map is not time-preserving.  The full Collatz time coordinate is
reconstructed by expanding each Syracuse step using
\(q_j=\nu_2(3M_j+1)\), which is determined by its odd source state \(M_j\).
\end{proof}

\begin{theorem}[repaired exact form of Tao's Theorem~1.3]
\label{thm:Tao-main-repaired}
Let \(f:\mathbb N+1\to\mathbb R\) satisfy \(f(N)\to+\infty\).  Then
\begin{equation}
\label{eq:Tao-main-collatz-strict}
 \operatorname{Col}_{\min}(N)<f(N)
\end{equation}
for a set of positive integers of logarithmic density one.
\end{theorem}

\begin{proof}
For each fixed \(a\in\mathbb N\), define
\[
 g_a(M)=f(2^aM)\qquad(M\in\mathcal O).
\]
Then \(g_a(M)\to+\infty\), so
Theorem~\ref{thm:Tao-main-syr-repaired} says that
\[
 B_a=\{M\in\mathcal O:
       \operatorname{Syr}_{\min}(M)\geq g_a(M)\}
\]
has harmonic mass \(o(\log T)\) below \(T\).  By
Proposition~\ref{prop:Tao-main-dyadic-morphism}, the bad set for
\eqref{eq:Tao-main-collatz-strict} in \(\mathcal S_a\) is exactly
\(\iota_a(B_a)\), and its mass below \(X\) is
\[
 2^{-a}\sum_{\substack{M\in B_a\\M\leq X/2^a}}\frac1M=o(\log X)
\]
for this fixed \(a\).

Fix \(A\in\mathbb N\).  Splitting the global bad mass into the first
\(A+1\) strata and the remaining tail gives
\begin{align}
 \sum_{\substack{N\leq X\\
   \operatorname{Col}_{\min}(N)\geq f(N)}}\frac1N
 &\leq\sum_{a=0}^A2^{-a}
   \sum_{\substack{M\in B_a\\M\leq X/2^a}}\frac1M
 +\sum_{\substack{N\leq X\\\nu_2(N)>A}}\frac1N,\notag\\
 \sum_{\substack{N\leq X\\\nu_2(N)>A}}\frac1N
 &=2^{-(A+1)}H(X/2^{A+1}).
\label{eq:Tao-main-dyadic-tail-limit}
\end{align}
The second identity follows from the unique formula
\(N=2^{A+1}Q\), with \(Q\) now an arbitrary positive integer.  Divide by
\(H(X)\) and let \(X\to+\infty\) at fixed \(A\).  The finite first sum
vanishes in the limit and
\eqref{eq:Tao-main-harmonic-asymptotics} gives
\[
 \limsup_{X\to\infty}
 \frac{\sum_{N\leq X,\ \operatorname{Col}_{\min}(N)\geq f(N)}1/N}
      {H(X)}
 \leq2^{-(A+1)}.
\]
Only now send \(A\to\infty\).  The bad logarithmic density is zero, which
proves the strict inequality \eqref{eq:Tao-main-collatz-strict}.  No
uniformity in \(a\), and no interchange of the two limits, has been used.
\end{proof}

As an independent cross-check, one may apply the fixed-\(K\) argument in
\eqref{eq:Tao-main-fixed-K-tail} directly to the Collatz estimate
\eqref{eq:Tao-main-alt-collatz}.  This gives
Theorem~\ref{thm:Tao-main-repaired} without passing through the dyadic
strata, and yields the same strict bad-set complement.

\begin{nonclaim}
Theorems~\ref{thm:Tao-main-alt-repaired},
\ref{thm:Tao-main-syr-repaired}, and~\ref{thm:Tao-main-repaired} prove the
exact repaired chain
\[
 \text{Proposition~1.11}\Longrightarrow\text{Theorem~3.1}
 \Longrightarrow\text{Theorem~1.6}
 \Longrightarrow\text{Theorem~1.3}.
\]
They prove logarithmic-density statements, not natural-density statements;
they do not give an absolute bound valid for almost all starting values,
force an individual orbit through a threshold, exclude divergent or
periodic exceptional orbits, prove convergence to one, construct an
 invariant logarithmic measure, or produce a coupling.  The dyadic bijection
preserves all starting-number information---it is a bijection on the
displayed starting coordinates---and preserves the orbit minimum at the
displayed scope.  It is not itself a time-preserving conjugacy; Collatz time
is recovered by the explicit valuation expansion above.
This closes the main-theorem reduction, not the durable whole-paper,
whole-corpus, formalisation, or release gates.
\end{nonclaim}


\subsection{Computational verification: exact coordinates, certificates, and chronology}
\label{subsec:computational-verification}

The computational literature does not use one unnamed Collatz map.  In this
subsection
\[
 U(n)=
 \begin{cases}
  3n+1,&n\text{ odd},\\
  n/2,&n\text{ even},
 \end{cases}
 \qquad
 T(n)=
 \begin{cases}
  (3n+1)/2,&n\text{ odd},\\
  n/2,&n\text{ even}
 \end{cases}
\]
denote respectively the unshortened and shortened maps.  One odd $T$-step is
two $U$-steps and one even $T$-step is one $U$-step.  Barina also uses the
auxiliary chart
\[
 T_1(n)=
 \begin{cases}
  (n+1)/2,&n\text{ odd},\\
  3n/2,&n\text{ even}.
 \end{cases}
\]
Its exact branchwise bridges are
\[
 T(n)=
 \begin{cases}T_1(n+1)-1,&n\text{ odd},\\ n/2,&n\text{ even},\end{cases}
 \quad
 T_1(n)=
 \begin{cases}T(n-1)+1,&n\text{ odd},\\3n/2,&n\text{ even}.
 \end{cases}
\]
They move the additive term between coordinate charts; they do not identify
the maps or their clocks.  In particular, the 2025 paper's page-4
\emph{trajectory on $n+1$}, $(7,4,6,9,5,3,2)$, is the shifted-coordinate
presentation associated with the trajectory from $12$, not the $T$-orbit
of $13$ \cite[pp.~2, 4]{Barina2025}.

\begin{proposition}[Barina's block map, with fibres]
\label{prop:Barina-block-map}
For $x\in\Npos$, put
\[
 \alpha=\nu_2(x+1),\qquad u=\frac{x+1}{2^\alpha},\qquad
 y=3^\alpha u-1,\qquad \beta=\nu_2(y),
\]
and $B(x)=y/2^\beta$.  Then $B(x)$ is a positive odd integer and
\begin{equation}
\stepcounter{equation}
\label{eq:Barina-block}
 B(x)=T^{\alpha+\beta}(x).
\tag{CV.1}
\end{equation}
For every odd $z>0$, the fibre $B^{-1}(z)$ is exactly the set of
\begin{equation}
\stepcounter{equation}
\label{eq:Barina-fibre}
 x=2^\alpha\frac{2^\beta z+1}{3^\alpha}-1
\tag{CV.2}
\end{equation}
over integers $\alpha\geq0,\beta\geq1$ satisfying
$3^\alpha\mid2^\beta z+1$.  Hence $B$ is surjective, has infinite fibres,
and has no left or two-sided inverse after the block boundary and elapsed
clock have been discarded.
\end{proposition}

\begin{proof}
The integer $u$ is odd.  For $0\leq j\leq\alpha$, direct induction gives
\[
 T^j(x)=3^j\frac{x+1}{2^j}-1.
\]
The state at $j=\alpha$ is $y$; the following $\beta$ steps divide it by
two, proving \eqref{eq:Barina-block}.  Conversely,
\eqref{eq:Barina-fibre} gives $x+1=2^\alpha u$, with
$u=(2^\beta z+1)/3^\alpha$ odd, and
$3^\alpha u-1=2^\beta z$.  Thus the two valuations are exactly
$\alpha,\beta$ and $B(x)=z$.  The section $z\mapsto2z$ proves
surjectivity.  Taking $\alpha=0$ in \eqref{eq:Barina-fibre} gives the
distinct points $2^\beta z$, $\beta\geq1$, in the same fibre.
\end{proof}

On the odd $2$-adic integers with normalized Haar measure,
$\Pr(\alpha=a)=2^{-a}$ and
$\Pr(\beta=b\mid\alpha=a)=2^{-b}$ for $a,b\geq1$.  Multiplication by the
odd unit $3^a$ proves the conditional statement, so the two valuations are
independent and
\[
 \mathbb E(\alpha+\beta)=4.
\]
This is the exact measure-theoretic content of the articles' phrase
\emph{average four iterations}; it is not an assertion that every finite
interval has raw arithmetic mean four \cite[pp.~3--5]{Barina2021}.

\begin{sourceissue}[the exceptional start $1$ in the printed algorithms]
\label{sourceissue:Barina-one-endpoints}
Equation~\eqref{eq:Barina-block} has $B(1)=1$, representing the shortened
cycle $1\mapsto2\mapsto1$.  The printed endpoints differ by version.  The
2021 Algorithm~1 uses strict \texttt{until n<n0} and loops at $n_0=1$.
Its Algorithm~2 uses \texttt{until n=1}: it terminates after one complete
$1\mapsto1$ macro, but a literal delay accumulator gives $2$ rather than
the conventional stopping time $0$.  Both 2025 Algorithms~1 and~2 use the
strict endpoint and loop at $1$ \cite[pp.~4--5]{Barina2021}
\cite[pp.~4--5]{Barina2025}.  The exact repairs are a separate base case, or
the domain $n_0>1$ where only descent is required.  The production inner
loop tests starts $3\bmod4$, so this exceptional input is not encountered
there.
\end{sourceissue}

\begin{proposition}[the dyadic affine-prefix coordinate]
\label{prop:computational-dyadic-prefix}
Fix $k\geq0$, $h\geq0$, and $0\leq r<2^k$, and let
\[
 a_j(r)=\#\{0\leq i<j:T^i(r)\text{ is odd}\}.
\]
For every $0\leq j\leq k$,
\begin{equation}
\stepcounter{equation}
\label{eq:computational-dyadic-prefix}
 T^j(2^kh+r)=3^{a_j(r)}2^{k-j}h+T^j(r).
\tag{CV.3}
\end{equation}
In particular,
\begin{equation}
\stepcounter{equation}
\label{eq:computational-dyadic-endpoint}
 T^k(2^kh+r)=3^{a_k(r)}h+T^k(r).
\tag{CV.4}
\end{equation}
\end{proposition}

\begin{proof}
For $j<k$, the high term in \eqref{eq:computational-dyadic-prefix} is even.
The full state and the low state $T^j(r)$ therefore have the same parity and
take the same next branch.  Applying that branch proves the induction.  This
is simultaneously Oliveira e Silva's Proposition~1 and Barina's equation~(6),
after their maps and variables have been typed rather than identified
\cite[p.~191]{OliveiraSilva2010}\cite[p.~5]{Barina2021}.
\end{proof}

Write $a=a_k(r)$, $d=T^k(r)$, $L=3^a-2^k$, and $R=r-d$.  The exact
necessary-and-sufficient condition for
\[
 T^k(2^kh+r)<2^kh+r\qquad(h=1,2,\ldots)
\]
is
\begin{equation}
\stepcounter{equation}
\label{eq:computational-uniform-descent}
 (L<0\text{ and }L<R)\quad\text{or}\quad(L=0\text{ and }R>0).
\tag{CV.5}
\end{equation}
Indeed the inequality is $Lh<R$, whose most restrictive positive $h$ is
$h=1$ when $L<0$; it is constant when $L=0$, and eventually fails when
$L>0$.  The generator's \texttt{is\_convergent} implements only the first
disjunct.  It is conservative in the equality case; it does not falsely
certify it.

If $r'<r$ and
\[
 a_k(r')=a_k(r),\qquad T^k(r')=T^k(r),
\]
then \eqref{eq:computational-dyadic-endpoint} proves, for every $h$,
\begin{equation}
\stepcounter{equation}
\label{eq:computational-endpoint-coalescence}
 T^k(2^kh+r')=T^k(2^kh+r),\qquad2^kh+r'<2^kh+r.
\tag{CV.6}
\end{equation}
Thus a dead sieve bit has a finite descent, coalescence, inherited-prefix, or
restricted-predecessor certificate.  A live bit is merely the complement of
the implemented certificates.  It is not a residue class proved not to
converge; the 2025 page-6 wording \emph{do not converge} must be read at this
finite algorithmic scope \cite{Barina2025}.

\begin{sourceissue}[cardinality and generator-history defects]
\label{sourceissue:Barina-sieve-cardinality}
Within one half-open work unit of width $2^{40}$, a fixed class modulo
$2^k$ contains $2^{40-k}$ starts, while there are $2^k$ classes.  The 2025
Section~3.4 reverses these quantities twice.  The 2021 page-6 example and
\path{src/worker/worker.c:336--369} use the correct $2^{40-34}$.

The separate \texttt{collatz-sieve} history has a dated implementation
  defect.  Commit
  \nolinkurl{49c5524d8d42867c866e863030a8e5d86237bf17}
(10 April 2020) passed the odd-step count \texttt{c\_} as both the third and
fourth arguments of \texttt{is\_convergent}; the fourth argument had to be
the endpoint \texttt{d\_}.  Commit
  \nolinkurl{b3d2f61d89eec55008d61c32577a9e8c8b3874ab}
(6 April 2023) repaired it.  The baseline committed maps predate the faulty
interval, while the \texttt{hesieve} and \texttt{h2esieve} maps postdate the
repair.  The defect governs regeneration from intermediate source history;
it is not silently attributed to blobs outside that interval.
\end{sourceissue}

The first exact lower-start reductions are
\[
 T(2q+1)=3q+2,
\]
and
\begin{equation}
\label{eq:computational-four-transports}
 T^3(8q+3)=9q+4,\qquad
 T^6(64q+7)=81q+10,\qquad
 T^7(128q+95)=243q+182.
\end{equation}
Their odd-step counts are respectively $1,2,4,5$.  Every source on the left
is strictly smaller than its target.  The first three families omit $37$ of
$81$ initial residues: the first two give
$\{2,4,5,8\}\bmod9$, their lifts modulo $27$ are
\[
 \{2,4,5,8,11,13,14,17,20,22,23,26\},
\]
and the $81q+10$ class is disjoint from them.  The fourth target is already
$2\bmod3$ as an initial residue, but is not redundant when it occurs as an
intermediate low-coordinate state.

\begin{proposition}[the Hercher coalescence lift]
\label{prop:Hercher-computational-coalescence}
Let $x$ be odd and put
\[
 \alpha=\nu_2(x+1),\quad u=(x+1)/2^\alpha,
 \quad y=3^\alpha u-1,\quad\beta=\nu_2(y).
\]
If $\beta\geq2$ and $m=(x-1)/2$, then
\begin{equation}
\stepcounter{equation}
\label{eq:Hercher-local-coalescence}
 T^{\alpha+\beta-1}(m)=T^{\alpha+\beta}(x)=y/2^\beta.
\tag{CV.7}
\end{equation}
If $x=T^s(r)$, the prefix contains $c$ odd steps, and $m<r$, then for every
$N=2^kh+r$ the lower lifted start
\[
 Q=3^c2^{k-s-1}h+m
\]
satisfies $Q<N$ and its trajectory coalesces with that of $N$.
\end{proposition}

\begin{proof}
Since $y\equiv0\pmod4$, direct iteration of $m$ gives
\eqref{eq:Hercher-local-coalescence}.  The affine prefix has nonnegative
intercept, positive when $c>0$.  From $m<r$ one has $x<2r$, hence
$3^c/2^s<2$.  This is exactly the strict high-coefficient comparison for
$Q<N$; \eqref{eq:computational-dyadic-prefix} and
\eqref{eq:Hercher-local-coalescence} then prove coalescence.
\end{proof}

The additional Oliveira--Roosendaal example is likewise an identity:
\begin{equation}
\stepcounter{equation}
\label{eq:Oliveira-Roosendaal-coalescence}
 T^6(64q+15)=T^6(64q+14)=T^5(32q+7)=81q+20.
\tag{CV.8}
\end{equation}

\begin{theorem}[finite closure of the recursive high-coefficient check]
\label{thm:Barina-recursive-k34}
Consider \texttt{collatz-sieve}'s restricted recursive closure generated by
the four inverse transports above.  For every invocation arising in its
packed implementation domain $k\leq34$, a returned lower representative
$P<r$ has a uniformly lower full-width lift.  This assertion is finite in
$k$; it does not prove an unbounded recursive invariant.
\end{theorem}

\begin{proof}
A recursive word with $d$ odd steps and $t$ total shortened steps gives
\[
 T^t(P)=L.
\]
For the four primitive words, $(d,t,H)$ is
\[
 (1,1,1),\ (2,3,5),\ (4,6,73),\ (5,7,211),
\]
and the inverse is the exact floor map
\[
 \phi(L)=\frac{2^tL-H}{3^d}
         =\left\lfloor\frac{2^tL}{3^d}\right\rfloor.
\]
A nonempty composite word has
\begin{equation}
\stepcounter{equation}
\label{eq:recursive-rounded-word}
 P=\frac{2^tL-H_w}{3^d}=\frac{2^t}{3^d}L-h_w,
 \qquad0<h_w<\ell(w)\leq d.
\tag{CV.9}
\end{equation}

Suppose $L=T^s(r)$ after $c$ odd steps.  Substitution in
\eqref{eq:computational-dyadic-prefix} shows that the unique
coefficient-correct full lift is
\[
 Q=3^{c-d}2^{k-s+t}h+P,
\]
and $Q<2^kh+r$ for every $h>0$ precisely when $P<r$ and
\begin{equation}
\stepcounter{equation}
\label{eq:recursive-current-coefficient}
 2^t3^{c-d}\leq2^s.
\tag{CV.10}
\end{equation}
For the recursive branch applied after
Proposition~\ref{prop:Hercher-computational-coalescence}, the corresponding
condition, at the preceding prefix $(s_0,c_0)$, is
\begin{equation}
\stepcounter{equation}
\label{eq:recursive-coalescence-coefficient}
 2^t3^{c_0-d}\leq2^{s_0+1}.
\tag{CV.11}
\end{equation}
The C function returns only $P$, retains neither $d$ nor $t$, and tests
neither inequality.

An exact finite reduction supplies the missing check.  In the relevant
exponent boxes the smallest coefficient strictly above one is
\begin{equation}
\stepcounter{equation}
\label{eq:recursive-minimum-unsafe}
 C_{\min}=\frac{3^{12}}{2^{19}}
 =\frac{531441}{524288},
 \qquad C_{\min}-1=\frac{7153}{524288}.
\tag{CV.12}
\end{equation}
If \eqref{eq:recursive-current-coefficient} failed while $P<r$,
equation~\eqref{eq:recursive-rounded-word} and the affine prefix imply
$P(C-1)<h_w<34$ and $r-P\leq33$; hence $P\leq2492$.  For
\eqref{eq:recursive-coalescence-coefficient}, the extra half-step contributes
less than $1/2$, so an unsafe witness has $P\leq2528$ and $r-P\leq34$.

The deterministic certificate exhausts these proved supersets.  In the
current-state branch it checks $7105$ restricted word states, $1557$
identity endpoint pairs, $2983$ nonidentity common-endpoint pairs, and all
$2585$ budget-admissible pairs.  In the coalescence branch it checks $7216$
word states, $1003$ common-endpoint pairs, including $480$ identity and $523$
nonidentity pairs, and all $487$ budget-admissible nonidentity pairs.  Every
required comparison satisfies \eqref{eq:recursive-current-coefficient} or
\eqref{eq:recursive-coalescence-coefficient}.  The finite minimum
\eqref{eq:recursive-minimum-unsafe}, rather than an asymptotic claim, is what
closes the implementation domain.
\end{proof}

\begin{sourceissue}[remaining code and witness boundaries]
\label{sourceissue:Barina-code-boundaries}
The name \texttt{smallest\_predecessor} means the minimum only in the
four-map recursive closure, not among all Collatz predecessors.  At residue
zero, the unsigned expression $(L_{\rm last}-1)/2$ wraps; it does not clear
zero but has no positive-integer predecessor meaning.  The effective
$k\leq34$ domain is the packed-field assertion $b<2^{34}$, not a function
precondition; disabling assertions outside it permits the 34-bit residue
field to alias.  The \texttt{ctzu64} helper assumes that
\texttt{unsigned long} has 64 bits and traps on LLP64.  Finally, a dead-map
bit stores no predecessor word, step counts, or elimination-reason witness,
so the bit alone cannot reconstruct its v3 certificate.
\end{sourceissue}

\begin{proposition}[finite convergence from certified descent]
\label{prop:finite-descent-induction}
If, for every $2\leq n\leq X$, a computation exhibits a finite iterate
$T^{j(n)}(n)<n$, then every positive start through $X$ reaches $1$.
\end{proposition}

\begin{proof}
Strong induction on $n$.  The displayed iterate is a smaller positive start,
so its orbit reaches $1$ by the induction hypothesis; concatenate the two
finite orbit segments.  The base $n=1$ is immediate.
\end{proof}

This proposition explains why sieve elimination and stopping-time descent
can establish a finite range without individually iterating every start to
one.  It does not establish that a reported machine run occurred.

\begin{center}
\begin{tabular}{@{}lll@{}}
\toprule
Source & exact reported range & status in this edition\\
\midrule
Oliveira e Silva (2010) & $n\leq20\cdot2^{58}$ & sourced finite result\\
Barina (2021) & $n<2^{68}$ & sourced finite result\\
Barina (2025) & $n<2^{71}$, 15 Jan.\ 2025 & sourced finite result\\
Project API, 28 Aug.\ 2026 & $2075\cdot2^{60}$ & later dynamic assertion\\
\bottomrule
\end{tabular}
\end{center}

Oliveira e Silva's printed inclusive endpoint is
\[
 20\cdot2^{58}=5764607523034234880.
\]
The twenty directly processed intervals are naturally half-open, but the
boundary itself is even and maps to the already covered $10\cdot2^{58}$;
hence the inclusive theorem follows.  The chapter reports a depth-46 pruned
parity tree with $269949796982$ surviving leaves, $2^{12}$-start blocks,
variable-step tables with $\Delta=14$ and tolerance $\tau=10$, two
independent workers per item, and exact agreement of candidate data and a
cyclic-redundancy checksum.  One mismatch in about 81 CPU-years was
invalidated and recomputed \cite[pp.~189--198]{OliveiraSilva2010}.  The
underlying worker source, completed-work ledger, and duplicate outputs were
not acquired, so the range is not relabelled an F031 reproduction.

The 2021 article reports the $2^{68}$ run from September 2019 to May 2020
\cite[p.~7]{Barina2021}.  The 2025 Table~10 records the strict boundary
$n<2^{71}$ and completion date; its surrounding phrase \emph{up to} does not
change the half-open work units
\[
 [j2^{40},(j+1)2^{40}).
\]
The boundary $2^{71}$ itself is a power of two and independently descends,
but was not in the first $2^{31}$ units.  Reported CPU/GPU years, throughput,
and speedups are measurements on named configurations, not
hardware-independent complexity theorems.  The page-2 throughput convention
also counts starts eliminated by sieve certificates whose trajectories were
not individually run to one \cite{Barina2025}.

The dated API snapshot reports
\[
 2075\cdot2^{60}=2392312122059207475200
 \approx2^{71.0188956211},
\]
with the first incomplete unit beginning at
\[
 2175975546\cdot2^{40}=2392510414583230365696.
\]
It is later than the paper, is not a completed-work ledger, and is not
backdated into the 2025 theorem \cite{BarinaStatus2026}.

\begin{sourceissue}[the public-code validation boundary]
\label{sourceissue:Barina-public-code-validation}
The 2025 paper pins the repository \nolinkurl{xbarin02/collatz} at commit
\linebreak[3]\nolinkurl{53c2a0608075d6fe3f10cc6eeeaf50e400c86338}.
It does not pin the
separate \texttt{collatz-sieve} dependency.  The bootstrap performs unpinned
clone-or-pull operations.  The edition's sieve commit and map hashes are
therefore separately attributed reconstruction pins, not historical facts
invented for the run \cite{BarinaCode2025}.

The default GPU kernel is the Git symlink \texttt{kernel.cl} to
\texttt{kernel32.cl}.  A GPU overflow zeroes a checksum lane; the host aborts
the unit and mclient attempts a CPU rerun.  GMP continuation exists only when
that CPU worker was compiled with \texttt{\_USE\_GMP}.  The generic bootstrap
does so, whereas the Lumi and Meta GPU-only scripts build no CPU worker.
Thus GPU--CPU--GMP is a conditional deployment chain, not GPU
multiprecision or a universal local fallback.

The optional/test \texttt{kernel32-precalc.cl} has a sieve-dependent
\texttt{continue} immediately before a work-group barrier.  A multi-lane
group with differing sieve bits violates barrier convergence.  This defect
is scoped to that variant and is not transferred to the default kernel.

The server first rejects stale returns with the wrong stored client ID and
rejects zero or non-whitelisted checksum prefixes.  For an accepted current
return, however, it marks the unit complete before comparing a previous
checksum, then logs and overwrites a mismatch; maximum offsets are handled in
the same order.  Mclient transmits the $\alpha$-sum and discards the
$\beta$-sum.  The verification utility recomputes only the trajectory maximum
at the stored offset, not the work unit.  These are consistency signals, not
Oliveira e Silva's independent duplicate-run protocol or a cryptographic
certificate.

No acquired immutable ledger binds every unit to its code tree, map hash,
binary, compiler, OpenCL environment, and return value.  Selected regression
tests and public source permit code-level audit; they do not independently
authenticate the full historical $2^{71}$ computation.
\end{sourceissue}

\begin{proposition}[shortened and unshortened maximum excursions]
\label{prop:shortened-unshortened-maxima}
For an odd $n>1$, if the orbit maxima are finite, then
\begin{equation}
\stepcounter{equation}
\label{eq:shortened-unshortened-maxima}
 M_U(n)=2M_T(n).
\tag{CV.13}
\end{equation}
The nontrivial record-start sequence is the same under both conventions,
although its maximum values differ by the factor two.
\end{proposition}

\begin{proof}
An odd shortened state $x>1$ cannot be the shortened maximum because
$(3x+1)/2>x$.  A maximal shortened state cannot have an even predecessor,
which would be twice as large.  It is therefore the output $y=(3x+1)/2$ of
an odd predecessor, and the unshortened orbit contains the intermediate
state $2y$.  Every state omitted by shortening is twice a shortened state,
proving \eqref{eq:shortened-unshortened-maxima}.  No even $n>2$ begins a new
unshortened record, because the lower odd start $n-1$ has first image
$3n-2>n$.  Thus record starts are odd and multiplication of every maximum by
two preserves their order.
\end{proof}

The 2025 abstract says four new records, while Section~6 lists five and calls
them new.  The first listed start, $274133054632352106267$, already appears
in the 2021 paper.  The exact reconciliation is four records new after 2021
and five cumulative project records; Section~6's phrase \emph{five new} is
imprecise rather than an unexplained numerical contradiction
\cite[p.~7]{Barina2021}\cite[pp.~12--13]{Barina2025}.  Finite record data are
consistent with the Lagarias--Weiss prediction but do not prove its
asymptotic form.

\begin{proposition}[Roosendaal's modular maximum theorem]
\label{prop:Roosendaal-mod36}
Let $N>2$ be odd.  If its unshortened maximum $q=M_U(N)$ is finite and
$q>3N+1$, then
\[
 q\equiv16\pmod{36}.
\]
\end{proposition}

\begin{proof}
The maximum $q$ is even and cannot be $2\bmod4$, since then the next two
states $q/2$ and $3q/2+1$ exceed it.  Because $q>3N+1$, it is neither the
start nor its first odd output; halving cannot create a maximum, so
$q=3p+1$ for an odd predecessor and $q\equiv1\bmod3$.  The remaining classes
modulo $36$ are $4,16,28$.  If $q=36a+4$, its forced preceding states include
$48a+4>q$.  If $q=36a+28$, they include $48a+36>q$.  Only $16\bmod36$
remains \cite{Roosendaal2026}.
\end{proof}

The page's following sentence applies this congruence to $P_i$, but the
theorem concerns $M_U(P_i)$, not the starting values.  Its Conjecture~3 also
prints $N<1$ against the surrounding $N>1$ domain.  Neither defect is
silently normalized.  Theorems~2--3 each begin with a residue-bound
hypothesis that is not proved on the page; this edition does not establish
that hypothesis and therefore does not import either conclusion.

Oliveira e Silva additionally studies the one-branch-per-iterate map
\begin{equation}
\stepcounter{equation}
\label{eq:Oliveira-generalized-five}
 C(n)=
 \begin{cases}
  n/2,&n\equiv0\pmod2,\\
  n/3,&n\not\equiv0\pmod2,\ n\equiv0\pmod3,\\
  5n+1,&\gcd(n,6)=1.
 \end{cases}
\tag{CV.14}
\end{equation}
Its clock is not the $T$- or $U$-clock.

\begin{sourceissue}[Oliveira e Silva's Proposition~2 is false as printed]
\label{sourceissue:Oliveira-proposition-two}
The proposition asks the first $i$ iterates of $m_i$ to determine counts
$k_2,k_3,k_5$, puts
\[
 r_i=2^{k_2}3^{k_3},\qquad
 m_i=n_0\bmod r_i,\qquad n_i=[n_0/r_i],
\]
and asserts
\[
 C^i(n_0)=5^{k_5}n_i+C^i(m_i).
\]
At $i=1,n_0=1$, $r_i=1$ forces $m_i=0$, whose division-by-two branch demands
$r_i=2$; $r_i=2$ or $3$ forces $m_i=1$, whose $5n+1$ branch demands
$r_i=1$.  Thus the circular definition has no fixed point.  If one instead
takes the branch of $n_0$ first, then $r_i=1,m_i=0,n_i=1$ and the asserted
equation reads $6=5$.  No domain qualification repairs it
\cite[p.~200]{OliveiraSilva2010}.

The chapter defines $[x]$ as the largest integer not smaller than $x$, which
is ceiling language.  Proposition~1's unique dyadic decomposition requires
$\lfloor n_0/2^i\rfloor$.  Replacing the bracket by a floor repairs that use;
it does not repair Proposition~2.
\end{sourceissue}

\begin{proposition}[the corrected $6^i$ branch-cylinder coordinate]
\label{prop:Oliveira-branch-cylinder}
For $i\geq0$, put
\[
 a_i=n_0\bmod6^i,\qquad
 Q_i=\left\lfloor\frac{n_0}{6^i}\right\rfloor.
\]
Let $k_2,k_3,k_5$ be the branch counts in the first $i$ iterates of $a_i$,
set $R_i=2^{k_2}3^{k_3}$, and
\[
 N_i=\frac{n_0-a_i}{R_i}=\frac{6^i}{R_i}Q_i.
\]
Then $N_i$ is an integer, $n_0$ and $a_i$ have the same first $i$ branches,
and
\begin{equation}
\stepcounter{equation}
\label{eq:Oliveira-branch-cylinder}
 C^i(n_0)=5^{k_5}N_i+C^i(a_i).
\tag{CV.15}
\end{equation}
The representative $a_i$ is a branch-cylinder representative, generally not
the canonical remainder modulo $R_i$.
\end{proposition}

\begin{proof}
If $x\equiv y\pmod{6^j}$, the two points take the same branch.  After
division by two their difference is divisible by $2^{j-1}3^j$, after
division by three by $2^j3^{j-1}$, and after $5x+1$ by $6^j$; every image
pair is congruent modulo $6^{j-1}$.  Induction gives the common length-$i$
branch word.

For a fixed word $w$, induction on its branches gives an integer $B_w$ with
\[
 C^i(x)=\frac{5^{k_5}x+B_w}{R_i}
\]
throughout its cylinder.  A division branch multiplies $R_i$ by two or
three, while a $5x+1$ branch replaces $B_w$ by $5B_w+R_i$.  Since
$R_i\mid6^i$, subtracting the formula at $n_0$ and $a_i$ proves
\eqref{eq:Oliveira-branch-cylinder}.  Reducing $a_i$ modulo $R_i$ is valid
only after a separate proof that the reduced point retains the same branch
word.
\end{proof}

The source reports exhaustive computations through $10^{16}$ for
\eqref{eq:Oliveira-generalized-five} and $10^{14}$ for the analogous
$C_{7,\{2,3,5\}}$, taking about five and two CPU-months respectively on one
450 MHz Pentium~III \cite[pp.~200--204]{OliveiraSilva2010}.  It does not say
that the replicated-worker protocol of the $3x+1$ run was used here.  Its
finite pruning counts, affine tables, transition matrix, and stationary
vector are reproducible finite calculations.  Section~4 expressly calls the
passage from a Chernoff upper bound to an asymptotic equality non-rigorous;
the printed values
\[
 \rho_5\approx1.442762198279,\qquad
 \rho_7\approx2.215508001593
\]
remain conjectured record-growth exponents, not proved rates.

\begin{nonclaim}
The deterministic F031 certificate pins the exact manifestations, code
trees, bundles, selected maps, and route artifacts; reproduces the basic
$2^{16}$ sieve byte for byte; proves
Theorem~\ref{thm:Barina-recursive-k34}; and checks the map, fibre, affine,
transport, maximum, modular, and corrected branch-cylinder identities above.
It does not reproduce a historical distributed run, authenticate a
completed-work database, independently regenerate the full $2^{24}$ or
$2^{34}$ maps, certify hardware or compilers, justify the Markov asymptotics,
or prove universal convergence. The eight legacy unresolved-dependency records
are audit-schema records rather than mathematical claims; this section does not
use them to preserve or conditionalize an unproved conclusion.
\end{nonclaim}


\begin{nonclaim}
Crandall's Theorem~6.1 proves only that
$\pi(X)>X^c$ for some $c>0$ and all sufficiently large $X$, where $\pi$
counts positive odd $m>1$ known to reach $1$ in accelerated time; the
repaired source argument permits the edition witness $c=1/100$
(Theorem~\ref{thm:Crandall-power-law-count}).  It is not a positive-density
theorem.  Tao's
theorem gives logarithmic-density-one descent below every prescribed
divergent threshold, but does not say those orbits reach $1$.  Neither result,
nor their conjunction, yields logarithmic-density-one convergence to $1$.
\end{nonclaim}
